% ============================================================
%  preamble.tex --- styling, environments, macros
%
%  Shared by BOTH editions. main-en.tex and main-pl.tex each define \booklang
%  before % ============================================================
%  preamble.tex --- styling, environments, macros
%
%  Shared by BOTH editions. main-en.tex and main-pl.tex each define \booklang
%  before % ============================================================
%  preamble.tex --- styling, environments, macros
%
%  Shared by BOTH editions. main-en.tex and main-pl.tex each define \booklang
%  before % ============================================================
%  preamble.tex --- styling, environments, macros
%
%  Shared by BOTH editions. main-en.tex and main-pl.tex each define \booklang
%  before \input{preamble}; nothing else differs between them at this level.
%  Every user-visible string lives in lang/en.tex or lang/pl.tex, so a label
%  that appears in one edition and not the other is a missing macro rather
%  than a silent divergence.
%
%  Lineage. This book is the third in a family and takes one thing from each
%  of the other two:
%
%    * from Mathematics from Zero for the AI Engineer: the bilingual single
%      source (\booklang, lang/, one body.tex), the rule that every number is
%      computed and reaches the page as \val{}, and the programmed-learning
%      method -- outcomes on entry, a question before every answer, a
%      misconception elicited before it is named, a self-check generated from
%      the outcomes, and checkpoint answers collected at the back;
%    * from Python for .NET Engineers: every listing is a FILE that CI runs,
%      every lab is a starter that must fail and a solution that must pass,
%      and a console transcript is a computed number that happens to be
%      verbatim.
%
%  What is new here is \plotfig. A book about chaos cannot be taught in boxes
%  and arrows: the bifurcation diagram and the Lorenz attractor ARE the
%  argument, so the plots are generated by code/plots/ from the same library
%  the listings use, per language, and never drawn by hand.
% ============================================================

% \booklang must be set by the main file. Default to English rather than
% failing, so a stray \input{preamble} still compiles and says what it did.
\providecommand{\booklang}{en}

\usepackage{etoolbox}   % loaded first: the language switch below needs it

% ---------- Page geometry ----------
% ONE format: A4, 12pt, ONE-SIDED. The book is read on a screen, usually with
% an editor open beside it, so there is no recto, no verso, no gutter and no
% blank leaf before a chapter. The measure is about seventy-five characters
% of prose. At \footnotesize a listing line of 79 characters fits without
% wrapping, and 79 is the width every file under code/ is held to by
% tools/check_structure.py --lines.
\usepackage[
  a4paper,
  left=3.1cm, right=3.1cm, top=2.6cm, bottom=3.0cm,
  head=26pt, headsep=0.8cm, footskip=1.4cm
]{geometry}

% ---------- Encoding & fonts ----------
\usepackage[T1]{fontenc}
\usepackage[utf8]{inputenc}
\IfFileExists{lmodern.sty}{\usepackage{lmodern}}{}
% newtxtext takes its TS1 symbols from TeX Gyre Termes (Debian: tex-gyre).
% A machine with newtx and without tex-gyre dies on the copyright page. CI's
% full TeX Live is the reference; locally install tex-gyre as well.
\IfFileExists{newtxtext.sty}{\usepackage{newtxtext}}{}
% amsmath before newtxmath; amssymb only when newtxmath is absent, because
% loading both is a fatal clash (\Bbbk already defined) that is invisible on
% a machine without newtx.
\usepackage{amsmath}
\IfFileExists{newtxmath.sty}{\usepackage{newtxmath}}{\usepackage{amssymb}}
\IfFileExists{inconsolata.sty}{\usepackage[varqu,scaled=0.95]{inconsolata}}{}
% Without lmodern, microtype runs with font expansion off, which is what
% resolves a marginal overfull line -- so a container short of lmodern
% disagrees with CI about line breaks. Install lmodern locally too.
\IfFileExists{lmodern.sty}{\usepackage{microtype}}{\usepackage[expansion=false]{microtype}}

% upquote: in T1 the input ' is a right single quote, and a listing that sets
% f('x') with curly quotes is a SyntaxError when typed back in.
\IfFileExists{upquote.sty}{\usepackage{upquote}}{}
% And its twin: in T1 the pair -- is an en-dash ligature in the typewriter
% family too, so \code{--flag} would print a flag nobody can type. Scoped to
% tt*, so prose keeps its own dashes.
\DisableLigatures[-]{encoding = T1, family = tt*}

% ---------- Language ----------
% babel with a missing language file is FATAL, so each .ldf is probed first.
\IfFileExists{babel.sty}{%
  \IfFileExists{polish.ldf}{%
    \IfFileExists{british.ldf}{%
      \ifdefstring{\booklang}{pl}%
        {\usepackage[british,main=polish]{babel}}%
        {\usepackage[polish,main=british]{babel}}%
    }{\usepackage[polish]{babel}}%
  }{%
    \IfFileExists{british.ldf}{\usepackage[british]{babel}}{}%
  }%
}{}
% babel-polish makes " an active shorthand, which is a category-code accident
% waiting to happen inside listings. Polish quotation marks come from
% \enquote, which csquotes sets per language.
\ifdefstring{\booklang}{pl}{\AtBeginDocument{\shorthandoff{"}}}{}

% ---------- Core packages ----------
\usepackage{graphicx}
\usepackage{xcolor}
\usepackage{booktabs}
\usepackage{tabularx}
\usepackage{longtable}
\usepackage{array}
\usepackage{enumitem}
\usepackage{listings}
\usepackage[most]{tcolorbox}
\usepackage{fancyhdr}
\usepackage{titlesec}
\usepackage{caption}
\usepackage{imakeidx}
\IfFileExists{csquotes.sty}{\usepackage[autostyle=true]{csquotes}}{%
  \providecommand{\enquote}[1]{``##1''}}
\usepackage[hidelinks]{hyperref}

% ---------- Palette ----------
% The family's palette, so the three books read as a set. Two colours are
% this book's own: the lab (green) and the stop-and-think box (violet).
\definecolor{mblue}{HTML}{1F4E79}
\definecolor{mteal}{HTML}{0E7C7B}
\definecolor{mamber}{HTML}{B26A00}
\definecolor{mred}{HTML}{9B2C2C}
\definecolor{mviolet}{HTML}{7B4397}
\definecolor{mgreen}{HTML}{3B7A3B}
\definecolor{mgrey}{HTML}{5A5A5A}
\definecolor{mlight}{HTML}{F4F6F8}
\definecolor{mfaint}{HTML}{FBFCFD}
\definecolor{codebg}{HTML}{FAFAFA}
\definecolor{codeframe}{HTML}{D8DEE4}
\definecolor{kw}{HTML}{0B5394}
\definecolor{str}{HTML}{9B2C2C}
\definecolor{cmt}{HTML}{4F7A4F}
\definecolor{typ}{HTML}{0E7C7B}
\definecolor{dec}{HTML}{7B4397}

\hypersetup{
  colorlinks=true,
  linkcolor=mblue,
  urlcolor=mteal,
  citecolor=mblue,
  pdfauthor={Konrad Cinkusz}
}

% \ifpl{Polish}{English} -- for the handful of places where a string is
% structural rather than content. Robust, because a fragile \ifpl inside a
% moving argument fails on the next run in exactly one of the two editions.
\DeclareRobustCommand{\ifpl}[2]{\ifdefstring{\booklang}{pl}{#1}{#2}}
\makeatletter
\pdfstringdefDisableCommands{%
  \ifdefstring{\booklang}{pl}{\let\ifpl\@firstoftwo}{\let\ifpl\@secondoftwo}%
  % hyperref drops \color from a bookmark and keeps its ARGUMENT, so a title
  % carrying \code{} would bookmark as "mtealx". Gobble the colour name.
  \def\color#1{}%
}
\makeatother

% ---------- Language strings ----------
\input{lang/\booklang}

% The PDF's subject, keywords and language read \lblPdf... from the language
% file, so they cannot sit in the colour block above, which runs first. Two
% blocks is the ordering, not an oversight.
\hypersetup{
  pdfsubject={\lblPdfSubject},
  pdfkeywords={\lblPdfKeywords},
  pdflang={\lblPdfLang}
}

% ---------- Headings ----------
\titleformat{\chapter}[display]
  {\normalfont\huge\bfseries\color{mblue}\raggedright}
  {\normalfont\normalsize\color{mgrey}\MakeUppercase{\chaptertitlename}\ \thechapter}
  {8pt}{\Huge\raggedright}
\titlespacing*{\chapter}{0pt}{10pt}{28pt}
\titleformat{\section}{\normalfont\Large\bfseries\color{mblue}\raggedright}{\thesection}{0.8em}{}
\titleformat{\subsection}{\normalfont\large\bfseries\color{mgrey}\raggedright}{\thesubsection}{0.7em}{}

\setcounter{tocdepth}{1}
\setcounter{secnumdepth}{2}

\pagestyle{fancy}
\fancyhf{}
\fancyhead[L]{\small\nouppercase{\leftmark}}
\fancyhead[R]{\small\thepage}
\renewcommand{\headrulewidth}{0.4pt}
\fancypagestyle{plain}{\fancyhf{}\fancyfoot[R]{\small\thepage}%
  \renewcommand{\headrulewidth}{0pt}}
% After \pagestyle{fancy}, which installs its own \chaptermark.
\renewcommand{\chaptermark}[1]{\markboth{\thechapter.\ #1}{}}

% ---------- Mathematics ----------
% The handful of notations the book uses, as macros so both editions set the
% same symbol. The book writes \ln for the natural logarithm and \log_{10}
% where the base is ten; a bare \log is ambiguous between the two readerships
% and is not used.
\newcommand{\R}{\mathbb{R}}
\newcommand{\dd}{\,\mathrm{d}}
\newcommand{\abs}[1]{\left\lvert #1\right\rvert}
\newcommand{\iu}{\mathrm{i}}

% ============================================================
%  CODE LISTINGS
%  Every listing in a chapter is a FILE under code/, pulled in with \pyfile or
%  \pyregion, and CI runs it on the pinned interpreter. An in-source
%  \begin{python} block is allowed for a fragment of three or four lines that
%  shows a shape rather than runs.
% ============================================================
\lstdefinelanguage{pythonx}{
  language=Python,
  morekeywords={async,await,match,case,nonlocal,yield,type},
  morekeywords=[2]{
    orbit,iterate,logistic,rk4_step,euler_step,integrate,lorenz,henon,
    lyapunov,separation,escape_time,box_count
  },
  sensitive=true
}

\lstdefinestyle{bookcode}{
  basicstyle=\ttfamily\footnotesize,
  keywordstyle=\color{kw}\bfseries,
  keywordstyle=[2]\color{typ},
  stringstyle=\color{str},
  % Colour only: inconsolata has no italic, and an italic comment style is
  % silently substituted and logged on every build.
  commentstyle=\color{cmt},
  numbers=left,
  numberstyle=\ttfamily\tiny\color{mgrey},
  numbersep=8pt,
  backgroundcolor=\color{codebg},
  frame=single,
  rulecolor=\color{codeframe},
  framesep=6pt,
  breaklines=true,
  breakatwhitespace=false,
  postbreak=\mbox{\textcolor{mgrey}{$\hookrightarrow$}\space},
  showstringspaces=false,
  tabsize=4,
  columns=fullflexible,
  keepspaces=true,
  captionpos=b,
  aboveskip=1.2em,
  belowskip=1.2em,
  % Do NOT add a mapping for U+00A0: it aborts the run with a message that
  % names neither the character nor this line. Listings are ASCII, and
  % parity's C13 checks it.
  literate={ł}{{\l}}1 {Ł}{{\L}}1 {ą}{{\k a}}1 {ę}{{\k e}}1
           {ś}{{\'s}}1 {ć}{{\'c}}1 {ż}{{\.z}}1 {ź}{{\'z}}1 {ó}{{\'o}}1 {ń}{{\'n}}1
           {Ą}{{\k A}}1 {Ę}{{\k E}}1 {Ś}{{\'S}}1 {Ć}{{\'C}}1
           {Ż}{{\.Z}}1 {Ź}{{\'Z}}1 {Ó}{{\'O}}1 {Ń}{{\'N}}1
           {—}{{-{}-}}2 {–}{{-}}1 {‘}{{'}}1 {’}{{'}}1
           {“}{{"}}1 {”}{{"}}1 {°}{{\textdegree}}1
}
\lstset{style=bookcode}

\lstnewenvironment{python}[1][]
  {\lstset{language=pythonx,style=bookcode,#1}}
  {}
\lstnewenvironment{shellcmd}[1][]
  {\lstset{language=bash,style=bookcode,numbers=none,keywordstyle=\color{mteal}\bfseries,#1}}
  {}

% The path is printed above every file listing, in the form the reader opens
% in the editor. \nobreak keeps it on the page of the listing it names.
% The path must not be stranded at the foot of a page with its listing
% overleaf: \nobreak alone does not hold, because listings opens with its own
% vertical material. \needspace asks for room for the path and the first few
% lines of code, and turns the page first when there is not.
% \Needspace (capital N) measures \pagegoal-\pagetotal exactly; lower-case
% \needspace works through glue and penalties and was measured, by the
% Chapter 11 pass, turning pages early and leaving gaps. It must be used
% between paragraphs, which \listingpath's own \par guarantees. Chapters use
% \Needspace* too, so the fallback for a machine without the package defines
% both forms (the star swallowed with its argument).
\makeatletter
\IfFileExists{needspace.sty}{\usepackage{needspace}}{%
  \providecommand{\needspace}[1]{}%
  \providecommand{\Needspace}{\@ifstar\@gobble\@gobble}}
\makeatother
\newcommand{\listingpath}[1]{%
  \par\medskip\Needspace{6\baselineskip}\noindent
  {\ttfamily\footnotesize\color{mgrey}\detokenize{#1}}%
  \par\nobreak}

% Usage: \pyfile{code/ch05/cobweb.py}{Caption}{lst:label}
\newcommand{\pyfile}[3]{%
  \listingpath{#1}%
  \lstinputlisting[language=pythonx,style=bookcode,aboveskip=0.3em,
    caption={#2},label={#3}]{#1}%
}

% A named region of a file, delimited by `# --8<-- [start:name]` and
% `# --8<-- [end:name]`. Matched through listings' range prefix and suffix
% keys -- the other obvious form matches nothing and prints the WHOLE file
% without a warning (measured in the sibling book). A region name is a word,
% never a number. The line numbers printed are the file's own.
% Usage: \pyregion{code/ch05/cobweb.py}{step}{Caption}{lst:label}
\lstset{
  rangebeginprefix=\#\ --8<--\ [start:, rangebeginsuffix=],
  rangeendprefix=\#\ --8<--\ [end:, rangeendsuffix=],
  includerangemarker=false}
\newcommand{\pyregion}[4]{%
  \listingpath{#1}%
  \lstinputlisting[language=pythonx,style=bookcode,aboveskip=0.3em,
    linerange={#2-#2},caption={#3},label={#4}]{#1}%
}

% A program's output, written by code/measure/transcripts.py. There is no
% typed form and there must not be: \val{} does not expand in a verbatim
% body, and a transcript nobody ran looks exactly like one that was.
\newcommand{\transcript}[1]{%
  \par\smallskip
  \IfFileExists{figures/transcripts/#1.txt}{%
    \lstinputlisting[style=bookcode,numbers=none,basicstyle=\ttfamily\footnotesize,
      keywordstyle=\color{black},backgroundcolor=\color{white},
      language={}]{figures/transcripts/#1.txt}%
  }{%
    \begin{tcolorbox}[enhanced, sharp corners, width=\linewidth,
      colback=mlight, colframe=mgrey, boxrule=0.6pt,
      left=8pt, right=8pt, top=6pt, bottom=6pt]
      {\bfseries\color{mgrey}\small \lblTranscriptMissing}\\[2pt]
      {\ttfamily\scriptsize\color{mgrey} figures/transcripts/#1.txt}
    \end{tcolorbox}%
  }%
  \par\smallskip
}

% Inline literals. Underscores inside these must be written \_.
\newcommand{\code}[1]{\texttt{\small #1}}
\newcommand{\pkg}[1]{\texttt{\small\color{mblue}#1}}
\newcommand{\api}[1]{\texttt{\small\color{mteal}#1}}

% ============================================================
%  ADMONITIONS
%  A small, fixed vocabulary. `tinted' may be read in passing; `outlined' is
%  a block the reader must stop at. A third treatment would mean the first two
%  are not carrying their meaning.
% ============================================================
\tcbset{
  cfz tinted/.style={
    enhanced, breakable, sharp corners,
    lines before break=4,
    colback=mlight, boxrule=0pt, leftrule=3pt,
    left=8pt, right=8pt, top=6pt, bottom=6pt,
    before skip=13pt, after skip=13pt,
    fonttitle=\bfseries\small,
    coltitle=black, colbacktitle=mlight,
    attach title to upper={\par\smallskip},
  },
  cfz outlined/.style={
    enhanced, breakable, sharp corners,
    % A box broken straight after its title is a heading stranded at the foot
    % of a page with its question overleaf. Four lines is title plus three.
    lines before break=4,
    colback=white, boxrule=0.6pt,
    left=8pt, right=8pt, top=6pt, bottom=6pt,
    before skip=13pt, after skip=13pt,
    fonttitle=\bfseries\small,
    coltitle=black, colbacktitle=white,
    attach title to upper={\par\smallskip},
  },
}

% Read in passing.
\newtcolorbox{note}[1][]{cfz tinted, colframe=mblue, title={\lblNote}, #1}
\newtcolorbox{warning}[1][]{cfz tinted, colframe=mred, title={\lblWarning}, #1}
% Where you meet this in the world. If it cannot name a specific system,
% instrument, paper or event, it is prose and not this box.
\newtcolorbox{worldbox}[1][]{cfz tinted, colframe=mteal, title={\lblWorld}, #1}
% What is stated and not proved here, and where the proof lives.
\newtcolorbox{rigourbox}[1][]{cfz tinted, colframe=mgrey, title={\lblRigour}, #1}
% The companion library the reader builds, under code/src/chaoslab/.
\newtcolorbox{projectbox}[1][]{cfz tinted, colframe=mamber, title={\lblProject}, #1}

% Stop and read.
% The misconception, named AFTER the reader has been asked the question it
% gets wrong. Appendix B is the catalogue; a trap box is one entry in place.
\newtcolorbox{trapbox}[1][]{cfz outlined, colframe=mred, title={\lblTrap}, #1}
% A listing that has not been run on the pinned versions. Counted by
% `make debt'; nothing ships to a reader with one attached.
\newtcolorbox{verifybox}[1][]{cfz outlined, colframe=mgrey, title={\lblVerify}, #1}
\newtcolorbox{labbox}[1][]{cfz outlined, colframe=mgreen, title={\lblLab}, #1}

% ============================================================
%  THE METHOD --- questions before answers
%  ---------------------------------------------------------
%  The unit of teaching is a question the reader answers BEFORE reading the
%  answer. That is retrieval practice and errorful generation, the two
%  findings the programmed-learning format is defended by, and it is the one
%  thing the sibling mathematics book does that this one keeps whole.
%
%  \begin{think} ... \end{think}   the question, numbered per chapter, ending
%                                  with the instruction to answer before
%                                  reading on;
%  \reveal{short answer}           a centred box with the answer, which is the
%                                  first thing after the question -- the edge
%                                  the reader covers with a hand, or scrolls
%                                  to;
%  \begin{revealblock} ... \end{revealblock}  the same, for a worked answer.
%
%  tools/check_structure.py --think fails when a think box is not followed by
%  a reveal before the next think box or section, and when a reveal answers
%  nothing. That is a property of the source, so it is checked mechanically.
%
%  The instruction says "before you read on", never "before you turn over":
%  the two editions paginate differently, so the second is true in one of them.
% ============================================================
\newcounter{think}[chapter]
\renewcommand{\thethink}{\thechapter.\arabic{think}}
\newtcolorbox{thinkbox}[1][]{cfz outlined, colframe=mviolet, #1}
\NewDocumentEnvironment{think}{}{%
  \refstepcounter{think}%
  \begin{thinkbox}[title={\lblThink~\thethink}]%
}{%
  \par\smallskip
  {\raggedleft\itshape\small\color{mviolet}\lblThinkCue\par}%
  \end{thinkbox}%
}
\tcbset{
  cfz answer/.style={
    enhanced, sharp corners,
    colback=mfaint, colframe=mviolet!55, boxrule=0.5pt,
    before skip=4pt, after skip=10pt,
  },
}
\newcommand{\reveal}[1]{%
  \par\nobreak
  \begin{tcolorbox}[cfz answer, unbreakable,
    left=8pt, right=8pt, top=5pt, bottom=5pt,
    width=0.86\linewidth, center, halign=flush center]
    {\small\bfseries\color{mviolet}\lblAnswerMark}\enspace #1
  \end{tcolorbox}%
  \nobreak\noindent\ignorespaces
}
\NewDocumentEnvironment{revealblock}{}{%
  \par\nopagebreak
  \begin{tcolorbox}[cfz answer, breakable,
    left=10pt, right=10pt, top=6pt, bottom=6pt]
  {\small\bfseries\color{mviolet}\lblAnswerMark}\par\smallskip
}{%
  \end{tcolorbox}%
  \nopagebreak\par\noindent\ignorespaces
}

% ---- Outcomes on entry, and "Can you?" generated from them ----------------
% The self-check at the end of a chapter IS the list of outcomes at its
% start, one for one, because it is built from the same tokens: each
% \outcome both prints its item and appends a row to \cfz@canrows. A \loop
% inside a tabularx body does not survive alignment scanning, which is why
% the rows are accumulated at declaration time (the sibling book's finding).
\makeatletter
\newcommand{\cfz@canrows}{}
\newcommand{\cfz@ratebox}{%
  \fbox{\rule{0pt}{0.9ex}\rule{0.9ex}{0pt}}}
\newcommand{\cfz@rates}{%
  {\footnotesize\color{mgrey}%
   1\,\cfz@ratebox\ 2\,\cfz@ratebox\ 3\,\cfz@ratebox\ 4\,\cfz@ratebox\ 5\,\cfz@ratebox}}
\NewDocumentEnvironment{outcomes}{}{%
  \gdef\cfz@canrows{}%
  \begin{tcolorbox}[cfz tinted, colframe=mblue, title={\lblOutcomes}]%
  \lblOutcomesLead
  \begin{itemize}[leftmargin=1.4em, itemsep=2pt, topsep=4pt]%
}{%
  \end{itemize}%
  \end{tcolorbox}%
}
\newcommand{\outcome}[1]{%
  \item #1%
  \gappto\cfz@canrows{#1 & \cfz@rates \\ \addlinespace[3pt]}%
}
\newcommand{\canyou}{%
  \section*{\lblCanYou}%
  \phantomsection\addcontentsline{toc}{section}{\lblCanYou}%
  \noindent\lblCanYouIntro\par\medskip
  \ifdefempty{\cfz@canrows}{%
    \noindent{\color{mred}\lblNoOutcomes}\par}{%
  {\small
  \noindent\begin{tabularx}{\linewidth}{@{}>{\raggedright\arraybackslash}X r@{}}
    \toprule
    \cfz@canrows
    \bottomrule
  \end{tabularx}\par}}%
}
\makeatother

% ---- Summary -------------------------------------------------------------
\newtcolorbox{summarybox}[1][]{cfz tinted, colframe=mblue, title={\lblSummary}, #1}

% ---- Checkpoint questions, answered at the back --------------------------
% Answers are written AT THE QUESTION with \answerto{} and printed in
% Appendix A. They are carried there by a global token list rather than an
% auxiliary file, because an answer contains mathematics and \val{}, and
% material written to an .aux is expanded at shipout, where fragile commands
% break with errors naming neither the answer nor the chapter.
% \include is sequential within a run, so the list is complete by the time
% the appendix is typeset.
\makeatletter
\newcommand{\cfz@answers}{}
\newcommand{\cfz@lastansch}{0}
\NewDocumentEnvironment{checkpoint}{}{%
  \section*{\lblCheckpoint}%
  \phantomsection\addcontentsline{toc}{section}{\lblCheckpoint}%
  \noindent\lblCheckpointIntro\par
  \begin{enumerate}[label=\textbf{\arabic*.}, leftmargin=2em, itemsep=4pt]%
}{%
  \end{enumerate}%
}
\newcommand{\answerto}[1]{%
  \ifnum\value{chapter}=\cfz@lastansch\relax\else
    \xdef\cfz@lastansch{\arabic{chapter}}%
    \xappto\cfz@answers{\noexpand\cfz@anschapter{\thechapter}}%
  \fi
  \xappto\cfz@answers{\noexpand\cfz@ansitem{\arabic{enumi}}}%
  \gappto\cfz@answers{{#1}}%
}
\newcommand{\cfz@anschapter}[1]{%
  \par\medskip\noindent{\bfseries\color{mblue}\chaptername~#1}\par\nobreak\smallskip}
\newcommand{\cfz@ansitem}[2]{%
  \par\noindent\hangindent=2em\hangafter=1
  \makebox[2em][l]{\textbf{#1.}}#2\par}
\newcommand{\answersbody}{%
  \ifdefempty{\cfz@answers}{\noindent{\color{mred}\lblNoAnswers}\par}{%
    {\raggedright\cfz@answers}}%
}
\makeatother

% ============================================================
%  LABS
%  \begin{lab}{key}{title} ... \end{lab}
%
%  A lab is not a paragraph but a FILE the reader opens: a starter under
%  code/labs/chNN/<key>.py that must FAIL its test, a reference solution under
%  code/labs/chNN/solutions/, and test_<key>.py beside it. The box prints the
%  paths and the one command that runs the test. An environment rather than a
%  macro, because a macro argument cannot hold a verbatim fragment.
%  tools/check_structure.py --labs fails when a file is missing or the key's
%  chapter number disagrees with its chapter.
% ============================================================
\newcounter{lab}[chapter]
\renewcommand{\thelab}{\thechapter.\arabic{lab}}
\makeatletter
\newcommand{\labdir}{ch\two@digits{\value{chapter}}}
\newcommand{\listoflabs}{%
  \section*{\lblLabManifest}%
  \phantomsection\addcontentsline{toc}{section}{\lblLabManifest}%
  \begingroup\c@tocdepth=2\relax\@starttoc{lab}\endgroup%
}
\makeatother
\NewDocumentEnvironment{lab}{mm}{%
  \refstepcounter{lab}%
  \addcontentsline{lab}{subsection}{\protect\numberline{\thelab}\texttt{\protect\detokenize{#1}} --- #2}%
  \begin{labbox}[title={\lblLab~\thelab\ \dash{} #2}]%
}{%
  \par\medskip
  {\small\setlength{\parindent}{0pt}%
  \lblLabStarter: \texttt{code/labs/\labdir/\detokenize{#1}.py}\\
  \lblLabTest: \texttt{code/labs/\labdir/test\_\detokenize{#1}.py}\\
  \lblLabRun: \texttt{cd code \&\& uv run pytest -k \detokenize{#1}}\par}%
  \end{labbox}%
}

% ============================================================
%  FIGURES --- plots and diagrams
%  ---------------------------------------------------------
%  \plotfig{key}{caption}{one-line manifest description}
%    A PLOT, generated by code/plots/<chapter>.py from the same library the
%    listings use, into figures/plots/<lang>/<key>.pdf. Per language, because
%    an axis label in the wrong language is the half-translation that makes a
%    bilingual book look machine-made, and because the Polish edition owes a
%    decimal comma on its tick labels too. Build output, gitignored.
%
%  \mermaidfig{key}{caption}{one-line manifest description}
%    A DIAGRAM of an idea rather than of data, from
%    figures/mermaid/<lang>/<key>.mmd, rendered by `make diagrams'.
%
%  If the rendered PDF is absent either macro prints a visible marker in the
%  flow rather than a float, so a missing figure cannot be mistaken for a
%  layout choice. \phantomsection before \addcontentsline is load-bearing:
%  without it the manifest entry records whatever anchor preceded the figure,
%  which after a listing is a line that was never printed.
% ============================================================
\makeatletter
\newcommand{\listofplots}{%
  \section*{\lblPlotManifest}%
  \phantomsection\addcontentsline{toc}{section}{\lblPlotManifest}%
  \begingroup\c@tocdepth=2\relax\@starttoc{plt}\endgroup%
}
\newcommand{\listofdiagrams}{%
  \section*{\lblDiagramManifest}%
  \phantomsection\addcontentsline{toc}{section}{\lblDiagramManifest}%
  \begingroup\c@tocdepth=2\relax\@starttoc{dgm}\endgroup%
}
\makeatother

\newcommand{\figmissing}[3]{% #1 message, #2 path, #3 caption
  \par\medskip
  \begin{tcolorbox}[enhanced, breakable, width=\linewidth, sharp corners,
    colback=white, colframe=mgrey, boxrule=0.8pt,
    left=8pt, right=8pt, top=8pt, bottom=8pt]
    {\bfseries\color{mgrey}\small #1}\\[4pt]
    {\ttfamily\scriptsize\color{mgrey} #2}
  \end{tcolorbox}%
  \captionof{figure}{#3}%
  \par\medskip
}

\newcommand{\plotfig}[3]{%
  \phantomsection%
  \addcontentsline{plt}{subsection}{\protect\numberline{}\texttt{#1} --- #3}%
  \IfFileExists{figures/plots/\booklang/#1.pdf}{%
    \begin{figure}[htbp]
    \centering
    \includegraphics[width=\linewidth,height=0.45\textheight,keepaspectratio]{figures/plots/\booklang/#1.pdf}%
    \caption{#2}\label{fig:#1}
    \end{figure}%
  }{%
    \figmissing{\lblPlotNotRendered}{figures/plots/\booklang/#1.pdf}{#2}\label{fig:#1}%
  }%
}

\newcommand{\mermaidfig}[3]{%
  \phantomsection%
  \addcontentsline{dgm}{subsection}{\protect\numberline{}\texttt{#1.mmd} --- #3}%
  \IfFileExists{figures/diagrams/\booklang/#1.pdf}{%
    \begin{figure}[htbp]
    \centering
    \includegraphics[width=0.95\linewidth,height=0.42\textheight,keepaspectratio]{figures/diagrams/\booklang/#1.pdf}%
    \caption{#2}\label{fig:#1}
    \end{figure}%
  }{%
    \figmissing{\lblNotRendered}{figures/mermaid/\booklang/#1.mmd}{#2}\label{fig:#1}%
  }%
}

% ============================================================
%  APPENDIX B: the misreadings, generated by tools/gen_traps.py
%
%  \trapchapter{Chapter~\ref{ch:x}}{title}    a heading per chapter
%  \trapentry{n}{habit}{truth}{\ref{sec:x}}  one misreading
%
%  Set ragged right: an entry is mostly short sentences with \code{} spans in
%  them, and a justified line full of unbreakable runs is an overfull box
%  waiting for the edition with the longer words.
% ============================================================
\newcommand{\trapchapter}[2]{%
  \par\medskip
  {\large\bfseries\color{mblue}#1\quad\normalfont\itshape\color{mgrey}#2\par}%
  \nopagebreak\smallskip}
\newcommand{\trapentry}[4]{%
  \par\noindent
  \begingroup\raggedright\hangindent=2.4em\hangafter=1
  \makebox[2.4em][l]{\textbf{\color{mred}#1}}%
  \textbf{\enquote{#2}}\enspace #3\enspace{\color{mgrey}\S\,#4}\par
  \endgroup\smallskip}

% ============================================================
%  CHAPTER STUBS
% ============================================================
\newcommand{\chapterstub}[1]{%
  \begin{tcolorbox}[
    enhanced, breakable, width=\linewidth, sharp corners,
    colback=mlight, colframe=mred, boxrule=1pt,
    borderline={0.6pt}{3pt}{mred, dashed},
    left=12pt, right=12pt, top=12pt, bottom=12pt]
    {\bfseries\color{mred}\large \lblNotWritten}\\[8pt]
    \small\raggedright #1
  \end{tcolorbox}%
}

% ============================================================
%  COMPUTED VALUES
%  A number the reader cannot do in their head is computed by a script under
%  code/measure/, written to figures/values/<chapter>.tex as
%      \pyval{ch06.steps.apart}{34}
%  and reaches the page as \val{ch06.steps.apart}. The book contains
%  references, not digits; `make verify' fails when a committed value no
%  longer matches its script. \val applies the edition's decimal marker, which
%  is how the Polish edition gets its comma without anybody remembering.
% ============================================================
% The comma is BRACED: siunitx sets numbers in maths mode, where a bare comma
% is punctuation and gets a thin space after it ("9, 81"). {,} is an ordinary
% symbol, which is what a decimal marker is.
\newcommand{\bookdecimalmarker}{\ifpl{{,}}{.}}
\IfFileExists{siunitx.sty}{%
  \usepackage{siunitx}%
  \sisetup{
    group-digits = integer,
    group-minimum-digits = 5,
    group-separator = {\,},
    output-decimal-marker = {\bookdecimalmarker},
    exponent-product = \times,
    print-unity-mantissa = false
  }%
  \newcommand{\booknum}[1]{\num{##1}}%
}{%
  \newcommand{\booknum}[1]{##1}%
  \providecommand{\num}[1]{##1}%
}

\makeatletter
\newcommand{\pyvalues}{}
\newcommand{\pyval}[2]{\csgdef{pyval@#1}{#2}\listxadd{\pyvalues}{#1}}
\newcommand{\pyvaltext}[2]{%
  \csgdef{pyval@#1}{#2}\csgdef{pyvaltext@#1}{}\listxadd{\pyvalues}{#1}}
% \num on a non-numeric body is a FATAL siunitx error, so \val refuses a text
% value by name rather than letting siunitx fail on it.
\newcommand{\val}[1]{%
  \ifcsdef{pyval@#1}{%
    \ifcsdef{pyvaltext@#1}{%
      \PackageError{chaosbook}{Value `#1' is not a number}%
        {Use \string\valtext{#1} instead.}%
    }{}%
    \booknum{\csuse{pyval@#1}}%
  }{\textcolor{mred}{\textbf{[\lblValueMissing: #1]}}}%
}
\newcommand{\valtext}[1]{%
  \ifcsdef{pyval@#1}{\csuse{pyval@#1}}%
    {\textcolor{mred}{\textbf{[\lblValueMissing: #1]}}}%
}
\makeatother

% figures/values/all.tex is generated by `make numbers' and \inputs each
% value file. A build with no values still compiles; every \val{} then prints
% a visible marker. \valuesindex lets the one-chapter build point elsewhere.
\providecommand{\valuesindex}{figures/values/all}
\IfFileExists{\valuesindex.tex}{\input{\valuesindex}}{%
  \AtBeginDocument{\PackageWarning{chaosbook}{No computed values found. Run `make numbers'.}}}

% ============================================================
%  PINNED VERSIONS --- single source of truth
%  Change these here AND in code/pyproject.toml; tools/check_structure.py
%  --pins fails when the two disagree. A chapter never types a version.
% ============================================================
\newcommand{\bookedition}{0.1}
\newcommand{\bookyear}{2026}
\newcommand{\pyver}{3.14}              % the minor the book targets
\newcommand{\pypatch}{3.14.7}          % the exact interpreter the code ran on
\newcommand{\uvver}{0.12.13}           % uv
\newcommand{\numpyver}{2.5.3}          % numpy
\newcommand{\matplotlibver}{3.11.2}    % matplotlib
\newcommand{\pytestver}{9.1.1}         % pytest
\newcommand{\ruffver}{0.16.9}          % ruff
% The prose licence is a name with a version in it, not a quantity: it is a
% macro so that neither edition typesets it as a decimal to be localised.
\newcommand{\proselicence}{CC~BY-NC-SA~4.0}

% ============================================================
%  INDEX
% ============================================================
\apptocmd{\theindex}{\raggedcolumns\raggedright}{}{%
  \PackageWarning{chaosbook}{Could not patch theindex for ragged setting}}
\providecommand*\see[2]{}
\renewcommand*\see[2]{\emph{\lblIndexSee} #1}
\providecommand*\seealso[2]{}
\renewcommand*\seealso[2]{\emph{\lblIndexSeeAlso} #1}

\makeindex
; nothing else differs between them at this level.
%  Every user-visible string lives in lang/en.tex or lang/pl.tex, so a label
%  that appears in one edition and not the other is a missing macro rather
%  than a silent divergence.
%
%  Lineage. This book is the third in a family and takes one thing from each
%  of the other two:
%
%    * from Mathematics from Zero for the AI Engineer: the bilingual single
%      source (\booklang, lang/, one body.tex), the rule that every number is
%      computed and reaches the page as \val{}, and the programmed-learning
%      method -- outcomes on entry, a question before every answer, a
%      misconception elicited before it is named, a self-check generated from
%      the outcomes, and checkpoint answers collected at the back;
%    * from Python for .NET Engineers: every listing is a FILE that CI runs,
%      every lab is a starter that must fail and a solution that must pass,
%      and a console transcript is a computed number that happens to be
%      verbatim.
%
%  What is new here is \plotfig. A book about chaos cannot be taught in boxes
%  and arrows: the bifurcation diagram and the Lorenz attractor ARE the
%  argument, so the plots are generated by code/plots/ from the same library
%  the listings use, per language, and never drawn by hand.
% ============================================================

% \booklang must be set by the main file. Default to English rather than
% failing, so a stray % ============================================================
%  preamble.tex --- styling, environments, macros
%
%  Shared by BOTH editions. main-en.tex and main-pl.tex each define \booklang
%  before \input{preamble}; nothing else differs between them at this level.
%  Every user-visible string lives in lang/en.tex or lang/pl.tex, so a label
%  that appears in one edition and not the other is a missing macro rather
%  than a silent divergence.
%
%  Lineage. This book is the third in a family and takes one thing from each
%  of the other two:
%
%    * from Mathematics from Zero for the AI Engineer: the bilingual single
%      source (\booklang, lang/, one body.tex), the rule that every number is
%      computed and reaches the page as \val{}, and the programmed-learning
%      method -- outcomes on entry, a question before every answer, a
%      misconception elicited before it is named, a self-check generated from
%      the outcomes, and checkpoint answers collected at the back;
%    * from Python for .NET Engineers: every listing is a FILE that CI runs,
%      every lab is a starter that must fail and a solution that must pass,
%      and a console transcript is a computed number that happens to be
%      verbatim.
%
%  What is new here is \plotfig. A book about chaos cannot be taught in boxes
%  and arrows: the bifurcation diagram and the Lorenz attractor ARE the
%  argument, so the plots are generated by code/plots/ from the same library
%  the listings use, per language, and never drawn by hand.
% ============================================================

% \booklang must be set by the main file. Default to English rather than
% failing, so a stray \input{preamble} still compiles and says what it did.
\providecommand{\booklang}{en}

\usepackage{etoolbox}   % loaded first: the language switch below needs it

% ---------- Page geometry ----------
% ONE format: A4, 12pt, ONE-SIDED. The book is read on a screen, usually with
% an editor open beside it, so there is no recto, no verso, no gutter and no
% blank leaf before a chapter. The measure is about seventy-five characters
% of prose. At \footnotesize a listing line of 79 characters fits without
% wrapping, and 79 is the width every file under code/ is held to by
% tools/check_structure.py --lines.
\usepackage[
  a4paper,
  left=3.1cm, right=3.1cm, top=2.6cm, bottom=3.0cm,
  head=26pt, headsep=0.8cm, footskip=1.4cm
]{geometry}

% ---------- Encoding & fonts ----------
\usepackage[T1]{fontenc}
\usepackage[utf8]{inputenc}
\IfFileExists{lmodern.sty}{\usepackage{lmodern}}{}
% newtxtext takes its TS1 symbols from TeX Gyre Termes (Debian: tex-gyre).
% A machine with newtx and without tex-gyre dies on the copyright page. CI's
% full TeX Live is the reference; locally install tex-gyre as well.
\IfFileExists{newtxtext.sty}{\usepackage{newtxtext}}{}
% amsmath before newtxmath; amssymb only when newtxmath is absent, because
% loading both is a fatal clash (\Bbbk already defined) that is invisible on
% a machine without newtx.
\usepackage{amsmath}
\IfFileExists{newtxmath.sty}{\usepackage{newtxmath}}{\usepackage{amssymb}}
\IfFileExists{inconsolata.sty}{\usepackage[varqu,scaled=0.95]{inconsolata}}{}
% Without lmodern, microtype runs with font expansion off, which is what
% resolves a marginal overfull line -- so a container short of lmodern
% disagrees with CI about line breaks. Install lmodern locally too.
\IfFileExists{lmodern.sty}{\usepackage{microtype}}{\usepackage[expansion=false]{microtype}}

% upquote: in T1 the input ' is a right single quote, and a listing that sets
% f('x') with curly quotes is a SyntaxError when typed back in.
\IfFileExists{upquote.sty}{\usepackage{upquote}}{}
% And its twin: in T1 the pair -- is an en-dash ligature in the typewriter
% family too, so \code{--flag} would print a flag nobody can type. Scoped to
% tt*, so prose keeps its own dashes.
\DisableLigatures[-]{encoding = T1, family = tt*}

% ---------- Language ----------
% babel with a missing language file is FATAL, so each .ldf is probed first.
\IfFileExists{babel.sty}{%
  \IfFileExists{polish.ldf}{%
    \IfFileExists{british.ldf}{%
      \ifdefstring{\booklang}{pl}%
        {\usepackage[british,main=polish]{babel}}%
        {\usepackage[polish,main=british]{babel}}%
    }{\usepackage[polish]{babel}}%
  }{%
    \IfFileExists{british.ldf}{\usepackage[british]{babel}}{}%
  }%
}{}
% babel-polish makes " an active shorthand, which is a category-code accident
% waiting to happen inside listings. Polish quotation marks come from
% \enquote, which csquotes sets per language.
\ifdefstring{\booklang}{pl}{\AtBeginDocument{\shorthandoff{"}}}{}

% ---------- Core packages ----------
\usepackage{graphicx}
\usepackage{xcolor}
\usepackage{booktabs}
\usepackage{tabularx}
\usepackage{longtable}
\usepackage{array}
\usepackage{enumitem}
\usepackage{listings}
\usepackage[most]{tcolorbox}
\usepackage{fancyhdr}
\usepackage{titlesec}
\usepackage{caption}
\usepackage{imakeidx}
\IfFileExists{csquotes.sty}{\usepackage[autostyle=true]{csquotes}}{%
  \providecommand{\enquote}[1]{``##1''}}
\usepackage[hidelinks]{hyperref}

% ---------- Palette ----------
% The family's palette, so the three books read as a set. Two colours are
% this book's own: the lab (green) and the stop-and-think box (violet).
\definecolor{mblue}{HTML}{1F4E79}
\definecolor{mteal}{HTML}{0E7C7B}
\definecolor{mamber}{HTML}{B26A00}
\definecolor{mred}{HTML}{9B2C2C}
\definecolor{mviolet}{HTML}{7B4397}
\definecolor{mgreen}{HTML}{3B7A3B}
\definecolor{mgrey}{HTML}{5A5A5A}
\definecolor{mlight}{HTML}{F4F6F8}
\definecolor{mfaint}{HTML}{FBFCFD}
\definecolor{codebg}{HTML}{FAFAFA}
\definecolor{codeframe}{HTML}{D8DEE4}
\definecolor{kw}{HTML}{0B5394}
\definecolor{str}{HTML}{9B2C2C}
\definecolor{cmt}{HTML}{4F7A4F}
\definecolor{typ}{HTML}{0E7C7B}
\definecolor{dec}{HTML}{7B4397}

\hypersetup{
  colorlinks=true,
  linkcolor=mblue,
  urlcolor=mteal,
  citecolor=mblue,
  pdfauthor={Konrad Cinkusz}
}

% \ifpl{Polish}{English} -- for the handful of places where a string is
% structural rather than content. Robust, because a fragile \ifpl inside a
% moving argument fails on the next run in exactly one of the two editions.
\DeclareRobustCommand{\ifpl}[2]{\ifdefstring{\booklang}{pl}{#1}{#2}}
\makeatletter
\pdfstringdefDisableCommands{%
  \ifdefstring{\booklang}{pl}{\let\ifpl\@firstoftwo}{\let\ifpl\@secondoftwo}%
  % hyperref drops \color from a bookmark and keeps its ARGUMENT, so a title
  % carrying \code{} would bookmark as "mtealx". Gobble the colour name.
  \def\color#1{}%
}
\makeatother

% ---------- Language strings ----------
\input{lang/\booklang}

% The PDF's subject, keywords and language read \lblPdf... from the language
% file, so they cannot sit in the colour block above, which runs first. Two
% blocks is the ordering, not an oversight.
\hypersetup{
  pdfsubject={\lblPdfSubject},
  pdfkeywords={\lblPdfKeywords},
  pdflang={\lblPdfLang}
}

% ---------- Headings ----------
\titleformat{\chapter}[display]
  {\normalfont\huge\bfseries\color{mblue}\raggedright}
  {\normalfont\normalsize\color{mgrey}\MakeUppercase{\chaptertitlename}\ \thechapter}
  {8pt}{\Huge\raggedright}
\titlespacing*{\chapter}{0pt}{10pt}{28pt}
\titleformat{\section}{\normalfont\Large\bfseries\color{mblue}\raggedright}{\thesection}{0.8em}{}
\titleformat{\subsection}{\normalfont\large\bfseries\color{mgrey}\raggedright}{\thesubsection}{0.7em}{}

\setcounter{tocdepth}{1}
\setcounter{secnumdepth}{2}

\pagestyle{fancy}
\fancyhf{}
\fancyhead[L]{\small\nouppercase{\leftmark}}
\fancyhead[R]{\small\thepage}
\renewcommand{\headrulewidth}{0.4pt}
\fancypagestyle{plain}{\fancyhf{}\fancyfoot[R]{\small\thepage}%
  \renewcommand{\headrulewidth}{0pt}}
% After \pagestyle{fancy}, which installs its own \chaptermark.
\renewcommand{\chaptermark}[1]{\markboth{\thechapter.\ #1}{}}

% ---------- Mathematics ----------
% The handful of notations the book uses, as macros so both editions set the
% same symbol. The book writes \ln for the natural logarithm and \log_{10}
% where the base is ten; a bare \log is ambiguous between the two readerships
% and is not used.
\newcommand{\R}{\mathbb{R}}
\newcommand{\dd}{\,\mathrm{d}}
\newcommand{\abs}[1]{\left\lvert #1\right\rvert}
\newcommand{\iu}{\mathrm{i}}

% ============================================================
%  CODE LISTINGS
%  Every listing in a chapter is a FILE under code/, pulled in with \pyfile or
%  \pyregion, and CI runs it on the pinned interpreter. An in-source
%  \begin{python} block is allowed for a fragment of three or four lines that
%  shows a shape rather than runs.
% ============================================================
\lstdefinelanguage{pythonx}{
  language=Python,
  morekeywords={async,await,match,case,nonlocal,yield,type},
  morekeywords=[2]{
    orbit,iterate,logistic,rk4_step,euler_step,integrate,lorenz,henon,
    lyapunov,separation,escape_time,box_count
  },
  sensitive=true
}

\lstdefinestyle{bookcode}{
  basicstyle=\ttfamily\footnotesize,
  keywordstyle=\color{kw}\bfseries,
  keywordstyle=[2]\color{typ},
  stringstyle=\color{str},
  % Colour only: inconsolata has no italic, and an italic comment style is
  % silently substituted and logged on every build.
  commentstyle=\color{cmt},
  numbers=left,
  numberstyle=\ttfamily\tiny\color{mgrey},
  numbersep=8pt,
  backgroundcolor=\color{codebg},
  frame=single,
  rulecolor=\color{codeframe},
  framesep=6pt,
  breaklines=true,
  breakatwhitespace=false,
  postbreak=\mbox{\textcolor{mgrey}{$\hookrightarrow$}\space},
  showstringspaces=false,
  tabsize=4,
  columns=fullflexible,
  keepspaces=true,
  captionpos=b,
  aboveskip=1.2em,
  belowskip=1.2em,
  % Do NOT add a mapping for U+00A0: it aborts the run with a message that
  % names neither the character nor this line. Listings are ASCII, and
  % parity's C13 checks it.
  literate={ł}{{\l}}1 {Ł}{{\L}}1 {ą}{{\k a}}1 {ę}{{\k e}}1
           {ś}{{\'s}}1 {ć}{{\'c}}1 {ż}{{\.z}}1 {ź}{{\'z}}1 {ó}{{\'o}}1 {ń}{{\'n}}1
           {Ą}{{\k A}}1 {Ę}{{\k E}}1 {Ś}{{\'S}}1 {Ć}{{\'C}}1
           {Ż}{{\.Z}}1 {Ź}{{\'Z}}1 {Ó}{{\'O}}1 {Ń}{{\'N}}1
           {—}{{-{}-}}2 {–}{{-}}1 {‘}{{'}}1 {’}{{'}}1
           {“}{{"}}1 {”}{{"}}1 {°}{{\textdegree}}1
}
\lstset{style=bookcode}

\lstnewenvironment{python}[1][]
  {\lstset{language=pythonx,style=bookcode,#1}}
  {}
\lstnewenvironment{shellcmd}[1][]
  {\lstset{language=bash,style=bookcode,numbers=none,keywordstyle=\color{mteal}\bfseries,#1}}
  {}

% The path is printed above every file listing, in the form the reader opens
% in the editor. \nobreak keeps it on the page of the listing it names.
% The path must not be stranded at the foot of a page with its listing
% overleaf: \nobreak alone does not hold, because listings opens with its own
% vertical material. \needspace asks for room for the path and the first few
% lines of code, and turns the page first when there is not.
% \Needspace (capital N) measures \pagegoal-\pagetotal exactly; lower-case
% \needspace works through glue and penalties and was measured, by the
% Chapter 11 pass, turning pages early and leaving gaps. It must be used
% between paragraphs, which \listingpath's own \par guarantees. Chapters use
% \Needspace* too, so the fallback for a machine without the package defines
% both forms (the star swallowed with its argument).
\makeatletter
\IfFileExists{needspace.sty}{\usepackage{needspace}}{%
  \providecommand{\needspace}[1]{}%
  \providecommand{\Needspace}{\@ifstar\@gobble\@gobble}}
\makeatother
\newcommand{\listingpath}[1]{%
  \par\medskip\Needspace{6\baselineskip}\noindent
  {\ttfamily\footnotesize\color{mgrey}\detokenize{#1}}%
  \par\nobreak}

% Usage: \pyfile{code/ch05/cobweb.py}{Caption}{lst:label}
\newcommand{\pyfile}[3]{%
  \listingpath{#1}%
  \lstinputlisting[language=pythonx,style=bookcode,aboveskip=0.3em,
    caption={#2},label={#3}]{#1}%
}

% A named region of a file, delimited by `# --8<-- [start:name]` and
% `# --8<-- [end:name]`. Matched through listings' range prefix and suffix
% keys -- the other obvious form matches nothing and prints the WHOLE file
% without a warning (measured in the sibling book). A region name is a word,
% never a number. The line numbers printed are the file's own.
% Usage: \pyregion{code/ch05/cobweb.py}{step}{Caption}{lst:label}
\lstset{
  rangebeginprefix=\#\ --8<--\ [start:, rangebeginsuffix=],
  rangeendprefix=\#\ --8<--\ [end:, rangeendsuffix=],
  includerangemarker=false}
\newcommand{\pyregion}[4]{%
  \listingpath{#1}%
  \lstinputlisting[language=pythonx,style=bookcode,aboveskip=0.3em,
    linerange={#2-#2},caption={#3},label={#4}]{#1}%
}

% A program's output, written by code/measure/transcripts.py. There is no
% typed form and there must not be: \val{} does not expand in a verbatim
% body, and a transcript nobody ran looks exactly like one that was.
\newcommand{\transcript}[1]{%
  \par\smallskip
  \IfFileExists{figures/transcripts/#1.txt}{%
    \lstinputlisting[style=bookcode,numbers=none,basicstyle=\ttfamily\footnotesize,
      keywordstyle=\color{black},backgroundcolor=\color{white},
      language={}]{figures/transcripts/#1.txt}%
  }{%
    \begin{tcolorbox}[enhanced, sharp corners, width=\linewidth,
      colback=mlight, colframe=mgrey, boxrule=0.6pt,
      left=8pt, right=8pt, top=6pt, bottom=6pt]
      {\bfseries\color{mgrey}\small \lblTranscriptMissing}\\[2pt]
      {\ttfamily\scriptsize\color{mgrey} figures/transcripts/#1.txt}
    \end{tcolorbox}%
  }%
  \par\smallskip
}

% Inline literals. Underscores inside these must be written \_.
\newcommand{\code}[1]{\texttt{\small #1}}
\newcommand{\pkg}[1]{\texttt{\small\color{mblue}#1}}
\newcommand{\api}[1]{\texttt{\small\color{mteal}#1}}

% ============================================================
%  ADMONITIONS
%  A small, fixed vocabulary. `tinted' may be read in passing; `outlined' is
%  a block the reader must stop at. A third treatment would mean the first two
%  are not carrying their meaning.
% ============================================================
\tcbset{
  cfz tinted/.style={
    enhanced, breakable, sharp corners,
    lines before break=4,
    colback=mlight, boxrule=0pt, leftrule=3pt,
    left=8pt, right=8pt, top=6pt, bottom=6pt,
    before skip=13pt, after skip=13pt,
    fonttitle=\bfseries\small,
    coltitle=black, colbacktitle=mlight,
    attach title to upper={\par\smallskip},
  },
  cfz outlined/.style={
    enhanced, breakable, sharp corners,
    % A box broken straight after its title is a heading stranded at the foot
    % of a page with its question overleaf. Four lines is title plus three.
    lines before break=4,
    colback=white, boxrule=0.6pt,
    left=8pt, right=8pt, top=6pt, bottom=6pt,
    before skip=13pt, after skip=13pt,
    fonttitle=\bfseries\small,
    coltitle=black, colbacktitle=white,
    attach title to upper={\par\smallskip},
  },
}

% Read in passing.
\newtcolorbox{note}[1][]{cfz tinted, colframe=mblue, title={\lblNote}, #1}
\newtcolorbox{warning}[1][]{cfz tinted, colframe=mred, title={\lblWarning}, #1}
% Where you meet this in the world. If it cannot name a specific system,
% instrument, paper or event, it is prose and not this box.
\newtcolorbox{worldbox}[1][]{cfz tinted, colframe=mteal, title={\lblWorld}, #1}
% What is stated and not proved here, and where the proof lives.
\newtcolorbox{rigourbox}[1][]{cfz tinted, colframe=mgrey, title={\lblRigour}, #1}
% The companion library the reader builds, under code/src/chaoslab/.
\newtcolorbox{projectbox}[1][]{cfz tinted, colframe=mamber, title={\lblProject}, #1}

% Stop and read.
% The misconception, named AFTER the reader has been asked the question it
% gets wrong. Appendix B is the catalogue; a trap box is one entry in place.
\newtcolorbox{trapbox}[1][]{cfz outlined, colframe=mred, title={\lblTrap}, #1}
% A listing that has not been run on the pinned versions. Counted by
% `make debt'; nothing ships to a reader with one attached.
\newtcolorbox{verifybox}[1][]{cfz outlined, colframe=mgrey, title={\lblVerify}, #1}
\newtcolorbox{labbox}[1][]{cfz outlined, colframe=mgreen, title={\lblLab}, #1}

% ============================================================
%  THE METHOD --- questions before answers
%  ---------------------------------------------------------
%  The unit of teaching is a question the reader answers BEFORE reading the
%  answer. That is retrieval practice and errorful generation, the two
%  findings the programmed-learning format is defended by, and it is the one
%  thing the sibling mathematics book does that this one keeps whole.
%
%  \begin{think} ... \end{think}   the question, numbered per chapter, ending
%                                  with the instruction to answer before
%                                  reading on;
%  \reveal{short answer}           a centred box with the answer, which is the
%                                  first thing after the question -- the edge
%                                  the reader covers with a hand, or scrolls
%                                  to;
%  \begin{revealblock} ... \end{revealblock}  the same, for a worked answer.
%
%  tools/check_structure.py --think fails when a think box is not followed by
%  a reveal before the next think box or section, and when a reveal answers
%  nothing. That is a property of the source, so it is checked mechanically.
%
%  The instruction says "before you read on", never "before you turn over":
%  the two editions paginate differently, so the second is true in one of them.
% ============================================================
\newcounter{think}[chapter]
\renewcommand{\thethink}{\thechapter.\arabic{think}}
\newtcolorbox{thinkbox}[1][]{cfz outlined, colframe=mviolet, #1}
\NewDocumentEnvironment{think}{}{%
  \refstepcounter{think}%
  \begin{thinkbox}[title={\lblThink~\thethink}]%
}{%
  \par\smallskip
  {\raggedleft\itshape\small\color{mviolet}\lblThinkCue\par}%
  \end{thinkbox}%
}
\tcbset{
  cfz answer/.style={
    enhanced, sharp corners,
    colback=mfaint, colframe=mviolet!55, boxrule=0.5pt,
    before skip=4pt, after skip=10pt,
  },
}
\newcommand{\reveal}[1]{%
  \par\nobreak
  \begin{tcolorbox}[cfz answer, unbreakable,
    left=8pt, right=8pt, top=5pt, bottom=5pt,
    width=0.86\linewidth, center, halign=flush center]
    {\small\bfseries\color{mviolet}\lblAnswerMark}\enspace #1
  \end{tcolorbox}%
  \nobreak\noindent\ignorespaces
}
\NewDocumentEnvironment{revealblock}{}{%
  \par\nopagebreak
  \begin{tcolorbox}[cfz answer, breakable,
    left=10pt, right=10pt, top=6pt, bottom=6pt]
  {\small\bfseries\color{mviolet}\lblAnswerMark}\par\smallskip
}{%
  \end{tcolorbox}%
  \nopagebreak\par\noindent\ignorespaces
}

% ---- Outcomes on entry, and "Can you?" generated from them ----------------
% The self-check at the end of a chapter IS the list of outcomes at its
% start, one for one, because it is built from the same tokens: each
% \outcome both prints its item and appends a row to \cfz@canrows. A \loop
% inside a tabularx body does not survive alignment scanning, which is why
% the rows are accumulated at declaration time (the sibling book's finding).
\makeatletter
\newcommand{\cfz@canrows}{}
\newcommand{\cfz@ratebox}{%
  \fbox{\rule{0pt}{0.9ex}\rule{0.9ex}{0pt}}}
\newcommand{\cfz@rates}{%
  {\footnotesize\color{mgrey}%
   1\,\cfz@ratebox\ 2\,\cfz@ratebox\ 3\,\cfz@ratebox\ 4\,\cfz@ratebox\ 5\,\cfz@ratebox}}
\NewDocumentEnvironment{outcomes}{}{%
  \gdef\cfz@canrows{}%
  \begin{tcolorbox}[cfz tinted, colframe=mblue, title={\lblOutcomes}]%
  \lblOutcomesLead
  \begin{itemize}[leftmargin=1.4em, itemsep=2pt, topsep=4pt]%
}{%
  \end{itemize}%
  \end{tcolorbox}%
}
\newcommand{\outcome}[1]{%
  \item #1%
  \gappto\cfz@canrows{#1 & \cfz@rates \\ \addlinespace[3pt]}%
}
\newcommand{\canyou}{%
  \section*{\lblCanYou}%
  \phantomsection\addcontentsline{toc}{section}{\lblCanYou}%
  \noindent\lblCanYouIntro\par\medskip
  \ifdefempty{\cfz@canrows}{%
    \noindent{\color{mred}\lblNoOutcomes}\par}{%
  {\small
  \noindent\begin{tabularx}{\linewidth}{@{}>{\raggedright\arraybackslash}X r@{}}
    \toprule
    \cfz@canrows
    \bottomrule
  \end{tabularx}\par}}%
}
\makeatother

% ---- Summary -------------------------------------------------------------
\newtcolorbox{summarybox}[1][]{cfz tinted, colframe=mblue, title={\lblSummary}, #1}

% ---- Checkpoint questions, answered at the back --------------------------
% Answers are written AT THE QUESTION with \answerto{} and printed in
% Appendix A. They are carried there by a global token list rather than an
% auxiliary file, because an answer contains mathematics and \val{}, and
% material written to an .aux is expanded at shipout, where fragile commands
% break with errors naming neither the answer nor the chapter.
% \include is sequential within a run, so the list is complete by the time
% the appendix is typeset.
\makeatletter
\newcommand{\cfz@answers}{}
\newcommand{\cfz@lastansch}{0}
\NewDocumentEnvironment{checkpoint}{}{%
  \section*{\lblCheckpoint}%
  \phantomsection\addcontentsline{toc}{section}{\lblCheckpoint}%
  \noindent\lblCheckpointIntro\par
  \begin{enumerate}[label=\textbf{\arabic*.}, leftmargin=2em, itemsep=4pt]%
}{%
  \end{enumerate}%
}
\newcommand{\answerto}[1]{%
  \ifnum\value{chapter}=\cfz@lastansch\relax\else
    \xdef\cfz@lastansch{\arabic{chapter}}%
    \xappto\cfz@answers{\noexpand\cfz@anschapter{\thechapter}}%
  \fi
  \xappto\cfz@answers{\noexpand\cfz@ansitem{\arabic{enumi}}}%
  \gappto\cfz@answers{{#1}}%
}
\newcommand{\cfz@anschapter}[1]{%
  \par\medskip\noindent{\bfseries\color{mblue}\chaptername~#1}\par\nobreak\smallskip}
\newcommand{\cfz@ansitem}[2]{%
  \par\noindent\hangindent=2em\hangafter=1
  \makebox[2em][l]{\textbf{#1.}}#2\par}
\newcommand{\answersbody}{%
  \ifdefempty{\cfz@answers}{\noindent{\color{mred}\lblNoAnswers}\par}{%
    {\raggedright\cfz@answers}}%
}
\makeatother

% ============================================================
%  LABS
%  \begin{lab}{key}{title} ... \end{lab}
%
%  A lab is not a paragraph but a FILE the reader opens: a starter under
%  code/labs/chNN/<key>.py that must FAIL its test, a reference solution under
%  code/labs/chNN/solutions/, and test_<key>.py beside it. The box prints the
%  paths and the one command that runs the test. An environment rather than a
%  macro, because a macro argument cannot hold a verbatim fragment.
%  tools/check_structure.py --labs fails when a file is missing or the key's
%  chapter number disagrees with its chapter.
% ============================================================
\newcounter{lab}[chapter]
\renewcommand{\thelab}{\thechapter.\arabic{lab}}
\makeatletter
\newcommand{\labdir}{ch\two@digits{\value{chapter}}}
\newcommand{\listoflabs}{%
  \section*{\lblLabManifest}%
  \phantomsection\addcontentsline{toc}{section}{\lblLabManifest}%
  \begingroup\c@tocdepth=2\relax\@starttoc{lab}\endgroup%
}
\makeatother
\NewDocumentEnvironment{lab}{mm}{%
  \refstepcounter{lab}%
  \addcontentsline{lab}{subsection}{\protect\numberline{\thelab}\texttt{\protect\detokenize{#1}} --- #2}%
  \begin{labbox}[title={\lblLab~\thelab\ \dash{} #2}]%
}{%
  \par\medskip
  {\small\setlength{\parindent}{0pt}%
  \lblLabStarter: \texttt{code/labs/\labdir/\detokenize{#1}.py}\\
  \lblLabTest: \texttt{code/labs/\labdir/test\_\detokenize{#1}.py}\\
  \lblLabRun: \texttt{cd code \&\& uv run pytest -k \detokenize{#1}}\par}%
  \end{labbox}%
}

% ============================================================
%  FIGURES --- plots and diagrams
%  ---------------------------------------------------------
%  \plotfig{key}{caption}{one-line manifest description}
%    A PLOT, generated by code/plots/<chapter>.py from the same library the
%    listings use, into figures/plots/<lang>/<key>.pdf. Per language, because
%    an axis label in the wrong language is the half-translation that makes a
%    bilingual book look machine-made, and because the Polish edition owes a
%    decimal comma on its tick labels too. Build output, gitignored.
%
%  \mermaidfig{key}{caption}{one-line manifest description}
%    A DIAGRAM of an idea rather than of data, from
%    figures/mermaid/<lang>/<key>.mmd, rendered by `make diagrams'.
%
%  If the rendered PDF is absent either macro prints a visible marker in the
%  flow rather than a float, so a missing figure cannot be mistaken for a
%  layout choice. \phantomsection before \addcontentsline is load-bearing:
%  without it the manifest entry records whatever anchor preceded the figure,
%  which after a listing is a line that was never printed.
% ============================================================
\makeatletter
\newcommand{\listofplots}{%
  \section*{\lblPlotManifest}%
  \phantomsection\addcontentsline{toc}{section}{\lblPlotManifest}%
  \begingroup\c@tocdepth=2\relax\@starttoc{plt}\endgroup%
}
\newcommand{\listofdiagrams}{%
  \section*{\lblDiagramManifest}%
  \phantomsection\addcontentsline{toc}{section}{\lblDiagramManifest}%
  \begingroup\c@tocdepth=2\relax\@starttoc{dgm}\endgroup%
}
\makeatother

\newcommand{\figmissing}[3]{% #1 message, #2 path, #3 caption
  \par\medskip
  \begin{tcolorbox}[enhanced, breakable, width=\linewidth, sharp corners,
    colback=white, colframe=mgrey, boxrule=0.8pt,
    left=8pt, right=8pt, top=8pt, bottom=8pt]
    {\bfseries\color{mgrey}\small #1}\\[4pt]
    {\ttfamily\scriptsize\color{mgrey} #2}
  \end{tcolorbox}%
  \captionof{figure}{#3}%
  \par\medskip
}

\newcommand{\plotfig}[3]{%
  \phantomsection%
  \addcontentsline{plt}{subsection}{\protect\numberline{}\texttt{#1} --- #3}%
  \IfFileExists{figures/plots/\booklang/#1.pdf}{%
    \begin{figure}[htbp]
    \centering
    \includegraphics[width=\linewidth,height=0.45\textheight,keepaspectratio]{figures/plots/\booklang/#1.pdf}%
    \caption{#2}\label{fig:#1}
    \end{figure}%
  }{%
    \figmissing{\lblPlotNotRendered}{figures/plots/\booklang/#1.pdf}{#2}\label{fig:#1}%
  }%
}

\newcommand{\mermaidfig}[3]{%
  \phantomsection%
  \addcontentsline{dgm}{subsection}{\protect\numberline{}\texttt{#1.mmd} --- #3}%
  \IfFileExists{figures/diagrams/\booklang/#1.pdf}{%
    \begin{figure}[htbp]
    \centering
    \includegraphics[width=0.95\linewidth,height=0.42\textheight,keepaspectratio]{figures/diagrams/\booklang/#1.pdf}%
    \caption{#2}\label{fig:#1}
    \end{figure}%
  }{%
    \figmissing{\lblNotRendered}{figures/mermaid/\booklang/#1.mmd}{#2}\label{fig:#1}%
  }%
}

% ============================================================
%  APPENDIX B: the misreadings, generated by tools/gen_traps.py
%
%  \trapchapter{Chapter~\ref{ch:x}}{title}    a heading per chapter
%  \trapentry{n}{habit}{truth}{\ref{sec:x}}  one misreading
%
%  Set ragged right: an entry is mostly short sentences with \code{} spans in
%  them, and a justified line full of unbreakable runs is an overfull box
%  waiting for the edition with the longer words.
% ============================================================
\newcommand{\trapchapter}[2]{%
  \par\medskip
  {\large\bfseries\color{mblue}#1\quad\normalfont\itshape\color{mgrey}#2\par}%
  \nopagebreak\smallskip}
\newcommand{\trapentry}[4]{%
  \par\noindent
  \begingroup\raggedright\hangindent=2.4em\hangafter=1
  \makebox[2.4em][l]{\textbf{\color{mred}#1}}%
  \textbf{\enquote{#2}}\enspace #3\enspace{\color{mgrey}\S\,#4}\par
  \endgroup\smallskip}

% ============================================================
%  CHAPTER STUBS
% ============================================================
\newcommand{\chapterstub}[1]{%
  \begin{tcolorbox}[
    enhanced, breakable, width=\linewidth, sharp corners,
    colback=mlight, colframe=mred, boxrule=1pt,
    borderline={0.6pt}{3pt}{mred, dashed},
    left=12pt, right=12pt, top=12pt, bottom=12pt]
    {\bfseries\color{mred}\large \lblNotWritten}\\[8pt]
    \small\raggedright #1
  \end{tcolorbox}%
}

% ============================================================
%  COMPUTED VALUES
%  A number the reader cannot do in their head is computed by a script under
%  code/measure/, written to figures/values/<chapter>.tex as
%      \pyval{ch06.steps.apart}{34}
%  and reaches the page as \val{ch06.steps.apart}. The book contains
%  references, not digits; `make verify' fails when a committed value no
%  longer matches its script. \val applies the edition's decimal marker, which
%  is how the Polish edition gets its comma without anybody remembering.
% ============================================================
% The comma is BRACED: siunitx sets numbers in maths mode, where a bare comma
% is punctuation and gets a thin space after it ("9, 81"). {,} is an ordinary
% symbol, which is what a decimal marker is.
\newcommand{\bookdecimalmarker}{\ifpl{{,}}{.}}
\IfFileExists{siunitx.sty}{%
  \usepackage{siunitx}%
  \sisetup{
    group-digits = integer,
    group-minimum-digits = 5,
    group-separator = {\,},
    output-decimal-marker = {\bookdecimalmarker},
    exponent-product = \times,
    print-unity-mantissa = false
  }%
  \newcommand{\booknum}[1]{\num{##1}}%
}{%
  \newcommand{\booknum}[1]{##1}%
  \providecommand{\num}[1]{##1}%
}

\makeatletter
\newcommand{\pyvalues}{}
\newcommand{\pyval}[2]{\csgdef{pyval@#1}{#2}\listxadd{\pyvalues}{#1}}
\newcommand{\pyvaltext}[2]{%
  \csgdef{pyval@#1}{#2}\csgdef{pyvaltext@#1}{}\listxadd{\pyvalues}{#1}}
% \num on a non-numeric body is a FATAL siunitx error, so \val refuses a text
% value by name rather than letting siunitx fail on it.
\newcommand{\val}[1]{%
  \ifcsdef{pyval@#1}{%
    \ifcsdef{pyvaltext@#1}{%
      \PackageError{chaosbook}{Value `#1' is not a number}%
        {Use \string\valtext{#1} instead.}%
    }{}%
    \booknum{\csuse{pyval@#1}}%
  }{\textcolor{mred}{\textbf{[\lblValueMissing: #1]}}}%
}
\newcommand{\valtext}[1]{%
  \ifcsdef{pyval@#1}{\csuse{pyval@#1}}%
    {\textcolor{mred}{\textbf{[\lblValueMissing: #1]}}}%
}
\makeatother

% figures/values/all.tex is generated by `make numbers' and \inputs each
% value file. A build with no values still compiles; every \val{} then prints
% a visible marker. \valuesindex lets the one-chapter build point elsewhere.
\providecommand{\valuesindex}{figures/values/all}
\IfFileExists{\valuesindex.tex}{\input{\valuesindex}}{%
  \AtBeginDocument{\PackageWarning{chaosbook}{No computed values found. Run `make numbers'.}}}

% ============================================================
%  PINNED VERSIONS --- single source of truth
%  Change these here AND in code/pyproject.toml; tools/check_structure.py
%  --pins fails when the two disagree. A chapter never types a version.
% ============================================================
\newcommand{\bookedition}{0.1}
\newcommand{\bookyear}{2026}
\newcommand{\pyver}{3.14}              % the minor the book targets
\newcommand{\pypatch}{3.14.7}          % the exact interpreter the code ran on
\newcommand{\uvver}{0.12.13}           % uv
\newcommand{\numpyver}{2.5.3}          % numpy
\newcommand{\matplotlibver}{3.11.2}    % matplotlib
\newcommand{\pytestver}{9.1.1}         % pytest
\newcommand{\ruffver}{0.16.9}          % ruff
% The prose licence is a name with a version in it, not a quantity: it is a
% macro so that neither edition typesets it as a decimal to be localised.
\newcommand{\proselicence}{CC~BY-NC-SA~4.0}

% ============================================================
%  INDEX
% ============================================================
\apptocmd{\theindex}{\raggedcolumns\raggedright}{}{%
  \PackageWarning{chaosbook}{Could not patch theindex for ragged setting}}
\providecommand*\see[2]{}
\renewcommand*\see[2]{\emph{\lblIndexSee} #1}
\providecommand*\seealso[2]{}
\renewcommand*\seealso[2]{\emph{\lblIndexSeeAlso} #1}

\makeindex
 still compiles and says what it did.
\providecommand{\booklang}{en}

\usepackage{etoolbox}   % loaded first: the language switch below needs it

% ---------- Page geometry ----------
% ONE format: A4, 12pt, ONE-SIDED. The book is read on a screen, usually with
% an editor open beside it, so there is no recto, no verso, no gutter and no
% blank leaf before a chapter. The measure is about seventy-five characters
% of prose. At \footnotesize a listing line of 79 characters fits without
% wrapping, and 79 is the width every file under code/ is held to by
% tools/check_structure.py --lines.
\usepackage[
  a4paper,
  left=3.1cm, right=3.1cm, top=2.6cm, bottom=3.0cm,
  head=26pt, headsep=0.8cm, footskip=1.4cm
]{geometry}

% ---------- Encoding & fonts ----------
\usepackage[T1]{fontenc}
\usepackage[utf8]{inputenc}
\IfFileExists{lmodern.sty}{\usepackage{lmodern}}{}
% newtxtext takes its TS1 symbols from TeX Gyre Termes (Debian: tex-gyre).
% A machine with newtx and without tex-gyre dies on the copyright page. CI's
% full TeX Live is the reference; locally install tex-gyre as well.
\IfFileExists{newtxtext.sty}{\usepackage{newtxtext}}{}
% amsmath before newtxmath; amssymb only when newtxmath is absent, because
% loading both is a fatal clash (\Bbbk already defined) that is invisible on
% a machine without newtx.
\usepackage{amsmath}
\IfFileExists{newtxmath.sty}{\usepackage{newtxmath}}{\usepackage{amssymb}}
\IfFileExists{inconsolata.sty}{\usepackage[varqu,scaled=0.95]{inconsolata}}{}
% Without lmodern, microtype runs with font expansion off, which is what
% resolves a marginal overfull line -- so a container short of lmodern
% disagrees with CI about line breaks. Install lmodern locally too.
\IfFileExists{lmodern.sty}{\usepackage{microtype}}{\usepackage[expansion=false]{microtype}}

% upquote: in T1 the input ' is a right single quote, and a listing that sets
% f('x') with curly quotes is a SyntaxError when typed back in.
\IfFileExists{upquote.sty}{\usepackage{upquote}}{}
% And its twin: in T1 the pair -- is an en-dash ligature in the typewriter
% family too, so \code{--flag} would print a flag nobody can type. Scoped to
% tt*, so prose keeps its own dashes.
\DisableLigatures[-]{encoding = T1, family = tt*}

% ---------- Language ----------
% babel with a missing language file is FATAL, so each .ldf is probed first.
\IfFileExists{babel.sty}{%
  \IfFileExists{polish.ldf}{%
    \IfFileExists{british.ldf}{%
      \ifdefstring{\booklang}{pl}%
        {\usepackage[british,main=polish]{babel}}%
        {\usepackage[polish,main=british]{babel}}%
    }{\usepackage[polish]{babel}}%
  }{%
    \IfFileExists{british.ldf}{\usepackage[british]{babel}}{}%
  }%
}{}
% babel-polish makes " an active shorthand, which is a category-code accident
% waiting to happen inside listings. Polish quotation marks come from
% \enquote, which csquotes sets per language.
\ifdefstring{\booklang}{pl}{\AtBeginDocument{\shorthandoff{"}}}{}

% ---------- Core packages ----------
\usepackage{graphicx}
\usepackage{xcolor}
\usepackage{booktabs}
\usepackage{tabularx}
\usepackage{longtable}
\usepackage{array}
\usepackage{enumitem}
\usepackage{listings}
\usepackage[most]{tcolorbox}
\usepackage{fancyhdr}
\usepackage{titlesec}
\usepackage{caption}
\usepackage{imakeidx}
\IfFileExists{csquotes.sty}{\usepackage[autostyle=true]{csquotes}}{%
  \providecommand{\enquote}[1]{``##1''}}
\usepackage[hidelinks]{hyperref}

% ---------- Palette ----------
% The family's palette, so the three books read as a set. Two colours are
% this book's own: the lab (green) and the stop-and-think box (violet).
\definecolor{mblue}{HTML}{1F4E79}
\definecolor{mteal}{HTML}{0E7C7B}
\definecolor{mamber}{HTML}{B26A00}
\definecolor{mred}{HTML}{9B2C2C}
\definecolor{mviolet}{HTML}{7B4397}
\definecolor{mgreen}{HTML}{3B7A3B}
\definecolor{mgrey}{HTML}{5A5A5A}
\definecolor{mlight}{HTML}{F4F6F8}
\definecolor{mfaint}{HTML}{FBFCFD}
\definecolor{codebg}{HTML}{FAFAFA}
\definecolor{codeframe}{HTML}{D8DEE4}
\definecolor{kw}{HTML}{0B5394}
\definecolor{str}{HTML}{9B2C2C}
\definecolor{cmt}{HTML}{4F7A4F}
\definecolor{typ}{HTML}{0E7C7B}
\definecolor{dec}{HTML}{7B4397}

\hypersetup{
  colorlinks=true,
  linkcolor=mblue,
  urlcolor=mteal,
  citecolor=mblue,
  pdfauthor={Konrad Cinkusz}
}

% \ifpl{Polish}{English} -- for the handful of places where a string is
% structural rather than content. Robust, because a fragile \ifpl inside a
% moving argument fails on the next run in exactly one of the two editions.
\DeclareRobustCommand{\ifpl}[2]{\ifdefstring{\booklang}{pl}{#1}{#2}}
\makeatletter
\pdfstringdefDisableCommands{%
  \ifdefstring{\booklang}{pl}{\let\ifpl\@firstoftwo}{\let\ifpl\@secondoftwo}%
  % hyperref drops \color from a bookmark and keeps its ARGUMENT, so a title
  % carrying \code{} would bookmark as "mtealx". Gobble the colour name.
  \def\color#1{}%
}
\makeatother

% ---------- Language strings ----------
\input{lang/\booklang}

% The PDF's subject, keywords and language read \lblPdf... from the language
% file, so they cannot sit in the colour block above, which runs first. Two
% blocks is the ordering, not an oversight.
\hypersetup{
  pdfsubject={\lblPdfSubject},
  pdfkeywords={\lblPdfKeywords},
  pdflang={\lblPdfLang}
}

% ---------- Headings ----------
\titleformat{\chapter}[display]
  {\normalfont\huge\bfseries\color{mblue}\raggedright}
  {\normalfont\normalsize\color{mgrey}\MakeUppercase{\chaptertitlename}\ \thechapter}
  {8pt}{\Huge\raggedright}
\titlespacing*{\chapter}{0pt}{10pt}{28pt}
\titleformat{\section}{\normalfont\Large\bfseries\color{mblue}\raggedright}{\thesection}{0.8em}{}
\titleformat{\subsection}{\normalfont\large\bfseries\color{mgrey}\raggedright}{\thesubsection}{0.7em}{}

\setcounter{tocdepth}{1}
\setcounter{secnumdepth}{2}

\pagestyle{fancy}
\fancyhf{}
\fancyhead[L]{\small\nouppercase{\leftmark}}
\fancyhead[R]{\small\thepage}
\renewcommand{\headrulewidth}{0.4pt}
\fancypagestyle{plain}{\fancyhf{}\fancyfoot[R]{\small\thepage}%
  \renewcommand{\headrulewidth}{0pt}}
% After \pagestyle{fancy}, which installs its own \chaptermark.
\renewcommand{\chaptermark}[1]{\markboth{\thechapter.\ #1}{}}

% ---------- Mathematics ----------
% The handful of notations the book uses, as macros so both editions set the
% same symbol. The book writes \ln for the natural logarithm and \log_{10}
% where the base is ten; a bare \log is ambiguous between the two readerships
% and is not used.
\newcommand{\R}{\mathbb{R}}
\newcommand{\dd}{\,\mathrm{d}}
\newcommand{\abs}[1]{\left\lvert #1\right\rvert}
\newcommand{\iu}{\mathrm{i}}

% ============================================================
%  CODE LISTINGS
%  Every listing in a chapter is a FILE under code/, pulled in with \pyfile or
%  \pyregion, and CI runs it on the pinned interpreter. An in-source
%  \begin{python} block is allowed for a fragment of three or four lines that
%  shows a shape rather than runs.
% ============================================================
\lstdefinelanguage{pythonx}{
  language=Python,
  morekeywords={async,await,match,case,nonlocal,yield,type},
  morekeywords=[2]{
    orbit,iterate,logistic,rk4_step,euler_step,integrate,lorenz,henon,
    lyapunov,separation,escape_time,box_count
  },
  sensitive=true
}

\lstdefinestyle{bookcode}{
  basicstyle=\ttfamily\footnotesize,
  keywordstyle=\color{kw}\bfseries,
  keywordstyle=[2]\color{typ},
  stringstyle=\color{str},
  % Colour only: inconsolata has no italic, and an italic comment style is
  % silently substituted and logged on every build.
  commentstyle=\color{cmt},
  numbers=left,
  numberstyle=\ttfamily\tiny\color{mgrey},
  numbersep=8pt,
  backgroundcolor=\color{codebg},
  frame=single,
  rulecolor=\color{codeframe},
  framesep=6pt,
  breaklines=true,
  breakatwhitespace=false,
  postbreak=\mbox{\textcolor{mgrey}{$\hookrightarrow$}\space},
  showstringspaces=false,
  tabsize=4,
  columns=fullflexible,
  keepspaces=true,
  captionpos=b,
  aboveskip=1.2em,
  belowskip=1.2em,
  % Do NOT add a mapping for U+00A0: it aborts the run with a message that
  % names neither the character nor this line. Listings are ASCII, and
  % parity's C13 checks it.
  literate={ł}{{\l}}1 {Ł}{{\L}}1 {ą}{{\k a}}1 {ę}{{\k e}}1
           {ś}{{\'s}}1 {ć}{{\'c}}1 {ż}{{\.z}}1 {ź}{{\'z}}1 {ó}{{\'o}}1 {ń}{{\'n}}1
           {Ą}{{\k A}}1 {Ę}{{\k E}}1 {Ś}{{\'S}}1 {Ć}{{\'C}}1
           {Ż}{{\.Z}}1 {Ź}{{\'Z}}1 {Ó}{{\'O}}1 {Ń}{{\'N}}1
           {—}{{-{}-}}2 {–}{{-}}1 {‘}{{'}}1 {’}{{'}}1
           {“}{{"}}1 {”}{{"}}1 {°}{{\textdegree}}1
}
\lstset{style=bookcode}

\lstnewenvironment{python}[1][]
  {\lstset{language=pythonx,style=bookcode,#1}}
  {}
\lstnewenvironment{shellcmd}[1][]
  {\lstset{language=bash,style=bookcode,numbers=none,keywordstyle=\color{mteal}\bfseries,#1}}
  {}

% The path is printed above every file listing, in the form the reader opens
% in the editor. \nobreak keeps it on the page of the listing it names.
% The path must not be stranded at the foot of a page with its listing
% overleaf: \nobreak alone does not hold, because listings opens with its own
% vertical material. \needspace asks for room for the path and the first few
% lines of code, and turns the page first when there is not.
% \Needspace (capital N) measures \pagegoal-\pagetotal exactly; lower-case
% \needspace works through glue and penalties and was measured, by the
% Chapter 11 pass, turning pages early and leaving gaps. It must be used
% between paragraphs, which \listingpath's own \par guarantees. Chapters use
% \Needspace* too, so the fallback for a machine without the package defines
% both forms (the star swallowed with its argument).
\makeatletter
\IfFileExists{needspace.sty}{\usepackage{needspace}}{%
  \providecommand{\needspace}[1]{}%
  \providecommand{\Needspace}{\@ifstar\@gobble\@gobble}}
\makeatother
\newcommand{\listingpath}[1]{%
  \par\medskip\Needspace{6\baselineskip}\noindent
  {\ttfamily\footnotesize\color{mgrey}\detokenize{#1}}%
  \par\nobreak}

% Usage: \pyfile{code/ch05/cobweb.py}{Caption}{lst:label}
\newcommand{\pyfile}[3]{%
  \listingpath{#1}%
  \lstinputlisting[language=pythonx,style=bookcode,aboveskip=0.3em,
    caption={#2},label={#3}]{#1}%
}

% A named region of a file, delimited by `# --8<-- [start:name]` and
% `# --8<-- [end:name]`. Matched through listings' range prefix and suffix
% keys -- the other obvious form matches nothing and prints the WHOLE file
% without a warning (measured in the sibling book). A region name is a word,
% never a number. The line numbers printed are the file's own.
% Usage: \pyregion{code/ch05/cobweb.py}{step}{Caption}{lst:label}
\lstset{
  rangebeginprefix=\#\ --8<--\ [start:, rangebeginsuffix=],
  rangeendprefix=\#\ --8<--\ [end:, rangeendsuffix=],
  includerangemarker=false}
\newcommand{\pyregion}[4]{%
  \listingpath{#1}%
  \lstinputlisting[language=pythonx,style=bookcode,aboveskip=0.3em,
    linerange={#2-#2},caption={#3},label={#4}]{#1}%
}

% A program's output, written by code/measure/transcripts.py. There is no
% typed form and there must not be: \val{} does not expand in a verbatim
% body, and a transcript nobody ran looks exactly like one that was.
\newcommand{\transcript}[1]{%
  \par\smallskip
  \IfFileExists{figures/transcripts/#1.txt}{%
    \lstinputlisting[style=bookcode,numbers=none,basicstyle=\ttfamily\footnotesize,
      keywordstyle=\color{black},backgroundcolor=\color{white},
      language={}]{figures/transcripts/#1.txt}%
  }{%
    \begin{tcolorbox}[enhanced, sharp corners, width=\linewidth,
      colback=mlight, colframe=mgrey, boxrule=0.6pt,
      left=8pt, right=8pt, top=6pt, bottom=6pt]
      {\bfseries\color{mgrey}\small \lblTranscriptMissing}\\[2pt]
      {\ttfamily\scriptsize\color{mgrey} figures/transcripts/#1.txt}
    \end{tcolorbox}%
  }%
  \par\smallskip
}

% Inline literals. Underscores inside these must be written \_.
\newcommand{\code}[1]{\texttt{\small #1}}
\newcommand{\pkg}[1]{\texttt{\small\color{mblue}#1}}
\newcommand{\api}[1]{\texttt{\small\color{mteal}#1}}

% ============================================================
%  ADMONITIONS
%  A small, fixed vocabulary. `tinted' may be read in passing; `outlined' is
%  a block the reader must stop at. A third treatment would mean the first two
%  are not carrying their meaning.
% ============================================================
\tcbset{
  cfz tinted/.style={
    enhanced, breakable, sharp corners,
    lines before break=4,
    colback=mlight, boxrule=0pt, leftrule=3pt,
    left=8pt, right=8pt, top=6pt, bottom=6pt,
    before skip=13pt, after skip=13pt,
    fonttitle=\bfseries\small,
    coltitle=black, colbacktitle=mlight,
    attach title to upper={\par\smallskip},
  },
  cfz outlined/.style={
    enhanced, breakable, sharp corners,
    % A box broken straight after its title is a heading stranded at the foot
    % of a page with its question overleaf. Four lines is title plus three.
    lines before break=4,
    colback=white, boxrule=0.6pt,
    left=8pt, right=8pt, top=6pt, bottom=6pt,
    before skip=13pt, after skip=13pt,
    fonttitle=\bfseries\small,
    coltitle=black, colbacktitle=white,
    attach title to upper={\par\smallskip},
  },
}

% Read in passing.
\newtcolorbox{note}[1][]{cfz tinted, colframe=mblue, title={\lblNote}, #1}
\newtcolorbox{warning}[1][]{cfz tinted, colframe=mred, title={\lblWarning}, #1}
% Where you meet this in the world. If it cannot name a specific system,
% instrument, paper or event, it is prose and not this box.
\newtcolorbox{worldbox}[1][]{cfz tinted, colframe=mteal, title={\lblWorld}, #1}
% What is stated and not proved here, and where the proof lives.
\newtcolorbox{rigourbox}[1][]{cfz tinted, colframe=mgrey, title={\lblRigour}, #1}
% The companion library the reader builds, under code/src/chaoslab/.
\newtcolorbox{projectbox}[1][]{cfz tinted, colframe=mamber, title={\lblProject}, #1}

% Stop and read.
% The misconception, named AFTER the reader has been asked the question it
% gets wrong. Appendix B is the catalogue; a trap box is one entry in place.
\newtcolorbox{trapbox}[1][]{cfz outlined, colframe=mred, title={\lblTrap}, #1}
% A listing that has not been run on the pinned versions. Counted by
% `make debt'; nothing ships to a reader with one attached.
\newtcolorbox{verifybox}[1][]{cfz outlined, colframe=mgrey, title={\lblVerify}, #1}
\newtcolorbox{labbox}[1][]{cfz outlined, colframe=mgreen, title={\lblLab}, #1}

% ============================================================
%  THE METHOD --- questions before answers
%  ---------------------------------------------------------
%  The unit of teaching is a question the reader answers BEFORE reading the
%  answer. That is retrieval practice and errorful generation, the two
%  findings the programmed-learning format is defended by, and it is the one
%  thing the sibling mathematics book does that this one keeps whole.
%
%  \begin{think} ... \end{think}   the question, numbered per chapter, ending
%                                  with the instruction to answer before
%                                  reading on;
%  \reveal{short answer}           a centred box with the answer, which is the
%                                  first thing after the question -- the edge
%                                  the reader covers with a hand, or scrolls
%                                  to;
%  \begin{revealblock} ... \end{revealblock}  the same, for a worked answer.
%
%  tools/check_structure.py --think fails when a think box is not followed by
%  a reveal before the next think box or section, and when a reveal answers
%  nothing. That is a property of the source, so it is checked mechanically.
%
%  The instruction says "before you read on", never "before you turn over":
%  the two editions paginate differently, so the second is true in one of them.
% ============================================================
\newcounter{think}[chapter]
\renewcommand{\thethink}{\thechapter.\arabic{think}}
\newtcolorbox{thinkbox}[1][]{cfz outlined, colframe=mviolet, #1}
\NewDocumentEnvironment{think}{}{%
  \refstepcounter{think}%
  \begin{thinkbox}[title={\lblThink~\thethink}]%
}{%
  \par\smallskip
  {\raggedleft\itshape\small\color{mviolet}\lblThinkCue\par}%
  \end{thinkbox}%
}
\tcbset{
  cfz answer/.style={
    enhanced, sharp corners,
    colback=mfaint, colframe=mviolet!55, boxrule=0.5pt,
    before skip=4pt, after skip=10pt,
  },
}
\newcommand{\reveal}[1]{%
  \par\nobreak
  \begin{tcolorbox}[cfz answer, unbreakable,
    left=8pt, right=8pt, top=5pt, bottom=5pt,
    width=0.86\linewidth, center, halign=flush center]
    {\small\bfseries\color{mviolet}\lblAnswerMark}\enspace #1
  \end{tcolorbox}%
  \nobreak\noindent\ignorespaces
}
\NewDocumentEnvironment{revealblock}{}{%
  \par\nopagebreak
  \begin{tcolorbox}[cfz answer, breakable,
    left=10pt, right=10pt, top=6pt, bottom=6pt]
  {\small\bfseries\color{mviolet}\lblAnswerMark}\par\smallskip
}{%
  \end{tcolorbox}%
  \nopagebreak\par\noindent\ignorespaces
}

% ---- Outcomes on entry, and "Can you?" generated from them ----------------
% The self-check at the end of a chapter IS the list of outcomes at its
% start, one for one, because it is built from the same tokens: each
% \outcome both prints its item and appends a row to \cfz@canrows. A \loop
% inside a tabularx body does not survive alignment scanning, which is why
% the rows are accumulated at declaration time (the sibling book's finding).
\makeatletter
\newcommand{\cfz@canrows}{}
\newcommand{\cfz@ratebox}{%
  \fbox{\rule{0pt}{0.9ex}\rule{0.9ex}{0pt}}}
\newcommand{\cfz@rates}{%
  {\footnotesize\color{mgrey}%
   1\,\cfz@ratebox\ 2\,\cfz@ratebox\ 3\,\cfz@ratebox\ 4\,\cfz@ratebox\ 5\,\cfz@ratebox}}
\NewDocumentEnvironment{outcomes}{}{%
  \gdef\cfz@canrows{}%
  \begin{tcolorbox}[cfz tinted, colframe=mblue, title={\lblOutcomes}]%
  \lblOutcomesLead
  \begin{itemize}[leftmargin=1.4em, itemsep=2pt, topsep=4pt]%
}{%
  \end{itemize}%
  \end{tcolorbox}%
}
\newcommand{\outcome}[1]{%
  \item #1%
  \gappto\cfz@canrows{#1 & \cfz@rates \\ \addlinespace[3pt]}%
}
\newcommand{\canyou}{%
  \section*{\lblCanYou}%
  \phantomsection\addcontentsline{toc}{section}{\lblCanYou}%
  \noindent\lblCanYouIntro\par\medskip
  \ifdefempty{\cfz@canrows}{%
    \noindent{\color{mred}\lblNoOutcomes}\par}{%
  {\small
  \noindent\begin{tabularx}{\linewidth}{@{}>{\raggedright\arraybackslash}X r@{}}
    \toprule
    \cfz@canrows
    \bottomrule
  \end{tabularx}\par}}%
}
\makeatother

% ---- Summary -------------------------------------------------------------
\newtcolorbox{summarybox}[1][]{cfz tinted, colframe=mblue, title={\lblSummary}, #1}

% ---- Checkpoint questions, answered at the back --------------------------
% Answers are written AT THE QUESTION with \answerto{} and printed in
% Appendix A. They are carried there by a global token list rather than an
% auxiliary file, because an answer contains mathematics and \val{}, and
% material written to an .aux is expanded at shipout, where fragile commands
% break with errors naming neither the answer nor the chapter.
% \include is sequential within a run, so the list is complete by the time
% the appendix is typeset.
\makeatletter
\newcommand{\cfz@answers}{}
\newcommand{\cfz@lastansch}{0}
\NewDocumentEnvironment{checkpoint}{}{%
  \section*{\lblCheckpoint}%
  \phantomsection\addcontentsline{toc}{section}{\lblCheckpoint}%
  \noindent\lblCheckpointIntro\par
  \begin{enumerate}[label=\textbf{\arabic*.}, leftmargin=2em, itemsep=4pt]%
}{%
  \end{enumerate}%
}
\newcommand{\answerto}[1]{%
  \ifnum\value{chapter}=\cfz@lastansch\relax\else
    \xdef\cfz@lastansch{\arabic{chapter}}%
    \xappto\cfz@answers{\noexpand\cfz@anschapter{\thechapter}}%
  \fi
  \xappto\cfz@answers{\noexpand\cfz@ansitem{\arabic{enumi}}}%
  \gappto\cfz@answers{{#1}}%
}
\newcommand{\cfz@anschapter}[1]{%
  \par\medskip\noindent{\bfseries\color{mblue}\chaptername~#1}\par\nobreak\smallskip}
\newcommand{\cfz@ansitem}[2]{%
  \par\noindent\hangindent=2em\hangafter=1
  \makebox[2em][l]{\textbf{#1.}}#2\par}
\newcommand{\answersbody}{%
  \ifdefempty{\cfz@answers}{\noindent{\color{mred}\lblNoAnswers}\par}{%
    {\raggedright\cfz@answers}}%
}
\makeatother

% ============================================================
%  LABS
%  \begin{lab}{key}{title} ... \end{lab}
%
%  A lab is not a paragraph but a FILE the reader opens: a starter under
%  code/labs/chNN/<key>.py that must FAIL its test, a reference solution under
%  code/labs/chNN/solutions/, and test_<key>.py beside it. The box prints the
%  paths and the one command that runs the test. An environment rather than a
%  macro, because a macro argument cannot hold a verbatim fragment.
%  tools/check_structure.py --labs fails when a file is missing or the key's
%  chapter number disagrees with its chapter.
% ============================================================
\newcounter{lab}[chapter]
\renewcommand{\thelab}{\thechapter.\arabic{lab}}
\makeatletter
\newcommand{\labdir}{ch\two@digits{\value{chapter}}}
\newcommand{\listoflabs}{%
  \section*{\lblLabManifest}%
  \phantomsection\addcontentsline{toc}{section}{\lblLabManifest}%
  \begingroup\c@tocdepth=2\relax\@starttoc{lab}\endgroup%
}
\makeatother
\NewDocumentEnvironment{lab}{mm}{%
  \refstepcounter{lab}%
  \addcontentsline{lab}{subsection}{\protect\numberline{\thelab}\texttt{\protect\detokenize{#1}} --- #2}%
  \begin{labbox}[title={\lblLab~\thelab\ \dash{} #2}]%
}{%
  \par\medskip
  {\small\setlength{\parindent}{0pt}%
  \lblLabStarter: \texttt{code/labs/\labdir/\detokenize{#1}.py}\\
  \lblLabTest: \texttt{code/labs/\labdir/test\_\detokenize{#1}.py}\\
  \lblLabRun: \texttt{cd code \&\& uv run pytest -k \detokenize{#1}}\par}%
  \end{labbox}%
}

% ============================================================
%  FIGURES --- plots and diagrams
%  ---------------------------------------------------------
%  \plotfig{key}{caption}{one-line manifest description}
%    A PLOT, generated by code/plots/<chapter>.py from the same library the
%    listings use, into figures/plots/<lang>/<key>.pdf. Per language, because
%    an axis label in the wrong language is the half-translation that makes a
%    bilingual book look machine-made, and because the Polish edition owes a
%    decimal comma on its tick labels too. Build output, gitignored.
%
%  \mermaidfig{key}{caption}{one-line manifest description}
%    A DIAGRAM of an idea rather than of data, from
%    figures/mermaid/<lang>/<key>.mmd, rendered by `make diagrams'.
%
%  If the rendered PDF is absent either macro prints a visible marker in the
%  flow rather than a float, so a missing figure cannot be mistaken for a
%  layout choice. \phantomsection before \addcontentsline is load-bearing:
%  without it the manifest entry records whatever anchor preceded the figure,
%  which after a listing is a line that was never printed.
% ============================================================
\makeatletter
\newcommand{\listofplots}{%
  \section*{\lblPlotManifest}%
  \phantomsection\addcontentsline{toc}{section}{\lblPlotManifest}%
  \begingroup\c@tocdepth=2\relax\@starttoc{plt}\endgroup%
}
\newcommand{\listofdiagrams}{%
  \section*{\lblDiagramManifest}%
  \phantomsection\addcontentsline{toc}{section}{\lblDiagramManifest}%
  \begingroup\c@tocdepth=2\relax\@starttoc{dgm}\endgroup%
}
\makeatother

\newcommand{\figmissing}[3]{% #1 message, #2 path, #3 caption
  \par\medskip
  \begin{tcolorbox}[enhanced, breakable, width=\linewidth, sharp corners,
    colback=white, colframe=mgrey, boxrule=0.8pt,
    left=8pt, right=8pt, top=8pt, bottom=8pt]
    {\bfseries\color{mgrey}\small #1}\\[4pt]
    {\ttfamily\scriptsize\color{mgrey} #2}
  \end{tcolorbox}%
  \captionof{figure}{#3}%
  \par\medskip
}

\newcommand{\plotfig}[3]{%
  \phantomsection%
  \addcontentsline{plt}{subsection}{\protect\numberline{}\texttt{#1} --- #3}%
  \IfFileExists{figures/plots/\booklang/#1.pdf}{%
    \begin{figure}[htbp]
    \centering
    \includegraphics[width=\linewidth,height=0.45\textheight,keepaspectratio]{figures/plots/\booklang/#1.pdf}%
    \caption{#2}\label{fig:#1}
    \end{figure}%
  }{%
    \figmissing{\lblPlotNotRendered}{figures/plots/\booklang/#1.pdf}{#2}\label{fig:#1}%
  }%
}

\newcommand{\mermaidfig}[3]{%
  \phantomsection%
  \addcontentsline{dgm}{subsection}{\protect\numberline{}\texttt{#1.mmd} --- #3}%
  \IfFileExists{figures/diagrams/\booklang/#1.pdf}{%
    \begin{figure}[htbp]
    \centering
    \includegraphics[width=0.95\linewidth,height=0.42\textheight,keepaspectratio]{figures/diagrams/\booklang/#1.pdf}%
    \caption{#2}\label{fig:#1}
    \end{figure}%
  }{%
    \figmissing{\lblNotRendered}{figures/mermaid/\booklang/#1.mmd}{#2}\label{fig:#1}%
  }%
}

% ============================================================
%  APPENDIX B: the misreadings, generated by tools/gen_traps.py
%
%  \trapchapter{Chapter~\ref{ch:x}}{title}    a heading per chapter
%  \trapentry{n}{habit}{truth}{\ref{sec:x}}  one misreading
%
%  Set ragged right: an entry is mostly short sentences with \code{} spans in
%  them, and a justified line full of unbreakable runs is an overfull box
%  waiting for the edition with the longer words.
% ============================================================
\newcommand{\trapchapter}[2]{%
  \par\medskip
  {\large\bfseries\color{mblue}#1\quad\normalfont\itshape\color{mgrey}#2\par}%
  \nopagebreak\smallskip}
\newcommand{\trapentry}[4]{%
  \par\noindent
  \begingroup\raggedright\hangindent=2.4em\hangafter=1
  \makebox[2.4em][l]{\textbf{\color{mred}#1}}%
  \textbf{\enquote{#2}}\enspace #3\enspace{\color{mgrey}\S\,#4}\par
  \endgroup\smallskip}

% ============================================================
%  CHAPTER STUBS
% ============================================================
\newcommand{\chapterstub}[1]{%
  \begin{tcolorbox}[
    enhanced, breakable, width=\linewidth, sharp corners,
    colback=mlight, colframe=mred, boxrule=1pt,
    borderline={0.6pt}{3pt}{mred, dashed},
    left=12pt, right=12pt, top=12pt, bottom=12pt]
    {\bfseries\color{mred}\large \lblNotWritten}\\[8pt]
    \small\raggedright #1
  \end{tcolorbox}%
}

% ============================================================
%  COMPUTED VALUES
%  A number the reader cannot do in their head is computed by a script under
%  code/measure/, written to figures/values/<chapter>.tex as
%      \pyval{ch06.steps.apart}{34}
%  and reaches the page as \val{ch06.steps.apart}. The book contains
%  references, not digits; `make verify' fails when a committed value no
%  longer matches its script. \val applies the edition's decimal marker, which
%  is how the Polish edition gets its comma without anybody remembering.
% ============================================================
% The comma is BRACED: siunitx sets numbers in maths mode, where a bare comma
% is punctuation and gets a thin space after it ("9, 81"). {,} is an ordinary
% symbol, which is what a decimal marker is.
\newcommand{\bookdecimalmarker}{\ifpl{{,}}{.}}
\IfFileExists{siunitx.sty}{%
  \usepackage{siunitx}%
  \sisetup{
    group-digits = integer,
    group-minimum-digits = 5,
    group-separator = {\,},
    output-decimal-marker = {\bookdecimalmarker},
    exponent-product = \times,
    print-unity-mantissa = false
  }%
  \newcommand{\booknum}[1]{\num{##1}}%
}{%
  \newcommand{\booknum}[1]{##1}%
  \providecommand{\num}[1]{##1}%
}

\makeatletter
\newcommand{\pyvalues}{}
\newcommand{\pyval}[2]{\csgdef{pyval@#1}{#2}\listxadd{\pyvalues}{#1}}
\newcommand{\pyvaltext}[2]{%
  \csgdef{pyval@#1}{#2}\csgdef{pyvaltext@#1}{}\listxadd{\pyvalues}{#1}}
% \num on a non-numeric body is a FATAL siunitx error, so \val refuses a text
% value by name rather than letting siunitx fail on it.
\newcommand{\val}[1]{%
  \ifcsdef{pyval@#1}{%
    \ifcsdef{pyvaltext@#1}{%
      \PackageError{chaosbook}{Value `#1' is not a number}%
        {Use \string\valtext{#1} instead.}%
    }{}%
    \booknum{\csuse{pyval@#1}}%
  }{\textcolor{mred}{\textbf{[\lblValueMissing: #1]}}}%
}
\newcommand{\valtext}[1]{%
  \ifcsdef{pyval@#1}{\csuse{pyval@#1}}%
    {\textcolor{mred}{\textbf{[\lblValueMissing: #1]}}}%
}
\makeatother

% figures/values/all.tex is generated by `make numbers' and \inputs each
% value file. A build with no values still compiles; every \val{} then prints
% a visible marker. \valuesindex lets the one-chapter build point elsewhere.
\providecommand{\valuesindex}{figures/values/all}
\IfFileExists{\valuesindex.tex}{\input{\valuesindex}}{%
  \AtBeginDocument{\PackageWarning{chaosbook}{No computed values found. Run `make numbers'.}}}

% ============================================================
%  PINNED VERSIONS --- single source of truth
%  Change these here AND in code/pyproject.toml; tools/check_structure.py
%  --pins fails when the two disagree. A chapter never types a version.
% ============================================================
\newcommand{\bookedition}{0.1}
\newcommand{\bookyear}{2026}
\newcommand{\pyver}{3.14}              % the minor the book targets
\newcommand{\pypatch}{3.14.7}          % the exact interpreter the code ran on
\newcommand{\uvver}{0.12.13}           % uv
\newcommand{\numpyver}{2.5.3}          % numpy
\newcommand{\matplotlibver}{3.11.2}    % matplotlib
\newcommand{\pytestver}{9.1.1}         % pytest
\newcommand{\ruffver}{0.16.9}          % ruff
% The prose licence is a name with a version in it, not a quantity: it is a
% macro so that neither edition typesets it as a decimal to be localised.
\newcommand{\proselicence}{CC~BY-NC-SA~4.0}

% ============================================================
%  INDEX
% ============================================================
\apptocmd{\theindex}{\raggedcolumns\raggedright}{}{%
  \PackageWarning{chaosbook}{Could not patch theindex for ragged setting}}
\providecommand*\see[2]{}
\renewcommand*\see[2]{\emph{\lblIndexSee} #1}
\providecommand*\seealso[2]{}
\renewcommand*\seealso[2]{\emph{\lblIndexSeeAlso} #1}

\makeindex
; nothing else differs between them at this level.
%  Every user-visible string lives in lang/en.tex or lang/pl.tex, so a label
%  that appears in one edition and not the other is a missing macro rather
%  than a silent divergence.
%
%  Lineage. This book is the third in a family and takes one thing from each
%  of the other two:
%
%    * from Mathematics from Zero for the AI Engineer: the bilingual single
%      source (\booklang, lang/, one body.tex), the rule that every number is
%      computed and reaches the page as \val{}, and the programmed-learning
%      method -- outcomes on entry, a question before every answer, a
%      misconception elicited before it is named, a self-check generated from
%      the outcomes, and checkpoint answers collected at the back;
%    * from Python for .NET Engineers: every listing is a FILE that CI runs,
%      every lab is a starter that must fail and a solution that must pass,
%      and a console transcript is a computed number that happens to be
%      verbatim.
%
%  What is new here is \plotfig. A book about chaos cannot be taught in boxes
%  and arrows: the bifurcation diagram and the Lorenz attractor ARE the
%  argument, so the plots are generated by code/plots/ from the same library
%  the listings use, per language, and never drawn by hand.
% ============================================================

% \booklang must be set by the main file. Default to English rather than
% failing, so a stray % ============================================================
%  preamble.tex --- styling, environments, macros
%
%  Shared by BOTH editions. main-en.tex and main-pl.tex each define \booklang
%  before % ============================================================
%  preamble.tex --- styling, environments, macros
%
%  Shared by BOTH editions. main-en.tex and main-pl.tex each define \booklang
%  before \input{preamble}; nothing else differs between them at this level.
%  Every user-visible string lives in lang/en.tex or lang/pl.tex, so a label
%  that appears in one edition and not the other is a missing macro rather
%  than a silent divergence.
%
%  Lineage. This book is the third in a family and takes one thing from each
%  of the other two:
%
%    * from Mathematics from Zero for the AI Engineer: the bilingual single
%      source (\booklang, lang/, one body.tex), the rule that every number is
%      computed and reaches the page as \val{}, and the programmed-learning
%      method -- outcomes on entry, a question before every answer, a
%      misconception elicited before it is named, a self-check generated from
%      the outcomes, and checkpoint answers collected at the back;
%    * from Python for .NET Engineers: every listing is a FILE that CI runs,
%      every lab is a starter that must fail and a solution that must pass,
%      and a console transcript is a computed number that happens to be
%      verbatim.
%
%  What is new here is \plotfig. A book about chaos cannot be taught in boxes
%  and arrows: the bifurcation diagram and the Lorenz attractor ARE the
%  argument, so the plots are generated by code/plots/ from the same library
%  the listings use, per language, and never drawn by hand.
% ============================================================

% \booklang must be set by the main file. Default to English rather than
% failing, so a stray \input{preamble} still compiles and says what it did.
\providecommand{\booklang}{en}

\usepackage{etoolbox}   % loaded first: the language switch below needs it

% ---------- Page geometry ----------
% ONE format: A4, 12pt, ONE-SIDED. The book is read on a screen, usually with
% an editor open beside it, so there is no recto, no verso, no gutter and no
% blank leaf before a chapter. The measure is about seventy-five characters
% of prose. At \footnotesize a listing line of 79 characters fits without
% wrapping, and 79 is the width every file under code/ is held to by
% tools/check_structure.py --lines.
\usepackage[
  a4paper,
  left=3.1cm, right=3.1cm, top=2.6cm, bottom=3.0cm,
  head=26pt, headsep=0.8cm, footskip=1.4cm
]{geometry}

% ---------- Encoding & fonts ----------
\usepackage[T1]{fontenc}
\usepackage[utf8]{inputenc}
\IfFileExists{lmodern.sty}{\usepackage{lmodern}}{}
% newtxtext takes its TS1 symbols from TeX Gyre Termes (Debian: tex-gyre).
% A machine with newtx and without tex-gyre dies on the copyright page. CI's
% full TeX Live is the reference; locally install tex-gyre as well.
\IfFileExists{newtxtext.sty}{\usepackage{newtxtext}}{}
% amsmath before newtxmath; amssymb only when newtxmath is absent, because
% loading both is a fatal clash (\Bbbk already defined) that is invisible on
% a machine without newtx.
\usepackage{amsmath}
\IfFileExists{newtxmath.sty}{\usepackage{newtxmath}}{\usepackage{amssymb}}
\IfFileExists{inconsolata.sty}{\usepackage[varqu,scaled=0.95]{inconsolata}}{}
% Without lmodern, microtype runs with font expansion off, which is what
% resolves a marginal overfull line -- so a container short of lmodern
% disagrees with CI about line breaks. Install lmodern locally too.
\IfFileExists{lmodern.sty}{\usepackage{microtype}}{\usepackage[expansion=false]{microtype}}

% upquote: in T1 the input ' is a right single quote, and a listing that sets
% f('x') with curly quotes is a SyntaxError when typed back in.
\IfFileExists{upquote.sty}{\usepackage{upquote}}{}
% And its twin: in T1 the pair -- is an en-dash ligature in the typewriter
% family too, so \code{--flag} would print a flag nobody can type. Scoped to
% tt*, so prose keeps its own dashes.
\DisableLigatures[-]{encoding = T1, family = tt*}

% ---------- Language ----------
% babel with a missing language file is FATAL, so each .ldf is probed first.
\IfFileExists{babel.sty}{%
  \IfFileExists{polish.ldf}{%
    \IfFileExists{british.ldf}{%
      \ifdefstring{\booklang}{pl}%
        {\usepackage[british,main=polish]{babel}}%
        {\usepackage[polish,main=british]{babel}}%
    }{\usepackage[polish]{babel}}%
  }{%
    \IfFileExists{british.ldf}{\usepackage[british]{babel}}{}%
  }%
}{}
% babel-polish makes " an active shorthand, which is a category-code accident
% waiting to happen inside listings. Polish quotation marks come from
% \enquote, which csquotes sets per language.
\ifdefstring{\booklang}{pl}{\AtBeginDocument{\shorthandoff{"}}}{}

% ---------- Core packages ----------
\usepackage{graphicx}
\usepackage{xcolor}
\usepackage{booktabs}
\usepackage{tabularx}
\usepackage{longtable}
\usepackage{array}
\usepackage{enumitem}
\usepackage{listings}
\usepackage[most]{tcolorbox}
\usepackage{fancyhdr}
\usepackage{titlesec}
\usepackage{caption}
\usepackage{imakeidx}
\IfFileExists{csquotes.sty}{\usepackage[autostyle=true]{csquotes}}{%
  \providecommand{\enquote}[1]{``##1''}}
\usepackage[hidelinks]{hyperref}

% ---------- Palette ----------
% The family's palette, so the three books read as a set. Two colours are
% this book's own: the lab (green) and the stop-and-think box (violet).
\definecolor{mblue}{HTML}{1F4E79}
\definecolor{mteal}{HTML}{0E7C7B}
\definecolor{mamber}{HTML}{B26A00}
\definecolor{mred}{HTML}{9B2C2C}
\definecolor{mviolet}{HTML}{7B4397}
\definecolor{mgreen}{HTML}{3B7A3B}
\definecolor{mgrey}{HTML}{5A5A5A}
\definecolor{mlight}{HTML}{F4F6F8}
\definecolor{mfaint}{HTML}{FBFCFD}
\definecolor{codebg}{HTML}{FAFAFA}
\definecolor{codeframe}{HTML}{D8DEE4}
\definecolor{kw}{HTML}{0B5394}
\definecolor{str}{HTML}{9B2C2C}
\definecolor{cmt}{HTML}{4F7A4F}
\definecolor{typ}{HTML}{0E7C7B}
\definecolor{dec}{HTML}{7B4397}

\hypersetup{
  colorlinks=true,
  linkcolor=mblue,
  urlcolor=mteal,
  citecolor=mblue,
  pdfauthor={Konrad Cinkusz}
}

% \ifpl{Polish}{English} -- for the handful of places where a string is
% structural rather than content. Robust, because a fragile \ifpl inside a
% moving argument fails on the next run in exactly one of the two editions.
\DeclareRobustCommand{\ifpl}[2]{\ifdefstring{\booklang}{pl}{#1}{#2}}
\makeatletter
\pdfstringdefDisableCommands{%
  \ifdefstring{\booklang}{pl}{\let\ifpl\@firstoftwo}{\let\ifpl\@secondoftwo}%
  % hyperref drops \color from a bookmark and keeps its ARGUMENT, so a title
  % carrying \code{} would bookmark as "mtealx". Gobble the colour name.
  \def\color#1{}%
}
\makeatother

% ---------- Language strings ----------
\input{lang/\booklang}

% The PDF's subject, keywords and language read \lblPdf... from the language
% file, so they cannot sit in the colour block above, which runs first. Two
% blocks is the ordering, not an oversight.
\hypersetup{
  pdfsubject={\lblPdfSubject},
  pdfkeywords={\lblPdfKeywords},
  pdflang={\lblPdfLang}
}

% ---------- Headings ----------
\titleformat{\chapter}[display]
  {\normalfont\huge\bfseries\color{mblue}\raggedright}
  {\normalfont\normalsize\color{mgrey}\MakeUppercase{\chaptertitlename}\ \thechapter}
  {8pt}{\Huge\raggedright}
\titlespacing*{\chapter}{0pt}{10pt}{28pt}
\titleformat{\section}{\normalfont\Large\bfseries\color{mblue}\raggedright}{\thesection}{0.8em}{}
\titleformat{\subsection}{\normalfont\large\bfseries\color{mgrey}\raggedright}{\thesubsection}{0.7em}{}

\setcounter{tocdepth}{1}
\setcounter{secnumdepth}{2}

\pagestyle{fancy}
\fancyhf{}
\fancyhead[L]{\small\nouppercase{\leftmark}}
\fancyhead[R]{\small\thepage}
\renewcommand{\headrulewidth}{0.4pt}
\fancypagestyle{plain}{\fancyhf{}\fancyfoot[R]{\small\thepage}%
  \renewcommand{\headrulewidth}{0pt}}
% After \pagestyle{fancy}, which installs its own \chaptermark.
\renewcommand{\chaptermark}[1]{\markboth{\thechapter.\ #1}{}}

% ---------- Mathematics ----------
% The handful of notations the book uses, as macros so both editions set the
% same symbol. The book writes \ln for the natural logarithm and \log_{10}
% where the base is ten; a bare \log is ambiguous between the two readerships
% and is not used.
\newcommand{\R}{\mathbb{R}}
\newcommand{\dd}{\,\mathrm{d}}
\newcommand{\abs}[1]{\left\lvert #1\right\rvert}
\newcommand{\iu}{\mathrm{i}}

% ============================================================
%  CODE LISTINGS
%  Every listing in a chapter is a FILE under code/, pulled in with \pyfile or
%  \pyregion, and CI runs it on the pinned interpreter. An in-source
%  \begin{python} block is allowed for a fragment of three or four lines that
%  shows a shape rather than runs.
% ============================================================
\lstdefinelanguage{pythonx}{
  language=Python,
  morekeywords={async,await,match,case,nonlocal,yield,type},
  morekeywords=[2]{
    orbit,iterate,logistic,rk4_step,euler_step,integrate,lorenz,henon,
    lyapunov,separation,escape_time,box_count
  },
  sensitive=true
}

\lstdefinestyle{bookcode}{
  basicstyle=\ttfamily\footnotesize,
  keywordstyle=\color{kw}\bfseries,
  keywordstyle=[2]\color{typ},
  stringstyle=\color{str},
  % Colour only: inconsolata has no italic, and an italic comment style is
  % silently substituted and logged on every build.
  commentstyle=\color{cmt},
  numbers=left,
  numberstyle=\ttfamily\tiny\color{mgrey},
  numbersep=8pt,
  backgroundcolor=\color{codebg},
  frame=single,
  rulecolor=\color{codeframe},
  framesep=6pt,
  breaklines=true,
  breakatwhitespace=false,
  postbreak=\mbox{\textcolor{mgrey}{$\hookrightarrow$}\space},
  showstringspaces=false,
  tabsize=4,
  columns=fullflexible,
  keepspaces=true,
  captionpos=b,
  aboveskip=1.2em,
  belowskip=1.2em,
  % Do NOT add a mapping for U+00A0: it aborts the run with a message that
  % names neither the character nor this line. Listings are ASCII, and
  % parity's C13 checks it.
  literate={ł}{{\l}}1 {Ł}{{\L}}1 {ą}{{\k a}}1 {ę}{{\k e}}1
           {ś}{{\'s}}1 {ć}{{\'c}}1 {ż}{{\.z}}1 {ź}{{\'z}}1 {ó}{{\'o}}1 {ń}{{\'n}}1
           {Ą}{{\k A}}1 {Ę}{{\k E}}1 {Ś}{{\'S}}1 {Ć}{{\'C}}1
           {Ż}{{\.Z}}1 {Ź}{{\'Z}}1 {Ó}{{\'O}}1 {Ń}{{\'N}}1
           {—}{{-{}-}}2 {–}{{-}}1 {‘}{{'}}1 {’}{{'}}1
           {“}{{"}}1 {”}{{"}}1 {°}{{\textdegree}}1
}
\lstset{style=bookcode}

\lstnewenvironment{python}[1][]
  {\lstset{language=pythonx,style=bookcode,#1}}
  {}
\lstnewenvironment{shellcmd}[1][]
  {\lstset{language=bash,style=bookcode,numbers=none,keywordstyle=\color{mteal}\bfseries,#1}}
  {}

% The path is printed above every file listing, in the form the reader opens
% in the editor. \nobreak keeps it on the page of the listing it names.
% The path must not be stranded at the foot of a page with its listing
% overleaf: \nobreak alone does not hold, because listings opens with its own
% vertical material. \needspace asks for room for the path and the first few
% lines of code, and turns the page first when there is not.
% \Needspace (capital N) measures \pagegoal-\pagetotal exactly; lower-case
% \needspace works through glue and penalties and was measured, by the
% Chapter 11 pass, turning pages early and leaving gaps. It must be used
% between paragraphs, which \listingpath's own \par guarantees. Chapters use
% \Needspace* too, so the fallback for a machine without the package defines
% both forms (the star swallowed with its argument).
\makeatletter
\IfFileExists{needspace.sty}{\usepackage{needspace}}{%
  \providecommand{\needspace}[1]{}%
  \providecommand{\Needspace}{\@ifstar\@gobble\@gobble}}
\makeatother
\newcommand{\listingpath}[1]{%
  \par\medskip\Needspace{6\baselineskip}\noindent
  {\ttfamily\footnotesize\color{mgrey}\detokenize{#1}}%
  \par\nobreak}

% Usage: \pyfile{code/ch05/cobweb.py}{Caption}{lst:label}
\newcommand{\pyfile}[3]{%
  \listingpath{#1}%
  \lstinputlisting[language=pythonx,style=bookcode,aboveskip=0.3em,
    caption={#2},label={#3}]{#1}%
}

% A named region of a file, delimited by `# --8<-- [start:name]` and
% `# --8<-- [end:name]`. Matched through listings' range prefix and suffix
% keys -- the other obvious form matches nothing and prints the WHOLE file
% without a warning (measured in the sibling book). A region name is a word,
% never a number. The line numbers printed are the file's own.
% Usage: \pyregion{code/ch05/cobweb.py}{step}{Caption}{lst:label}
\lstset{
  rangebeginprefix=\#\ --8<--\ [start:, rangebeginsuffix=],
  rangeendprefix=\#\ --8<--\ [end:, rangeendsuffix=],
  includerangemarker=false}
\newcommand{\pyregion}[4]{%
  \listingpath{#1}%
  \lstinputlisting[language=pythonx,style=bookcode,aboveskip=0.3em,
    linerange={#2-#2},caption={#3},label={#4}]{#1}%
}

% A program's output, written by code/measure/transcripts.py. There is no
% typed form and there must not be: \val{} does not expand in a verbatim
% body, and a transcript nobody ran looks exactly like one that was.
\newcommand{\transcript}[1]{%
  \par\smallskip
  \IfFileExists{figures/transcripts/#1.txt}{%
    \lstinputlisting[style=bookcode,numbers=none,basicstyle=\ttfamily\footnotesize,
      keywordstyle=\color{black},backgroundcolor=\color{white},
      language={}]{figures/transcripts/#1.txt}%
  }{%
    \begin{tcolorbox}[enhanced, sharp corners, width=\linewidth,
      colback=mlight, colframe=mgrey, boxrule=0.6pt,
      left=8pt, right=8pt, top=6pt, bottom=6pt]
      {\bfseries\color{mgrey}\small \lblTranscriptMissing}\\[2pt]
      {\ttfamily\scriptsize\color{mgrey} figures/transcripts/#1.txt}
    \end{tcolorbox}%
  }%
  \par\smallskip
}

% Inline literals. Underscores inside these must be written \_.
\newcommand{\code}[1]{\texttt{\small #1}}
\newcommand{\pkg}[1]{\texttt{\small\color{mblue}#1}}
\newcommand{\api}[1]{\texttt{\small\color{mteal}#1}}

% ============================================================
%  ADMONITIONS
%  A small, fixed vocabulary. `tinted' may be read in passing; `outlined' is
%  a block the reader must stop at. A third treatment would mean the first two
%  are not carrying their meaning.
% ============================================================
\tcbset{
  cfz tinted/.style={
    enhanced, breakable, sharp corners,
    lines before break=4,
    colback=mlight, boxrule=0pt, leftrule=3pt,
    left=8pt, right=8pt, top=6pt, bottom=6pt,
    before skip=13pt, after skip=13pt,
    fonttitle=\bfseries\small,
    coltitle=black, colbacktitle=mlight,
    attach title to upper={\par\smallskip},
  },
  cfz outlined/.style={
    enhanced, breakable, sharp corners,
    % A box broken straight after its title is a heading stranded at the foot
    % of a page with its question overleaf. Four lines is title plus three.
    lines before break=4,
    colback=white, boxrule=0.6pt,
    left=8pt, right=8pt, top=6pt, bottom=6pt,
    before skip=13pt, after skip=13pt,
    fonttitle=\bfseries\small,
    coltitle=black, colbacktitle=white,
    attach title to upper={\par\smallskip},
  },
}

% Read in passing.
\newtcolorbox{note}[1][]{cfz tinted, colframe=mblue, title={\lblNote}, #1}
\newtcolorbox{warning}[1][]{cfz tinted, colframe=mred, title={\lblWarning}, #1}
% Where you meet this in the world. If it cannot name a specific system,
% instrument, paper or event, it is prose and not this box.
\newtcolorbox{worldbox}[1][]{cfz tinted, colframe=mteal, title={\lblWorld}, #1}
% What is stated and not proved here, and where the proof lives.
\newtcolorbox{rigourbox}[1][]{cfz tinted, colframe=mgrey, title={\lblRigour}, #1}
% The companion library the reader builds, under code/src/chaoslab/.
\newtcolorbox{projectbox}[1][]{cfz tinted, colframe=mamber, title={\lblProject}, #1}

% Stop and read.
% The misconception, named AFTER the reader has been asked the question it
% gets wrong. Appendix B is the catalogue; a trap box is one entry in place.
\newtcolorbox{trapbox}[1][]{cfz outlined, colframe=mred, title={\lblTrap}, #1}
% A listing that has not been run on the pinned versions. Counted by
% `make debt'; nothing ships to a reader with one attached.
\newtcolorbox{verifybox}[1][]{cfz outlined, colframe=mgrey, title={\lblVerify}, #1}
\newtcolorbox{labbox}[1][]{cfz outlined, colframe=mgreen, title={\lblLab}, #1}

% ============================================================
%  THE METHOD --- questions before answers
%  ---------------------------------------------------------
%  The unit of teaching is a question the reader answers BEFORE reading the
%  answer. That is retrieval practice and errorful generation, the two
%  findings the programmed-learning format is defended by, and it is the one
%  thing the sibling mathematics book does that this one keeps whole.
%
%  \begin{think} ... \end{think}   the question, numbered per chapter, ending
%                                  with the instruction to answer before
%                                  reading on;
%  \reveal{short answer}           a centred box with the answer, which is the
%                                  first thing after the question -- the edge
%                                  the reader covers with a hand, or scrolls
%                                  to;
%  \begin{revealblock} ... \end{revealblock}  the same, for a worked answer.
%
%  tools/check_structure.py --think fails when a think box is not followed by
%  a reveal before the next think box or section, and when a reveal answers
%  nothing. That is a property of the source, so it is checked mechanically.
%
%  The instruction says "before you read on", never "before you turn over":
%  the two editions paginate differently, so the second is true in one of them.
% ============================================================
\newcounter{think}[chapter]
\renewcommand{\thethink}{\thechapter.\arabic{think}}
\newtcolorbox{thinkbox}[1][]{cfz outlined, colframe=mviolet, #1}
\NewDocumentEnvironment{think}{}{%
  \refstepcounter{think}%
  \begin{thinkbox}[title={\lblThink~\thethink}]%
}{%
  \par\smallskip
  {\raggedleft\itshape\small\color{mviolet}\lblThinkCue\par}%
  \end{thinkbox}%
}
\tcbset{
  cfz answer/.style={
    enhanced, sharp corners,
    colback=mfaint, colframe=mviolet!55, boxrule=0.5pt,
    before skip=4pt, after skip=10pt,
  },
}
\newcommand{\reveal}[1]{%
  \par\nobreak
  \begin{tcolorbox}[cfz answer, unbreakable,
    left=8pt, right=8pt, top=5pt, bottom=5pt,
    width=0.86\linewidth, center, halign=flush center]
    {\small\bfseries\color{mviolet}\lblAnswerMark}\enspace #1
  \end{tcolorbox}%
  \nobreak\noindent\ignorespaces
}
\NewDocumentEnvironment{revealblock}{}{%
  \par\nopagebreak
  \begin{tcolorbox}[cfz answer, breakable,
    left=10pt, right=10pt, top=6pt, bottom=6pt]
  {\small\bfseries\color{mviolet}\lblAnswerMark}\par\smallskip
}{%
  \end{tcolorbox}%
  \nopagebreak\par\noindent\ignorespaces
}

% ---- Outcomes on entry, and "Can you?" generated from them ----------------
% The self-check at the end of a chapter IS the list of outcomes at its
% start, one for one, because it is built from the same tokens: each
% \outcome both prints its item and appends a row to \cfz@canrows. A \loop
% inside a tabularx body does not survive alignment scanning, which is why
% the rows are accumulated at declaration time (the sibling book's finding).
\makeatletter
\newcommand{\cfz@canrows}{}
\newcommand{\cfz@ratebox}{%
  \fbox{\rule{0pt}{0.9ex}\rule{0.9ex}{0pt}}}
\newcommand{\cfz@rates}{%
  {\footnotesize\color{mgrey}%
   1\,\cfz@ratebox\ 2\,\cfz@ratebox\ 3\,\cfz@ratebox\ 4\,\cfz@ratebox\ 5\,\cfz@ratebox}}
\NewDocumentEnvironment{outcomes}{}{%
  \gdef\cfz@canrows{}%
  \begin{tcolorbox}[cfz tinted, colframe=mblue, title={\lblOutcomes}]%
  \lblOutcomesLead
  \begin{itemize}[leftmargin=1.4em, itemsep=2pt, topsep=4pt]%
}{%
  \end{itemize}%
  \end{tcolorbox}%
}
\newcommand{\outcome}[1]{%
  \item #1%
  \gappto\cfz@canrows{#1 & \cfz@rates \\ \addlinespace[3pt]}%
}
\newcommand{\canyou}{%
  \section*{\lblCanYou}%
  \phantomsection\addcontentsline{toc}{section}{\lblCanYou}%
  \noindent\lblCanYouIntro\par\medskip
  \ifdefempty{\cfz@canrows}{%
    \noindent{\color{mred}\lblNoOutcomes}\par}{%
  {\small
  \noindent\begin{tabularx}{\linewidth}{@{}>{\raggedright\arraybackslash}X r@{}}
    \toprule
    \cfz@canrows
    \bottomrule
  \end{tabularx}\par}}%
}
\makeatother

% ---- Summary -------------------------------------------------------------
\newtcolorbox{summarybox}[1][]{cfz tinted, colframe=mblue, title={\lblSummary}, #1}

% ---- Checkpoint questions, answered at the back --------------------------
% Answers are written AT THE QUESTION with \answerto{} and printed in
% Appendix A. They are carried there by a global token list rather than an
% auxiliary file, because an answer contains mathematics and \val{}, and
% material written to an .aux is expanded at shipout, where fragile commands
% break with errors naming neither the answer nor the chapter.
% \include is sequential within a run, so the list is complete by the time
% the appendix is typeset.
\makeatletter
\newcommand{\cfz@answers}{}
\newcommand{\cfz@lastansch}{0}
\NewDocumentEnvironment{checkpoint}{}{%
  \section*{\lblCheckpoint}%
  \phantomsection\addcontentsline{toc}{section}{\lblCheckpoint}%
  \noindent\lblCheckpointIntro\par
  \begin{enumerate}[label=\textbf{\arabic*.}, leftmargin=2em, itemsep=4pt]%
}{%
  \end{enumerate}%
}
\newcommand{\answerto}[1]{%
  \ifnum\value{chapter}=\cfz@lastansch\relax\else
    \xdef\cfz@lastansch{\arabic{chapter}}%
    \xappto\cfz@answers{\noexpand\cfz@anschapter{\thechapter}}%
  \fi
  \xappto\cfz@answers{\noexpand\cfz@ansitem{\arabic{enumi}}}%
  \gappto\cfz@answers{{#1}}%
}
\newcommand{\cfz@anschapter}[1]{%
  \par\medskip\noindent{\bfseries\color{mblue}\chaptername~#1}\par\nobreak\smallskip}
\newcommand{\cfz@ansitem}[2]{%
  \par\noindent\hangindent=2em\hangafter=1
  \makebox[2em][l]{\textbf{#1.}}#2\par}
\newcommand{\answersbody}{%
  \ifdefempty{\cfz@answers}{\noindent{\color{mred}\lblNoAnswers}\par}{%
    {\raggedright\cfz@answers}}%
}
\makeatother

% ============================================================
%  LABS
%  \begin{lab}{key}{title} ... \end{lab}
%
%  A lab is not a paragraph but a FILE the reader opens: a starter under
%  code/labs/chNN/<key>.py that must FAIL its test, a reference solution under
%  code/labs/chNN/solutions/, and test_<key>.py beside it. The box prints the
%  paths and the one command that runs the test. An environment rather than a
%  macro, because a macro argument cannot hold a verbatim fragment.
%  tools/check_structure.py --labs fails when a file is missing or the key's
%  chapter number disagrees with its chapter.
% ============================================================
\newcounter{lab}[chapter]
\renewcommand{\thelab}{\thechapter.\arabic{lab}}
\makeatletter
\newcommand{\labdir}{ch\two@digits{\value{chapter}}}
\newcommand{\listoflabs}{%
  \section*{\lblLabManifest}%
  \phantomsection\addcontentsline{toc}{section}{\lblLabManifest}%
  \begingroup\c@tocdepth=2\relax\@starttoc{lab}\endgroup%
}
\makeatother
\NewDocumentEnvironment{lab}{mm}{%
  \refstepcounter{lab}%
  \addcontentsline{lab}{subsection}{\protect\numberline{\thelab}\texttt{\protect\detokenize{#1}} --- #2}%
  \begin{labbox}[title={\lblLab~\thelab\ \dash{} #2}]%
}{%
  \par\medskip
  {\small\setlength{\parindent}{0pt}%
  \lblLabStarter: \texttt{code/labs/\labdir/\detokenize{#1}.py}\\
  \lblLabTest: \texttt{code/labs/\labdir/test\_\detokenize{#1}.py}\\
  \lblLabRun: \texttt{cd code \&\& uv run pytest -k \detokenize{#1}}\par}%
  \end{labbox}%
}

% ============================================================
%  FIGURES --- plots and diagrams
%  ---------------------------------------------------------
%  \plotfig{key}{caption}{one-line manifest description}
%    A PLOT, generated by code/plots/<chapter>.py from the same library the
%    listings use, into figures/plots/<lang>/<key>.pdf. Per language, because
%    an axis label in the wrong language is the half-translation that makes a
%    bilingual book look machine-made, and because the Polish edition owes a
%    decimal comma on its tick labels too. Build output, gitignored.
%
%  \mermaidfig{key}{caption}{one-line manifest description}
%    A DIAGRAM of an idea rather than of data, from
%    figures/mermaid/<lang>/<key>.mmd, rendered by `make diagrams'.
%
%  If the rendered PDF is absent either macro prints a visible marker in the
%  flow rather than a float, so a missing figure cannot be mistaken for a
%  layout choice. \phantomsection before \addcontentsline is load-bearing:
%  without it the manifest entry records whatever anchor preceded the figure,
%  which after a listing is a line that was never printed.
% ============================================================
\makeatletter
\newcommand{\listofplots}{%
  \section*{\lblPlotManifest}%
  \phantomsection\addcontentsline{toc}{section}{\lblPlotManifest}%
  \begingroup\c@tocdepth=2\relax\@starttoc{plt}\endgroup%
}
\newcommand{\listofdiagrams}{%
  \section*{\lblDiagramManifest}%
  \phantomsection\addcontentsline{toc}{section}{\lblDiagramManifest}%
  \begingroup\c@tocdepth=2\relax\@starttoc{dgm}\endgroup%
}
\makeatother

\newcommand{\figmissing}[3]{% #1 message, #2 path, #3 caption
  \par\medskip
  \begin{tcolorbox}[enhanced, breakable, width=\linewidth, sharp corners,
    colback=white, colframe=mgrey, boxrule=0.8pt,
    left=8pt, right=8pt, top=8pt, bottom=8pt]
    {\bfseries\color{mgrey}\small #1}\\[4pt]
    {\ttfamily\scriptsize\color{mgrey} #2}
  \end{tcolorbox}%
  \captionof{figure}{#3}%
  \par\medskip
}

\newcommand{\plotfig}[3]{%
  \phantomsection%
  \addcontentsline{plt}{subsection}{\protect\numberline{}\texttt{#1} --- #3}%
  \IfFileExists{figures/plots/\booklang/#1.pdf}{%
    \begin{figure}[htbp]
    \centering
    \includegraphics[width=\linewidth,height=0.45\textheight,keepaspectratio]{figures/plots/\booklang/#1.pdf}%
    \caption{#2}\label{fig:#1}
    \end{figure}%
  }{%
    \figmissing{\lblPlotNotRendered}{figures/plots/\booklang/#1.pdf}{#2}\label{fig:#1}%
  }%
}

\newcommand{\mermaidfig}[3]{%
  \phantomsection%
  \addcontentsline{dgm}{subsection}{\protect\numberline{}\texttt{#1.mmd} --- #3}%
  \IfFileExists{figures/diagrams/\booklang/#1.pdf}{%
    \begin{figure}[htbp]
    \centering
    \includegraphics[width=0.95\linewidth,height=0.42\textheight,keepaspectratio]{figures/diagrams/\booklang/#1.pdf}%
    \caption{#2}\label{fig:#1}
    \end{figure}%
  }{%
    \figmissing{\lblNotRendered}{figures/mermaid/\booklang/#1.mmd}{#2}\label{fig:#1}%
  }%
}

% ============================================================
%  APPENDIX B: the misreadings, generated by tools/gen_traps.py
%
%  \trapchapter{Chapter~\ref{ch:x}}{title}    a heading per chapter
%  \trapentry{n}{habit}{truth}{\ref{sec:x}}  one misreading
%
%  Set ragged right: an entry is mostly short sentences with \code{} spans in
%  them, and a justified line full of unbreakable runs is an overfull box
%  waiting for the edition with the longer words.
% ============================================================
\newcommand{\trapchapter}[2]{%
  \par\medskip
  {\large\bfseries\color{mblue}#1\quad\normalfont\itshape\color{mgrey}#2\par}%
  \nopagebreak\smallskip}
\newcommand{\trapentry}[4]{%
  \par\noindent
  \begingroup\raggedright\hangindent=2.4em\hangafter=1
  \makebox[2.4em][l]{\textbf{\color{mred}#1}}%
  \textbf{\enquote{#2}}\enspace #3\enspace{\color{mgrey}\S\,#4}\par
  \endgroup\smallskip}

% ============================================================
%  CHAPTER STUBS
% ============================================================
\newcommand{\chapterstub}[1]{%
  \begin{tcolorbox}[
    enhanced, breakable, width=\linewidth, sharp corners,
    colback=mlight, colframe=mred, boxrule=1pt,
    borderline={0.6pt}{3pt}{mred, dashed},
    left=12pt, right=12pt, top=12pt, bottom=12pt]
    {\bfseries\color{mred}\large \lblNotWritten}\\[8pt]
    \small\raggedright #1
  \end{tcolorbox}%
}

% ============================================================
%  COMPUTED VALUES
%  A number the reader cannot do in their head is computed by a script under
%  code/measure/, written to figures/values/<chapter>.tex as
%      \pyval{ch06.steps.apart}{34}
%  and reaches the page as \val{ch06.steps.apart}. The book contains
%  references, not digits; `make verify' fails when a committed value no
%  longer matches its script. \val applies the edition's decimal marker, which
%  is how the Polish edition gets its comma without anybody remembering.
% ============================================================
% The comma is BRACED: siunitx sets numbers in maths mode, where a bare comma
% is punctuation and gets a thin space after it ("9, 81"). {,} is an ordinary
% symbol, which is what a decimal marker is.
\newcommand{\bookdecimalmarker}{\ifpl{{,}}{.}}
\IfFileExists{siunitx.sty}{%
  \usepackage{siunitx}%
  \sisetup{
    group-digits = integer,
    group-minimum-digits = 5,
    group-separator = {\,},
    output-decimal-marker = {\bookdecimalmarker},
    exponent-product = \times,
    print-unity-mantissa = false
  }%
  \newcommand{\booknum}[1]{\num{##1}}%
}{%
  \newcommand{\booknum}[1]{##1}%
  \providecommand{\num}[1]{##1}%
}

\makeatletter
\newcommand{\pyvalues}{}
\newcommand{\pyval}[2]{\csgdef{pyval@#1}{#2}\listxadd{\pyvalues}{#1}}
\newcommand{\pyvaltext}[2]{%
  \csgdef{pyval@#1}{#2}\csgdef{pyvaltext@#1}{}\listxadd{\pyvalues}{#1}}
% \num on a non-numeric body is a FATAL siunitx error, so \val refuses a text
% value by name rather than letting siunitx fail on it.
\newcommand{\val}[1]{%
  \ifcsdef{pyval@#1}{%
    \ifcsdef{pyvaltext@#1}{%
      \PackageError{chaosbook}{Value `#1' is not a number}%
        {Use \string\valtext{#1} instead.}%
    }{}%
    \booknum{\csuse{pyval@#1}}%
  }{\textcolor{mred}{\textbf{[\lblValueMissing: #1]}}}%
}
\newcommand{\valtext}[1]{%
  \ifcsdef{pyval@#1}{\csuse{pyval@#1}}%
    {\textcolor{mred}{\textbf{[\lblValueMissing: #1]}}}%
}
\makeatother

% figures/values/all.tex is generated by `make numbers' and \inputs each
% value file. A build with no values still compiles; every \val{} then prints
% a visible marker. \valuesindex lets the one-chapter build point elsewhere.
\providecommand{\valuesindex}{figures/values/all}
\IfFileExists{\valuesindex.tex}{\input{\valuesindex}}{%
  \AtBeginDocument{\PackageWarning{chaosbook}{No computed values found. Run `make numbers'.}}}

% ============================================================
%  PINNED VERSIONS --- single source of truth
%  Change these here AND in code/pyproject.toml; tools/check_structure.py
%  --pins fails when the two disagree. A chapter never types a version.
% ============================================================
\newcommand{\bookedition}{0.1}
\newcommand{\bookyear}{2026}
\newcommand{\pyver}{3.14}              % the minor the book targets
\newcommand{\pypatch}{3.14.7}          % the exact interpreter the code ran on
\newcommand{\uvver}{0.12.13}           % uv
\newcommand{\numpyver}{2.5.3}          % numpy
\newcommand{\matplotlibver}{3.11.2}    % matplotlib
\newcommand{\pytestver}{9.1.1}         % pytest
\newcommand{\ruffver}{0.16.9}          % ruff
% The prose licence is a name with a version in it, not a quantity: it is a
% macro so that neither edition typesets it as a decimal to be localised.
\newcommand{\proselicence}{CC~BY-NC-SA~4.0}

% ============================================================
%  INDEX
% ============================================================
\apptocmd{\theindex}{\raggedcolumns\raggedright}{}{%
  \PackageWarning{chaosbook}{Could not patch theindex for ragged setting}}
\providecommand*\see[2]{}
\renewcommand*\see[2]{\emph{\lblIndexSee} #1}
\providecommand*\seealso[2]{}
\renewcommand*\seealso[2]{\emph{\lblIndexSeeAlso} #1}

\makeindex
; nothing else differs between them at this level.
%  Every user-visible string lives in lang/en.tex or lang/pl.tex, so a label
%  that appears in one edition and not the other is a missing macro rather
%  than a silent divergence.
%
%  Lineage. This book is the third in a family and takes one thing from each
%  of the other two:
%
%    * from Mathematics from Zero for the AI Engineer: the bilingual single
%      source (\booklang, lang/, one body.tex), the rule that every number is
%      computed and reaches the page as \val{}, and the programmed-learning
%      method -- outcomes on entry, a question before every answer, a
%      misconception elicited before it is named, a self-check generated from
%      the outcomes, and checkpoint answers collected at the back;
%    * from Python for .NET Engineers: every listing is a FILE that CI runs,
%      every lab is a starter that must fail and a solution that must pass,
%      and a console transcript is a computed number that happens to be
%      verbatim.
%
%  What is new here is \plotfig. A book about chaos cannot be taught in boxes
%  and arrows: the bifurcation diagram and the Lorenz attractor ARE the
%  argument, so the plots are generated by code/plots/ from the same library
%  the listings use, per language, and never drawn by hand.
% ============================================================

% \booklang must be set by the main file. Default to English rather than
% failing, so a stray % ============================================================
%  preamble.tex --- styling, environments, macros
%
%  Shared by BOTH editions. main-en.tex and main-pl.tex each define \booklang
%  before \input{preamble}; nothing else differs between them at this level.
%  Every user-visible string lives in lang/en.tex or lang/pl.tex, so a label
%  that appears in one edition and not the other is a missing macro rather
%  than a silent divergence.
%
%  Lineage. This book is the third in a family and takes one thing from each
%  of the other two:
%
%    * from Mathematics from Zero for the AI Engineer: the bilingual single
%      source (\booklang, lang/, one body.tex), the rule that every number is
%      computed and reaches the page as \val{}, and the programmed-learning
%      method -- outcomes on entry, a question before every answer, a
%      misconception elicited before it is named, a self-check generated from
%      the outcomes, and checkpoint answers collected at the back;
%    * from Python for .NET Engineers: every listing is a FILE that CI runs,
%      every lab is a starter that must fail and a solution that must pass,
%      and a console transcript is a computed number that happens to be
%      verbatim.
%
%  What is new here is \plotfig. A book about chaos cannot be taught in boxes
%  and arrows: the bifurcation diagram and the Lorenz attractor ARE the
%  argument, so the plots are generated by code/plots/ from the same library
%  the listings use, per language, and never drawn by hand.
% ============================================================

% \booklang must be set by the main file. Default to English rather than
% failing, so a stray \input{preamble} still compiles and says what it did.
\providecommand{\booklang}{en}

\usepackage{etoolbox}   % loaded first: the language switch below needs it

% ---------- Page geometry ----------
% ONE format: A4, 12pt, ONE-SIDED. The book is read on a screen, usually with
% an editor open beside it, so there is no recto, no verso, no gutter and no
% blank leaf before a chapter. The measure is about seventy-five characters
% of prose. At \footnotesize a listing line of 79 characters fits without
% wrapping, and 79 is the width every file under code/ is held to by
% tools/check_structure.py --lines.
\usepackage[
  a4paper,
  left=3.1cm, right=3.1cm, top=2.6cm, bottom=3.0cm,
  head=26pt, headsep=0.8cm, footskip=1.4cm
]{geometry}

% ---------- Encoding & fonts ----------
\usepackage[T1]{fontenc}
\usepackage[utf8]{inputenc}
\IfFileExists{lmodern.sty}{\usepackage{lmodern}}{}
% newtxtext takes its TS1 symbols from TeX Gyre Termes (Debian: tex-gyre).
% A machine with newtx and without tex-gyre dies on the copyright page. CI's
% full TeX Live is the reference; locally install tex-gyre as well.
\IfFileExists{newtxtext.sty}{\usepackage{newtxtext}}{}
% amsmath before newtxmath; amssymb only when newtxmath is absent, because
% loading both is a fatal clash (\Bbbk already defined) that is invisible on
% a machine without newtx.
\usepackage{amsmath}
\IfFileExists{newtxmath.sty}{\usepackage{newtxmath}}{\usepackage{amssymb}}
\IfFileExists{inconsolata.sty}{\usepackage[varqu,scaled=0.95]{inconsolata}}{}
% Without lmodern, microtype runs with font expansion off, which is what
% resolves a marginal overfull line -- so a container short of lmodern
% disagrees with CI about line breaks. Install lmodern locally too.
\IfFileExists{lmodern.sty}{\usepackage{microtype}}{\usepackage[expansion=false]{microtype}}

% upquote: in T1 the input ' is a right single quote, and a listing that sets
% f('x') with curly quotes is a SyntaxError when typed back in.
\IfFileExists{upquote.sty}{\usepackage{upquote}}{}
% And its twin: in T1 the pair -- is an en-dash ligature in the typewriter
% family too, so \code{--flag} would print a flag nobody can type. Scoped to
% tt*, so prose keeps its own dashes.
\DisableLigatures[-]{encoding = T1, family = tt*}

% ---------- Language ----------
% babel with a missing language file is FATAL, so each .ldf is probed first.
\IfFileExists{babel.sty}{%
  \IfFileExists{polish.ldf}{%
    \IfFileExists{british.ldf}{%
      \ifdefstring{\booklang}{pl}%
        {\usepackage[british,main=polish]{babel}}%
        {\usepackage[polish,main=british]{babel}}%
    }{\usepackage[polish]{babel}}%
  }{%
    \IfFileExists{british.ldf}{\usepackage[british]{babel}}{}%
  }%
}{}
% babel-polish makes " an active shorthand, which is a category-code accident
% waiting to happen inside listings. Polish quotation marks come from
% \enquote, which csquotes sets per language.
\ifdefstring{\booklang}{pl}{\AtBeginDocument{\shorthandoff{"}}}{}

% ---------- Core packages ----------
\usepackage{graphicx}
\usepackage{xcolor}
\usepackage{booktabs}
\usepackage{tabularx}
\usepackage{longtable}
\usepackage{array}
\usepackage{enumitem}
\usepackage{listings}
\usepackage[most]{tcolorbox}
\usepackage{fancyhdr}
\usepackage{titlesec}
\usepackage{caption}
\usepackage{imakeidx}
\IfFileExists{csquotes.sty}{\usepackage[autostyle=true]{csquotes}}{%
  \providecommand{\enquote}[1]{``##1''}}
\usepackage[hidelinks]{hyperref}

% ---------- Palette ----------
% The family's palette, so the three books read as a set. Two colours are
% this book's own: the lab (green) and the stop-and-think box (violet).
\definecolor{mblue}{HTML}{1F4E79}
\definecolor{mteal}{HTML}{0E7C7B}
\definecolor{mamber}{HTML}{B26A00}
\definecolor{mred}{HTML}{9B2C2C}
\definecolor{mviolet}{HTML}{7B4397}
\definecolor{mgreen}{HTML}{3B7A3B}
\definecolor{mgrey}{HTML}{5A5A5A}
\definecolor{mlight}{HTML}{F4F6F8}
\definecolor{mfaint}{HTML}{FBFCFD}
\definecolor{codebg}{HTML}{FAFAFA}
\definecolor{codeframe}{HTML}{D8DEE4}
\definecolor{kw}{HTML}{0B5394}
\definecolor{str}{HTML}{9B2C2C}
\definecolor{cmt}{HTML}{4F7A4F}
\definecolor{typ}{HTML}{0E7C7B}
\definecolor{dec}{HTML}{7B4397}

\hypersetup{
  colorlinks=true,
  linkcolor=mblue,
  urlcolor=mteal,
  citecolor=mblue,
  pdfauthor={Konrad Cinkusz}
}

% \ifpl{Polish}{English} -- for the handful of places where a string is
% structural rather than content. Robust, because a fragile \ifpl inside a
% moving argument fails on the next run in exactly one of the two editions.
\DeclareRobustCommand{\ifpl}[2]{\ifdefstring{\booklang}{pl}{#1}{#2}}
\makeatletter
\pdfstringdefDisableCommands{%
  \ifdefstring{\booklang}{pl}{\let\ifpl\@firstoftwo}{\let\ifpl\@secondoftwo}%
  % hyperref drops \color from a bookmark and keeps its ARGUMENT, so a title
  % carrying \code{} would bookmark as "mtealx". Gobble the colour name.
  \def\color#1{}%
}
\makeatother

% ---------- Language strings ----------
\input{lang/\booklang}

% The PDF's subject, keywords and language read \lblPdf... from the language
% file, so they cannot sit in the colour block above, which runs first. Two
% blocks is the ordering, not an oversight.
\hypersetup{
  pdfsubject={\lblPdfSubject},
  pdfkeywords={\lblPdfKeywords},
  pdflang={\lblPdfLang}
}

% ---------- Headings ----------
\titleformat{\chapter}[display]
  {\normalfont\huge\bfseries\color{mblue}\raggedright}
  {\normalfont\normalsize\color{mgrey}\MakeUppercase{\chaptertitlename}\ \thechapter}
  {8pt}{\Huge\raggedright}
\titlespacing*{\chapter}{0pt}{10pt}{28pt}
\titleformat{\section}{\normalfont\Large\bfseries\color{mblue}\raggedright}{\thesection}{0.8em}{}
\titleformat{\subsection}{\normalfont\large\bfseries\color{mgrey}\raggedright}{\thesubsection}{0.7em}{}

\setcounter{tocdepth}{1}
\setcounter{secnumdepth}{2}

\pagestyle{fancy}
\fancyhf{}
\fancyhead[L]{\small\nouppercase{\leftmark}}
\fancyhead[R]{\small\thepage}
\renewcommand{\headrulewidth}{0.4pt}
\fancypagestyle{plain}{\fancyhf{}\fancyfoot[R]{\small\thepage}%
  \renewcommand{\headrulewidth}{0pt}}
% After \pagestyle{fancy}, which installs its own \chaptermark.
\renewcommand{\chaptermark}[1]{\markboth{\thechapter.\ #1}{}}

% ---------- Mathematics ----------
% The handful of notations the book uses, as macros so both editions set the
% same symbol. The book writes \ln for the natural logarithm and \log_{10}
% where the base is ten; a bare \log is ambiguous between the two readerships
% and is not used.
\newcommand{\R}{\mathbb{R}}
\newcommand{\dd}{\,\mathrm{d}}
\newcommand{\abs}[1]{\left\lvert #1\right\rvert}
\newcommand{\iu}{\mathrm{i}}

% ============================================================
%  CODE LISTINGS
%  Every listing in a chapter is a FILE under code/, pulled in with \pyfile or
%  \pyregion, and CI runs it on the pinned interpreter. An in-source
%  \begin{python} block is allowed for a fragment of three or four lines that
%  shows a shape rather than runs.
% ============================================================
\lstdefinelanguage{pythonx}{
  language=Python,
  morekeywords={async,await,match,case,nonlocal,yield,type},
  morekeywords=[2]{
    orbit,iterate,logistic,rk4_step,euler_step,integrate,lorenz,henon,
    lyapunov,separation,escape_time,box_count
  },
  sensitive=true
}

\lstdefinestyle{bookcode}{
  basicstyle=\ttfamily\footnotesize,
  keywordstyle=\color{kw}\bfseries,
  keywordstyle=[2]\color{typ},
  stringstyle=\color{str},
  % Colour only: inconsolata has no italic, and an italic comment style is
  % silently substituted and logged on every build.
  commentstyle=\color{cmt},
  numbers=left,
  numberstyle=\ttfamily\tiny\color{mgrey},
  numbersep=8pt,
  backgroundcolor=\color{codebg},
  frame=single,
  rulecolor=\color{codeframe},
  framesep=6pt,
  breaklines=true,
  breakatwhitespace=false,
  postbreak=\mbox{\textcolor{mgrey}{$\hookrightarrow$}\space},
  showstringspaces=false,
  tabsize=4,
  columns=fullflexible,
  keepspaces=true,
  captionpos=b,
  aboveskip=1.2em,
  belowskip=1.2em,
  % Do NOT add a mapping for U+00A0: it aborts the run with a message that
  % names neither the character nor this line. Listings are ASCII, and
  % parity's C13 checks it.
  literate={ł}{{\l}}1 {Ł}{{\L}}1 {ą}{{\k a}}1 {ę}{{\k e}}1
           {ś}{{\'s}}1 {ć}{{\'c}}1 {ż}{{\.z}}1 {ź}{{\'z}}1 {ó}{{\'o}}1 {ń}{{\'n}}1
           {Ą}{{\k A}}1 {Ę}{{\k E}}1 {Ś}{{\'S}}1 {Ć}{{\'C}}1
           {Ż}{{\.Z}}1 {Ź}{{\'Z}}1 {Ó}{{\'O}}1 {Ń}{{\'N}}1
           {—}{{-{}-}}2 {–}{{-}}1 {‘}{{'}}1 {’}{{'}}1
           {“}{{"}}1 {”}{{"}}1 {°}{{\textdegree}}1
}
\lstset{style=bookcode}

\lstnewenvironment{python}[1][]
  {\lstset{language=pythonx,style=bookcode,#1}}
  {}
\lstnewenvironment{shellcmd}[1][]
  {\lstset{language=bash,style=bookcode,numbers=none,keywordstyle=\color{mteal}\bfseries,#1}}
  {}

% The path is printed above every file listing, in the form the reader opens
% in the editor. \nobreak keeps it on the page of the listing it names.
% The path must not be stranded at the foot of a page with its listing
% overleaf: \nobreak alone does not hold, because listings opens with its own
% vertical material. \needspace asks for room for the path and the first few
% lines of code, and turns the page first when there is not.
% \Needspace (capital N) measures \pagegoal-\pagetotal exactly; lower-case
% \needspace works through glue and penalties and was measured, by the
% Chapter 11 pass, turning pages early and leaving gaps. It must be used
% between paragraphs, which \listingpath's own \par guarantees. Chapters use
% \Needspace* too, so the fallback for a machine without the package defines
% both forms (the star swallowed with its argument).
\makeatletter
\IfFileExists{needspace.sty}{\usepackage{needspace}}{%
  \providecommand{\needspace}[1]{}%
  \providecommand{\Needspace}{\@ifstar\@gobble\@gobble}}
\makeatother
\newcommand{\listingpath}[1]{%
  \par\medskip\Needspace{6\baselineskip}\noindent
  {\ttfamily\footnotesize\color{mgrey}\detokenize{#1}}%
  \par\nobreak}

% Usage: \pyfile{code/ch05/cobweb.py}{Caption}{lst:label}
\newcommand{\pyfile}[3]{%
  \listingpath{#1}%
  \lstinputlisting[language=pythonx,style=bookcode,aboveskip=0.3em,
    caption={#2},label={#3}]{#1}%
}

% A named region of a file, delimited by `# --8<-- [start:name]` and
% `# --8<-- [end:name]`. Matched through listings' range prefix and suffix
% keys -- the other obvious form matches nothing and prints the WHOLE file
% without a warning (measured in the sibling book). A region name is a word,
% never a number. The line numbers printed are the file's own.
% Usage: \pyregion{code/ch05/cobweb.py}{step}{Caption}{lst:label}
\lstset{
  rangebeginprefix=\#\ --8<--\ [start:, rangebeginsuffix=],
  rangeendprefix=\#\ --8<--\ [end:, rangeendsuffix=],
  includerangemarker=false}
\newcommand{\pyregion}[4]{%
  \listingpath{#1}%
  \lstinputlisting[language=pythonx,style=bookcode,aboveskip=0.3em,
    linerange={#2-#2},caption={#3},label={#4}]{#1}%
}

% A program's output, written by code/measure/transcripts.py. There is no
% typed form and there must not be: \val{} does not expand in a verbatim
% body, and a transcript nobody ran looks exactly like one that was.
\newcommand{\transcript}[1]{%
  \par\smallskip
  \IfFileExists{figures/transcripts/#1.txt}{%
    \lstinputlisting[style=bookcode,numbers=none,basicstyle=\ttfamily\footnotesize,
      keywordstyle=\color{black},backgroundcolor=\color{white},
      language={}]{figures/transcripts/#1.txt}%
  }{%
    \begin{tcolorbox}[enhanced, sharp corners, width=\linewidth,
      colback=mlight, colframe=mgrey, boxrule=0.6pt,
      left=8pt, right=8pt, top=6pt, bottom=6pt]
      {\bfseries\color{mgrey}\small \lblTranscriptMissing}\\[2pt]
      {\ttfamily\scriptsize\color{mgrey} figures/transcripts/#1.txt}
    \end{tcolorbox}%
  }%
  \par\smallskip
}

% Inline literals. Underscores inside these must be written \_.
\newcommand{\code}[1]{\texttt{\small #1}}
\newcommand{\pkg}[1]{\texttt{\small\color{mblue}#1}}
\newcommand{\api}[1]{\texttt{\small\color{mteal}#1}}

% ============================================================
%  ADMONITIONS
%  A small, fixed vocabulary. `tinted' may be read in passing; `outlined' is
%  a block the reader must stop at. A third treatment would mean the first two
%  are not carrying their meaning.
% ============================================================
\tcbset{
  cfz tinted/.style={
    enhanced, breakable, sharp corners,
    lines before break=4,
    colback=mlight, boxrule=0pt, leftrule=3pt,
    left=8pt, right=8pt, top=6pt, bottom=6pt,
    before skip=13pt, after skip=13pt,
    fonttitle=\bfseries\small,
    coltitle=black, colbacktitle=mlight,
    attach title to upper={\par\smallskip},
  },
  cfz outlined/.style={
    enhanced, breakable, sharp corners,
    % A box broken straight after its title is a heading stranded at the foot
    % of a page with its question overleaf. Four lines is title plus three.
    lines before break=4,
    colback=white, boxrule=0.6pt,
    left=8pt, right=8pt, top=6pt, bottom=6pt,
    before skip=13pt, after skip=13pt,
    fonttitle=\bfseries\small,
    coltitle=black, colbacktitle=white,
    attach title to upper={\par\smallskip},
  },
}

% Read in passing.
\newtcolorbox{note}[1][]{cfz tinted, colframe=mblue, title={\lblNote}, #1}
\newtcolorbox{warning}[1][]{cfz tinted, colframe=mred, title={\lblWarning}, #1}
% Where you meet this in the world. If it cannot name a specific system,
% instrument, paper or event, it is prose and not this box.
\newtcolorbox{worldbox}[1][]{cfz tinted, colframe=mteal, title={\lblWorld}, #1}
% What is stated and not proved here, and where the proof lives.
\newtcolorbox{rigourbox}[1][]{cfz tinted, colframe=mgrey, title={\lblRigour}, #1}
% The companion library the reader builds, under code/src/chaoslab/.
\newtcolorbox{projectbox}[1][]{cfz tinted, colframe=mamber, title={\lblProject}, #1}

% Stop and read.
% The misconception, named AFTER the reader has been asked the question it
% gets wrong. Appendix B is the catalogue; a trap box is one entry in place.
\newtcolorbox{trapbox}[1][]{cfz outlined, colframe=mred, title={\lblTrap}, #1}
% A listing that has not been run on the pinned versions. Counted by
% `make debt'; nothing ships to a reader with one attached.
\newtcolorbox{verifybox}[1][]{cfz outlined, colframe=mgrey, title={\lblVerify}, #1}
\newtcolorbox{labbox}[1][]{cfz outlined, colframe=mgreen, title={\lblLab}, #1}

% ============================================================
%  THE METHOD --- questions before answers
%  ---------------------------------------------------------
%  The unit of teaching is a question the reader answers BEFORE reading the
%  answer. That is retrieval practice and errorful generation, the two
%  findings the programmed-learning format is defended by, and it is the one
%  thing the sibling mathematics book does that this one keeps whole.
%
%  \begin{think} ... \end{think}   the question, numbered per chapter, ending
%                                  with the instruction to answer before
%                                  reading on;
%  \reveal{short answer}           a centred box with the answer, which is the
%                                  first thing after the question -- the edge
%                                  the reader covers with a hand, or scrolls
%                                  to;
%  \begin{revealblock} ... \end{revealblock}  the same, for a worked answer.
%
%  tools/check_structure.py --think fails when a think box is not followed by
%  a reveal before the next think box or section, and when a reveal answers
%  nothing. That is a property of the source, so it is checked mechanically.
%
%  The instruction says "before you read on", never "before you turn over":
%  the two editions paginate differently, so the second is true in one of them.
% ============================================================
\newcounter{think}[chapter]
\renewcommand{\thethink}{\thechapter.\arabic{think}}
\newtcolorbox{thinkbox}[1][]{cfz outlined, colframe=mviolet, #1}
\NewDocumentEnvironment{think}{}{%
  \refstepcounter{think}%
  \begin{thinkbox}[title={\lblThink~\thethink}]%
}{%
  \par\smallskip
  {\raggedleft\itshape\small\color{mviolet}\lblThinkCue\par}%
  \end{thinkbox}%
}
\tcbset{
  cfz answer/.style={
    enhanced, sharp corners,
    colback=mfaint, colframe=mviolet!55, boxrule=0.5pt,
    before skip=4pt, after skip=10pt,
  },
}
\newcommand{\reveal}[1]{%
  \par\nobreak
  \begin{tcolorbox}[cfz answer, unbreakable,
    left=8pt, right=8pt, top=5pt, bottom=5pt,
    width=0.86\linewidth, center, halign=flush center]
    {\small\bfseries\color{mviolet}\lblAnswerMark}\enspace #1
  \end{tcolorbox}%
  \nobreak\noindent\ignorespaces
}
\NewDocumentEnvironment{revealblock}{}{%
  \par\nopagebreak
  \begin{tcolorbox}[cfz answer, breakable,
    left=10pt, right=10pt, top=6pt, bottom=6pt]
  {\small\bfseries\color{mviolet}\lblAnswerMark}\par\smallskip
}{%
  \end{tcolorbox}%
  \nopagebreak\par\noindent\ignorespaces
}

% ---- Outcomes on entry, and "Can you?" generated from them ----------------
% The self-check at the end of a chapter IS the list of outcomes at its
% start, one for one, because it is built from the same tokens: each
% \outcome both prints its item and appends a row to \cfz@canrows. A \loop
% inside a tabularx body does not survive alignment scanning, which is why
% the rows are accumulated at declaration time (the sibling book's finding).
\makeatletter
\newcommand{\cfz@canrows}{}
\newcommand{\cfz@ratebox}{%
  \fbox{\rule{0pt}{0.9ex}\rule{0.9ex}{0pt}}}
\newcommand{\cfz@rates}{%
  {\footnotesize\color{mgrey}%
   1\,\cfz@ratebox\ 2\,\cfz@ratebox\ 3\,\cfz@ratebox\ 4\,\cfz@ratebox\ 5\,\cfz@ratebox}}
\NewDocumentEnvironment{outcomes}{}{%
  \gdef\cfz@canrows{}%
  \begin{tcolorbox}[cfz tinted, colframe=mblue, title={\lblOutcomes}]%
  \lblOutcomesLead
  \begin{itemize}[leftmargin=1.4em, itemsep=2pt, topsep=4pt]%
}{%
  \end{itemize}%
  \end{tcolorbox}%
}
\newcommand{\outcome}[1]{%
  \item #1%
  \gappto\cfz@canrows{#1 & \cfz@rates \\ \addlinespace[3pt]}%
}
\newcommand{\canyou}{%
  \section*{\lblCanYou}%
  \phantomsection\addcontentsline{toc}{section}{\lblCanYou}%
  \noindent\lblCanYouIntro\par\medskip
  \ifdefempty{\cfz@canrows}{%
    \noindent{\color{mred}\lblNoOutcomes}\par}{%
  {\small
  \noindent\begin{tabularx}{\linewidth}{@{}>{\raggedright\arraybackslash}X r@{}}
    \toprule
    \cfz@canrows
    \bottomrule
  \end{tabularx}\par}}%
}
\makeatother

% ---- Summary -------------------------------------------------------------
\newtcolorbox{summarybox}[1][]{cfz tinted, colframe=mblue, title={\lblSummary}, #1}

% ---- Checkpoint questions, answered at the back --------------------------
% Answers are written AT THE QUESTION with \answerto{} and printed in
% Appendix A. They are carried there by a global token list rather than an
% auxiliary file, because an answer contains mathematics and \val{}, and
% material written to an .aux is expanded at shipout, where fragile commands
% break with errors naming neither the answer nor the chapter.
% \include is sequential within a run, so the list is complete by the time
% the appendix is typeset.
\makeatletter
\newcommand{\cfz@answers}{}
\newcommand{\cfz@lastansch}{0}
\NewDocumentEnvironment{checkpoint}{}{%
  \section*{\lblCheckpoint}%
  \phantomsection\addcontentsline{toc}{section}{\lblCheckpoint}%
  \noindent\lblCheckpointIntro\par
  \begin{enumerate}[label=\textbf{\arabic*.}, leftmargin=2em, itemsep=4pt]%
}{%
  \end{enumerate}%
}
\newcommand{\answerto}[1]{%
  \ifnum\value{chapter}=\cfz@lastansch\relax\else
    \xdef\cfz@lastansch{\arabic{chapter}}%
    \xappto\cfz@answers{\noexpand\cfz@anschapter{\thechapter}}%
  \fi
  \xappto\cfz@answers{\noexpand\cfz@ansitem{\arabic{enumi}}}%
  \gappto\cfz@answers{{#1}}%
}
\newcommand{\cfz@anschapter}[1]{%
  \par\medskip\noindent{\bfseries\color{mblue}\chaptername~#1}\par\nobreak\smallskip}
\newcommand{\cfz@ansitem}[2]{%
  \par\noindent\hangindent=2em\hangafter=1
  \makebox[2em][l]{\textbf{#1.}}#2\par}
\newcommand{\answersbody}{%
  \ifdefempty{\cfz@answers}{\noindent{\color{mred}\lblNoAnswers}\par}{%
    {\raggedright\cfz@answers}}%
}
\makeatother

% ============================================================
%  LABS
%  \begin{lab}{key}{title} ... \end{lab}
%
%  A lab is not a paragraph but a FILE the reader opens: a starter under
%  code/labs/chNN/<key>.py that must FAIL its test, a reference solution under
%  code/labs/chNN/solutions/, and test_<key>.py beside it. The box prints the
%  paths and the one command that runs the test. An environment rather than a
%  macro, because a macro argument cannot hold a verbatim fragment.
%  tools/check_structure.py --labs fails when a file is missing or the key's
%  chapter number disagrees with its chapter.
% ============================================================
\newcounter{lab}[chapter]
\renewcommand{\thelab}{\thechapter.\arabic{lab}}
\makeatletter
\newcommand{\labdir}{ch\two@digits{\value{chapter}}}
\newcommand{\listoflabs}{%
  \section*{\lblLabManifest}%
  \phantomsection\addcontentsline{toc}{section}{\lblLabManifest}%
  \begingroup\c@tocdepth=2\relax\@starttoc{lab}\endgroup%
}
\makeatother
\NewDocumentEnvironment{lab}{mm}{%
  \refstepcounter{lab}%
  \addcontentsline{lab}{subsection}{\protect\numberline{\thelab}\texttt{\protect\detokenize{#1}} --- #2}%
  \begin{labbox}[title={\lblLab~\thelab\ \dash{} #2}]%
}{%
  \par\medskip
  {\small\setlength{\parindent}{0pt}%
  \lblLabStarter: \texttt{code/labs/\labdir/\detokenize{#1}.py}\\
  \lblLabTest: \texttt{code/labs/\labdir/test\_\detokenize{#1}.py}\\
  \lblLabRun: \texttt{cd code \&\& uv run pytest -k \detokenize{#1}}\par}%
  \end{labbox}%
}

% ============================================================
%  FIGURES --- plots and diagrams
%  ---------------------------------------------------------
%  \plotfig{key}{caption}{one-line manifest description}
%    A PLOT, generated by code/plots/<chapter>.py from the same library the
%    listings use, into figures/plots/<lang>/<key>.pdf. Per language, because
%    an axis label in the wrong language is the half-translation that makes a
%    bilingual book look machine-made, and because the Polish edition owes a
%    decimal comma on its tick labels too. Build output, gitignored.
%
%  \mermaidfig{key}{caption}{one-line manifest description}
%    A DIAGRAM of an idea rather than of data, from
%    figures/mermaid/<lang>/<key>.mmd, rendered by `make diagrams'.
%
%  If the rendered PDF is absent either macro prints a visible marker in the
%  flow rather than a float, so a missing figure cannot be mistaken for a
%  layout choice. \phantomsection before \addcontentsline is load-bearing:
%  without it the manifest entry records whatever anchor preceded the figure,
%  which after a listing is a line that was never printed.
% ============================================================
\makeatletter
\newcommand{\listofplots}{%
  \section*{\lblPlotManifest}%
  \phantomsection\addcontentsline{toc}{section}{\lblPlotManifest}%
  \begingroup\c@tocdepth=2\relax\@starttoc{plt}\endgroup%
}
\newcommand{\listofdiagrams}{%
  \section*{\lblDiagramManifest}%
  \phantomsection\addcontentsline{toc}{section}{\lblDiagramManifest}%
  \begingroup\c@tocdepth=2\relax\@starttoc{dgm}\endgroup%
}
\makeatother

\newcommand{\figmissing}[3]{% #1 message, #2 path, #3 caption
  \par\medskip
  \begin{tcolorbox}[enhanced, breakable, width=\linewidth, sharp corners,
    colback=white, colframe=mgrey, boxrule=0.8pt,
    left=8pt, right=8pt, top=8pt, bottom=8pt]
    {\bfseries\color{mgrey}\small #1}\\[4pt]
    {\ttfamily\scriptsize\color{mgrey} #2}
  \end{tcolorbox}%
  \captionof{figure}{#3}%
  \par\medskip
}

\newcommand{\plotfig}[3]{%
  \phantomsection%
  \addcontentsline{plt}{subsection}{\protect\numberline{}\texttt{#1} --- #3}%
  \IfFileExists{figures/plots/\booklang/#1.pdf}{%
    \begin{figure}[htbp]
    \centering
    \includegraphics[width=\linewidth,height=0.45\textheight,keepaspectratio]{figures/plots/\booklang/#1.pdf}%
    \caption{#2}\label{fig:#1}
    \end{figure}%
  }{%
    \figmissing{\lblPlotNotRendered}{figures/plots/\booklang/#1.pdf}{#2}\label{fig:#1}%
  }%
}

\newcommand{\mermaidfig}[3]{%
  \phantomsection%
  \addcontentsline{dgm}{subsection}{\protect\numberline{}\texttt{#1.mmd} --- #3}%
  \IfFileExists{figures/diagrams/\booklang/#1.pdf}{%
    \begin{figure}[htbp]
    \centering
    \includegraphics[width=0.95\linewidth,height=0.42\textheight,keepaspectratio]{figures/diagrams/\booklang/#1.pdf}%
    \caption{#2}\label{fig:#1}
    \end{figure}%
  }{%
    \figmissing{\lblNotRendered}{figures/mermaid/\booklang/#1.mmd}{#2}\label{fig:#1}%
  }%
}

% ============================================================
%  APPENDIX B: the misreadings, generated by tools/gen_traps.py
%
%  \trapchapter{Chapter~\ref{ch:x}}{title}    a heading per chapter
%  \trapentry{n}{habit}{truth}{\ref{sec:x}}  one misreading
%
%  Set ragged right: an entry is mostly short sentences with \code{} spans in
%  them, and a justified line full of unbreakable runs is an overfull box
%  waiting for the edition with the longer words.
% ============================================================
\newcommand{\trapchapter}[2]{%
  \par\medskip
  {\large\bfseries\color{mblue}#1\quad\normalfont\itshape\color{mgrey}#2\par}%
  \nopagebreak\smallskip}
\newcommand{\trapentry}[4]{%
  \par\noindent
  \begingroup\raggedright\hangindent=2.4em\hangafter=1
  \makebox[2.4em][l]{\textbf{\color{mred}#1}}%
  \textbf{\enquote{#2}}\enspace #3\enspace{\color{mgrey}\S\,#4}\par
  \endgroup\smallskip}

% ============================================================
%  CHAPTER STUBS
% ============================================================
\newcommand{\chapterstub}[1]{%
  \begin{tcolorbox}[
    enhanced, breakable, width=\linewidth, sharp corners,
    colback=mlight, colframe=mred, boxrule=1pt,
    borderline={0.6pt}{3pt}{mred, dashed},
    left=12pt, right=12pt, top=12pt, bottom=12pt]
    {\bfseries\color{mred}\large \lblNotWritten}\\[8pt]
    \small\raggedright #1
  \end{tcolorbox}%
}

% ============================================================
%  COMPUTED VALUES
%  A number the reader cannot do in their head is computed by a script under
%  code/measure/, written to figures/values/<chapter>.tex as
%      \pyval{ch06.steps.apart}{34}
%  and reaches the page as \val{ch06.steps.apart}. The book contains
%  references, not digits; `make verify' fails when a committed value no
%  longer matches its script. \val applies the edition's decimal marker, which
%  is how the Polish edition gets its comma without anybody remembering.
% ============================================================
% The comma is BRACED: siunitx sets numbers in maths mode, where a bare comma
% is punctuation and gets a thin space after it ("9, 81"). {,} is an ordinary
% symbol, which is what a decimal marker is.
\newcommand{\bookdecimalmarker}{\ifpl{{,}}{.}}
\IfFileExists{siunitx.sty}{%
  \usepackage{siunitx}%
  \sisetup{
    group-digits = integer,
    group-minimum-digits = 5,
    group-separator = {\,},
    output-decimal-marker = {\bookdecimalmarker},
    exponent-product = \times,
    print-unity-mantissa = false
  }%
  \newcommand{\booknum}[1]{\num{##1}}%
}{%
  \newcommand{\booknum}[1]{##1}%
  \providecommand{\num}[1]{##1}%
}

\makeatletter
\newcommand{\pyvalues}{}
\newcommand{\pyval}[2]{\csgdef{pyval@#1}{#2}\listxadd{\pyvalues}{#1}}
\newcommand{\pyvaltext}[2]{%
  \csgdef{pyval@#1}{#2}\csgdef{pyvaltext@#1}{}\listxadd{\pyvalues}{#1}}
% \num on a non-numeric body is a FATAL siunitx error, so \val refuses a text
% value by name rather than letting siunitx fail on it.
\newcommand{\val}[1]{%
  \ifcsdef{pyval@#1}{%
    \ifcsdef{pyvaltext@#1}{%
      \PackageError{chaosbook}{Value `#1' is not a number}%
        {Use \string\valtext{#1} instead.}%
    }{}%
    \booknum{\csuse{pyval@#1}}%
  }{\textcolor{mred}{\textbf{[\lblValueMissing: #1]}}}%
}
\newcommand{\valtext}[1]{%
  \ifcsdef{pyval@#1}{\csuse{pyval@#1}}%
    {\textcolor{mred}{\textbf{[\lblValueMissing: #1]}}}%
}
\makeatother

% figures/values/all.tex is generated by `make numbers' and \inputs each
% value file. A build with no values still compiles; every \val{} then prints
% a visible marker. \valuesindex lets the one-chapter build point elsewhere.
\providecommand{\valuesindex}{figures/values/all}
\IfFileExists{\valuesindex.tex}{\input{\valuesindex}}{%
  \AtBeginDocument{\PackageWarning{chaosbook}{No computed values found. Run `make numbers'.}}}

% ============================================================
%  PINNED VERSIONS --- single source of truth
%  Change these here AND in code/pyproject.toml; tools/check_structure.py
%  --pins fails when the two disagree. A chapter never types a version.
% ============================================================
\newcommand{\bookedition}{0.1}
\newcommand{\bookyear}{2026}
\newcommand{\pyver}{3.14}              % the minor the book targets
\newcommand{\pypatch}{3.14.7}          % the exact interpreter the code ran on
\newcommand{\uvver}{0.12.13}           % uv
\newcommand{\numpyver}{2.5.3}          % numpy
\newcommand{\matplotlibver}{3.11.2}    % matplotlib
\newcommand{\pytestver}{9.1.1}         % pytest
\newcommand{\ruffver}{0.16.9}          % ruff
% The prose licence is a name with a version in it, not a quantity: it is a
% macro so that neither edition typesets it as a decimal to be localised.
\newcommand{\proselicence}{CC~BY-NC-SA~4.0}

% ============================================================
%  INDEX
% ============================================================
\apptocmd{\theindex}{\raggedcolumns\raggedright}{}{%
  \PackageWarning{chaosbook}{Could not patch theindex for ragged setting}}
\providecommand*\see[2]{}
\renewcommand*\see[2]{\emph{\lblIndexSee} #1}
\providecommand*\seealso[2]{}
\renewcommand*\seealso[2]{\emph{\lblIndexSeeAlso} #1}

\makeindex
 still compiles and says what it did.
\providecommand{\booklang}{en}

\usepackage{etoolbox}   % loaded first: the language switch below needs it

% ---------- Page geometry ----------
% ONE format: A4, 12pt, ONE-SIDED. The book is read on a screen, usually with
% an editor open beside it, so there is no recto, no verso, no gutter and no
% blank leaf before a chapter. The measure is about seventy-five characters
% of prose. At \footnotesize a listing line of 79 characters fits without
% wrapping, and 79 is the width every file under code/ is held to by
% tools/check_structure.py --lines.
\usepackage[
  a4paper,
  left=3.1cm, right=3.1cm, top=2.6cm, bottom=3.0cm,
  head=26pt, headsep=0.8cm, footskip=1.4cm
]{geometry}

% ---------- Encoding & fonts ----------
\usepackage[T1]{fontenc}
\usepackage[utf8]{inputenc}
\IfFileExists{lmodern.sty}{\usepackage{lmodern}}{}
% newtxtext takes its TS1 symbols from TeX Gyre Termes (Debian: tex-gyre).
% A machine with newtx and without tex-gyre dies on the copyright page. CI's
% full TeX Live is the reference; locally install tex-gyre as well.
\IfFileExists{newtxtext.sty}{\usepackage{newtxtext}}{}
% amsmath before newtxmath; amssymb only when newtxmath is absent, because
% loading both is a fatal clash (\Bbbk already defined) that is invisible on
% a machine without newtx.
\usepackage{amsmath}
\IfFileExists{newtxmath.sty}{\usepackage{newtxmath}}{\usepackage{amssymb}}
\IfFileExists{inconsolata.sty}{\usepackage[varqu,scaled=0.95]{inconsolata}}{}
% Without lmodern, microtype runs with font expansion off, which is what
% resolves a marginal overfull line -- so a container short of lmodern
% disagrees with CI about line breaks. Install lmodern locally too.
\IfFileExists{lmodern.sty}{\usepackage{microtype}}{\usepackage[expansion=false]{microtype}}

% upquote: in T1 the input ' is a right single quote, and a listing that sets
% f('x') with curly quotes is a SyntaxError when typed back in.
\IfFileExists{upquote.sty}{\usepackage{upquote}}{}
% And its twin: in T1 the pair -- is an en-dash ligature in the typewriter
% family too, so \code{--flag} would print a flag nobody can type. Scoped to
% tt*, so prose keeps its own dashes.
\DisableLigatures[-]{encoding = T1, family = tt*}

% ---------- Language ----------
% babel with a missing language file is FATAL, so each .ldf is probed first.
\IfFileExists{babel.sty}{%
  \IfFileExists{polish.ldf}{%
    \IfFileExists{british.ldf}{%
      \ifdefstring{\booklang}{pl}%
        {\usepackage[british,main=polish]{babel}}%
        {\usepackage[polish,main=british]{babel}}%
    }{\usepackage[polish]{babel}}%
  }{%
    \IfFileExists{british.ldf}{\usepackage[british]{babel}}{}%
  }%
}{}
% babel-polish makes " an active shorthand, which is a category-code accident
% waiting to happen inside listings. Polish quotation marks come from
% \enquote, which csquotes sets per language.
\ifdefstring{\booklang}{pl}{\AtBeginDocument{\shorthandoff{"}}}{}

% ---------- Core packages ----------
\usepackage{graphicx}
\usepackage{xcolor}
\usepackage{booktabs}
\usepackage{tabularx}
\usepackage{longtable}
\usepackage{array}
\usepackage{enumitem}
\usepackage{listings}
\usepackage[most]{tcolorbox}
\usepackage{fancyhdr}
\usepackage{titlesec}
\usepackage{caption}
\usepackage{imakeidx}
\IfFileExists{csquotes.sty}{\usepackage[autostyle=true]{csquotes}}{%
  \providecommand{\enquote}[1]{``##1''}}
\usepackage[hidelinks]{hyperref}

% ---------- Palette ----------
% The family's palette, so the three books read as a set. Two colours are
% this book's own: the lab (green) and the stop-and-think box (violet).
\definecolor{mblue}{HTML}{1F4E79}
\definecolor{mteal}{HTML}{0E7C7B}
\definecolor{mamber}{HTML}{B26A00}
\definecolor{mred}{HTML}{9B2C2C}
\definecolor{mviolet}{HTML}{7B4397}
\definecolor{mgreen}{HTML}{3B7A3B}
\definecolor{mgrey}{HTML}{5A5A5A}
\definecolor{mlight}{HTML}{F4F6F8}
\definecolor{mfaint}{HTML}{FBFCFD}
\definecolor{codebg}{HTML}{FAFAFA}
\definecolor{codeframe}{HTML}{D8DEE4}
\definecolor{kw}{HTML}{0B5394}
\definecolor{str}{HTML}{9B2C2C}
\definecolor{cmt}{HTML}{4F7A4F}
\definecolor{typ}{HTML}{0E7C7B}
\definecolor{dec}{HTML}{7B4397}

\hypersetup{
  colorlinks=true,
  linkcolor=mblue,
  urlcolor=mteal,
  citecolor=mblue,
  pdfauthor={Konrad Cinkusz}
}

% \ifpl{Polish}{English} -- for the handful of places where a string is
% structural rather than content. Robust, because a fragile \ifpl inside a
% moving argument fails on the next run in exactly one of the two editions.
\DeclareRobustCommand{\ifpl}[2]{\ifdefstring{\booklang}{pl}{#1}{#2}}
\makeatletter
\pdfstringdefDisableCommands{%
  \ifdefstring{\booklang}{pl}{\let\ifpl\@firstoftwo}{\let\ifpl\@secondoftwo}%
  % hyperref drops \color from a bookmark and keeps its ARGUMENT, so a title
  % carrying \code{} would bookmark as "mtealx". Gobble the colour name.
  \def\color#1{}%
}
\makeatother

% ---------- Language strings ----------
\input{lang/\booklang}

% The PDF's subject, keywords and language read \lblPdf... from the language
% file, so they cannot sit in the colour block above, which runs first. Two
% blocks is the ordering, not an oversight.
\hypersetup{
  pdfsubject={\lblPdfSubject},
  pdfkeywords={\lblPdfKeywords},
  pdflang={\lblPdfLang}
}

% ---------- Headings ----------
\titleformat{\chapter}[display]
  {\normalfont\huge\bfseries\color{mblue}\raggedright}
  {\normalfont\normalsize\color{mgrey}\MakeUppercase{\chaptertitlename}\ \thechapter}
  {8pt}{\Huge\raggedright}
\titlespacing*{\chapter}{0pt}{10pt}{28pt}
\titleformat{\section}{\normalfont\Large\bfseries\color{mblue}\raggedright}{\thesection}{0.8em}{}
\titleformat{\subsection}{\normalfont\large\bfseries\color{mgrey}\raggedright}{\thesubsection}{0.7em}{}

\setcounter{tocdepth}{1}
\setcounter{secnumdepth}{2}

\pagestyle{fancy}
\fancyhf{}
\fancyhead[L]{\small\nouppercase{\leftmark}}
\fancyhead[R]{\small\thepage}
\renewcommand{\headrulewidth}{0.4pt}
\fancypagestyle{plain}{\fancyhf{}\fancyfoot[R]{\small\thepage}%
  \renewcommand{\headrulewidth}{0pt}}
% After \pagestyle{fancy}, which installs its own \chaptermark.
\renewcommand{\chaptermark}[1]{\markboth{\thechapter.\ #1}{}}

% ---------- Mathematics ----------
% The handful of notations the book uses, as macros so both editions set the
% same symbol. The book writes \ln for the natural logarithm and \log_{10}
% where the base is ten; a bare \log is ambiguous between the two readerships
% and is not used.
\newcommand{\R}{\mathbb{R}}
\newcommand{\dd}{\,\mathrm{d}}
\newcommand{\abs}[1]{\left\lvert #1\right\rvert}
\newcommand{\iu}{\mathrm{i}}

% ============================================================
%  CODE LISTINGS
%  Every listing in a chapter is a FILE under code/, pulled in with \pyfile or
%  \pyregion, and CI runs it on the pinned interpreter. An in-source
%  \begin{python} block is allowed for a fragment of three or four lines that
%  shows a shape rather than runs.
% ============================================================
\lstdefinelanguage{pythonx}{
  language=Python,
  morekeywords={async,await,match,case,nonlocal,yield,type},
  morekeywords=[2]{
    orbit,iterate,logistic,rk4_step,euler_step,integrate,lorenz,henon,
    lyapunov,separation,escape_time,box_count
  },
  sensitive=true
}

\lstdefinestyle{bookcode}{
  basicstyle=\ttfamily\footnotesize,
  keywordstyle=\color{kw}\bfseries,
  keywordstyle=[2]\color{typ},
  stringstyle=\color{str},
  % Colour only: inconsolata has no italic, and an italic comment style is
  % silently substituted and logged on every build.
  commentstyle=\color{cmt},
  numbers=left,
  numberstyle=\ttfamily\tiny\color{mgrey},
  numbersep=8pt,
  backgroundcolor=\color{codebg},
  frame=single,
  rulecolor=\color{codeframe},
  framesep=6pt,
  breaklines=true,
  breakatwhitespace=false,
  postbreak=\mbox{\textcolor{mgrey}{$\hookrightarrow$}\space},
  showstringspaces=false,
  tabsize=4,
  columns=fullflexible,
  keepspaces=true,
  captionpos=b,
  aboveskip=1.2em,
  belowskip=1.2em,
  % Do NOT add a mapping for U+00A0: it aborts the run with a message that
  % names neither the character nor this line. Listings are ASCII, and
  % parity's C13 checks it.
  literate={ł}{{\l}}1 {Ł}{{\L}}1 {ą}{{\k a}}1 {ę}{{\k e}}1
           {ś}{{\'s}}1 {ć}{{\'c}}1 {ż}{{\.z}}1 {ź}{{\'z}}1 {ó}{{\'o}}1 {ń}{{\'n}}1
           {Ą}{{\k A}}1 {Ę}{{\k E}}1 {Ś}{{\'S}}1 {Ć}{{\'C}}1
           {Ż}{{\.Z}}1 {Ź}{{\'Z}}1 {Ó}{{\'O}}1 {Ń}{{\'N}}1
           {—}{{-{}-}}2 {–}{{-}}1 {‘}{{'}}1 {’}{{'}}1
           {“}{{"}}1 {”}{{"}}1 {°}{{\textdegree}}1
}
\lstset{style=bookcode}

\lstnewenvironment{python}[1][]
  {\lstset{language=pythonx,style=bookcode,#1}}
  {}
\lstnewenvironment{shellcmd}[1][]
  {\lstset{language=bash,style=bookcode,numbers=none,keywordstyle=\color{mteal}\bfseries,#1}}
  {}

% The path is printed above every file listing, in the form the reader opens
% in the editor. \nobreak keeps it on the page of the listing it names.
% The path must not be stranded at the foot of a page with its listing
% overleaf: \nobreak alone does not hold, because listings opens with its own
% vertical material. \needspace asks for room for the path and the first few
% lines of code, and turns the page first when there is not.
% \Needspace (capital N) measures \pagegoal-\pagetotal exactly; lower-case
% \needspace works through glue and penalties and was measured, by the
% Chapter 11 pass, turning pages early and leaving gaps. It must be used
% between paragraphs, which \listingpath's own \par guarantees. Chapters use
% \Needspace* too, so the fallback for a machine without the package defines
% both forms (the star swallowed with its argument).
\makeatletter
\IfFileExists{needspace.sty}{\usepackage{needspace}}{%
  \providecommand{\needspace}[1]{}%
  \providecommand{\Needspace}{\@ifstar\@gobble\@gobble}}
\makeatother
\newcommand{\listingpath}[1]{%
  \par\medskip\Needspace{6\baselineskip}\noindent
  {\ttfamily\footnotesize\color{mgrey}\detokenize{#1}}%
  \par\nobreak}

% Usage: \pyfile{code/ch05/cobweb.py}{Caption}{lst:label}
\newcommand{\pyfile}[3]{%
  \listingpath{#1}%
  \lstinputlisting[language=pythonx,style=bookcode,aboveskip=0.3em,
    caption={#2},label={#3}]{#1}%
}

% A named region of a file, delimited by `# --8<-- [start:name]` and
% `# --8<-- [end:name]`. Matched through listings' range prefix and suffix
% keys -- the other obvious form matches nothing and prints the WHOLE file
% without a warning (measured in the sibling book). A region name is a word,
% never a number. The line numbers printed are the file's own.
% Usage: \pyregion{code/ch05/cobweb.py}{step}{Caption}{lst:label}
\lstset{
  rangebeginprefix=\#\ --8<--\ [start:, rangebeginsuffix=],
  rangeendprefix=\#\ --8<--\ [end:, rangeendsuffix=],
  includerangemarker=false}
\newcommand{\pyregion}[4]{%
  \listingpath{#1}%
  \lstinputlisting[language=pythonx,style=bookcode,aboveskip=0.3em,
    linerange={#2-#2},caption={#3},label={#4}]{#1}%
}

% A program's output, written by code/measure/transcripts.py. There is no
% typed form and there must not be: \val{} does not expand in a verbatim
% body, and a transcript nobody ran looks exactly like one that was.
\newcommand{\transcript}[1]{%
  \par\smallskip
  \IfFileExists{figures/transcripts/#1.txt}{%
    \lstinputlisting[style=bookcode,numbers=none,basicstyle=\ttfamily\footnotesize,
      keywordstyle=\color{black},backgroundcolor=\color{white},
      language={}]{figures/transcripts/#1.txt}%
  }{%
    \begin{tcolorbox}[enhanced, sharp corners, width=\linewidth,
      colback=mlight, colframe=mgrey, boxrule=0.6pt,
      left=8pt, right=8pt, top=6pt, bottom=6pt]
      {\bfseries\color{mgrey}\small \lblTranscriptMissing}\\[2pt]
      {\ttfamily\scriptsize\color{mgrey} figures/transcripts/#1.txt}
    \end{tcolorbox}%
  }%
  \par\smallskip
}

% Inline literals. Underscores inside these must be written \_.
\newcommand{\code}[1]{\texttt{\small #1}}
\newcommand{\pkg}[1]{\texttt{\small\color{mblue}#1}}
\newcommand{\api}[1]{\texttt{\small\color{mteal}#1}}

% ============================================================
%  ADMONITIONS
%  A small, fixed vocabulary. `tinted' may be read in passing; `outlined' is
%  a block the reader must stop at. A third treatment would mean the first two
%  are not carrying their meaning.
% ============================================================
\tcbset{
  cfz tinted/.style={
    enhanced, breakable, sharp corners,
    lines before break=4,
    colback=mlight, boxrule=0pt, leftrule=3pt,
    left=8pt, right=8pt, top=6pt, bottom=6pt,
    before skip=13pt, after skip=13pt,
    fonttitle=\bfseries\small,
    coltitle=black, colbacktitle=mlight,
    attach title to upper={\par\smallskip},
  },
  cfz outlined/.style={
    enhanced, breakable, sharp corners,
    % A box broken straight after its title is a heading stranded at the foot
    % of a page with its question overleaf. Four lines is title plus three.
    lines before break=4,
    colback=white, boxrule=0.6pt,
    left=8pt, right=8pt, top=6pt, bottom=6pt,
    before skip=13pt, after skip=13pt,
    fonttitle=\bfseries\small,
    coltitle=black, colbacktitle=white,
    attach title to upper={\par\smallskip},
  },
}

% Read in passing.
\newtcolorbox{note}[1][]{cfz tinted, colframe=mblue, title={\lblNote}, #1}
\newtcolorbox{warning}[1][]{cfz tinted, colframe=mred, title={\lblWarning}, #1}
% Where you meet this in the world. If it cannot name a specific system,
% instrument, paper or event, it is prose and not this box.
\newtcolorbox{worldbox}[1][]{cfz tinted, colframe=mteal, title={\lblWorld}, #1}
% What is stated and not proved here, and where the proof lives.
\newtcolorbox{rigourbox}[1][]{cfz tinted, colframe=mgrey, title={\lblRigour}, #1}
% The companion library the reader builds, under code/src/chaoslab/.
\newtcolorbox{projectbox}[1][]{cfz tinted, colframe=mamber, title={\lblProject}, #1}

% Stop and read.
% The misconception, named AFTER the reader has been asked the question it
% gets wrong. Appendix B is the catalogue; a trap box is one entry in place.
\newtcolorbox{trapbox}[1][]{cfz outlined, colframe=mred, title={\lblTrap}, #1}
% A listing that has not been run on the pinned versions. Counted by
% `make debt'; nothing ships to a reader with one attached.
\newtcolorbox{verifybox}[1][]{cfz outlined, colframe=mgrey, title={\lblVerify}, #1}
\newtcolorbox{labbox}[1][]{cfz outlined, colframe=mgreen, title={\lblLab}, #1}

% ============================================================
%  THE METHOD --- questions before answers
%  ---------------------------------------------------------
%  The unit of teaching is a question the reader answers BEFORE reading the
%  answer. That is retrieval practice and errorful generation, the two
%  findings the programmed-learning format is defended by, and it is the one
%  thing the sibling mathematics book does that this one keeps whole.
%
%  \begin{think} ... \end{think}   the question, numbered per chapter, ending
%                                  with the instruction to answer before
%                                  reading on;
%  \reveal{short answer}           a centred box with the answer, which is the
%                                  first thing after the question -- the edge
%                                  the reader covers with a hand, or scrolls
%                                  to;
%  \begin{revealblock} ... \end{revealblock}  the same, for a worked answer.
%
%  tools/check_structure.py --think fails when a think box is not followed by
%  a reveal before the next think box or section, and when a reveal answers
%  nothing. That is a property of the source, so it is checked mechanically.
%
%  The instruction says "before you read on", never "before you turn over":
%  the two editions paginate differently, so the second is true in one of them.
% ============================================================
\newcounter{think}[chapter]
\renewcommand{\thethink}{\thechapter.\arabic{think}}
\newtcolorbox{thinkbox}[1][]{cfz outlined, colframe=mviolet, #1}
\NewDocumentEnvironment{think}{}{%
  \refstepcounter{think}%
  \begin{thinkbox}[title={\lblThink~\thethink}]%
}{%
  \par\smallskip
  {\raggedleft\itshape\small\color{mviolet}\lblThinkCue\par}%
  \end{thinkbox}%
}
\tcbset{
  cfz answer/.style={
    enhanced, sharp corners,
    colback=mfaint, colframe=mviolet!55, boxrule=0.5pt,
    before skip=4pt, after skip=10pt,
  },
}
\newcommand{\reveal}[1]{%
  \par\nobreak
  \begin{tcolorbox}[cfz answer, unbreakable,
    left=8pt, right=8pt, top=5pt, bottom=5pt,
    width=0.86\linewidth, center, halign=flush center]
    {\small\bfseries\color{mviolet}\lblAnswerMark}\enspace #1
  \end{tcolorbox}%
  \nobreak\noindent\ignorespaces
}
\NewDocumentEnvironment{revealblock}{}{%
  \par\nopagebreak
  \begin{tcolorbox}[cfz answer, breakable,
    left=10pt, right=10pt, top=6pt, bottom=6pt]
  {\small\bfseries\color{mviolet}\lblAnswerMark}\par\smallskip
}{%
  \end{tcolorbox}%
  \nopagebreak\par\noindent\ignorespaces
}

% ---- Outcomes on entry, and "Can you?" generated from them ----------------
% The self-check at the end of a chapter IS the list of outcomes at its
% start, one for one, because it is built from the same tokens: each
% \outcome both prints its item and appends a row to \cfz@canrows. A \loop
% inside a tabularx body does not survive alignment scanning, which is why
% the rows are accumulated at declaration time (the sibling book's finding).
\makeatletter
\newcommand{\cfz@canrows}{}
\newcommand{\cfz@ratebox}{%
  \fbox{\rule{0pt}{0.9ex}\rule{0.9ex}{0pt}}}
\newcommand{\cfz@rates}{%
  {\footnotesize\color{mgrey}%
   1\,\cfz@ratebox\ 2\,\cfz@ratebox\ 3\,\cfz@ratebox\ 4\,\cfz@ratebox\ 5\,\cfz@ratebox}}
\NewDocumentEnvironment{outcomes}{}{%
  \gdef\cfz@canrows{}%
  \begin{tcolorbox}[cfz tinted, colframe=mblue, title={\lblOutcomes}]%
  \lblOutcomesLead
  \begin{itemize}[leftmargin=1.4em, itemsep=2pt, topsep=4pt]%
}{%
  \end{itemize}%
  \end{tcolorbox}%
}
\newcommand{\outcome}[1]{%
  \item #1%
  \gappto\cfz@canrows{#1 & \cfz@rates \\ \addlinespace[3pt]}%
}
\newcommand{\canyou}{%
  \section*{\lblCanYou}%
  \phantomsection\addcontentsline{toc}{section}{\lblCanYou}%
  \noindent\lblCanYouIntro\par\medskip
  \ifdefempty{\cfz@canrows}{%
    \noindent{\color{mred}\lblNoOutcomes}\par}{%
  {\small
  \noindent\begin{tabularx}{\linewidth}{@{}>{\raggedright\arraybackslash}X r@{}}
    \toprule
    \cfz@canrows
    \bottomrule
  \end{tabularx}\par}}%
}
\makeatother

% ---- Summary -------------------------------------------------------------
\newtcolorbox{summarybox}[1][]{cfz tinted, colframe=mblue, title={\lblSummary}, #1}

% ---- Checkpoint questions, answered at the back --------------------------
% Answers are written AT THE QUESTION with \answerto{} and printed in
% Appendix A. They are carried there by a global token list rather than an
% auxiliary file, because an answer contains mathematics and \val{}, and
% material written to an .aux is expanded at shipout, where fragile commands
% break with errors naming neither the answer nor the chapter.
% \include is sequential within a run, so the list is complete by the time
% the appendix is typeset.
\makeatletter
\newcommand{\cfz@answers}{}
\newcommand{\cfz@lastansch}{0}
\NewDocumentEnvironment{checkpoint}{}{%
  \section*{\lblCheckpoint}%
  \phantomsection\addcontentsline{toc}{section}{\lblCheckpoint}%
  \noindent\lblCheckpointIntro\par
  \begin{enumerate}[label=\textbf{\arabic*.}, leftmargin=2em, itemsep=4pt]%
}{%
  \end{enumerate}%
}
\newcommand{\answerto}[1]{%
  \ifnum\value{chapter}=\cfz@lastansch\relax\else
    \xdef\cfz@lastansch{\arabic{chapter}}%
    \xappto\cfz@answers{\noexpand\cfz@anschapter{\thechapter}}%
  \fi
  \xappto\cfz@answers{\noexpand\cfz@ansitem{\arabic{enumi}}}%
  \gappto\cfz@answers{{#1}}%
}
\newcommand{\cfz@anschapter}[1]{%
  \par\medskip\noindent{\bfseries\color{mblue}\chaptername~#1}\par\nobreak\smallskip}
\newcommand{\cfz@ansitem}[2]{%
  \par\noindent\hangindent=2em\hangafter=1
  \makebox[2em][l]{\textbf{#1.}}#2\par}
\newcommand{\answersbody}{%
  \ifdefempty{\cfz@answers}{\noindent{\color{mred}\lblNoAnswers}\par}{%
    {\raggedright\cfz@answers}}%
}
\makeatother

% ============================================================
%  LABS
%  \begin{lab}{key}{title} ... \end{lab}
%
%  A lab is not a paragraph but a FILE the reader opens: a starter under
%  code/labs/chNN/<key>.py that must FAIL its test, a reference solution under
%  code/labs/chNN/solutions/, and test_<key>.py beside it. The box prints the
%  paths and the one command that runs the test. An environment rather than a
%  macro, because a macro argument cannot hold a verbatim fragment.
%  tools/check_structure.py --labs fails when a file is missing or the key's
%  chapter number disagrees with its chapter.
% ============================================================
\newcounter{lab}[chapter]
\renewcommand{\thelab}{\thechapter.\arabic{lab}}
\makeatletter
\newcommand{\labdir}{ch\two@digits{\value{chapter}}}
\newcommand{\listoflabs}{%
  \section*{\lblLabManifest}%
  \phantomsection\addcontentsline{toc}{section}{\lblLabManifest}%
  \begingroup\c@tocdepth=2\relax\@starttoc{lab}\endgroup%
}
\makeatother
\NewDocumentEnvironment{lab}{mm}{%
  \refstepcounter{lab}%
  \addcontentsline{lab}{subsection}{\protect\numberline{\thelab}\texttt{\protect\detokenize{#1}} --- #2}%
  \begin{labbox}[title={\lblLab~\thelab\ \dash{} #2}]%
}{%
  \par\medskip
  {\small\setlength{\parindent}{0pt}%
  \lblLabStarter: \texttt{code/labs/\labdir/\detokenize{#1}.py}\\
  \lblLabTest: \texttt{code/labs/\labdir/test\_\detokenize{#1}.py}\\
  \lblLabRun: \texttt{cd code \&\& uv run pytest -k \detokenize{#1}}\par}%
  \end{labbox}%
}

% ============================================================
%  FIGURES --- plots and diagrams
%  ---------------------------------------------------------
%  \plotfig{key}{caption}{one-line manifest description}
%    A PLOT, generated by code/plots/<chapter>.py from the same library the
%    listings use, into figures/plots/<lang>/<key>.pdf. Per language, because
%    an axis label in the wrong language is the half-translation that makes a
%    bilingual book look machine-made, and because the Polish edition owes a
%    decimal comma on its tick labels too. Build output, gitignored.
%
%  \mermaidfig{key}{caption}{one-line manifest description}
%    A DIAGRAM of an idea rather than of data, from
%    figures/mermaid/<lang>/<key>.mmd, rendered by `make diagrams'.
%
%  If the rendered PDF is absent either macro prints a visible marker in the
%  flow rather than a float, so a missing figure cannot be mistaken for a
%  layout choice. \phantomsection before \addcontentsline is load-bearing:
%  without it the manifest entry records whatever anchor preceded the figure,
%  which after a listing is a line that was never printed.
% ============================================================
\makeatletter
\newcommand{\listofplots}{%
  \section*{\lblPlotManifest}%
  \phantomsection\addcontentsline{toc}{section}{\lblPlotManifest}%
  \begingroup\c@tocdepth=2\relax\@starttoc{plt}\endgroup%
}
\newcommand{\listofdiagrams}{%
  \section*{\lblDiagramManifest}%
  \phantomsection\addcontentsline{toc}{section}{\lblDiagramManifest}%
  \begingroup\c@tocdepth=2\relax\@starttoc{dgm}\endgroup%
}
\makeatother

\newcommand{\figmissing}[3]{% #1 message, #2 path, #3 caption
  \par\medskip
  \begin{tcolorbox}[enhanced, breakable, width=\linewidth, sharp corners,
    colback=white, colframe=mgrey, boxrule=0.8pt,
    left=8pt, right=8pt, top=8pt, bottom=8pt]
    {\bfseries\color{mgrey}\small #1}\\[4pt]
    {\ttfamily\scriptsize\color{mgrey} #2}
  \end{tcolorbox}%
  \captionof{figure}{#3}%
  \par\medskip
}

\newcommand{\plotfig}[3]{%
  \phantomsection%
  \addcontentsline{plt}{subsection}{\protect\numberline{}\texttt{#1} --- #3}%
  \IfFileExists{figures/plots/\booklang/#1.pdf}{%
    \begin{figure}[htbp]
    \centering
    \includegraphics[width=\linewidth,height=0.45\textheight,keepaspectratio]{figures/plots/\booklang/#1.pdf}%
    \caption{#2}\label{fig:#1}
    \end{figure}%
  }{%
    \figmissing{\lblPlotNotRendered}{figures/plots/\booklang/#1.pdf}{#2}\label{fig:#1}%
  }%
}

\newcommand{\mermaidfig}[3]{%
  \phantomsection%
  \addcontentsline{dgm}{subsection}{\protect\numberline{}\texttt{#1.mmd} --- #3}%
  \IfFileExists{figures/diagrams/\booklang/#1.pdf}{%
    \begin{figure}[htbp]
    \centering
    \includegraphics[width=0.95\linewidth,height=0.42\textheight,keepaspectratio]{figures/diagrams/\booklang/#1.pdf}%
    \caption{#2}\label{fig:#1}
    \end{figure}%
  }{%
    \figmissing{\lblNotRendered}{figures/mermaid/\booklang/#1.mmd}{#2}\label{fig:#1}%
  }%
}

% ============================================================
%  APPENDIX B: the misreadings, generated by tools/gen_traps.py
%
%  \trapchapter{Chapter~\ref{ch:x}}{title}    a heading per chapter
%  \trapentry{n}{habit}{truth}{\ref{sec:x}}  one misreading
%
%  Set ragged right: an entry is mostly short sentences with \code{} spans in
%  them, and a justified line full of unbreakable runs is an overfull box
%  waiting for the edition with the longer words.
% ============================================================
\newcommand{\trapchapter}[2]{%
  \par\medskip
  {\large\bfseries\color{mblue}#1\quad\normalfont\itshape\color{mgrey}#2\par}%
  \nopagebreak\smallskip}
\newcommand{\trapentry}[4]{%
  \par\noindent
  \begingroup\raggedright\hangindent=2.4em\hangafter=1
  \makebox[2.4em][l]{\textbf{\color{mred}#1}}%
  \textbf{\enquote{#2}}\enspace #3\enspace{\color{mgrey}\S\,#4}\par
  \endgroup\smallskip}

% ============================================================
%  CHAPTER STUBS
% ============================================================
\newcommand{\chapterstub}[1]{%
  \begin{tcolorbox}[
    enhanced, breakable, width=\linewidth, sharp corners,
    colback=mlight, colframe=mred, boxrule=1pt,
    borderline={0.6pt}{3pt}{mred, dashed},
    left=12pt, right=12pt, top=12pt, bottom=12pt]
    {\bfseries\color{mred}\large \lblNotWritten}\\[8pt]
    \small\raggedright #1
  \end{tcolorbox}%
}

% ============================================================
%  COMPUTED VALUES
%  A number the reader cannot do in their head is computed by a script under
%  code/measure/, written to figures/values/<chapter>.tex as
%      \pyval{ch06.steps.apart}{34}
%  and reaches the page as \val{ch06.steps.apart}. The book contains
%  references, not digits; `make verify' fails when a committed value no
%  longer matches its script. \val applies the edition's decimal marker, which
%  is how the Polish edition gets its comma without anybody remembering.
% ============================================================
% The comma is BRACED: siunitx sets numbers in maths mode, where a bare comma
% is punctuation and gets a thin space after it ("9, 81"). {,} is an ordinary
% symbol, which is what a decimal marker is.
\newcommand{\bookdecimalmarker}{\ifpl{{,}}{.}}
\IfFileExists{siunitx.sty}{%
  \usepackage{siunitx}%
  \sisetup{
    group-digits = integer,
    group-minimum-digits = 5,
    group-separator = {\,},
    output-decimal-marker = {\bookdecimalmarker},
    exponent-product = \times,
    print-unity-mantissa = false
  }%
  \newcommand{\booknum}[1]{\num{##1}}%
}{%
  \newcommand{\booknum}[1]{##1}%
  \providecommand{\num}[1]{##1}%
}

\makeatletter
\newcommand{\pyvalues}{}
\newcommand{\pyval}[2]{\csgdef{pyval@#1}{#2}\listxadd{\pyvalues}{#1}}
\newcommand{\pyvaltext}[2]{%
  \csgdef{pyval@#1}{#2}\csgdef{pyvaltext@#1}{}\listxadd{\pyvalues}{#1}}
% \num on a non-numeric body is a FATAL siunitx error, so \val refuses a text
% value by name rather than letting siunitx fail on it.
\newcommand{\val}[1]{%
  \ifcsdef{pyval@#1}{%
    \ifcsdef{pyvaltext@#1}{%
      \PackageError{chaosbook}{Value `#1' is not a number}%
        {Use \string\valtext{#1} instead.}%
    }{}%
    \booknum{\csuse{pyval@#1}}%
  }{\textcolor{mred}{\textbf{[\lblValueMissing: #1]}}}%
}
\newcommand{\valtext}[1]{%
  \ifcsdef{pyval@#1}{\csuse{pyval@#1}}%
    {\textcolor{mred}{\textbf{[\lblValueMissing: #1]}}}%
}
\makeatother

% figures/values/all.tex is generated by `make numbers' and \inputs each
% value file. A build with no values still compiles; every \val{} then prints
% a visible marker. \valuesindex lets the one-chapter build point elsewhere.
\providecommand{\valuesindex}{figures/values/all}
\IfFileExists{\valuesindex.tex}{\input{\valuesindex}}{%
  \AtBeginDocument{\PackageWarning{chaosbook}{No computed values found. Run `make numbers'.}}}

% ============================================================
%  PINNED VERSIONS --- single source of truth
%  Change these here AND in code/pyproject.toml; tools/check_structure.py
%  --pins fails when the two disagree. A chapter never types a version.
% ============================================================
\newcommand{\bookedition}{0.1}
\newcommand{\bookyear}{2026}
\newcommand{\pyver}{3.14}              % the minor the book targets
\newcommand{\pypatch}{3.14.7}          % the exact interpreter the code ran on
\newcommand{\uvver}{0.12.13}           % uv
\newcommand{\numpyver}{2.5.3}          % numpy
\newcommand{\matplotlibver}{3.11.2}    % matplotlib
\newcommand{\pytestver}{9.1.1}         % pytest
\newcommand{\ruffver}{0.16.9}          % ruff
% The prose licence is a name with a version in it, not a quantity: it is a
% macro so that neither edition typesets it as a decimal to be localised.
\newcommand{\proselicence}{CC~BY-NC-SA~4.0}

% ============================================================
%  INDEX
% ============================================================
\apptocmd{\theindex}{\raggedcolumns\raggedright}{}{%
  \PackageWarning{chaosbook}{Could not patch theindex for ragged setting}}
\providecommand*\see[2]{}
\renewcommand*\see[2]{\emph{\lblIndexSee} #1}
\providecommand*\seealso[2]{}
\renewcommand*\seealso[2]{\emph{\lblIndexSeeAlso} #1}

\makeindex
 still compiles and says what it did.
\providecommand{\booklang}{en}

\usepackage{etoolbox}   % loaded first: the language switch below needs it

% ---------- Page geometry ----------
% ONE format: A4, 12pt, ONE-SIDED. The book is read on a screen, usually with
% an editor open beside it, so there is no recto, no verso, no gutter and no
% blank leaf before a chapter. The measure is about seventy-five characters
% of prose. At \footnotesize a listing line of 79 characters fits without
% wrapping, and 79 is the width every file under code/ is held to by
% tools/check_structure.py --lines.
\usepackage[
  a4paper,
  left=3.1cm, right=3.1cm, top=2.6cm, bottom=3.0cm,
  head=26pt, headsep=0.8cm, footskip=1.4cm
]{geometry}

% ---------- Encoding & fonts ----------
\usepackage[T1]{fontenc}
\usepackage[utf8]{inputenc}
\IfFileExists{lmodern.sty}{\usepackage{lmodern}}{}
% newtxtext takes its TS1 symbols from TeX Gyre Termes (Debian: tex-gyre).
% A machine with newtx and without tex-gyre dies on the copyright page. CI's
% full TeX Live is the reference; locally install tex-gyre as well.
\IfFileExists{newtxtext.sty}{\usepackage{newtxtext}}{}
% amsmath before newtxmath; amssymb only when newtxmath is absent, because
% loading both is a fatal clash (\Bbbk already defined) that is invisible on
% a machine without newtx.
\usepackage{amsmath}
\IfFileExists{newtxmath.sty}{\usepackage{newtxmath}}{\usepackage{amssymb}}
\IfFileExists{inconsolata.sty}{\usepackage[varqu,scaled=0.95]{inconsolata}}{}
% Without lmodern, microtype runs with font expansion off, which is what
% resolves a marginal overfull line -- so a container short of lmodern
% disagrees with CI about line breaks. Install lmodern locally too.
\IfFileExists{lmodern.sty}{\usepackage{microtype}}{\usepackage[expansion=false]{microtype}}

% upquote: in T1 the input ' is a right single quote, and a listing that sets
% f('x') with curly quotes is a SyntaxError when typed back in.
\IfFileExists{upquote.sty}{\usepackage{upquote}}{}
% And its twin: in T1 the pair -- is an en-dash ligature in the typewriter
% family too, so \code{--flag} would print a flag nobody can type. Scoped to
% tt*, so prose keeps its own dashes.
\DisableLigatures[-]{encoding = T1, family = tt*}

% ---------- Language ----------
% babel with a missing language file is FATAL, so each .ldf is probed first.
\IfFileExists{babel.sty}{%
  \IfFileExists{polish.ldf}{%
    \IfFileExists{british.ldf}{%
      \ifdefstring{\booklang}{pl}%
        {\usepackage[british,main=polish]{babel}}%
        {\usepackage[polish,main=british]{babel}}%
    }{\usepackage[polish]{babel}}%
  }{%
    \IfFileExists{british.ldf}{\usepackage[british]{babel}}{}%
  }%
}{}
% babel-polish makes " an active shorthand, which is a category-code accident
% waiting to happen inside listings. Polish quotation marks come from
% \enquote, which csquotes sets per language.
\ifdefstring{\booklang}{pl}{\AtBeginDocument{\shorthandoff{"}}}{}

% ---------- Core packages ----------
\usepackage{graphicx}
\usepackage{xcolor}
\usepackage{booktabs}
\usepackage{tabularx}
\usepackage{longtable}
\usepackage{array}
\usepackage{enumitem}
\usepackage{listings}
\usepackage[most]{tcolorbox}
\usepackage{fancyhdr}
\usepackage{titlesec}
\usepackage{caption}
\usepackage{imakeidx}
\IfFileExists{csquotes.sty}{\usepackage[autostyle=true]{csquotes}}{%
  \providecommand{\enquote}[1]{``##1''}}
\usepackage[hidelinks]{hyperref}

% ---------- Palette ----------
% The family's palette, so the three books read as a set. Two colours are
% this book's own: the lab (green) and the stop-and-think box (violet).
\definecolor{mblue}{HTML}{1F4E79}
\definecolor{mteal}{HTML}{0E7C7B}
\definecolor{mamber}{HTML}{B26A00}
\definecolor{mred}{HTML}{9B2C2C}
\definecolor{mviolet}{HTML}{7B4397}
\definecolor{mgreen}{HTML}{3B7A3B}
\definecolor{mgrey}{HTML}{5A5A5A}
\definecolor{mlight}{HTML}{F4F6F8}
\definecolor{mfaint}{HTML}{FBFCFD}
\definecolor{codebg}{HTML}{FAFAFA}
\definecolor{codeframe}{HTML}{D8DEE4}
\definecolor{kw}{HTML}{0B5394}
\definecolor{str}{HTML}{9B2C2C}
\definecolor{cmt}{HTML}{4F7A4F}
\definecolor{typ}{HTML}{0E7C7B}
\definecolor{dec}{HTML}{7B4397}

\hypersetup{
  colorlinks=true,
  linkcolor=mblue,
  urlcolor=mteal,
  citecolor=mblue,
  pdfauthor={Konrad Cinkusz}
}

% \ifpl{Polish}{English} -- for the handful of places where a string is
% structural rather than content. Robust, because a fragile \ifpl inside a
% moving argument fails on the next run in exactly one of the two editions.
\DeclareRobustCommand{\ifpl}[2]{\ifdefstring{\booklang}{pl}{#1}{#2}}
\makeatletter
\pdfstringdefDisableCommands{%
  \ifdefstring{\booklang}{pl}{\let\ifpl\@firstoftwo}{\let\ifpl\@secondoftwo}%
  % hyperref drops \color from a bookmark and keeps its ARGUMENT, so a title
  % carrying \code{} would bookmark as "mtealx". Gobble the colour name.
  \def\color#1{}%
}
\makeatother

% ---------- Language strings ----------
\input{lang/\booklang}

% The PDF's subject, keywords and language read \lblPdf... from the language
% file, so they cannot sit in the colour block above, which runs first. Two
% blocks is the ordering, not an oversight.
\hypersetup{
  pdfsubject={\lblPdfSubject},
  pdfkeywords={\lblPdfKeywords},
  pdflang={\lblPdfLang}
}

% ---------- Headings ----------
\titleformat{\chapter}[display]
  {\normalfont\huge\bfseries\color{mblue}\raggedright}
  {\normalfont\normalsize\color{mgrey}\MakeUppercase{\chaptertitlename}\ \thechapter}
  {8pt}{\Huge\raggedright}
\titlespacing*{\chapter}{0pt}{10pt}{28pt}
\titleformat{\section}{\normalfont\Large\bfseries\color{mblue}\raggedright}{\thesection}{0.8em}{}
\titleformat{\subsection}{\normalfont\large\bfseries\color{mgrey}\raggedright}{\thesubsection}{0.7em}{}

\setcounter{tocdepth}{1}
\setcounter{secnumdepth}{2}

\pagestyle{fancy}
\fancyhf{}
\fancyhead[L]{\small\nouppercase{\leftmark}}
\fancyhead[R]{\small\thepage}
\renewcommand{\headrulewidth}{0.4pt}
\fancypagestyle{plain}{\fancyhf{}\fancyfoot[R]{\small\thepage}%
  \renewcommand{\headrulewidth}{0pt}}
% After \pagestyle{fancy}, which installs its own \chaptermark.
\renewcommand{\chaptermark}[1]{\markboth{\thechapter.\ #1}{}}

% ---------- Mathematics ----------
% The handful of notations the book uses, as macros so both editions set the
% same symbol. The book writes \ln for the natural logarithm and \log_{10}
% where the base is ten; a bare \log is ambiguous between the two readerships
% and is not used.
\newcommand{\R}{\mathbb{R}}
\newcommand{\dd}{\,\mathrm{d}}
\newcommand{\abs}[1]{\left\lvert #1\right\rvert}
\newcommand{\iu}{\mathrm{i}}

% ============================================================
%  CODE LISTINGS
%  Every listing in a chapter is a FILE under code/, pulled in with \pyfile or
%  \pyregion, and CI runs it on the pinned interpreter. An in-source
%  \begin{python} block is allowed for a fragment of three or four lines that
%  shows a shape rather than runs.
% ============================================================
\lstdefinelanguage{pythonx}{
  language=Python,
  morekeywords={async,await,match,case,nonlocal,yield,type},
  morekeywords=[2]{
    orbit,iterate,logistic,rk4_step,euler_step,integrate,lorenz,henon,
    lyapunov,separation,escape_time,box_count
  },
  sensitive=true
}

\lstdefinestyle{bookcode}{
  basicstyle=\ttfamily\footnotesize,
  keywordstyle=\color{kw}\bfseries,
  keywordstyle=[2]\color{typ},
  stringstyle=\color{str},
  % Colour only: inconsolata has no italic, and an italic comment style is
  % silently substituted and logged on every build.
  commentstyle=\color{cmt},
  numbers=left,
  numberstyle=\ttfamily\tiny\color{mgrey},
  numbersep=8pt,
  backgroundcolor=\color{codebg},
  frame=single,
  rulecolor=\color{codeframe},
  framesep=6pt,
  breaklines=true,
  breakatwhitespace=false,
  postbreak=\mbox{\textcolor{mgrey}{$\hookrightarrow$}\space},
  showstringspaces=false,
  tabsize=4,
  columns=fullflexible,
  keepspaces=true,
  captionpos=b,
  aboveskip=1.2em,
  belowskip=1.2em,
  % Do NOT add a mapping for U+00A0: it aborts the run with a message that
  % names neither the character nor this line. Listings are ASCII, and
  % parity's C13 checks it.
  literate={ł}{{\l}}1 {Ł}{{\L}}1 {ą}{{\k a}}1 {ę}{{\k e}}1
           {ś}{{\'s}}1 {ć}{{\'c}}1 {ż}{{\.z}}1 {ź}{{\'z}}1 {ó}{{\'o}}1 {ń}{{\'n}}1
           {Ą}{{\k A}}1 {Ę}{{\k E}}1 {Ś}{{\'S}}1 {Ć}{{\'C}}1
           {Ż}{{\.Z}}1 {Ź}{{\'Z}}1 {Ó}{{\'O}}1 {Ń}{{\'N}}1
           {—}{{-{}-}}2 {–}{{-}}1 {‘}{{'}}1 {’}{{'}}1
           {“}{{"}}1 {”}{{"}}1 {°}{{\textdegree}}1
}
\lstset{style=bookcode}

\lstnewenvironment{python}[1][]
  {\lstset{language=pythonx,style=bookcode,#1}}
  {}
\lstnewenvironment{shellcmd}[1][]
  {\lstset{language=bash,style=bookcode,numbers=none,keywordstyle=\color{mteal}\bfseries,#1}}
  {}

% The path is printed above every file listing, in the form the reader opens
% in the editor. \nobreak keeps it on the page of the listing it names.
% The path must not be stranded at the foot of a page with its listing
% overleaf: \nobreak alone does not hold, because listings opens with its own
% vertical material. \needspace asks for room for the path and the first few
% lines of code, and turns the page first when there is not.
% \Needspace (capital N) measures \pagegoal-\pagetotal exactly; lower-case
% \needspace works through glue and penalties and was measured, by the
% Chapter 11 pass, turning pages early and leaving gaps. It must be used
% between paragraphs, which \listingpath's own \par guarantees. Chapters use
% \Needspace* too, so the fallback for a machine without the package defines
% both forms (the star swallowed with its argument).
\makeatletter
\IfFileExists{needspace.sty}{\usepackage{needspace}}{%
  \providecommand{\needspace}[1]{}%
  \providecommand{\Needspace}{\@ifstar\@gobble\@gobble}}
\makeatother
\newcommand{\listingpath}[1]{%
  \par\medskip\Needspace{6\baselineskip}\noindent
  {\ttfamily\footnotesize\color{mgrey}\detokenize{#1}}%
  \par\nobreak}

% Usage: \pyfile{code/ch05/cobweb.py}{Caption}{lst:label}
\newcommand{\pyfile}[3]{%
  \listingpath{#1}%
  \lstinputlisting[language=pythonx,style=bookcode,aboveskip=0.3em,
    caption={#2},label={#3}]{#1}%
}

% A named region of a file, delimited by `# --8<-- [start:name]` and
% `# --8<-- [end:name]`. Matched through listings' range prefix and suffix
% keys -- the other obvious form matches nothing and prints the WHOLE file
% without a warning (measured in the sibling book). A region name is a word,
% never a number. The line numbers printed are the file's own.
% Usage: \pyregion{code/ch05/cobweb.py}{step}{Caption}{lst:label}
\lstset{
  rangebeginprefix=\#\ --8<--\ [start:, rangebeginsuffix=],
  rangeendprefix=\#\ --8<--\ [end:, rangeendsuffix=],
  includerangemarker=false}
\newcommand{\pyregion}[4]{%
  \listingpath{#1}%
  \lstinputlisting[language=pythonx,style=bookcode,aboveskip=0.3em,
    linerange={#2-#2},caption={#3},label={#4}]{#1}%
}

% A program's output, written by code/measure/transcripts.py. There is no
% typed form and there must not be: \val{} does not expand in a verbatim
% body, and a transcript nobody ran looks exactly like one that was.
\newcommand{\transcript}[1]{%
  \par\smallskip
  \IfFileExists{figures/transcripts/#1.txt}{%
    \lstinputlisting[style=bookcode,numbers=none,basicstyle=\ttfamily\footnotesize,
      keywordstyle=\color{black},backgroundcolor=\color{white},
      language={}]{figures/transcripts/#1.txt}%
  }{%
    \begin{tcolorbox}[enhanced, sharp corners, width=\linewidth,
      colback=mlight, colframe=mgrey, boxrule=0.6pt,
      left=8pt, right=8pt, top=6pt, bottom=6pt]
      {\bfseries\color{mgrey}\small \lblTranscriptMissing}\\[2pt]
      {\ttfamily\scriptsize\color{mgrey} figures/transcripts/#1.txt}
    \end{tcolorbox}%
  }%
  \par\smallskip
}

% Inline literals. Underscores inside these must be written \_.
\newcommand{\code}[1]{\texttt{\small #1}}
\newcommand{\pkg}[1]{\texttt{\small\color{mblue}#1}}
\newcommand{\api}[1]{\texttt{\small\color{mteal}#1}}

% ============================================================
%  ADMONITIONS
%  A small, fixed vocabulary. `tinted' may be read in passing; `outlined' is
%  a block the reader must stop at. A third treatment would mean the first two
%  are not carrying their meaning.
% ============================================================
\tcbset{
  cfz tinted/.style={
    enhanced, breakable, sharp corners,
    lines before break=4,
    colback=mlight, boxrule=0pt, leftrule=3pt,
    left=8pt, right=8pt, top=6pt, bottom=6pt,
    before skip=13pt, after skip=13pt,
    fonttitle=\bfseries\small,
    coltitle=black, colbacktitle=mlight,
    attach title to upper={\par\smallskip},
  },
  cfz outlined/.style={
    enhanced, breakable, sharp corners,
    % A box broken straight after its title is a heading stranded at the foot
    % of a page with its question overleaf. Four lines is title plus three.
    lines before break=4,
    colback=white, boxrule=0.6pt,
    left=8pt, right=8pt, top=6pt, bottom=6pt,
    before skip=13pt, after skip=13pt,
    fonttitle=\bfseries\small,
    coltitle=black, colbacktitle=white,
    attach title to upper={\par\smallskip},
  },
}

% Read in passing.
\newtcolorbox{note}[1][]{cfz tinted, colframe=mblue, title={\lblNote}, #1}
\newtcolorbox{warning}[1][]{cfz tinted, colframe=mred, title={\lblWarning}, #1}
% Where you meet this in the world. If it cannot name a specific system,
% instrument, paper or event, it is prose and not this box.
\newtcolorbox{worldbox}[1][]{cfz tinted, colframe=mteal, title={\lblWorld}, #1}
% What is stated and not proved here, and where the proof lives.
\newtcolorbox{rigourbox}[1][]{cfz tinted, colframe=mgrey, title={\lblRigour}, #1}
% The companion library the reader builds, under code/src/chaoslab/.
\newtcolorbox{projectbox}[1][]{cfz tinted, colframe=mamber, title={\lblProject}, #1}

% Stop and read.
% The misconception, named AFTER the reader has been asked the question it
% gets wrong. Appendix B is the catalogue; a trap box is one entry in place.
\newtcolorbox{trapbox}[1][]{cfz outlined, colframe=mred, title={\lblTrap}, #1}
% A listing that has not been run on the pinned versions. Counted by
% `make debt'; nothing ships to a reader with one attached.
\newtcolorbox{verifybox}[1][]{cfz outlined, colframe=mgrey, title={\lblVerify}, #1}
\newtcolorbox{labbox}[1][]{cfz outlined, colframe=mgreen, title={\lblLab}, #1}

% ============================================================
%  THE METHOD --- questions before answers
%  ---------------------------------------------------------
%  The unit of teaching is a question the reader answers BEFORE reading the
%  answer. That is retrieval practice and errorful generation, the two
%  findings the programmed-learning format is defended by, and it is the one
%  thing the sibling mathematics book does that this one keeps whole.
%
%  \begin{think} ... \end{think}   the question, numbered per chapter, ending
%                                  with the instruction to answer before
%                                  reading on;
%  \reveal{short answer}           a centred box with the answer, which is the
%                                  first thing after the question -- the edge
%                                  the reader covers with a hand, or scrolls
%                                  to;
%  \begin{revealblock} ... \end{revealblock}  the same, for a worked answer.
%
%  tools/check_structure.py --think fails when a think box is not followed by
%  a reveal before the next think box or section, and when a reveal answers
%  nothing. That is a property of the source, so it is checked mechanically.
%
%  The instruction says "before you read on", never "before you turn over":
%  the two editions paginate differently, so the second is true in one of them.
% ============================================================
\newcounter{think}[chapter]
\renewcommand{\thethink}{\thechapter.\arabic{think}}
\newtcolorbox{thinkbox}[1][]{cfz outlined, colframe=mviolet, #1}
\NewDocumentEnvironment{think}{}{%
  \refstepcounter{think}%
  \begin{thinkbox}[title={\lblThink~\thethink}]%
}{%
  \par\smallskip
  {\raggedleft\itshape\small\color{mviolet}\lblThinkCue\par}%
  \end{thinkbox}%
}
\tcbset{
  cfz answer/.style={
    enhanced, sharp corners,
    colback=mfaint, colframe=mviolet!55, boxrule=0.5pt,
    before skip=4pt, after skip=10pt,
  },
}
\newcommand{\reveal}[1]{%
  \par\nobreak
  \begin{tcolorbox}[cfz answer, unbreakable,
    left=8pt, right=8pt, top=5pt, bottom=5pt,
    width=0.86\linewidth, center, halign=flush center]
    {\small\bfseries\color{mviolet}\lblAnswerMark}\enspace #1
  \end{tcolorbox}%
  \nobreak\noindent\ignorespaces
}
\NewDocumentEnvironment{revealblock}{}{%
  \par\nopagebreak
  \begin{tcolorbox}[cfz answer, breakable,
    left=10pt, right=10pt, top=6pt, bottom=6pt]
  {\small\bfseries\color{mviolet}\lblAnswerMark}\par\smallskip
}{%
  \end{tcolorbox}%
  \nopagebreak\par\noindent\ignorespaces
}

% ---- Outcomes on entry, and "Can you?" generated from them ----------------
% The self-check at the end of a chapter IS the list of outcomes at its
% start, one for one, because it is built from the same tokens: each
% \outcome both prints its item and appends a row to \cfz@canrows. A \loop
% inside a tabularx body does not survive alignment scanning, which is why
% the rows are accumulated at declaration time (the sibling book's finding).
\makeatletter
\newcommand{\cfz@canrows}{}
\newcommand{\cfz@ratebox}{%
  \fbox{\rule{0pt}{0.9ex}\rule{0.9ex}{0pt}}}
\newcommand{\cfz@rates}{%
  {\footnotesize\color{mgrey}%
   1\,\cfz@ratebox\ 2\,\cfz@ratebox\ 3\,\cfz@ratebox\ 4\,\cfz@ratebox\ 5\,\cfz@ratebox}}
\NewDocumentEnvironment{outcomes}{}{%
  \gdef\cfz@canrows{}%
  \begin{tcolorbox}[cfz tinted, colframe=mblue, title={\lblOutcomes}]%
  \lblOutcomesLead
  \begin{itemize}[leftmargin=1.4em, itemsep=2pt, topsep=4pt]%
}{%
  \end{itemize}%
  \end{tcolorbox}%
}
\newcommand{\outcome}[1]{%
  \item #1%
  \gappto\cfz@canrows{#1 & \cfz@rates \\ \addlinespace[3pt]}%
}
\newcommand{\canyou}{%
  \section*{\lblCanYou}%
  \phantomsection\addcontentsline{toc}{section}{\lblCanYou}%
  \noindent\lblCanYouIntro\par\medskip
  \ifdefempty{\cfz@canrows}{%
    \noindent{\color{mred}\lblNoOutcomes}\par}{%
  {\small
  \noindent\begin{tabularx}{\linewidth}{@{}>{\raggedright\arraybackslash}X r@{}}
    \toprule
    \cfz@canrows
    \bottomrule
  \end{tabularx}\par}}%
}
\makeatother

% ---- Summary -------------------------------------------------------------
\newtcolorbox{summarybox}[1][]{cfz tinted, colframe=mblue, title={\lblSummary}, #1}

% ---- Checkpoint questions, answered at the back --------------------------
% Answers are written AT THE QUESTION with \answerto{} and printed in
% Appendix A. They are carried there by a global token list rather than an
% auxiliary file, because an answer contains mathematics and \val{}, and
% material written to an .aux is expanded at shipout, where fragile commands
% break with errors naming neither the answer nor the chapter.
% \include is sequential within a run, so the list is complete by the time
% the appendix is typeset.
\makeatletter
\newcommand{\cfz@answers}{}
\newcommand{\cfz@lastansch}{0}
\NewDocumentEnvironment{checkpoint}{}{%
  \section*{\lblCheckpoint}%
  \phantomsection\addcontentsline{toc}{section}{\lblCheckpoint}%
  \noindent\lblCheckpointIntro\par
  \begin{enumerate}[label=\textbf{\arabic*.}, leftmargin=2em, itemsep=4pt]%
}{%
  \end{enumerate}%
}
\newcommand{\answerto}[1]{%
  \ifnum\value{chapter}=\cfz@lastansch\relax\else
    \xdef\cfz@lastansch{\arabic{chapter}}%
    \xappto\cfz@answers{\noexpand\cfz@anschapter{\thechapter}}%
  \fi
  \xappto\cfz@answers{\noexpand\cfz@ansitem{\arabic{enumi}}}%
  \gappto\cfz@answers{{#1}}%
}
\newcommand{\cfz@anschapter}[1]{%
  \par\medskip\noindent{\bfseries\color{mblue}\chaptername~#1}\par\nobreak\smallskip}
\newcommand{\cfz@ansitem}[2]{%
  \par\noindent\hangindent=2em\hangafter=1
  \makebox[2em][l]{\textbf{#1.}}#2\par}
\newcommand{\answersbody}{%
  \ifdefempty{\cfz@answers}{\noindent{\color{mred}\lblNoAnswers}\par}{%
    {\raggedright\cfz@answers}}%
}
\makeatother

% ============================================================
%  LABS
%  \begin{lab}{key}{title} ... \end{lab}
%
%  A lab is not a paragraph but a FILE the reader opens: a starter under
%  code/labs/chNN/<key>.py that must FAIL its test, a reference solution under
%  code/labs/chNN/solutions/, and test_<key>.py beside it. The box prints the
%  paths and the one command that runs the test. An environment rather than a
%  macro, because a macro argument cannot hold a verbatim fragment.
%  tools/check_structure.py --labs fails when a file is missing or the key's
%  chapter number disagrees with its chapter.
% ============================================================
\newcounter{lab}[chapter]
\renewcommand{\thelab}{\thechapter.\arabic{lab}}
\makeatletter
\newcommand{\labdir}{ch\two@digits{\value{chapter}}}
\newcommand{\listoflabs}{%
  \section*{\lblLabManifest}%
  \phantomsection\addcontentsline{toc}{section}{\lblLabManifest}%
  \begingroup\c@tocdepth=2\relax\@starttoc{lab}\endgroup%
}
\makeatother
\NewDocumentEnvironment{lab}{mm}{%
  \refstepcounter{lab}%
  \addcontentsline{lab}{subsection}{\protect\numberline{\thelab}\texttt{\protect\detokenize{#1}} --- #2}%
  \begin{labbox}[title={\lblLab~\thelab\ \dash{} #2}]%
}{%
  \par\medskip
  {\small\setlength{\parindent}{0pt}%
  \lblLabStarter: \texttt{code/labs/\labdir/\detokenize{#1}.py}\\
  \lblLabTest: \texttt{code/labs/\labdir/test\_\detokenize{#1}.py}\\
  \lblLabRun: \texttt{cd code \&\& uv run pytest -k \detokenize{#1}}\par}%
  \end{labbox}%
}

% ============================================================
%  FIGURES --- plots and diagrams
%  ---------------------------------------------------------
%  \plotfig{key}{caption}{one-line manifest description}
%    A PLOT, generated by code/plots/<chapter>.py from the same library the
%    listings use, into figures/plots/<lang>/<key>.pdf. Per language, because
%    an axis label in the wrong language is the half-translation that makes a
%    bilingual book look machine-made, and because the Polish edition owes a
%    decimal comma on its tick labels too. Build output, gitignored.
%
%  \mermaidfig{key}{caption}{one-line manifest description}
%    A DIAGRAM of an idea rather than of data, from
%    figures/mermaid/<lang>/<key>.mmd, rendered by `make diagrams'.
%
%  If the rendered PDF is absent either macro prints a visible marker in the
%  flow rather than a float, so a missing figure cannot be mistaken for a
%  layout choice. \phantomsection before \addcontentsline is load-bearing:
%  without it the manifest entry records whatever anchor preceded the figure,
%  which after a listing is a line that was never printed.
% ============================================================
\makeatletter
\newcommand{\listofplots}{%
  \section*{\lblPlotManifest}%
  \phantomsection\addcontentsline{toc}{section}{\lblPlotManifest}%
  \begingroup\c@tocdepth=2\relax\@starttoc{plt}\endgroup%
}
\newcommand{\listofdiagrams}{%
  \section*{\lblDiagramManifest}%
  \phantomsection\addcontentsline{toc}{section}{\lblDiagramManifest}%
  \begingroup\c@tocdepth=2\relax\@starttoc{dgm}\endgroup%
}
\makeatother

\newcommand{\figmissing}[3]{% #1 message, #2 path, #3 caption
  \par\medskip
  \begin{tcolorbox}[enhanced, breakable, width=\linewidth, sharp corners,
    colback=white, colframe=mgrey, boxrule=0.8pt,
    left=8pt, right=8pt, top=8pt, bottom=8pt]
    {\bfseries\color{mgrey}\small #1}\\[4pt]
    {\ttfamily\scriptsize\color{mgrey} #2}
  \end{tcolorbox}%
  \captionof{figure}{#3}%
  \par\medskip
}

\newcommand{\plotfig}[3]{%
  \phantomsection%
  \addcontentsline{plt}{subsection}{\protect\numberline{}\texttt{#1} --- #3}%
  \IfFileExists{figures/plots/\booklang/#1.pdf}{%
    \begin{figure}[htbp]
    \centering
    \includegraphics[width=\linewidth,height=0.45\textheight,keepaspectratio]{figures/plots/\booklang/#1.pdf}%
    \caption{#2}\label{fig:#1}
    \end{figure}%
  }{%
    \figmissing{\lblPlotNotRendered}{figures/plots/\booklang/#1.pdf}{#2}\label{fig:#1}%
  }%
}

\newcommand{\mermaidfig}[3]{%
  \phantomsection%
  \addcontentsline{dgm}{subsection}{\protect\numberline{}\texttt{#1.mmd} --- #3}%
  \IfFileExists{figures/diagrams/\booklang/#1.pdf}{%
    \begin{figure}[htbp]
    \centering
    \includegraphics[width=0.95\linewidth,height=0.42\textheight,keepaspectratio]{figures/diagrams/\booklang/#1.pdf}%
    \caption{#2}\label{fig:#1}
    \end{figure}%
  }{%
    \figmissing{\lblNotRendered}{figures/mermaid/\booklang/#1.mmd}{#2}\label{fig:#1}%
  }%
}

% ============================================================
%  APPENDIX B: the misreadings, generated by tools/gen_traps.py
%
%  \trapchapter{Chapter~\ref{ch:x}}{title}    a heading per chapter
%  \trapentry{n}{habit}{truth}{\ref{sec:x}}  one misreading
%
%  Set ragged right: an entry is mostly short sentences with \code{} spans in
%  them, and a justified line full of unbreakable runs is an overfull box
%  waiting for the edition with the longer words.
% ============================================================
\newcommand{\trapchapter}[2]{%
  \par\medskip
  {\large\bfseries\color{mblue}#1\quad\normalfont\itshape\color{mgrey}#2\par}%
  \nopagebreak\smallskip}
\newcommand{\trapentry}[4]{%
  \par\noindent
  \begingroup\raggedright\hangindent=2.4em\hangafter=1
  \makebox[2.4em][l]{\textbf{\color{mred}#1}}%
  \textbf{\enquote{#2}}\enspace #3\enspace{\color{mgrey}\S\,#4}\par
  \endgroup\smallskip}

% ============================================================
%  CHAPTER STUBS
% ============================================================
\newcommand{\chapterstub}[1]{%
  \begin{tcolorbox}[
    enhanced, breakable, width=\linewidth, sharp corners,
    colback=mlight, colframe=mred, boxrule=1pt,
    borderline={0.6pt}{3pt}{mred, dashed},
    left=12pt, right=12pt, top=12pt, bottom=12pt]
    {\bfseries\color{mred}\large \lblNotWritten}\\[8pt]
    \small\raggedright #1
  \end{tcolorbox}%
}

% ============================================================
%  COMPUTED VALUES
%  A number the reader cannot do in their head is computed by a script under
%  code/measure/, written to figures/values/<chapter>.tex as
%      \pyval{ch06.steps.apart}{34}
%  and reaches the page as \val{ch06.steps.apart}. The book contains
%  references, not digits; `make verify' fails when a committed value no
%  longer matches its script. \val applies the edition's decimal marker, which
%  is how the Polish edition gets its comma without anybody remembering.
% ============================================================
% The comma is BRACED: siunitx sets numbers in maths mode, where a bare comma
% is punctuation and gets a thin space after it ("9, 81"). {,} is an ordinary
% symbol, which is what a decimal marker is.
\newcommand{\bookdecimalmarker}{\ifpl{{,}}{.}}
\IfFileExists{siunitx.sty}{%
  \usepackage{siunitx}%
  \sisetup{
    group-digits = integer,
    group-minimum-digits = 5,
    group-separator = {\,},
    output-decimal-marker = {\bookdecimalmarker},
    exponent-product = \times,
    print-unity-mantissa = false
  }%
  \newcommand{\booknum}[1]{\num{##1}}%
}{%
  \newcommand{\booknum}[1]{##1}%
  \providecommand{\num}[1]{##1}%
}

\makeatletter
\newcommand{\pyvalues}{}
\newcommand{\pyval}[2]{\csgdef{pyval@#1}{#2}\listxadd{\pyvalues}{#1}}
\newcommand{\pyvaltext}[2]{%
  \csgdef{pyval@#1}{#2}\csgdef{pyvaltext@#1}{}\listxadd{\pyvalues}{#1}}
% \num on a non-numeric body is a FATAL siunitx error, so \val refuses a text
% value by name rather than letting siunitx fail on it.
\newcommand{\val}[1]{%
  \ifcsdef{pyval@#1}{%
    \ifcsdef{pyvaltext@#1}{%
      \PackageError{chaosbook}{Value `#1' is not a number}%
        {Use \string\valtext{#1} instead.}%
    }{}%
    \booknum{\csuse{pyval@#1}}%
  }{\textcolor{mred}{\textbf{[\lblValueMissing: #1]}}}%
}
\newcommand{\valtext}[1]{%
  \ifcsdef{pyval@#1}{\csuse{pyval@#1}}%
    {\textcolor{mred}{\textbf{[\lblValueMissing: #1]}}}%
}
\makeatother

% figures/values/all.tex is generated by `make numbers' and \inputs each
% value file. A build with no values still compiles; every \val{} then prints
% a visible marker. \valuesindex lets the one-chapter build point elsewhere.
\providecommand{\valuesindex}{figures/values/all}
\IfFileExists{\valuesindex.tex}{\input{\valuesindex}}{%
  \AtBeginDocument{\PackageWarning{chaosbook}{No computed values found. Run `make numbers'.}}}

% ============================================================
%  PINNED VERSIONS --- single source of truth
%  Change these here AND in code/pyproject.toml; tools/check_structure.py
%  --pins fails when the two disagree. A chapter never types a version.
% ============================================================
\newcommand{\bookedition}{0.1}
\newcommand{\bookyear}{2026}
\newcommand{\pyver}{3.14}              % the minor the book targets
\newcommand{\pypatch}{3.14.7}          % the exact interpreter the code ran on
\newcommand{\uvver}{0.12.13}           % uv
\newcommand{\numpyver}{2.5.3}          % numpy
\newcommand{\matplotlibver}{3.11.2}    % matplotlib
\newcommand{\pytestver}{9.1.1}         % pytest
\newcommand{\ruffver}{0.16.9}          % ruff
% The prose licence is a name with a version in it, not a quantity: it is a
% macro so that neither edition typesets it as a decimal to be localised.
\newcommand{\proselicence}{CC~BY-NC-SA~4.0}

% ============================================================
%  INDEX
% ============================================================
\apptocmd{\theindex}{\raggedcolumns\raggedright}{}{%
  \PackageWarning{chaosbook}{Could not patch theindex for ragged setting}}
\providecommand*\see[2]{}
\renewcommand*\see[2]{\emph{\lblIndexSee} #1}
\providecommand*\seealso[2]{}
\renewcommand*\seealso[2]{\emph{\lblIndexSeeAlso} #1}

\makeindex
; nothing else differs between them at this level.
%  Every user-visible string lives in lang/en.tex or lang/pl.tex, so a label
%  that appears in one edition and not the other is a missing macro rather
%  than a silent divergence.
%
%  Lineage. This book is the third in a family and takes one thing from each
%  of the other two:
%
%    * from Mathematics from Zero for the AI Engineer: the bilingual single
%      source (\booklang, lang/, one body.tex), the rule that every number is
%      computed and reaches the page as \val{}, and the programmed-learning
%      method -- outcomes on entry, a question before every answer, a
%      misconception elicited before it is named, a self-check generated from
%      the outcomes, and checkpoint answers collected at the back;
%    * from Python for .NET Engineers: every listing is a FILE that CI runs,
%      every lab is a starter that must fail and a solution that must pass,
%      and a console transcript is a computed number that happens to be
%      verbatim.
%
%  What is new here is \plotfig. A book about chaos cannot be taught in boxes
%  and arrows: the bifurcation diagram and the Lorenz attractor ARE the
%  argument, so the plots are generated by code/plots/ from the same library
%  the listings use, per language, and never drawn by hand.
% ============================================================

% \booklang must be set by the main file. Default to English rather than
% failing, so a stray % ============================================================
%  preamble.tex --- styling, environments, macros
%
%  Shared by BOTH editions. main-en.tex and main-pl.tex each define \booklang
%  before % ============================================================
%  preamble.tex --- styling, environments, macros
%
%  Shared by BOTH editions. main-en.tex and main-pl.tex each define \booklang
%  before % ============================================================
%  preamble.tex --- styling, environments, macros
%
%  Shared by BOTH editions. main-en.tex and main-pl.tex each define \booklang
%  before \input{preamble}; nothing else differs between them at this level.
%  Every user-visible string lives in lang/en.tex or lang/pl.tex, so a label
%  that appears in one edition and not the other is a missing macro rather
%  than a silent divergence.
%
%  Lineage. This book is the third in a family and takes one thing from each
%  of the other two:
%
%    * from Mathematics from Zero for the AI Engineer: the bilingual single
%      source (\booklang, lang/, one body.tex), the rule that every number is
%      computed and reaches the page as \val{}, and the programmed-learning
%      method -- outcomes on entry, a question before every answer, a
%      misconception elicited before it is named, a self-check generated from
%      the outcomes, and checkpoint answers collected at the back;
%    * from Python for .NET Engineers: every listing is a FILE that CI runs,
%      every lab is a starter that must fail and a solution that must pass,
%      and a console transcript is a computed number that happens to be
%      verbatim.
%
%  What is new here is \plotfig. A book about chaos cannot be taught in boxes
%  and arrows: the bifurcation diagram and the Lorenz attractor ARE the
%  argument, so the plots are generated by code/plots/ from the same library
%  the listings use, per language, and never drawn by hand.
% ============================================================

% \booklang must be set by the main file. Default to English rather than
% failing, so a stray \input{preamble} still compiles and says what it did.
\providecommand{\booklang}{en}

\usepackage{etoolbox}   % loaded first: the language switch below needs it

% ---------- Page geometry ----------
% ONE format: A4, 12pt, ONE-SIDED. The book is read on a screen, usually with
% an editor open beside it, so there is no recto, no verso, no gutter and no
% blank leaf before a chapter. The measure is about seventy-five characters
% of prose. At \footnotesize a listing line of 79 characters fits without
% wrapping, and 79 is the width every file under code/ is held to by
% tools/check_structure.py --lines.
\usepackage[
  a4paper,
  left=3.1cm, right=3.1cm, top=2.6cm, bottom=3.0cm,
  head=26pt, headsep=0.8cm, footskip=1.4cm
]{geometry}

% ---------- Encoding & fonts ----------
\usepackage[T1]{fontenc}
\usepackage[utf8]{inputenc}
\IfFileExists{lmodern.sty}{\usepackage{lmodern}}{}
% newtxtext takes its TS1 symbols from TeX Gyre Termes (Debian: tex-gyre).
% A machine with newtx and without tex-gyre dies on the copyright page. CI's
% full TeX Live is the reference; locally install tex-gyre as well.
\IfFileExists{newtxtext.sty}{\usepackage{newtxtext}}{}
% amsmath before newtxmath; amssymb only when newtxmath is absent, because
% loading both is a fatal clash (\Bbbk already defined) that is invisible on
% a machine without newtx.
\usepackage{amsmath}
\IfFileExists{newtxmath.sty}{\usepackage{newtxmath}}{\usepackage{amssymb}}
\IfFileExists{inconsolata.sty}{\usepackage[varqu,scaled=0.95]{inconsolata}}{}
% Without lmodern, microtype runs with font expansion off, which is what
% resolves a marginal overfull line -- so a container short of lmodern
% disagrees with CI about line breaks. Install lmodern locally too.
\IfFileExists{lmodern.sty}{\usepackage{microtype}}{\usepackage[expansion=false]{microtype}}

% upquote: in T1 the input ' is a right single quote, and a listing that sets
% f('x') with curly quotes is a SyntaxError when typed back in.
\IfFileExists{upquote.sty}{\usepackage{upquote}}{}
% And its twin: in T1 the pair -- is an en-dash ligature in the typewriter
% family too, so \code{--flag} would print a flag nobody can type. Scoped to
% tt*, so prose keeps its own dashes.
\DisableLigatures[-]{encoding = T1, family = tt*}

% ---------- Language ----------
% babel with a missing language file is FATAL, so each .ldf is probed first.
\IfFileExists{babel.sty}{%
  \IfFileExists{polish.ldf}{%
    \IfFileExists{british.ldf}{%
      \ifdefstring{\booklang}{pl}%
        {\usepackage[british,main=polish]{babel}}%
        {\usepackage[polish,main=british]{babel}}%
    }{\usepackage[polish]{babel}}%
  }{%
    \IfFileExists{british.ldf}{\usepackage[british]{babel}}{}%
  }%
}{}
% babel-polish makes " an active shorthand, which is a category-code accident
% waiting to happen inside listings. Polish quotation marks come from
% \enquote, which csquotes sets per language.
\ifdefstring{\booklang}{pl}{\AtBeginDocument{\shorthandoff{"}}}{}

% ---------- Core packages ----------
\usepackage{graphicx}
\usepackage{xcolor}
\usepackage{booktabs}
\usepackage{tabularx}
\usepackage{longtable}
\usepackage{array}
\usepackage{enumitem}
\usepackage{listings}
\usepackage[most]{tcolorbox}
\usepackage{fancyhdr}
\usepackage{titlesec}
\usepackage{caption}
\usepackage{imakeidx}
\IfFileExists{csquotes.sty}{\usepackage[autostyle=true]{csquotes}}{%
  \providecommand{\enquote}[1]{``##1''}}
\usepackage[hidelinks]{hyperref}

% ---------- Palette ----------
% The family's palette, so the three books read as a set. Two colours are
% this book's own: the lab (green) and the stop-and-think box (violet).
\definecolor{mblue}{HTML}{1F4E79}
\definecolor{mteal}{HTML}{0E7C7B}
\definecolor{mamber}{HTML}{B26A00}
\definecolor{mred}{HTML}{9B2C2C}
\definecolor{mviolet}{HTML}{7B4397}
\definecolor{mgreen}{HTML}{3B7A3B}
\definecolor{mgrey}{HTML}{5A5A5A}
\definecolor{mlight}{HTML}{F4F6F8}
\definecolor{mfaint}{HTML}{FBFCFD}
\definecolor{codebg}{HTML}{FAFAFA}
\definecolor{codeframe}{HTML}{D8DEE4}
\definecolor{kw}{HTML}{0B5394}
\definecolor{str}{HTML}{9B2C2C}
\definecolor{cmt}{HTML}{4F7A4F}
\definecolor{typ}{HTML}{0E7C7B}
\definecolor{dec}{HTML}{7B4397}

\hypersetup{
  colorlinks=true,
  linkcolor=mblue,
  urlcolor=mteal,
  citecolor=mblue,
  pdfauthor={Konrad Cinkusz}
}

% \ifpl{Polish}{English} -- for the handful of places where a string is
% structural rather than content. Robust, because a fragile \ifpl inside a
% moving argument fails on the next run in exactly one of the two editions.
\DeclareRobustCommand{\ifpl}[2]{\ifdefstring{\booklang}{pl}{#1}{#2}}
\makeatletter
\pdfstringdefDisableCommands{%
  \ifdefstring{\booklang}{pl}{\let\ifpl\@firstoftwo}{\let\ifpl\@secondoftwo}%
  % hyperref drops \color from a bookmark and keeps its ARGUMENT, so a title
  % carrying \code{} would bookmark as "mtealx". Gobble the colour name.
  \def\color#1{}%
}
\makeatother

% ---------- Language strings ----------
\input{lang/\booklang}

% The PDF's subject, keywords and language read \lblPdf... from the language
% file, so they cannot sit in the colour block above, which runs first. Two
% blocks is the ordering, not an oversight.
\hypersetup{
  pdfsubject={\lblPdfSubject},
  pdfkeywords={\lblPdfKeywords},
  pdflang={\lblPdfLang}
}

% ---------- Headings ----------
\titleformat{\chapter}[display]
  {\normalfont\huge\bfseries\color{mblue}\raggedright}
  {\normalfont\normalsize\color{mgrey}\MakeUppercase{\chaptertitlename}\ \thechapter}
  {8pt}{\Huge\raggedright}
\titlespacing*{\chapter}{0pt}{10pt}{28pt}
\titleformat{\section}{\normalfont\Large\bfseries\color{mblue}\raggedright}{\thesection}{0.8em}{}
\titleformat{\subsection}{\normalfont\large\bfseries\color{mgrey}\raggedright}{\thesubsection}{0.7em}{}

\setcounter{tocdepth}{1}
\setcounter{secnumdepth}{2}

\pagestyle{fancy}
\fancyhf{}
\fancyhead[L]{\small\nouppercase{\leftmark}}
\fancyhead[R]{\small\thepage}
\renewcommand{\headrulewidth}{0.4pt}
\fancypagestyle{plain}{\fancyhf{}\fancyfoot[R]{\small\thepage}%
  \renewcommand{\headrulewidth}{0pt}}
% After \pagestyle{fancy}, which installs its own \chaptermark.
\renewcommand{\chaptermark}[1]{\markboth{\thechapter.\ #1}{}}

% ---------- Mathematics ----------
% The handful of notations the book uses, as macros so both editions set the
% same symbol. The book writes \ln for the natural logarithm and \log_{10}
% where the base is ten; a bare \log is ambiguous between the two readerships
% and is not used.
\newcommand{\R}{\mathbb{R}}
\newcommand{\dd}{\,\mathrm{d}}
\newcommand{\abs}[1]{\left\lvert #1\right\rvert}
\newcommand{\iu}{\mathrm{i}}

% ============================================================
%  CODE LISTINGS
%  Every listing in a chapter is a FILE under code/, pulled in with \pyfile or
%  \pyregion, and CI runs it on the pinned interpreter. An in-source
%  \begin{python} block is allowed for a fragment of three or four lines that
%  shows a shape rather than runs.
% ============================================================
\lstdefinelanguage{pythonx}{
  language=Python,
  morekeywords={async,await,match,case,nonlocal,yield,type},
  morekeywords=[2]{
    orbit,iterate,logistic,rk4_step,euler_step,integrate,lorenz,henon,
    lyapunov,separation,escape_time,box_count
  },
  sensitive=true
}

\lstdefinestyle{bookcode}{
  basicstyle=\ttfamily\footnotesize,
  keywordstyle=\color{kw}\bfseries,
  keywordstyle=[2]\color{typ},
  stringstyle=\color{str},
  % Colour only: inconsolata has no italic, and an italic comment style is
  % silently substituted and logged on every build.
  commentstyle=\color{cmt},
  numbers=left,
  numberstyle=\ttfamily\tiny\color{mgrey},
  numbersep=8pt,
  backgroundcolor=\color{codebg},
  frame=single,
  rulecolor=\color{codeframe},
  framesep=6pt,
  breaklines=true,
  breakatwhitespace=false,
  postbreak=\mbox{\textcolor{mgrey}{$\hookrightarrow$}\space},
  showstringspaces=false,
  tabsize=4,
  columns=fullflexible,
  keepspaces=true,
  captionpos=b,
  aboveskip=1.2em,
  belowskip=1.2em,
  % Do NOT add a mapping for U+00A0: it aborts the run with a message that
  % names neither the character nor this line. Listings are ASCII, and
  % parity's C13 checks it.
  literate={ł}{{\l}}1 {Ł}{{\L}}1 {ą}{{\k a}}1 {ę}{{\k e}}1
           {ś}{{\'s}}1 {ć}{{\'c}}1 {ż}{{\.z}}1 {ź}{{\'z}}1 {ó}{{\'o}}1 {ń}{{\'n}}1
           {Ą}{{\k A}}1 {Ę}{{\k E}}1 {Ś}{{\'S}}1 {Ć}{{\'C}}1
           {Ż}{{\.Z}}1 {Ź}{{\'Z}}1 {Ó}{{\'O}}1 {Ń}{{\'N}}1
           {—}{{-{}-}}2 {–}{{-}}1 {‘}{{'}}1 {’}{{'}}1
           {“}{{"}}1 {”}{{"}}1 {°}{{\textdegree}}1
}
\lstset{style=bookcode}

\lstnewenvironment{python}[1][]
  {\lstset{language=pythonx,style=bookcode,#1}}
  {}
\lstnewenvironment{shellcmd}[1][]
  {\lstset{language=bash,style=bookcode,numbers=none,keywordstyle=\color{mteal}\bfseries,#1}}
  {}

% The path is printed above every file listing, in the form the reader opens
% in the editor. \nobreak keeps it on the page of the listing it names.
% The path must not be stranded at the foot of a page with its listing
% overleaf: \nobreak alone does not hold, because listings opens with its own
% vertical material. \needspace asks for room for the path and the first few
% lines of code, and turns the page first when there is not.
% \Needspace (capital N) measures \pagegoal-\pagetotal exactly; lower-case
% \needspace works through glue and penalties and was measured, by the
% Chapter 11 pass, turning pages early and leaving gaps. It must be used
% between paragraphs, which \listingpath's own \par guarantees. Chapters use
% \Needspace* too, so the fallback for a machine without the package defines
% both forms (the star swallowed with its argument).
\makeatletter
\IfFileExists{needspace.sty}{\usepackage{needspace}}{%
  \providecommand{\needspace}[1]{}%
  \providecommand{\Needspace}{\@ifstar\@gobble\@gobble}}
\makeatother
\newcommand{\listingpath}[1]{%
  \par\medskip\Needspace{6\baselineskip}\noindent
  {\ttfamily\footnotesize\color{mgrey}\detokenize{#1}}%
  \par\nobreak}

% Usage: \pyfile{code/ch05/cobweb.py}{Caption}{lst:label}
\newcommand{\pyfile}[3]{%
  \listingpath{#1}%
  \lstinputlisting[language=pythonx,style=bookcode,aboveskip=0.3em,
    caption={#2},label={#3}]{#1}%
}

% A named region of a file, delimited by `# --8<-- [start:name]` and
% `# --8<-- [end:name]`. Matched through listings' range prefix and suffix
% keys -- the other obvious form matches nothing and prints the WHOLE file
% without a warning (measured in the sibling book). A region name is a word,
% never a number. The line numbers printed are the file's own.
% Usage: \pyregion{code/ch05/cobweb.py}{step}{Caption}{lst:label}
\lstset{
  rangebeginprefix=\#\ --8<--\ [start:, rangebeginsuffix=],
  rangeendprefix=\#\ --8<--\ [end:, rangeendsuffix=],
  includerangemarker=false}
\newcommand{\pyregion}[4]{%
  \listingpath{#1}%
  \lstinputlisting[language=pythonx,style=bookcode,aboveskip=0.3em,
    linerange={#2-#2},caption={#3},label={#4}]{#1}%
}

% A program's output, written by code/measure/transcripts.py. There is no
% typed form and there must not be: \val{} does not expand in a verbatim
% body, and a transcript nobody ran looks exactly like one that was.
\newcommand{\transcript}[1]{%
  \par\smallskip
  \IfFileExists{figures/transcripts/#1.txt}{%
    \lstinputlisting[style=bookcode,numbers=none,basicstyle=\ttfamily\footnotesize,
      keywordstyle=\color{black},backgroundcolor=\color{white},
      language={}]{figures/transcripts/#1.txt}%
  }{%
    \begin{tcolorbox}[enhanced, sharp corners, width=\linewidth,
      colback=mlight, colframe=mgrey, boxrule=0.6pt,
      left=8pt, right=8pt, top=6pt, bottom=6pt]
      {\bfseries\color{mgrey}\small \lblTranscriptMissing}\\[2pt]
      {\ttfamily\scriptsize\color{mgrey} figures/transcripts/#1.txt}
    \end{tcolorbox}%
  }%
  \par\smallskip
}

% Inline literals. Underscores inside these must be written \_.
\newcommand{\code}[1]{\texttt{\small #1}}
\newcommand{\pkg}[1]{\texttt{\small\color{mblue}#1}}
\newcommand{\api}[1]{\texttt{\small\color{mteal}#1}}

% ============================================================
%  ADMONITIONS
%  A small, fixed vocabulary. `tinted' may be read in passing; `outlined' is
%  a block the reader must stop at. A third treatment would mean the first two
%  are not carrying their meaning.
% ============================================================
\tcbset{
  cfz tinted/.style={
    enhanced, breakable, sharp corners,
    lines before break=4,
    colback=mlight, boxrule=0pt, leftrule=3pt,
    left=8pt, right=8pt, top=6pt, bottom=6pt,
    before skip=13pt, after skip=13pt,
    fonttitle=\bfseries\small,
    coltitle=black, colbacktitle=mlight,
    attach title to upper={\par\smallskip},
  },
  cfz outlined/.style={
    enhanced, breakable, sharp corners,
    % A box broken straight after its title is a heading stranded at the foot
    % of a page with its question overleaf. Four lines is title plus three.
    lines before break=4,
    colback=white, boxrule=0.6pt,
    left=8pt, right=8pt, top=6pt, bottom=6pt,
    before skip=13pt, after skip=13pt,
    fonttitle=\bfseries\small,
    coltitle=black, colbacktitle=white,
    attach title to upper={\par\smallskip},
  },
}

% Read in passing.
\newtcolorbox{note}[1][]{cfz tinted, colframe=mblue, title={\lblNote}, #1}
\newtcolorbox{warning}[1][]{cfz tinted, colframe=mred, title={\lblWarning}, #1}
% Where you meet this in the world. If it cannot name a specific system,
% instrument, paper or event, it is prose and not this box.
\newtcolorbox{worldbox}[1][]{cfz tinted, colframe=mteal, title={\lblWorld}, #1}
% What is stated and not proved here, and where the proof lives.
\newtcolorbox{rigourbox}[1][]{cfz tinted, colframe=mgrey, title={\lblRigour}, #1}
% The companion library the reader builds, under code/src/chaoslab/.
\newtcolorbox{projectbox}[1][]{cfz tinted, colframe=mamber, title={\lblProject}, #1}

% Stop and read.
% The misconception, named AFTER the reader has been asked the question it
% gets wrong. Appendix B is the catalogue; a trap box is one entry in place.
\newtcolorbox{trapbox}[1][]{cfz outlined, colframe=mred, title={\lblTrap}, #1}
% A listing that has not been run on the pinned versions. Counted by
% `make debt'; nothing ships to a reader with one attached.
\newtcolorbox{verifybox}[1][]{cfz outlined, colframe=mgrey, title={\lblVerify}, #1}
\newtcolorbox{labbox}[1][]{cfz outlined, colframe=mgreen, title={\lblLab}, #1}

% ============================================================
%  THE METHOD --- questions before answers
%  ---------------------------------------------------------
%  The unit of teaching is a question the reader answers BEFORE reading the
%  answer. That is retrieval practice and errorful generation, the two
%  findings the programmed-learning format is defended by, and it is the one
%  thing the sibling mathematics book does that this one keeps whole.
%
%  \begin{think} ... \end{think}   the question, numbered per chapter, ending
%                                  with the instruction to answer before
%                                  reading on;
%  \reveal{short answer}           a centred box with the answer, which is the
%                                  first thing after the question -- the edge
%                                  the reader covers with a hand, or scrolls
%                                  to;
%  \begin{revealblock} ... \end{revealblock}  the same, for a worked answer.
%
%  tools/check_structure.py --think fails when a think box is not followed by
%  a reveal before the next think box or section, and when a reveal answers
%  nothing. That is a property of the source, so it is checked mechanically.
%
%  The instruction says "before you read on", never "before you turn over":
%  the two editions paginate differently, so the second is true in one of them.
% ============================================================
\newcounter{think}[chapter]
\renewcommand{\thethink}{\thechapter.\arabic{think}}
\newtcolorbox{thinkbox}[1][]{cfz outlined, colframe=mviolet, #1}
\NewDocumentEnvironment{think}{}{%
  \refstepcounter{think}%
  \begin{thinkbox}[title={\lblThink~\thethink}]%
}{%
  \par\smallskip
  {\raggedleft\itshape\small\color{mviolet}\lblThinkCue\par}%
  \end{thinkbox}%
}
\tcbset{
  cfz answer/.style={
    enhanced, sharp corners,
    colback=mfaint, colframe=mviolet!55, boxrule=0.5pt,
    before skip=4pt, after skip=10pt,
  },
}
\newcommand{\reveal}[1]{%
  \par\nobreak
  \begin{tcolorbox}[cfz answer, unbreakable,
    left=8pt, right=8pt, top=5pt, bottom=5pt,
    width=0.86\linewidth, center, halign=flush center]
    {\small\bfseries\color{mviolet}\lblAnswerMark}\enspace #1
  \end{tcolorbox}%
  \nobreak\noindent\ignorespaces
}
\NewDocumentEnvironment{revealblock}{}{%
  \par\nopagebreak
  \begin{tcolorbox}[cfz answer, breakable,
    left=10pt, right=10pt, top=6pt, bottom=6pt]
  {\small\bfseries\color{mviolet}\lblAnswerMark}\par\smallskip
}{%
  \end{tcolorbox}%
  \nopagebreak\par\noindent\ignorespaces
}

% ---- Outcomes on entry, and "Can you?" generated from them ----------------
% The self-check at the end of a chapter IS the list of outcomes at its
% start, one for one, because it is built from the same tokens: each
% \outcome both prints its item and appends a row to \cfz@canrows. A \loop
% inside a tabularx body does not survive alignment scanning, which is why
% the rows are accumulated at declaration time (the sibling book's finding).
\makeatletter
\newcommand{\cfz@canrows}{}
\newcommand{\cfz@ratebox}{%
  \fbox{\rule{0pt}{0.9ex}\rule{0.9ex}{0pt}}}
\newcommand{\cfz@rates}{%
  {\footnotesize\color{mgrey}%
   1\,\cfz@ratebox\ 2\,\cfz@ratebox\ 3\,\cfz@ratebox\ 4\,\cfz@ratebox\ 5\,\cfz@ratebox}}
\NewDocumentEnvironment{outcomes}{}{%
  \gdef\cfz@canrows{}%
  \begin{tcolorbox}[cfz tinted, colframe=mblue, title={\lblOutcomes}]%
  \lblOutcomesLead
  \begin{itemize}[leftmargin=1.4em, itemsep=2pt, topsep=4pt]%
}{%
  \end{itemize}%
  \end{tcolorbox}%
}
\newcommand{\outcome}[1]{%
  \item #1%
  \gappto\cfz@canrows{#1 & \cfz@rates \\ \addlinespace[3pt]}%
}
\newcommand{\canyou}{%
  \section*{\lblCanYou}%
  \phantomsection\addcontentsline{toc}{section}{\lblCanYou}%
  \noindent\lblCanYouIntro\par\medskip
  \ifdefempty{\cfz@canrows}{%
    \noindent{\color{mred}\lblNoOutcomes}\par}{%
  {\small
  \noindent\begin{tabularx}{\linewidth}{@{}>{\raggedright\arraybackslash}X r@{}}
    \toprule
    \cfz@canrows
    \bottomrule
  \end{tabularx}\par}}%
}
\makeatother

% ---- Summary -------------------------------------------------------------
\newtcolorbox{summarybox}[1][]{cfz tinted, colframe=mblue, title={\lblSummary}, #1}

% ---- Checkpoint questions, answered at the back --------------------------
% Answers are written AT THE QUESTION with \answerto{} and printed in
% Appendix A. They are carried there by a global token list rather than an
% auxiliary file, because an answer contains mathematics and \val{}, and
% material written to an .aux is expanded at shipout, where fragile commands
% break with errors naming neither the answer nor the chapter.
% \include is sequential within a run, so the list is complete by the time
% the appendix is typeset.
\makeatletter
\newcommand{\cfz@answers}{}
\newcommand{\cfz@lastansch}{0}
\NewDocumentEnvironment{checkpoint}{}{%
  \section*{\lblCheckpoint}%
  \phantomsection\addcontentsline{toc}{section}{\lblCheckpoint}%
  \noindent\lblCheckpointIntro\par
  \begin{enumerate}[label=\textbf{\arabic*.}, leftmargin=2em, itemsep=4pt]%
}{%
  \end{enumerate}%
}
\newcommand{\answerto}[1]{%
  \ifnum\value{chapter}=\cfz@lastansch\relax\else
    \xdef\cfz@lastansch{\arabic{chapter}}%
    \xappto\cfz@answers{\noexpand\cfz@anschapter{\thechapter}}%
  \fi
  \xappto\cfz@answers{\noexpand\cfz@ansitem{\arabic{enumi}}}%
  \gappto\cfz@answers{{#1}}%
}
\newcommand{\cfz@anschapter}[1]{%
  \par\medskip\noindent{\bfseries\color{mblue}\chaptername~#1}\par\nobreak\smallskip}
\newcommand{\cfz@ansitem}[2]{%
  \par\noindent\hangindent=2em\hangafter=1
  \makebox[2em][l]{\textbf{#1.}}#2\par}
\newcommand{\answersbody}{%
  \ifdefempty{\cfz@answers}{\noindent{\color{mred}\lblNoAnswers}\par}{%
    {\raggedright\cfz@answers}}%
}
\makeatother

% ============================================================
%  LABS
%  \begin{lab}{key}{title} ... \end{lab}
%
%  A lab is not a paragraph but a FILE the reader opens: a starter under
%  code/labs/chNN/<key>.py that must FAIL its test, a reference solution under
%  code/labs/chNN/solutions/, and test_<key>.py beside it. The box prints the
%  paths and the one command that runs the test. An environment rather than a
%  macro, because a macro argument cannot hold a verbatim fragment.
%  tools/check_structure.py --labs fails when a file is missing or the key's
%  chapter number disagrees with its chapter.
% ============================================================
\newcounter{lab}[chapter]
\renewcommand{\thelab}{\thechapter.\arabic{lab}}
\makeatletter
\newcommand{\labdir}{ch\two@digits{\value{chapter}}}
\newcommand{\listoflabs}{%
  \section*{\lblLabManifest}%
  \phantomsection\addcontentsline{toc}{section}{\lblLabManifest}%
  \begingroup\c@tocdepth=2\relax\@starttoc{lab}\endgroup%
}
\makeatother
\NewDocumentEnvironment{lab}{mm}{%
  \refstepcounter{lab}%
  \addcontentsline{lab}{subsection}{\protect\numberline{\thelab}\texttt{\protect\detokenize{#1}} --- #2}%
  \begin{labbox}[title={\lblLab~\thelab\ \dash{} #2}]%
}{%
  \par\medskip
  {\small\setlength{\parindent}{0pt}%
  \lblLabStarter: \texttt{code/labs/\labdir/\detokenize{#1}.py}\\
  \lblLabTest: \texttt{code/labs/\labdir/test\_\detokenize{#1}.py}\\
  \lblLabRun: \texttt{cd code \&\& uv run pytest -k \detokenize{#1}}\par}%
  \end{labbox}%
}

% ============================================================
%  FIGURES --- plots and diagrams
%  ---------------------------------------------------------
%  \plotfig{key}{caption}{one-line manifest description}
%    A PLOT, generated by code/plots/<chapter>.py from the same library the
%    listings use, into figures/plots/<lang>/<key>.pdf. Per language, because
%    an axis label in the wrong language is the half-translation that makes a
%    bilingual book look machine-made, and because the Polish edition owes a
%    decimal comma on its tick labels too. Build output, gitignored.
%
%  \mermaidfig{key}{caption}{one-line manifest description}
%    A DIAGRAM of an idea rather than of data, from
%    figures/mermaid/<lang>/<key>.mmd, rendered by `make diagrams'.
%
%  If the rendered PDF is absent either macro prints a visible marker in the
%  flow rather than a float, so a missing figure cannot be mistaken for a
%  layout choice. \phantomsection before \addcontentsline is load-bearing:
%  without it the manifest entry records whatever anchor preceded the figure,
%  which after a listing is a line that was never printed.
% ============================================================
\makeatletter
\newcommand{\listofplots}{%
  \section*{\lblPlotManifest}%
  \phantomsection\addcontentsline{toc}{section}{\lblPlotManifest}%
  \begingroup\c@tocdepth=2\relax\@starttoc{plt}\endgroup%
}
\newcommand{\listofdiagrams}{%
  \section*{\lblDiagramManifest}%
  \phantomsection\addcontentsline{toc}{section}{\lblDiagramManifest}%
  \begingroup\c@tocdepth=2\relax\@starttoc{dgm}\endgroup%
}
\makeatother

\newcommand{\figmissing}[3]{% #1 message, #2 path, #3 caption
  \par\medskip
  \begin{tcolorbox}[enhanced, breakable, width=\linewidth, sharp corners,
    colback=white, colframe=mgrey, boxrule=0.8pt,
    left=8pt, right=8pt, top=8pt, bottom=8pt]
    {\bfseries\color{mgrey}\small #1}\\[4pt]
    {\ttfamily\scriptsize\color{mgrey} #2}
  \end{tcolorbox}%
  \captionof{figure}{#3}%
  \par\medskip
}

\newcommand{\plotfig}[3]{%
  \phantomsection%
  \addcontentsline{plt}{subsection}{\protect\numberline{}\texttt{#1} --- #3}%
  \IfFileExists{figures/plots/\booklang/#1.pdf}{%
    \begin{figure}[htbp]
    \centering
    \includegraphics[width=\linewidth,height=0.45\textheight,keepaspectratio]{figures/plots/\booklang/#1.pdf}%
    \caption{#2}\label{fig:#1}
    \end{figure}%
  }{%
    \figmissing{\lblPlotNotRendered}{figures/plots/\booklang/#1.pdf}{#2}\label{fig:#1}%
  }%
}

\newcommand{\mermaidfig}[3]{%
  \phantomsection%
  \addcontentsline{dgm}{subsection}{\protect\numberline{}\texttt{#1.mmd} --- #3}%
  \IfFileExists{figures/diagrams/\booklang/#1.pdf}{%
    \begin{figure}[htbp]
    \centering
    \includegraphics[width=0.95\linewidth,height=0.42\textheight,keepaspectratio]{figures/diagrams/\booklang/#1.pdf}%
    \caption{#2}\label{fig:#1}
    \end{figure}%
  }{%
    \figmissing{\lblNotRendered}{figures/mermaid/\booklang/#1.mmd}{#2}\label{fig:#1}%
  }%
}

% ============================================================
%  APPENDIX B: the misreadings, generated by tools/gen_traps.py
%
%  \trapchapter{Chapter~\ref{ch:x}}{title}    a heading per chapter
%  \trapentry{n}{habit}{truth}{\ref{sec:x}}  one misreading
%
%  Set ragged right: an entry is mostly short sentences with \code{} spans in
%  them, and a justified line full of unbreakable runs is an overfull box
%  waiting for the edition with the longer words.
% ============================================================
\newcommand{\trapchapter}[2]{%
  \par\medskip
  {\large\bfseries\color{mblue}#1\quad\normalfont\itshape\color{mgrey}#2\par}%
  \nopagebreak\smallskip}
\newcommand{\trapentry}[4]{%
  \par\noindent
  \begingroup\raggedright\hangindent=2.4em\hangafter=1
  \makebox[2.4em][l]{\textbf{\color{mred}#1}}%
  \textbf{\enquote{#2}}\enspace #3\enspace{\color{mgrey}\S\,#4}\par
  \endgroup\smallskip}

% ============================================================
%  CHAPTER STUBS
% ============================================================
\newcommand{\chapterstub}[1]{%
  \begin{tcolorbox}[
    enhanced, breakable, width=\linewidth, sharp corners,
    colback=mlight, colframe=mred, boxrule=1pt,
    borderline={0.6pt}{3pt}{mred, dashed},
    left=12pt, right=12pt, top=12pt, bottom=12pt]
    {\bfseries\color{mred}\large \lblNotWritten}\\[8pt]
    \small\raggedright #1
  \end{tcolorbox}%
}

% ============================================================
%  COMPUTED VALUES
%  A number the reader cannot do in their head is computed by a script under
%  code/measure/, written to figures/values/<chapter>.tex as
%      \pyval{ch06.steps.apart}{34}
%  and reaches the page as \val{ch06.steps.apart}. The book contains
%  references, not digits; `make verify' fails when a committed value no
%  longer matches its script. \val applies the edition's decimal marker, which
%  is how the Polish edition gets its comma without anybody remembering.
% ============================================================
% The comma is BRACED: siunitx sets numbers in maths mode, where a bare comma
% is punctuation and gets a thin space after it ("9, 81"). {,} is an ordinary
% symbol, which is what a decimal marker is.
\newcommand{\bookdecimalmarker}{\ifpl{{,}}{.}}
\IfFileExists{siunitx.sty}{%
  \usepackage{siunitx}%
  \sisetup{
    group-digits = integer,
    group-minimum-digits = 5,
    group-separator = {\,},
    output-decimal-marker = {\bookdecimalmarker},
    exponent-product = \times,
    print-unity-mantissa = false
  }%
  \newcommand{\booknum}[1]{\num{##1}}%
}{%
  \newcommand{\booknum}[1]{##1}%
  \providecommand{\num}[1]{##1}%
}

\makeatletter
\newcommand{\pyvalues}{}
\newcommand{\pyval}[2]{\csgdef{pyval@#1}{#2}\listxadd{\pyvalues}{#1}}
\newcommand{\pyvaltext}[2]{%
  \csgdef{pyval@#1}{#2}\csgdef{pyvaltext@#1}{}\listxadd{\pyvalues}{#1}}
% \num on a non-numeric body is a FATAL siunitx error, so \val refuses a text
% value by name rather than letting siunitx fail on it.
\newcommand{\val}[1]{%
  \ifcsdef{pyval@#1}{%
    \ifcsdef{pyvaltext@#1}{%
      \PackageError{chaosbook}{Value `#1' is not a number}%
        {Use \string\valtext{#1} instead.}%
    }{}%
    \booknum{\csuse{pyval@#1}}%
  }{\textcolor{mred}{\textbf{[\lblValueMissing: #1]}}}%
}
\newcommand{\valtext}[1]{%
  \ifcsdef{pyval@#1}{\csuse{pyval@#1}}%
    {\textcolor{mred}{\textbf{[\lblValueMissing: #1]}}}%
}
\makeatother

% figures/values/all.tex is generated by `make numbers' and \inputs each
% value file. A build with no values still compiles; every \val{} then prints
% a visible marker. \valuesindex lets the one-chapter build point elsewhere.
\providecommand{\valuesindex}{figures/values/all}
\IfFileExists{\valuesindex.tex}{\input{\valuesindex}}{%
  \AtBeginDocument{\PackageWarning{chaosbook}{No computed values found. Run `make numbers'.}}}

% ============================================================
%  PINNED VERSIONS --- single source of truth
%  Change these here AND in code/pyproject.toml; tools/check_structure.py
%  --pins fails when the two disagree. A chapter never types a version.
% ============================================================
\newcommand{\bookedition}{0.1}
\newcommand{\bookyear}{2026}
\newcommand{\pyver}{3.14}              % the minor the book targets
\newcommand{\pypatch}{3.14.7}          % the exact interpreter the code ran on
\newcommand{\uvver}{0.12.13}           % uv
\newcommand{\numpyver}{2.5.3}          % numpy
\newcommand{\matplotlibver}{3.11.2}    % matplotlib
\newcommand{\pytestver}{9.1.1}         % pytest
\newcommand{\ruffver}{0.16.9}          % ruff
% The prose licence is a name with a version in it, not a quantity: it is a
% macro so that neither edition typesets it as a decimal to be localised.
\newcommand{\proselicence}{CC~BY-NC-SA~4.0}

% ============================================================
%  INDEX
% ============================================================
\apptocmd{\theindex}{\raggedcolumns\raggedright}{}{%
  \PackageWarning{chaosbook}{Could not patch theindex for ragged setting}}
\providecommand*\see[2]{}
\renewcommand*\see[2]{\emph{\lblIndexSee} #1}
\providecommand*\seealso[2]{}
\renewcommand*\seealso[2]{\emph{\lblIndexSeeAlso} #1}

\makeindex
; nothing else differs between them at this level.
%  Every user-visible string lives in lang/en.tex or lang/pl.tex, so a label
%  that appears in one edition and not the other is a missing macro rather
%  than a silent divergence.
%
%  Lineage. This book is the third in a family and takes one thing from each
%  of the other two:
%
%    * from Mathematics from Zero for the AI Engineer: the bilingual single
%      source (\booklang, lang/, one body.tex), the rule that every number is
%      computed and reaches the page as \val{}, and the programmed-learning
%      method -- outcomes on entry, a question before every answer, a
%      misconception elicited before it is named, a self-check generated from
%      the outcomes, and checkpoint answers collected at the back;
%    * from Python for .NET Engineers: every listing is a FILE that CI runs,
%      every lab is a starter that must fail and a solution that must pass,
%      and a console transcript is a computed number that happens to be
%      verbatim.
%
%  What is new here is \plotfig. A book about chaos cannot be taught in boxes
%  and arrows: the bifurcation diagram and the Lorenz attractor ARE the
%  argument, so the plots are generated by code/plots/ from the same library
%  the listings use, per language, and never drawn by hand.
% ============================================================

% \booklang must be set by the main file. Default to English rather than
% failing, so a stray % ============================================================
%  preamble.tex --- styling, environments, macros
%
%  Shared by BOTH editions. main-en.tex and main-pl.tex each define \booklang
%  before \input{preamble}; nothing else differs between them at this level.
%  Every user-visible string lives in lang/en.tex or lang/pl.tex, so a label
%  that appears in one edition and not the other is a missing macro rather
%  than a silent divergence.
%
%  Lineage. This book is the third in a family and takes one thing from each
%  of the other two:
%
%    * from Mathematics from Zero for the AI Engineer: the bilingual single
%      source (\booklang, lang/, one body.tex), the rule that every number is
%      computed and reaches the page as \val{}, and the programmed-learning
%      method -- outcomes on entry, a question before every answer, a
%      misconception elicited before it is named, a self-check generated from
%      the outcomes, and checkpoint answers collected at the back;
%    * from Python for .NET Engineers: every listing is a FILE that CI runs,
%      every lab is a starter that must fail and a solution that must pass,
%      and a console transcript is a computed number that happens to be
%      verbatim.
%
%  What is new here is \plotfig. A book about chaos cannot be taught in boxes
%  and arrows: the bifurcation diagram and the Lorenz attractor ARE the
%  argument, so the plots are generated by code/plots/ from the same library
%  the listings use, per language, and never drawn by hand.
% ============================================================

% \booklang must be set by the main file. Default to English rather than
% failing, so a stray \input{preamble} still compiles and says what it did.
\providecommand{\booklang}{en}

\usepackage{etoolbox}   % loaded first: the language switch below needs it

% ---------- Page geometry ----------
% ONE format: A4, 12pt, ONE-SIDED. The book is read on a screen, usually with
% an editor open beside it, so there is no recto, no verso, no gutter and no
% blank leaf before a chapter. The measure is about seventy-five characters
% of prose. At \footnotesize a listing line of 79 characters fits without
% wrapping, and 79 is the width every file under code/ is held to by
% tools/check_structure.py --lines.
\usepackage[
  a4paper,
  left=3.1cm, right=3.1cm, top=2.6cm, bottom=3.0cm,
  head=26pt, headsep=0.8cm, footskip=1.4cm
]{geometry}

% ---------- Encoding & fonts ----------
\usepackage[T1]{fontenc}
\usepackage[utf8]{inputenc}
\IfFileExists{lmodern.sty}{\usepackage{lmodern}}{}
% newtxtext takes its TS1 symbols from TeX Gyre Termes (Debian: tex-gyre).
% A machine with newtx and without tex-gyre dies on the copyright page. CI's
% full TeX Live is the reference; locally install tex-gyre as well.
\IfFileExists{newtxtext.sty}{\usepackage{newtxtext}}{}
% amsmath before newtxmath; amssymb only when newtxmath is absent, because
% loading both is a fatal clash (\Bbbk already defined) that is invisible on
% a machine without newtx.
\usepackage{amsmath}
\IfFileExists{newtxmath.sty}{\usepackage{newtxmath}}{\usepackage{amssymb}}
\IfFileExists{inconsolata.sty}{\usepackage[varqu,scaled=0.95]{inconsolata}}{}
% Without lmodern, microtype runs with font expansion off, which is what
% resolves a marginal overfull line -- so a container short of lmodern
% disagrees with CI about line breaks. Install lmodern locally too.
\IfFileExists{lmodern.sty}{\usepackage{microtype}}{\usepackage[expansion=false]{microtype}}

% upquote: in T1 the input ' is a right single quote, and a listing that sets
% f('x') with curly quotes is a SyntaxError when typed back in.
\IfFileExists{upquote.sty}{\usepackage{upquote}}{}
% And its twin: in T1 the pair -- is an en-dash ligature in the typewriter
% family too, so \code{--flag} would print a flag nobody can type. Scoped to
% tt*, so prose keeps its own dashes.
\DisableLigatures[-]{encoding = T1, family = tt*}

% ---------- Language ----------
% babel with a missing language file is FATAL, so each .ldf is probed first.
\IfFileExists{babel.sty}{%
  \IfFileExists{polish.ldf}{%
    \IfFileExists{british.ldf}{%
      \ifdefstring{\booklang}{pl}%
        {\usepackage[british,main=polish]{babel}}%
        {\usepackage[polish,main=british]{babel}}%
    }{\usepackage[polish]{babel}}%
  }{%
    \IfFileExists{british.ldf}{\usepackage[british]{babel}}{}%
  }%
}{}
% babel-polish makes " an active shorthand, which is a category-code accident
% waiting to happen inside listings. Polish quotation marks come from
% \enquote, which csquotes sets per language.
\ifdefstring{\booklang}{pl}{\AtBeginDocument{\shorthandoff{"}}}{}

% ---------- Core packages ----------
\usepackage{graphicx}
\usepackage{xcolor}
\usepackage{booktabs}
\usepackage{tabularx}
\usepackage{longtable}
\usepackage{array}
\usepackage{enumitem}
\usepackage{listings}
\usepackage[most]{tcolorbox}
\usepackage{fancyhdr}
\usepackage{titlesec}
\usepackage{caption}
\usepackage{imakeidx}
\IfFileExists{csquotes.sty}{\usepackage[autostyle=true]{csquotes}}{%
  \providecommand{\enquote}[1]{``##1''}}
\usepackage[hidelinks]{hyperref}

% ---------- Palette ----------
% The family's palette, so the three books read as a set. Two colours are
% this book's own: the lab (green) and the stop-and-think box (violet).
\definecolor{mblue}{HTML}{1F4E79}
\definecolor{mteal}{HTML}{0E7C7B}
\definecolor{mamber}{HTML}{B26A00}
\definecolor{mred}{HTML}{9B2C2C}
\definecolor{mviolet}{HTML}{7B4397}
\definecolor{mgreen}{HTML}{3B7A3B}
\definecolor{mgrey}{HTML}{5A5A5A}
\definecolor{mlight}{HTML}{F4F6F8}
\definecolor{mfaint}{HTML}{FBFCFD}
\definecolor{codebg}{HTML}{FAFAFA}
\definecolor{codeframe}{HTML}{D8DEE4}
\definecolor{kw}{HTML}{0B5394}
\definecolor{str}{HTML}{9B2C2C}
\definecolor{cmt}{HTML}{4F7A4F}
\definecolor{typ}{HTML}{0E7C7B}
\definecolor{dec}{HTML}{7B4397}

\hypersetup{
  colorlinks=true,
  linkcolor=mblue,
  urlcolor=mteal,
  citecolor=mblue,
  pdfauthor={Konrad Cinkusz}
}

% \ifpl{Polish}{English} -- for the handful of places where a string is
% structural rather than content. Robust, because a fragile \ifpl inside a
% moving argument fails on the next run in exactly one of the two editions.
\DeclareRobustCommand{\ifpl}[2]{\ifdefstring{\booklang}{pl}{#1}{#2}}
\makeatletter
\pdfstringdefDisableCommands{%
  \ifdefstring{\booklang}{pl}{\let\ifpl\@firstoftwo}{\let\ifpl\@secondoftwo}%
  % hyperref drops \color from a bookmark and keeps its ARGUMENT, so a title
  % carrying \code{} would bookmark as "mtealx". Gobble the colour name.
  \def\color#1{}%
}
\makeatother

% ---------- Language strings ----------
\input{lang/\booklang}

% The PDF's subject, keywords and language read \lblPdf... from the language
% file, so they cannot sit in the colour block above, which runs first. Two
% blocks is the ordering, not an oversight.
\hypersetup{
  pdfsubject={\lblPdfSubject},
  pdfkeywords={\lblPdfKeywords},
  pdflang={\lblPdfLang}
}

% ---------- Headings ----------
\titleformat{\chapter}[display]
  {\normalfont\huge\bfseries\color{mblue}\raggedright}
  {\normalfont\normalsize\color{mgrey}\MakeUppercase{\chaptertitlename}\ \thechapter}
  {8pt}{\Huge\raggedright}
\titlespacing*{\chapter}{0pt}{10pt}{28pt}
\titleformat{\section}{\normalfont\Large\bfseries\color{mblue}\raggedright}{\thesection}{0.8em}{}
\titleformat{\subsection}{\normalfont\large\bfseries\color{mgrey}\raggedright}{\thesubsection}{0.7em}{}

\setcounter{tocdepth}{1}
\setcounter{secnumdepth}{2}

\pagestyle{fancy}
\fancyhf{}
\fancyhead[L]{\small\nouppercase{\leftmark}}
\fancyhead[R]{\small\thepage}
\renewcommand{\headrulewidth}{0.4pt}
\fancypagestyle{plain}{\fancyhf{}\fancyfoot[R]{\small\thepage}%
  \renewcommand{\headrulewidth}{0pt}}
% After \pagestyle{fancy}, which installs its own \chaptermark.
\renewcommand{\chaptermark}[1]{\markboth{\thechapter.\ #1}{}}

% ---------- Mathematics ----------
% The handful of notations the book uses, as macros so both editions set the
% same symbol. The book writes \ln for the natural logarithm and \log_{10}
% where the base is ten; a bare \log is ambiguous between the two readerships
% and is not used.
\newcommand{\R}{\mathbb{R}}
\newcommand{\dd}{\,\mathrm{d}}
\newcommand{\abs}[1]{\left\lvert #1\right\rvert}
\newcommand{\iu}{\mathrm{i}}

% ============================================================
%  CODE LISTINGS
%  Every listing in a chapter is a FILE under code/, pulled in with \pyfile or
%  \pyregion, and CI runs it on the pinned interpreter. An in-source
%  \begin{python} block is allowed for a fragment of three or four lines that
%  shows a shape rather than runs.
% ============================================================
\lstdefinelanguage{pythonx}{
  language=Python,
  morekeywords={async,await,match,case,nonlocal,yield,type},
  morekeywords=[2]{
    orbit,iterate,logistic,rk4_step,euler_step,integrate,lorenz,henon,
    lyapunov,separation,escape_time,box_count
  },
  sensitive=true
}

\lstdefinestyle{bookcode}{
  basicstyle=\ttfamily\footnotesize,
  keywordstyle=\color{kw}\bfseries,
  keywordstyle=[2]\color{typ},
  stringstyle=\color{str},
  % Colour only: inconsolata has no italic, and an italic comment style is
  % silently substituted and logged on every build.
  commentstyle=\color{cmt},
  numbers=left,
  numberstyle=\ttfamily\tiny\color{mgrey},
  numbersep=8pt,
  backgroundcolor=\color{codebg},
  frame=single,
  rulecolor=\color{codeframe},
  framesep=6pt,
  breaklines=true,
  breakatwhitespace=false,
  postbreak=\mbox{\textcolor{mgrey}{$\hookrightarrow$}\space},
  showstringspaces=false,
  tabsize=4,
  columns=fullflexible,
  keepspaces=true,
  captionpos=b,
  aboveskip=1.2em,
  belowskip=1.2em,
  % Do NOT add a mapping for U+00A0: it aborts the run with a message that
  % names neither the character nor this line. Listings are ASCII, and
  % parity's C13 checks it.
  literate={ł}{{\l}}1 {Ł}{{\L}}1 {ą}{{\k a}}1 {ę}{{\k e}}1
           {ś}{{\'s}}1 {ć}{{\'c}}1 {ż}{{\.z}}1 {ź}{{\'z}}1 {ó}{{\'o}}1 {ń}{{\'n}}1
           {Ą}{{\k A}}1 {Ę}{{\k E}}1 {Ś}{{\'S}}1 {Ć}{{\'C}}1
           {Ż}{{\.Z}}1 {Ź}{{\'Z}}1 {Ó}{{\'O}}1 {Ń}{{\'N}}1
           {—}{{-{}-}}2 {–}{{-}}1 {‘}{{'}}1 {’}{{'}}1
           {“}{{"}}1 {”}{{"}}1 {°}{{\textdegree}}1
}
\lstset{style=bookcode}

\lstnewenvironment{python}[1][]
  {\lstset{language=pythonx,style=bookcode,#1}}
  {}
\lstnewenvironment{shellcmd}[1][]
  {\lstset{language=bash,style=bookcode,numbers=none,keywordstyle=\color{mteal}\bfseries,#1}}
  {}

% The path is printed above every file listing, in the form the reader opens
% in the editor. \nobreak keeps it on the page of the listing it names.
% The path must not be stranded at the foot of a page with its listing
% overleaf: \nobreak alone does not hold, because listings opens with its own
% vertical material. \needspace asks for room for the path and the first few
% lines of code, and turns the page first when there is not.
% \Needspace (capital N) measures \pagegoal-\pagetotal exactly; lower-case
% \needspace works through glue and penalties and was measured, by the
% Chapter 11 pass, turning pages early and leaving gaps. It must be used
% between paragraphs, which \listingpath's own \par guarantees. Chapters use
% \Needspace* too, so the fallback for a machine without the package defines
% both forms (the star swallowed with its argument).
\makeatletter
\IfFileExists{needspace.sty}{\usepackage{needspace}}{%
  \providecommand{\needspace}[1]{}%
  \providecommand{\Needspace}{\@ifstar\@gobble\@gobble}}
\makeatother
\newcommand{\listingpath}[1]{%
  \par\medskip\Needspace{6\baselineskip}\noindent
  {\ttfamily\footnotesize\color{mgrey}\detokenize{#1}}%
  \par\nobreak}

% Usage: \pyfile{code/ch05/cobweb.py}{Caption}{lst:label}
\newcommand{\pyfile}[3]{%
  \listingpath{#1}%
  \lstinputlisting[language=pythonx,style=bookcode,aboveskip=0.3em,
    caption={#2},label={#3}]{#1}%
}

% A named region of a file, delimited by `# --8<-- [start:name]` and
% `# --8<-- [end:name]`. Matched through listings' range prefix and suffix
% keys -- the other obvious form matches nothing and prints the WHOLE file
% without a warning (measured in the sibling book). A region name is a word,
% never a number. The line numbers printed are the file's own.
% Usage: \pyregion{code/ch05/cobweb.py}{step}{Caption}{lst:label}
\lstset{
  rangebeginprefix=\#\ --8<--\ [start:, rangebeginsuffix=],
  rangeendprefix=\#\ --8<--\ [end:, rangeendsuffix=],
  includerangemarker=false}
\newcommand{\pyregion}[4]{%
  \listingpath{#1}%
  \lstinputlisting[language=pythonx,style=bookcode,aboveskip=0.3em,
    linerange={#2-#2},caption={#3},label={#4}]{#1}%
}

% A program's output, written by code/measure/transcripts.py. There is no
% typed form and there must not be: \val{} does not expand in a verbatim
% body, and a transcript nobody ran looks exactly like one that was.
\newcommand{\transcript}[1]{%
  \par\smallskip
  \IfFileExists{figures/transcripts/#1.txt}{%
    \lstinputlisting[style=bookcode,numbers=none,basicstyle=\ttfamily\footnotesize,
      keywordstyle=\color{black},backgroundcolor=\color{white},
      language={}]{figures/transcripts/#1.txt}%
  }{%
    \begin{tcolorbox}[enhanced, sharp corners, width=\linewidth,
      colback=mlight, colframe=mgrey, boxrule=0.6pt,
      left=8pt, right=8pt, top=6pt, bottom=6pt]
      {\bfseries\color{mgrey}\small \lblTranscriptMissing}\\[2pt]
      {\ttfamily\scriptsize\color{mgrey} figures/transcripts/#1.txt}
    \end{tcolorbox}%
  }%
  \par\smallskip
}

% Inline literals. Underscores inside these must be written \_.
\newcommand{\code}[1]{\texttt{\small #1}}
\newcommand{\pkg}[1]{\texttt{\small\color{mblue}#1}}
\newcommand{\api}[1]{\texttt{\small\color{mteal}#1}}

% ============================================================
%  ADMONITIONS
%  A small, fixed vocabulary. `tinted' may be read in passing; `outlined' is
%  a block the reader must stop at. A third treatment would mean the first two
%  are not carrying their meaning.
% ============================================================
\tcbset{
  cfz tinted/.style={
    enhanced, breakable, sharp corners,
    lines before break=4,
    colback=mlight, boxrule=0pt, leftrule=3pt,
    left=8pt, right=8pt, top=6pt, bottom=6pt,
    before skip=13pt, after skip=13pt,
    fonttitle=\bfseries\small,
    coltitle=black, colbacktitle=mlight,
    attach title to upper={\par\smallskip},
  },
  cfz outlined/.style={
    enhanced, breakable, sharp corners,
    % A box broken straight after its title is a heading stranded at the foot
    % of a page with its question overleaf. Four lines is title plus three.
    lines before break=4,
    colback=white, boxrule=0.6pt,
    left=8pt, right=8pt, top=6pt, bottom=6pt,
    before skip=13pt, after skip=13pt,
    fonttitle=\bfseries\small,
    coltitle=black, colbacktitle=white,
    attach title to upper={\par\smallskip},
  },
}

% Read in passing.
\newtcolorbox{note}[1][]{cfz tinted, colframe=mblue, title={\lblNote}, #1}
\newtcolorbox{warning}[1][]{cfz tinted, colframe=mred, title={\lblWarning}, #1}
% Where you meet this in the world. If it cannot name a specific system,
% instrument, paper or event, it is prose and not this box.
\newtcolorbox{worldbox}[1][]{cfz tinted, colframe=mteal, title={\lblWorld}, #1}
% What is stated and not proved here, and where the proof lives.
\newtcolorbox{rigourbox}[1][]{cfz tinted, colframe=mgrey, title={\lblRigour}, #1}
% The companion library the reader builds, under code/src/chaoslab/.
\newtcolorbox{projectbox}[1][]{cfz tinted, colframe=mamber, title={\lblProject}, #1}

% Stop and read.
% The misconception, named AFTER the reader has been asked the question it
% gets wrong. Appendix B is the catalogue; a trap box is one entry in place.
\newtcolorbox{trapbox}[1][]{cfz outlined, colframe=mred, title={\lblTrap}, #1}
% A listing that has not been run on the pinned versions. Counted by
% `make debt'; nothing ships to a reader with one attached.
\newtcolorbox{verifybox}[1][]{cfz outlined, colframe=mgrey, title={\lblVerify}, #1}
\newtcolorbox{labbox}[1][]{cfz outlined, colframe=mgreen, title={\lblLab}, #1}

% ============================================================
%  THE METHOD --- questions before answers
%  ---------------------------------------------------------
%  The unit of teaching is a question the reader answers BEFORE reading the
%  answer. That is retrieval practice and errorful generation, the two
%  findings the programmed-learning format is defended by, and it is the one
%  thing the sibling mathematics book does that this one keeps whole.
%
%  \begin{think} ... \end{think}   the question, numbered per chapter, ending
%                                  with the instruction to answer before
%                                  reading on;
%  \reveal{short answer}           a centred box with the answer, which is the
%                                  first thing after the question -- the edge
%                                  the reader covers with a hand, or scrolls
%                                  to;
%  \begin{revealblock} ... \end{revealblock}  the same, for a worked answer.
%
%  tools/check_structure.py --think fails when a think box is not followed by
%  a reveal before the next think box or section, and when a reveal answers
%  nothing. That is a property of the source, so it is checked mechanically.
%
%  The instruction says "before you read on", never "before you turn over":
%  the two editions paginate differently, so the second is true in one of them.
% ============================================================
\newcounter{think}[chapter]
\renewcommand{\thethink}{\thechapter.\arabic{think}}
\newtcolorbox{thinkbox}[1][]{cfz outlined, colframe=mviolet, #1}
\NewDocumentEnvironment{think}{}{%
  \refstepcounter{think}%
  \begin{thinkbox}[title={\lblThink~\thethink}]%
}{%
  \par\smallskip
  {\raggedleft\itshape\small\color{mviolet}\lblThinkCue\par}%
  \end{thinkbox}%
}
\tcbset{
  cfz answer/.style={
    enhanced, sharp corners,
    colback=mfaint, colframe=mviolet!55, boxrule=0.5pt,
    before skip=4pt, after skip=10pt,
  },
}
\newcommand{\reveal}[1]{%
  \par\nobreak
  \begin{tcolorbox}[cfz answer, unbreakable,
    left=8pt, right=8pt, top=5pt, bottom=5pt,
    width=0.86\linewidth, center, halign=flush center]
    {\small\bfseries\color{mviolet}\lblAnswerMark}\enspace #1
  \end{tcolorbox}%
  \nobreak\noindent\ignorespaces
}
\NewDocumentEnvironment{revealblock}{}{%
  \par\nopagebreak
  \begin{tcolorbox}[cfz answer, breakable,
    left=10pt, right=10pt, top=6pt, bottom=6pt]
  {\small\bfseries\color{mviolet}\lblAnswerMark}\par\smallskip
}{%
  \end{tcolorbox}%
  \nopagebreak\par\noindent\ignorespaces
}

% ---- Outcomes on entry, and "Can you?" generated from them ----------------
% The self-check at the end of a chapter IS the list of outcomes at its
% start, one for one, because it is built from the same tokens: each
% \outcome both prints its item and appends a row to \cfz@canrows. A \loop
% inside a tabularx body does not survive alignment scanning, which is why
% the rows are accumulated at declaration time (the sibling book's finding).
\makeatletter
\newcommand{\cfz@canrows}{}
\newcommand{\cfz@ratebox}{%
  \fbox{\rule{0pt}{0.9ex}\rule{0.9ex}{0pt}}}
\newcommand{\cfz@rates}{%
  {\footnotesize\color{mgrey}%
   1\,\cfz@ratebox\ 2\,\cfz@ratebox\ 3\,\cfz@ratebox\ 4\,\cfz@ratebox\ 5\,\cfz@ratebox}}
\NewDocumentEnvironment{outcomes}{}{%
  \gdef\cfz@canrows{}%
  \begin{tcolorbox}[cfz tinted, colframe=mblue, title={\lblOutcomes}]%
  \lblOutcomesLead
  \begin{itemize}[leftmargin=1.4em, itemsep=2pt, topsep=4pt]%
}{%
  \end{itemize}%
  \end{tcolorbox}%
}
\newcommand{\outcome}[1]{%
  \item #1%
  \gappto\cfz@canrows{#1 & \cfz@rates \\ \addlinespace[3pt]}%
}
\newcommand{\canyou}{%
  \section*{\lblCanYou}%
  \phantomsection\addcontentsline{toc}{section}{\lblCanYou}%
  \noindent\lblCanYouIntro\par\medskip
  \ifdefempty{\cfz@canrows}{%
    \noindent{\color{mred}\lblNoOutcomes}\par}{%
  {\small
  \noindent\begin{tabularx}{\linewidth}{@{}>{\raggedright\arraybackslash}X r@{}}
    \toprule
    \cfz@canrows
    \bottomrule
  \end{tabularx}\par}}%
}
\makeatother

% ---- Summary -------------------------------------------------------------
\newtcolorbox{summarybox}[1][]{cfz tinted, colframe=mblue, title={\lblSummary}, #1}

% ---- Checkpoint questions, answered at the back --------------------------
% Answers are written AT THE QUESTION with \answerto{} and printed in
% Appendix A. They are carried there by a global token list rather than an
% auxiliary file, because an answer contains mathematics and \val{}, and
% material written to an .aux is expanded at shipout, where fragile commands
% break with errors naming neither the answer nor the chapter.
% \include is sequential within a run, so the list is complete by the time
% the appendix is typeset.
\makeatletter
\newcommand{\cfz@answers}{}
\newcommand{\cfz@lastansch}{0}
\NewDocumentEnvironment{checkpoint}{}{%
  \section*{\lblCheckpoint}%
  \phantomsection\addcontentsline{toc}{section}{\lblCheckpoint}%
  \noindent\lblCheckpointIntro\par
  \begin{enumerate}[label=\textbf{\arabic*.}, leftmargin=2em, itemsep=4pt]%
}{%
  \end{enumerate}%
}
\newcommand{\answerto}[1]{%
  \ifnum\value{chapter}=\cfz@lastansch\relax\else
    \xdef\cfz@lastansch{\arabic{chapter}}%
    \xappto\cfz@answers{\noexpand\cfz@anschapter{\thechapter}}%
  \fi
  \xappto\cfz@answers{\noexpand\cfz@ansitem{\arabic{enumi}}}%
  \gappto\cfz@answers{{#1}}%
}
\newcommand{\cfz@anschapter}[1]{%
  \par\medskip\noindent{\bfseries\color{mblue}\chaptername~#1}\par\nobreak\smallskip}
\newcommand{\cfz@ansitem}[2]{%
  \par\noindent\hangindent=2em\hangafter=1
  \makebox[2em][l]{\textbf{#1.}}#2\par}
\newcommand{\answersbody}{%
  \ifdefempty{\cfz@answers}{\noindent{\color{mred}\lblNoAnswers}\par}{%
    {\raggedright\cfz@answers}}%
}
\makeatother

% ============================================================
%  LABS
%  \begin{lab}{key}{title} ... \end{lab}
%
%  A lab is not a paragraph but a FILE the reader opens: a starter under
%  code/labs/chNN/<key>.py that must FAIL its test, a reference solution under
%  code/labs/chNN/solutions/, and test_<key>.py beside it. The box prints the
%  paths and the one command that runs the test. An environment rather than a
%  macro, because a macro argument cannot hold a verbatim fragment.
%  tools/check_structure.py --labs fails when a file is missing or the key's
%  chapter number disagrees with its chapter.
% ============================================================
\newcounter{lab}[chapter]
\renewcommand{\thelab}{\thechapter.\arabic{lab}}
\makeatletter
\newcommand{\labdir}{ch\two@digits{\value{chapter}}}
\newcommand{\listoflabs}{%
  \section*{\lblLabManifest}%
  \phantomsection\addcontentsline{toc}{section}{\lblLabManifest}%
  \begingroup\c@tocdepth=2\relax\@starttoc{lab}\endgroup%
}
\makeatother
\NewDocumentEnvironment{lab}{mm}{%
  \refstepcounter{lab}%
  \addcontentsline{lab}{subsection}{\protect\numberline{\thelab}\texttt{\protect\detokenize{#1}} --- #2}%
  \begin{labbox}[title={\lblLab~\thelab\ \dash{} #2}]%
}{%
  \par\medskip
  {\small\setlength{\parindent}{0pt}%
  \lblLabStarter: \texttt{code/labs/\labdir/\detokenize{#1}.py}\\
  \lblLabTest: \texttt{code/labs/\labdir/test\_\detokenize{#1}.py}\\
  \lblLabRun: \texttt{cd code \&\& uv run pytest -k \detokenize{#1}}\par}%
  \end{labbox}%
}

% ============================================================
%  FIGURES --- plots and diagrams
%  ---------------------------------------------------------
%  \plotfig{key}{caption}{one-line manifest description}
%    A PLOT, generated by code/plots/<chapter>.py from the same library the
%    listings use, into figures/plots/<lang>/<key>.pdf. Per language, because
%    an axis label in the wrong language is the half-translation that makes a
%    bilingual book look machine-made, and because the Polish edition owes a
%    decimal comma on its tick labels too. Build output, gitignored.
%
%  \mermaidfig{key}{caption}{one-line manifest description}
%    A DIAGRAM of an idea rather than of data, from
%    figures/mermaid/<lang>/<key>.mmd, rendered by `make diagrams'.
%
%  If the rendered PDF is absent either macro prints a visible marker in the
%  flow rather than a float, so a missing figure cannot be mistaken for a
%  layout choice. \phantomsection before \addcontentsline is load-bearing:
%  without it the manifest entry records whatever anchor preceded the figure,
%  which after a listing is a line that was never printed.
% ============================================================
\makeatletter
\newcommand{\listofplots}{%
  \section*{\lblPlotManifest}%
  \phantomsection\addcontentsline{toc}{section}{\lblPlotManifest}%
  \begingroup\c@tocdepth=2\relax\@starttoc{plt}\endgroup%
}
\newcommand{\listofdiagrams}{%
  \section*{\lblDiagramManifest}%
  \phantomsection\addcontentsline{toc}{section}{\lblDiagramManifest}%
  \begingroup\c@tocdepth=2\relax\@starttoc{dgm}\endgroup%
}
\makeatother

\newcommand{\figmissing}[3]{% #1 message, #2 path, #3 caption
  \par\medskip
  \begin{tcolorbox}[enhanced, breakable, width=\linewidth, sharp corners,
    colback=white, colframe=mgrey, boxrule=0.8pt,
    left=8pt, right=8pt, top=8pt, bottom=8pt]
    {\bfseries\color{mgrey}\small #1}\\[4pt]
    {\ttfamily\scriptsize\color{mgrey} #2}
  \end{tcolorbox}%
  \captionof{figure}{#3}%
  \par\medskip
}

\newcommand{\plotfig}[3]{%
  \phantomsection%
  \addcontentsline{plt}{subsection}{\protect\numberline{}\texttt{#1} --- #3}%
  \IfFileExists{figures/plots/\booklang/#1.pdf}{%
    \begin{figure}[htbp]
    \centering
    \includegraphics[width=\linewidth,height=0.45\textheight,keepaspectratio]{figures/plots/\booklang/#1.pdf}%
    \caption{#2}\label{fig:#1}
    \end{figure}%
  }{%
    \figmissing{\lblPlotNotRendered}{figures/plots/\booklang/#1.pdf}{#2}\label{fig:#1}%
  }%
}

\newcommand{\mermaidfig}[3]{%
  \phantomsection%
  \addcontentsline{dgm}{subsection}{\protect\numberline{}\texttt{#1.mmd} --- #3}%
  \IfFileExists{figures/diagrams/\booklang/#1.pdf}{%
    \begin{figure}[htbp]
    \centering
    \includegraphics[width=0.95\linewidth,height=0.42\textheight,keepaspectratio]{figures/diagrams/\booklang/#1.pdf}%
    \caption{#2}\label{fig:#1}
    \end{figure}%
  }{%
    \figmissing{\lblNotRendered}{figures/mermaid/\booklang/#1.mmd}{#2}\label{fig:#1}%
  }%
}

% ============================================================
%  APPENDIX B: the misreadings, generated by tools/gen_traps.py
%
%  \trapchapter{Chapter~\ref{ch:x}}{title}    a heading per chapter
%  \trapentry{n}{habit}{truth}{\ref{sec:x}}  one misreading
%
%  Set ragged right: an entry is mostly short sentences with \code{} spans in
%  them, and a justified line full of unbreakable runs is an overfull box
%  waiting for the edition with the longer words.
% ============================================================
\newcommand{\trapchapter}[2]{%
  \par\medskip
  {\large\bfseries\color{mblue}#1\quad\normalfont\itshape\color{mgrey}#2\par}%
  \nopagebreak\smallskip}
\newcommand{\trapentry}[4]{%
  \par\noindent
  \begingroup\raggedright\hangindent=2.4em\hangafter=1
  \makebox[2.4em][l]{\textbf{\color{mred}#1}}%
  \textbf{\enquote{#2}}\enspace #3\enspace{\color{mgrey}\S\,#4}\par
  \endgroup\smallskip}

% ============================================================
%  CHAPTER STUBS
% ============================================================
\newcommand{\chapterstub}[1]{%
  \begin{tcolorbox}[
    enhanced, breakable, width=\linewidth, sharp corners,
    colback=mlight, colframe=mred, boxrule=1pt,
    borderline={0.6pt}{3pt}{mred, dashed},
    left=12pt, right=12pt, top=12pt, bottom=12pt]
    {\bfseries\color{mred}\large \lblNotWritten}\\[8pt]
    \small\raggedright #1
  \end{tcolorbox}%
}

% ============================================================
%  COMPUTED VALUES
%  A number the reader cannot do in their head is computed by a script under
%  code/measure/, written to figures/values/<chapter>.tex as
%      \pyval{ch06.steps.apart}{34}
%  and reaches the page as \val{ch06.steps.apart}. The book contains
%  references, not digits; `make verify' fails when a committed value no
%  longer matches its script. \val applies the edition's decimal marker, which
%  is how the Polish edition gets its comma without anybody remembering.
% ============================================================
% The comma is BRACED: siunitx sets numbers in maths mode, where a bare comma
% is punctuation and gets a thin space after it ("9, 81"). {,} is an ordinary
% symbol, which is what a decimal marker is.
\newcommand{\bookdecimalmarker}{\ifpl{{,}}{.}}
\IfFileExists{siunitx.sty}{%
  \usepackage{siunitx}%
  \sisetup{
    group-digits = integer,
    group-minimum-digits = 5,
    group-separator = {\,},
    output-decimal-marker = {\bookdecimalmarker},
    exponent-product = \times,
    print-unity-mantissa = false
  }%
  \newcommand{\booknum}[1]{\num{##1}}%
}{%
  \newcommand{\booknum}[1]{##1}%
  \providecommand{\num}[1]{##1}%
}

\makeatletter
\newcommand{\pyvalues}{}
\newcommand{\pyval}[2]{\csgdef{pyval@#1}{#2}\listxadd{\pyvalues}{#1}}
\newcommand{\pyvaltext}[2]{%
  \csgdef{pyval@#1}{#2}\csgdef{pyvaltext@#1}{}\listxadd{\pyvalues}{#1}}
% \num on a non-numeric body is a FATAL siunitx error, so \val refuses a text
% value by name rather than letting siunitx fail on it.
\newcommand{\val}[1]{%
  \ifcsdef{pyval@#1}{%
    \ifcsdef{pyvaltext@#1}{%
      \PackageError{chaosbook}{Value `#1' is not a number}%
        {Use \string\valtext{#1} instead.}%
    }{}%
    \booknum{\csuse{pyval@#1}}%
  }{\textcolor{mred}{\textbf{[\lblValueMissing: #1]}}}%
}
\newcommand{\valtext}[1]{%
  \ifcsdef{pyval@#1}{\csuse{pyval@#1}}%
    {\textcolor{mred}{\textbf{[\lblValueMissing: #1]}}}%
}
\makeatother

% figures/values/all.tex is generated by `make numbers' and \inputs each
% value file. A build with no values still compiles; every \val{} then prints
% a visible marker. \valuesindex lets the one-chapter build point elsewhere.
\providecommand{\valuesindex}{figures/values/all}
\IfFileExists{\valuesindex.tex}{\input{\valuesindex}}{%
  \AtBeginDocument{\PackageWarning{chaosbook}{No computed values found. Run `make numbers'.}}}

% ============================================================
%  PINNED VERSIONS --- single source of truth
%  Change these here AND in code/pyproject.toml; tools/check_structure.py
%  --pins fails when the two disagree. A chapter never types a version.
% ============================================================
\newcommand{\bookedition}{0.1}
\newcommand{\bookyear}{2026}
\newcommand{\pyver}{3.14}              % the minor the book targets
\newcommand{\pypatch}{3.14.7}          % the exact interpreter the code ran on
\newcommand{\uvver}{0.12.13}           % uv
\newcommand{\numpyver}{2.5.3}          % numpy
\newcommand{\matplotlibver}{3.11.2}    % matplotlib
\newcommand{\pytestver}{9.1.1}         % pytest
\newcommand{\ruffver}{0.16.9}          % ruff
% The prose licence is a name with a version in it, not a quantity: it is a
% macro so that neither edition typesets it as a decimal to be localised.
\newcommand{\proselicence}{CC~BY-NC-SA~4.0}

% ============================================================
%  INDEX
% ============================================================
\apptocmd{\theindex}{\raggedcolumns\raggedright}{}{%
  \PackageWarning{chaosbook}{Could not patch theindex for ragged setting}}
\providecommand*\see[2]{}
\renewcommand*\see[2]{\emph{\lblIndexSee} #1}
\providecommand*\seealso[2]{}
\renewcommand*\seealso[2]{\emph{\lblIndexSeeAlso} #1}

\makeindex
 still compiles and says what it did.
\providecommand{\booklang}{en}

\usepackage{etoolbox}   % loaded first: the language switch below needs it

% ---------- Page geometry ----------
% ONE format: A4, 12pt, ONE-SIDED. The book is read on a screen, usually with
% an editor open beside it, so there is no recto, no verso, no gutter and no
% blank leaf before a chapter. The measure is about seventy-five characters
% of prose. At \footnotesize a listing line of 79 characters fits without
% wrapping, and 79 is the width every file under code/ is held to by
% tools/check_structure.py --lines.
\usepackage[
  a4paper,
  left=3.1cm, right=3.1cm, top=2.6cm, bottom=3.0cm,
  head=26pt, headsep=0.8cm, footskip=1.4cm
]{geometry}

% ---------- Encoding & fonts ----------
\usepackage[T1]{fontenc}
\usepackage[utf8]{inputenc}
\IfFileExists{lmodern.sty}{\usepackage{lmodern}}{}
% newtxtext takes its TS1 symbols from TeX Gyre Termes (Debian: tex-gyre).
% A machine with newtx and without tex-gyre dies on the copyright page. CI's
% full TeX Live is the reference; locally install tex-gyre as well.
\IfFileExists{newtxtext.sty}{\usepackage{newtxtext}}{}
% amsmath before newtxmath; amssymb only when newtxmath is absent, because
% loading both is a fatal clash (\Bbbk already defined) that is invisible on
% a machine without newtx.
\usepackage{amsmath}
\IfFileExists{newtxmath.sty}{\usepackage{newtxmath}}{\usepackage{amssymb}}
\IfFileExists{inconsolata.sty}{\usepackage[varqu,scaled=0.95]{inconsolata}}{}
% Without lmodern, microtype runs with font expansion off, which is what
% resolves a marginal overfull line -- so a container short of lmodern
% disagrees with CI about line breaks. Install lmodern locally too.
\IfFileExists{lmodern.sty}{\usepackage{microtype}}{\usepackage[expansion=false]{microtype}}

% upquote: in T1 the input ' is a right single quote, and a listing that sets
% f('x') with curly quotes is a SyntaxError when typed back in.
\IfFileExists{upquote.sty}{\usepackage{upquote}}{}
% And its twin: in T1 the pair -- is an en-dash ligature in the typewriter
% family too, so \code{--flag} would print a flag nobody can type. Scoped to
% tt*, so prose keeps its own dashes.
\DisableLigatures[-]{encoding = T1, family = tt*}

% ---------- Language ----------
% babel with a missing language file is FATAL, so each .ldf is probed first.
\IfFileExists{babel.sty}{%
  \IfFileExists{polish.ldf}{%
    \IfFileExists{british.ldf}{%
      \ifdefstring{\booklang}{pl}%
        {\usepackage[british,main=polish]{babel}}%
        {\usepackage[polish,main=british]{babel}}%
    }{\usepackage[polish]{babel}}%
  }{%
    \IfFileExists{british.ldf}{\usepackage[british]{babel}}{}%
  }%
}{}
% babel-polish makes " an active shorthand, which is a category-code accident
% waiting to happen inside listings. Polish quotation marks come from
% \enquote, which csquotes sets per language.
\ifdefstring{\booklang}{pl}{\AtBeginDocument{\shorthandoff{"}}}{}

% ---------- Core packages ----------
\usepackage{graphicx}
\usepackage{xcolor}
\usepackage{booktabs}
\usepackage{tabularx}
\usepackage{longtable}
\usepackage{array}
\usepackage{enumitem}
\usepackage{listings}
\usepackage[most]{tcolorbox}
\usepackage{fancyhdr}
\usepackage{titlesec}
\usepackage{caption}
\usepackage{imakeidx}
\IfFileExists{csquotes.sty}{\usepackage[autostyle=true]{csquotes}}{%
  \providecommand{\enquote}[1]{``##1''}}
\usepackage[hidelinks]{hyperref}

% ---------- Palette ----------
% The family's palette, so the three books read as a set. Two colours are
% this book's own: the lab (green) and the stop-and-think box (violet).
\definecolor{mblue}{HTML}{1F4E79}
\definecolor{mteal}{HTML}{0E7C7B}
\definecolor{mamber}{HTML}{B26A00}
\definecolor{mred}{HTML}{9B2C2C}
\definecolor{mviolet}{HTML}{7B4397}
\definecolor{mgreen}{HTML}{3B7A3B}
\definecolor{mgrey}{HTML}{5A5A5A}
\definecolor{mlight}{HTML}{F4F6F8}
\definecolor{mfaint}{HTML}{FBFCFD}
\definecolor{codebg}{HTML}{FAFAFA}
\definecolor{codeframe}{HTML}{D8DEE4}
\definecolor{kw}{HTML}{0B5394}
\definecolor{str}{HTML}{9B2C2C}
\definecolor{cmt}{HTML}{4F7A4F}
\definecolor{typ}{HTML}{0E7C7B}
\definecolor{dec}{HTML}{7B4397}

\hypersetup{
  colorlinks=true,
  linkcolor=mblue,
  urlcolor=mteal,
  citecolor=mblue,
  pdfauthor={Konrad Cinkusz}
}

% \ifpl{Polish}{English} -- for the handful of places where a string is
% structural rather than content. Robust, because a fragile \ifpl inside a
% moving argument fails on the next run in exactly one of the two editions.
\DeclareRobustCommand{\ifpl}[2]{\ifdefstring{\booklang}{pl}{#1}{#2}}
\makeatletter
\pdfstringdefDisableCommands{%
  \ifdefstring{\booklang}{pl}{\let\ifpl\@firstoftwo}{\let\ifpl\@secondoftwo}%
  % hyperref drops \color from a bookmark and keeps its ARGUMENT, so a title
  % carrying \code{} would bookmark as "mtealx". Gobble the colour name.
  \def\color#1{}%
}
\makeatother

% ---------- Language strings ----------
\input{lang/\booklang}

% The PDF's subject, keywords and language read \lblPdf... from the language
% file, so they cannot sit in the colour block above, which runs first. Two
% blocks is the ordering, not an oversight.
\hypersetup{
  pdfsubject={\lblPdfSubject},
  pdfkeywords={\lblPdfKeywords},
  pdflang={\lblPdfLang}
}

% ---------- Headings ----------
\titleformat{\chapter}[display]
  {\normalfont\huge\bfseries\color{mblue}\raggedright}
  {\normalfont\normalsize\color{mgrey}\MakeUppercase{\chaptertitlename}\ \thechapter}
  {8pt}{\Huge\raggedright}
\titlespacing*{\chapter}{0pt}{10pt}{28pt}
\titleformat{\section}{\normalfont\Large\bfseries\color{mblue}\raggedright}{\thesection}{0.8em}{}
\titleformat{\subsection}{\normalfont\large\bfseries\color{mgrey}\raggedright}{\thesubsection}{0.7em}{}

\setcounter{tocdepth}{1}
\setcounter{secnumdepth}{2}

\pagestyle{fancy}
\fancyhf{}
\fancyhead[L]{\small\nouppercase{\leftmark}}
\fancyhead[R]{\small\thepage}
\renewcommand{\headrulewidth}{0.4pt}
\fancypagestyle{plain}{\fancyhf{}\fancyfoot[R]{\small\thepage}%
  \renewcommand{\headrulewidth}{0pt}}
% After \pagestyle{fancy}, which installs its own \chaptermark.
\renewcommand{\chaptermark}[1]{\markboth{\thechapter.\ #1}{}}

% ---------- Mathematics ----------
% The handful of notations the book uses, as macros so both editions set the
% same symbol. The book writes \ln for the natural logarithm and \log_{10}
% where the base is ten; a bare \log is ambiguous between the two readerships
% and is not used.
\newcommand{\R}{\mathbb{R}}
\newcommand{\dd}{\,\mathrm{d}}
\newcommand{\abs}[1]{\left\lvert #1\right\rvert}
\newcommand{\iu}{\mathrm{i}}

% ============================================================
%  CODE LISTINGS
%  Every listing in a chapter is a FILE under code/, pulled in with \pyfile or
%  \pyregion, and CI runs it on the pinned interpreter. An in-source
%  \begin{python} block is allowed for a fragment of three or four lines that
%  shows a shape rather than runs.
% ============================================================
\lstdefinelanguage{pythonx}{
  language=Python,
  morekeywords={async,await,match,case,nonlocal,yield,type},
  morekeywords=[2]{
    orbit,iterate,logistic,rk4_step,euler_step,integrate,lorenz,henon,
    lyapunov,separation,escape_time,box_count
  },
  sensitive=true
}

\lstdefinestyle{bookcode}{
  basicstyle=\ttfamily\footnotesize,
  keywordstyle=\color{kw}\bfseries,
  keywordstyle=[2]\color{typ},
  stringstyle=\color{str},
  % Colour only: inconsolata has no italic, and an italic comment style is
  % silently substituted and logged on every build.
  commentstyle=\color{cmt},
  numbers=left,
  numberstyle=\ttfamily\tiny\color{mgrey},
  numbersep=8pt,
  backgroundcolor=\color{codebg},
  frame=single,
  rulecolor=\color{codeframe},
  framesep=6pt,
  breaklines=true,
  breakatwhitespace=false,
  postbreak=\mbox{\textcolor{mgrey}{$\hookrightarrow$}\space},
  showstringspaces=false,
  tabsize=4,
  columns=fullflexible,
  keepspaces=true,
  captionpos=b,
  aboveskip=1.2em,
  belowskip=1.2em,
  % Do NOT add a mapping for U+00A0: it aborts the run with a message that
  % names neither the character nor this line. Listings are ASCII, and
  % parity's C13 checks it.
  literate={ł}{{\l}}1 {Ł}{{\L}}1 {ą}{{\k a}}1 {ę}{{\k e}}1
           {ś}{{\'s}}1 {ć}{{\'c}}1 {ż}{{\.z}}1 {ź}{{\'z}}1 {ó}{{\'o}}1 {ń}{{\'n}}1
           {Ą}{{\k A}}1 {Ę}{{\k E}}1 {Ś}{{\'S}}1 {Ć}{{\'C}}1
           {Ż}{{\.Z}}1 {Ź}{{\'Z}}1 {Ó}{{\'O}}1 {Ń}{{\'N}}1
           {—}{{-{}-}}2 {–}{{-}}1 {‘}{{'}}1 {’}{{'}}1
           {“}{{"}}1 {”}{{"}}1 {°}{{\textdegree}}1
}
\lstset{style=bookcode}

\lstnewenvironment{python}[1][]
  {\lstset{language=pythonx,style=bookcode,#1}}
  {}
\lstnewenvironment{shellcmd}[1][]
  {\lstset{language=bash,style=bookcode,numbers=none,keywordstyle=\color{mteal}\bfseries,#1}}
  {}

% The path is printed above every file listing, in the form the reader opens
% in the editor. \nobreak keeps it on the page of the listing it names.
% The path must not be stranded at the foot of a page with its listing
% overleaf: \nobreak alone does not hold, because listings opens with its own
% vertical material. \needspace asks for room for the path and the first few
% lines of code, and turns the page first when there is not.
% \Needspace (capital N) measures \pagegoal-\pagetotal exactly; lower-case
% \needspace works through glue and penalties and was measured, by the
% Chapter 11 pass, turning pages early and leaving gaps. It must be used
% between paragraphs, which \listingpath's own \par guarantees. Chapters use
% \Needspace* too, so the fallback for a machine without the package defines
% both forms (the star swallowed with its argument).
\makeatletter
\IfFileExists{needspace.sty}{\usepackage{needspace}}{%
  \providecommand{\needspace}[1]{}%
  \providecommand{\Needspace}{\@ifstar\@gobble\@gobble}}
\makeatother
\newcommand{\listingpath}[1]{%
  \par\medskip\Needspace{6\baselineskip}\noindent
  {\ttfamily\footnotesize\color{mgrey}\detokenize{#1}}%
  \par\nobreak}

% Usage: \pyfile{code/ch05/cobweb.py}{Caption}{lst:label}
\newcommand{\pyfile}[3]{%
  \listingpath{#1}%
  \lstinputlisting[language=pythonx,style=bookcode,aboveskip=0.3em,
    caption={#2},label={#3}]{#1}%
}

% A named region of a file, delimited by `# --8<-- [start:name]` and
% `# --8<-- [end:name]`. Matched through listings' range prefix and suffix
% keys -- the other obvious form matches nothing and prints the WHOLE file
% without a warning (measured in the sibling book). A region name is a word,
% never a number. The line numbers printed are the file's own.
% Usage: \pyregion{code/ch05/cobweb.py}{step}{Caption}{lst:label}
\lstset{
  rangebeginprefix=\#\ --8<--\ [start:, rangebeginsuffix=],
  rangeendprefix=\#\ --8<--\ [end:, rangeendsuffix=],
  includerangemarker=false}
\newcommand{\pyregion}[4]{%
  \listingpath{#1}%
  \lstinputlisting[language=pythonx,style=bookcode,aboveskip=0.3em,
    linerange={#2-#2},caption={#3},label={#4}]{#1}%
}

% A program's output, written by code/measure/transcripts.py. There is no
% typed form and there must not be: \val{} does not expand in a verbatim
% body, and a transcript nobody ran looks exactly like one that was.
\newcommand{\transcript}[1]{%
  \par\smallskip
  \IfFileExists{figures/transcripts/#1.txt}{%
    \lstinputlisting[style=bookcode,numbers=none,basicstyle=\ttfamily\footnotesize,
      keywordstyle=\color{black},backgroundcolor=\color{white},
      language={}]{figures/transcripts/#1.txt}%
  }{%
    \begin{tcolorbox}[enhanced, sharp corners, width=\linewidth,
      colback=mlight, colframe=mgrey, boxrule=0.6pt,
      left=8pt, right=8pt, top=6pt, bottom=6pt]
      {\bfseries\color{mgrey}\small \lblTranscriptMissing}\\[2pt]
      {\ttfamily\scriptsize\color{mgrey} figures/transcripts/#1.txt}
    \end{tcolorbox}%
  }%
  \par\smallskip
}

% Inline literals. Underscores inside these must be written \_.
\newcommand{\code}[1]{\texttt{\small #1}}
\newcommand{\pkg}[1]{\texttt{\small\color{mblue}#1}}
\newcommand{\api}[1]{\texttt{\small\color{mteal}#1}}

% ============================================================
%  ADMONITIONS
%  A small, fixed vocabulary. `tinted' may be read in passing; `outlined' is
%  a block the reader must stop at. A third treatment would mean the first two
%  are not carrying their meaning.
% ============================================================
\tcbset{
  cfz tinted/.style={
    enhanced, breakable, sharp corners,
    lines before break=4,
    colback=mlight, boxrule=0pt, leftrule=3pt,
    left=8pt, right=8pt, top=6pt, bottom=6pt,
    before skip=13pt, after skip=13pt,
    fonttitle=\bfseries\small,
    coltitle=black, colbacktitle=mlight,
    attach title to upper={\par\smallskip},
  },
  cfz outlined/.style={
    enhanced, breakable, sharp corners,
    % A box broken straight after its title is a heading stranded at the foot
    % of a page with its question overleaf. Four lines is title plus three.
    lines before break=4,
    colback=white, boxrule=0.6pt,
    left=8pt, right=8pt, top=6pt, bottom=6pt,
    before skip=13pt, after skip=13pt,
    fonttitle=\bfseries\small,
    coltitle=black, colbacktitle=white,
    attach title to upper={\par\smallskip},
  },
}

% Read in passing.
\newtcolorbox{note}[1][]{cfz tinted, colframe=mblue, title={\lblNote}, #1}
\newtcolorbox{warning}[1][]{cfz tinted, colframe=mred, title={\lblWarning}, #1}
% Where you meet this in the world. If it cannot name a specific system,
% instrument, paper or event, it is prose and not this box.
\newtcolorbox{worldbox}[1][]{cfz tinted, colframe=mteal, title={\lblWorld}, #1}
% What is stated and not proved here, and where the proof lives.
\newtcolorbox{rigourbox}[1][]{cfz tinted, colframe=mgrey, title={\lblRigour}, #1}
% The companion library the reader builds, under code/src/chaoslab/.
\newtcolorbox{projectbox}[1][]{cfz tinted, colframe=mamber, title={\lblProject}, #1}

% Stop and read.
% The misconception, named AFTER the reader has been asked the question it
% gets wrong. Appendix B is the catalogue; a trap box is one entry in place.
\newtcolorbox{trapbox}[1][]{cfz outlined, colframe=mred, title={\lblTrap}, #1}
% A listing that has not been run on the pinned versions. Counted by
% `make debt'; nothing ships to a reader with one attached.
\newtcolorbox{verifybox}[1][]{cfz outlined, colframe=mgrey, title={\lblVerify}, #1}
\newtcolorbox{labbox}[1][]{cfz outlined, colframe=mgreen, title={\lblLab}, #1}

% ============================================================
%  THE METHOD --- questions before answers
%  ---------------------------------------------------------
%  The unit of teaching is a question the reader answers BEFORE reading the
%  answer. That is retrieval practice and errorful generation, the two
%  findings the programmed-learning format is defended by, and it is the one
%  thing the sibling mathematics book does that this one keeps whole.
%
%  \begin{think} ... \end{think}   the question, numbered per chapter, ending
%                                  with the instruction to answer before
%                                  reading on;
%  \reveal{short answer}           a centred box with the answer, which is the
%                                  first thing after the question -- the edge
%                                  the reader covers with a hand, or scrolls
%                                  to;
%  \begin{revealblock} ... \end{revealblock}  the same, for a worked answer.
%
%  tools/check_structure.py --think fails when a think box is not followed by
%  a reveal before the next think box or section, and when a reveal answers
%  nothing. That is a property of the source, so it is checked mechanically.
%
%  The instruction says "before you read on", never "before you turn over":
%  the two editions paginate differently, so the second is true in one of them.
% ============================================================
\newcounter{think}[chapter]
\renewcommand{\thethink}{\thechapter.\arabic{think}}
\newtcolorbox{thinkbox}[1][]{cfz outlined, colframe=mviolet, #1}
\NewDocumentEnvironment{think}{}{%
  \refstepcounter{think}%
  \begin{thinkbox}[title={\lblThink~\thethink}]%
}{%
  \par\smallskip
  {\raggedleft\itshape\small\color{mviolet}\lblThinkCue\par}%
  \end{thinkbox}%
}
\tcbset{
  cfz answer/.style={
    enhanced, sharp corners,
    colback=mfaint, colframe=mviolet!55, boxrule=0.5pt,
    before skip=4pt, after skip=10pt,
  },
}
\newcommand{\reveal}[1]{%
  \par\nobreak
  \begin{tcolorbox}[cfz answer, unbreakable,
    left=8pt, right=8pt, top=5pt, bottom=5pt,
    width=0.86\linewidth, center, halign=flush center]
    {\small\bfseries\color{mviolet}\lblAnswerMark}\enspace #1
  \end{tcolorbox}%
  \nobreak\noindent\ignorespaces
}
\NewDocumentEnvironment{revealblock}{}{%
  \par\nopagebreak
  \begin{tcolorbox}[cfz answer, breakable,
    left=10pt, right=10pt, top=6pt, bottom=6pt]
  {\small\bfseries\color{mviolet}\lblAnswerMark}\par\smallskip
}{%
  \end{tcolorbox}%
  \nopagebreak\par\noindent\ignorespaces
}

% ---- Outcomes on entry, and "Can you?" generated from them ----------------
% The self-check at the end of a chapter IS the list of outcomes at its
% start, one for one, because it is built from the same tokens: each
% \outcome both prints its item and appends a row to \cfz@canrows. A \loop
% inside a tabularx body does not survive alignment scanning, which is why
% the rows are accumulated at declaration time (the sibling book's finding).
\makeatletter
\newcommand{\cfz@canrows}{}
\newcommand{\cfz@ratebox}{%
  \fbox{\rule{0pt}{0.9ex}\rule{0.9ex}{0pt}}}
\newcommand{\cfz@rates}{%
  {\footnotesize\color{mgrey}%
   1\,\cfz@ratebox\ 2\,\cfz@ratebox\ 3\,\cfz@ratebox\ 4\,\cfz@ratebox\ 5\,\cfz@ratebox}}
\NewDocumentEnvironment{outcomes}{}{%
  \gdef\cfz@canrows{}%
  \begin{tcolorbox}[cfz tinted, colframe=mblue, title={\lblOutcomes}]%
  \lblOutcomesLead
  \begin{itemize}[leftmargin=1.4em, itemsep=2pt, topsep=4pt]%
}{%
  \end{itemize}%
  \end{tcolorbox}%
}
\newcommand{\outcome}[1]{%
  \item #1%
  \gappto\cfz@canrows{#1 & \cfz@rates \\ \addlinespace[3pt]}%
}
\newcommand{\canyou}{%
  \section*{\lblCanYou}%
  \phantomsection\addcontentsline{toc}{section}{\lblCanYou}%
  \noindent\lblCanYouIntro\par\medskip
  \ifdefempty{\cfz@canrows}{%
    \noindent{\color{mred}\lblNoOutcomes}\par}{%
  {\small
  \noindent\begin{tabularx}{\linewidth}{@{}>{\raggedright\arraybackslash}X r@{}}
    \toprule
    \cfz@canrows
    \bottomrule
  \end{tabularx}\par}}%
}
\makeatother

% ---- Summary -------------------------------------------------------------
\newtcolorbox{summarybox}[1][]{cfz tinted, colframe=mblue, title={\lblSummary}, #1}

% ---- Checkpoint questions, answered at the back --------------------------
% Answers are written AT THE QUESTION with \answerto{} and printed in
% Appendix A. They are carried there by a global token list rather than an
% auxiliary file, because an answer contains mathematics and \val{}, and
% material written to an .aux is expanded at shipout, where fragile commands
% break with errors naming neither the answer nor the chapter.
% \include is sequential within a run, so the list is complete by the time
% the appendix is typeset.
\makeatletter
\newcommand{\cfz@answers}{}
\newcommand{\cfz@lastansch}{0}
\NewDocumentEnvironment{checkpoint}{}{%
  \section*{\lblCheckpoint}%
  \phantomsection\addcontentsline{toc}{section}{\lblCheckpoint}%
  \noindent\lblCheckpointIntro\par
  \begin{enumerate}[label=\textbf{\arabic*.}, leftmargin=2em, itemsep=4pt]%
}{%
  \end{enumerate}%
}
\newcommand{\answerto}[1]{%
  \ifnum\value{chapter}=\cfz@lastansch\relax\else
    \xdef\cfz@lastansch{\arabic{chapter}}%
    \xappto\cfz@answers{\noexpand\cfz@anschapter{\thechapter}}%
  \fi
  \xappto\cfz@answers{\noexpand\cfz@ansitem{\arabic{enumi}}}%
  \gappto\cfz@answers{{#1}}%
}
\newcommand{\cfz@anschapter}[1]{%
  \par\medskip\noindent{\bfseries\color{mblue}\chaptername~#1}\par\nobreak\smallskip}
\newcommand{\cfz@ansitem}[2]{%
  \par\noindent\hangindent=2em\hangafter=1
  \makebox[2em][l]{\textbf{#1.}}#2\par}
\newcommand{\answersbody}{%
  \ifdefempty{\cfz@answers}{\noindent{\color{mred}\lblNoAnswers}\par}{%
    {\raggedright\cfz@answers}}%
}
\makeatother

% ============================================================
%  LABS
%  \begin{lab}{key}{title} ... \end{lab}
%
%  A lab is not a paragraph but a FILE the reader opens: a starter under
%  code/labs/chNN/<key>.py that must FAIL its test, a reference solution under
%  code/labs/chNN/solutions/, and test_<key>.py beside it. The box prints the
%  paths and the one command that runs the test. An environment rather than a
%  macro, because a macro argument cannot hold a verbatim fragment.
%  tools/check_structure.py --labs fails when a file is missing or the key's
%  chapter number disagrees with its chapter.
% ============================================================
\newcounter{lab}[chapter]
\renewcommand{\thelab}{\thechapter.\arabic{lab}}
\makeatletter
\newcommand{\labdir}{ch\two@digits{\value{chapter}}}
\newcommand{\listoflabs}{%
  \section*{\lblLabManifest}%
  \phantomsection\addcontentsline{toc}{section}{\lblLabManifest}%
  \begingroup\c@tocdepth=2\relax\@starttoc{lab}\endgroup%
}
\makeatother
\NewDocumentEnvironment{lab}{mm}{%
  \refstepcounter{lab}%
  \addcontentsline{lab}{subsection}{\protect\numberline{\thelab}\texttt{\protect\detokenize{#1}} --- #2}%
  \begin{labbox}[title={\lblLab~\thelab\ \dash{} #2}]%
}{%
  \par\medskip
  {\small\setlength{\parindent}{0pt}%
  \lblLabStarter: \texttt{code/labs/\labdir/\detokenize{#1}.py}\\
  \lblLabTest: \texttt{code/labs/\labdir/test\_\detokenize{#1}.py}\\
  \lblLabRun: \texttt{cd code \&\& uv run pytest -k \detokenize{#1}}\par}%
  \end{labbox}%
}

% ============================================================
%  FIGURES --- plots and diagrams
%  ---------------------------------------------------------
%  \plotfig{key}{caption}{one-line manifest description}
%    A PLOT, generated by code/plots/<chapter>.py from the same library the
%    listings use, into figures/plots/<lang>/<key>.pdf. Per language, because
%    an axis label in the wrong language is the half-translation that makes a
%    bilingual book look machine-made, and because the Polish edition owes a
%    decimal comma on its tick labels too. Build output, gitignored.
%
%  \mermaidfig{key}{caption}{one-line manifest description}
%    A DIAGRAM of an idea rather than of data, from
%    figures/mermaid/<lang>/<key>.mmd, rendered by `make diagrams'.
%
%  If the rendered PDF is absent either macro prints a visible marker in the
%  flow rather than a float, so a missing figure cannot be mistaken for a
%  layout choice. \phantomsection before \addcontentsline is load-bearing:
%  without it the manifest entry records whatever anchor preceded the figure,
%  which after a listing is a line that was never printed.
% ============================================================
\makeatletter
\newcommand{\listofplots}{%
  \section*{\lblPlotManifest}%
  \phantomsection\addcontentsline{toc}{section}{\lblPlotManifest}%
  \begingroup\c@tocdepth=2\relax\@starttoc{plt}\endgroup%
}
\newcommand{\listofdiagrams}{%
  \section*{\lblDiagramManifest}%
  \phantomsection\addcontentsline{toc}{section}{\lblDiagramManifest}%
  \begingroup\c@tocdepth=2\relax\@starttoc{dgm}\endgroup%
}
\makeatother

\newcommand{\figmissing}[3]{% #1 message, #2 path, #3 caption
  \par\medskip
  \begin{tcolorbox}[enhanced, breakable, width=\linewidth, sharp corners,
    colback=white, colframe=mgrey, boxrule=0.8pt,
    left=8pt, right=8pt, top=8pt, bottom=8pt]
    {\bfseries\color{mgrey}\small #1}\\[4pt]
    {\ttfamily\scriptsize\color{mgrey} #2}
  \end{tcolorbox}%
  \captionof{figure}{#3}%
  \par\medskip
}

\newcommand{\plotfig}[3]{%
  \phantomsection%
  \addcontentsline{plt}{subsection}{\protect\numberline{}\texttt{#1} --- #3}%
  \IfFileExists{figures/plots/\booklang/#1.pdf}{%
    \begin{figure}[htbp]
    \centering
    \includegraphics[width=\linewidth,height=0.45\textheight,keepaspectratio]{figures/plots/\booklang/#1.pdf}%
    \caption{#2}\label{fig:#1}
    \end{figure}%
  }{%
    \figmissing{\lblPlotNotRendered}{figures/plots/\booklang/#1.pdf}{#2}\label{fig:#1}%
  }%
}

\newcommand{\mermaidfig}[3]{%
  \phantomsection%
  \addcontentsline{dgm}{subsection}{\protect\numberline{}\texttt{#1.mmd} --- #3}%
  \IfFileExists{figures/diagrams/\booklang/#1.pdf}{%
    \begin{figure}[htbp]
    \centering
    \includegraphics[width=0.95\linewidth,height=0.42\textheight,keepaspectratio]{figures/diagrams/\booklang/#1.pdf}%
    \caption{#2}\label{fig:#1}
    \end{figure}%
  }{%
    \figmissing{\lblNotRendered}{figures/mermaid/\booklang/#1.mmd}{#2}\label{fig:#1}%
  }%
}

% ============================================================
%  APPENDIX B: the misreadings, generated by tools/gen_traps.py
%
%  \trapchapter{Chapter~\ref{ch:x}}{title}    a heading per chapter
%  \trapentry{n}{habit}{truth}{\ref{sec:x}}  one misreading
%
%  Set ragged right: an entry is mostly short sentences with \code{} spans in
%  them, and a justified line full of unbreakable runs is an overfull box
%  waiting for the edition with the longer words.
% ============================================================
\newcommand{\trapchapter}[2]{%
  \par\medskip
  {\large\bfseries\color{mblue}#1\quad\normalfont\itshape\color{mgrey}#2\par}%
  \nopagebreak\smallskip}
\newcommand{\trapentry}[4]{%
  \par\noindent
  \begingroup\raggedright\hangindent=2.4em\hangafter=1
  \makebox[2.4em][l]{\textbf{\color{mred}#1}}%
  \textbf{\enquote{#2}}\enspace #3\enspace{\color{mgrey}\S\,#4}\par
  \endgroup\smallskip}

% ============================================================
%  CHAPTER STUBS
% ============================================================
\newcommand{\chapterstub}[1]{%
  \begin{tcolorbox}[
    enhanced, breakable, width=\linewidth, sharp corners,
    colback=mlight, colframe=mred, boxrule=1pt,
    borderline={0.6pt}{3pt}{mred, dashed},
    left=12pt, right=12pt, top=12pt, bottom=12pt]
    {\bfseries\color{mred}\large \lblNotWritten}\\[8pt]
    \small\raggedright #1
  \end{tcolorbox}%
}

% ============================================================
%  COMPUTED VALUES
%  A number the reader cannot do in their head is computed by a script under
%  code/measure/, written to figures/values/<chapter>.tex as
%      \pyval{ch06.steps.apart}{34}
%  and reaches the page as \val{ch06.steps.apart}. The book contains
%  references, not digits; `make verify' fails when a committed value no
%  longer matches its script. \val applies the edition's decimal marker, which
%  is how the Polish edition gets its comma without anybody remembering.
% ============================================================
% The comma is BRACED: siunitx sets numbers in maths mode, where a bare comma
% is punctuation and gets a thin space after it ("9, 81"). {,} is an ordinary
% symbol, which is what a decimal marker is.
\newcommand{\bookdecimalmarker}{\ifpl{{,}}{.}}
\IfFileExists{siunitx.sty}{%
  \usepackage{siunitx}%
  \sisetup{
    group-digits = integer,
    group-minimum-digits = 5,
    group-separator = {\,},
    output-decimal-marker = {\bookdecimalmarker},
    exponent-product = \times,
    print-unity-mantissa = false
  }%
  \newcommand{\booknum}[1]{\num{##1}}%
}{%
  \newcommand{\booknum}[1]{##1}%
  \providecommand{\num}[1]{##1}%
}

\makeatletter
\newcommand{\pyvalues}{}
\newcommand{\pyval}[2]{\csgdef{pyval@#1}{#2}\listxadd{\pyvalues}{#1}}
\newcommand{\pyvaltext}[2]{%
  \csgdef{pyval@#1}{#2}\csgdef{pyvaltext@#1}{}\listxadd{\pyvalues}{#1}}
% \num on a non-numeric body is a FATAL siunitx error, so \val refuses a text
% value by name rather than letting siunitx fail on it.
\newcommand{\val}[1]{%
  \ifcsdef{pyval@#1}{%
    \ifcsdef{pyvaltext@#1}{%
      \PackageError{chaosbook}{Value `#1' is not a number}%
        {Use \string\valtext{#1} instead.}%
    }{}%
    \booknum{\csuse{pyval@#1}}%
  }{\textcolor{mred}{\textbf{[\lblValueMissing: #1]}}}%
}
\newcommand{\valtext}[1]{%
  \ifcsdef{pyval@#1}{\csuse{pyval@#1}}%
    {\textcolor{mred}{\textbf{[\lblValueMissing: #1]}}}%
}
\makeatother

% figures/values/all.tex is generated by `make numbers' and \inputs each
% value file. A build with no values still compiles; every \val{} then prints
% a visible marker. \valuesindex lets the one-chapter build point elsewhere.
\providecommand{\valuesindex}{figures/values/all}
\IfFileExists{\valuesindex.tex}{\input{\valuesindex}}{%
  \AtBeginDocument{\PackageWarning{chaosbook}{No computed values found. Run `make numbers'.}}}

% ============================================================
%  PINNED VERSIONS --- single source of truth
%  Change these here AND in code/pyproject.toml; tools/check_structure.py
%  --pins fails when the two disagree. A chapter never types a version.
% ============================================================
\newcommand{\bookedition}{0.1}
\newcommand{\bookyear}{2026}
\newcommand{\pyver}{3.14}              % the minor the book targets
\newcommand{\pypatch}{3.14.7}          % the exact interpreter the code ran on
\newcommand{\uvver}{0.12.13}           % uv
\newcommand{\numpyver}{2.5.3}          % numpy
\newcommand{\matplotlibver}{3.11.2}    % matplotlib
\newcommand{\pytestver}{9.1.1}         % pytest
\newcommand{\ruffver}{0.16.9}          % ruff
% The prose licence is a name with a version in it, not a quantity: it is a
% macro so that neither edition typesets it as a decimal to be localised.
\newcommand{\proselicence}{CC~BY-NC-SA~4.0}

% ============================================================
%  INDEX
% ============================================================
\apptocmd{\theindex}{\raggedcolumns\raggedright}{}{%
  \PackageWarning{chaosbook}{Could not patch theindex for ragged setting}}
\providecommand*\see[2]{}
\renewcommand*\see[2]{\emph{\lblIndexSee} #1}
\providecommand*\seealso[2]{}
\renewcommand*\seealso[2]{\emph{\lblIndexSeeAlso} #1}

\makeindex
; nothing else differs between them at this level.
%  Every user-visible string lives in lang/en.tex or lang/pl.tex, so a label
%  that appears in one edition and not the other is a missing macro rather
%  than a silent divergence.
%
%  Lineage. This book is the third in a family and takes one thing from each
%  of the other two:
%
%    * from Mathematics from Zero for the AI Engineer: the bilingual single
%      source (\booklang, lang/, one body.tex), the rule that every number is
%      computed and reaches the page as \val{}, and the programmed-learning
%      method -- outcomes on entry, a question before every answer, a
%      misconception elicited before it is named, a self-check generated from
%      the outcomes, and checkpoint answers collected at the back;
%    * from Python for .NET Engineers: every listing is a FILE that CI runs,
%      every lab is a starter that must fail and a solution that must pass,
%      and a console transcript is a computed number that happens to be
%      verbatim.
%
%  What is new here is \plotfig. A book about chaos cannot be taught in boxes
%  and arrows: the bifurcation diagram and the Lorenz attractor ARE the
%  argument, so the plots are generated by code/plots/ from the same library
%  the listings use, per language, and never drawn by hand.
% ============================================================

% \booklang must be set by the main file. Default to English rather than
% failing, so a stray % ============================================================
%  preamble.tex --- styling, environments, macros
%
%  Shared by BOTH editions. main-en.tex and main-pl.tex each define \booklang
%  before % ============================================================
%  preamble.tex --- styling, environments, macros
%
%  Shared by BOTH editions. main-en.tex and main-pl.tex each define \booklang
%  before \input{preamble}; nothing else differs between them at this level.
%  Every user-visible string lives in lang/en.tex or lang/pl.tex, so a label
%  that appears in one edition and not the other is a missing macro rather
%  than a silent divergence.
%
%  Lineage. This book is the third in a family and takes one thing from each
%  of the other two:
%
%    * from Mathematics from Zero for the AI Engineer: the bilingual single
%      source (\booklang, lang/, one body.tex), the rule that every number is
%      computed and reaches the page as \val{}, and the programmed-learning
%      method -- outcomes on entry, a question before every answer, a
%      misconception elicited before it is named, a self-check generated from
%      the outcomes, and checkpoint answers collected at the back;
%    * from Python for .NET Engineers: every listing is a FILE that CI runs,
%      every lab is a starter that must fail and a solution that must pass,
%      and a console transcript is a computed number that happens to be
%      verbatim.
%
%  What is new here is \plotfig. A book about chaos cannot be taught in boxes
%  and arrows: the bifurcation diagram and the Lorenz attractor ARE the
%  argument, so the plots are generated by code/plots/ from the same library
%  the listings use, per language, and never drawn by hand.
% ============================================================

% \booklang must be set by the main file. Default to English rather than
% failing, so a stray \input{preamble} still compiles and says what it did.
\providecommand{\booklang}{en}

\usepackage{etoolbox}   % loaded first: the language switch below needs it

% ---------- Page geometry ----------
% ONE format: A4, 12pt, ONE-SIDED. The book is read on a screen, usually with
% an editor open beside it, so there is no recto, no verso, no gutter and no
% blank leaf before a chapter. The measure is about seventy-five characters
% of prose. At \footnotesize a listing line of 79 characters fits without
% wrapping, and 79 is the width every file under code/ is held to by
% tools/check_structure.py --lines.
\usepackage[
  a4paper,
  left=3.1cm, right=3.1cm, top=2.6cm, bottom=3.0cm,
  head=26pt, headsep=0.8cm, footskip=1.4cm
]{geometry}

% ---------- Encoding & fonts ----------
\usepackage[T1]{fontenc}
\usepackage[utf8]{inputenc}
\IfFileExists{lmodern.sty}{\usepackage{lmodern}}{}
% newtxtext takes its TS1 symbols from TeX Gyre Termes (Debian: tex-gyre).
% A machine with newtx and without tex-gyre dies on the copyright page. CI's
% full TeX Live is the reference; locally install tex-gyre as well.
\IfFileExists{newtxtext.sty}{\usepackage{newtxtext}}{}
% amsmath before newtxmath; amssymb only when newtxmath is absent, because
% loading both is a fatal clash (\Bbbk already defined) that is invisible on
% a machine without newtx.
\usepackage{amsmath}
\IfFileExists{newtxmath.sty}{\usepackage{newtxmath}}{\usepackage{amssymb}}
\IfFileExists{inconsolata.sty}{\usepackage[varqu,scaled=0.95]{inconsolata}}{}
% Without lmodern, microtype runs with font expansion off, which is what
% resolves a marginal overfull line -- so a container short of lmodern
% disagrees with CI about line breaks. Install lmodern locally too.
\IfFileExists{lmodern.sty}{\usepackage{microtype}}{\usepackage[expansion=false]{microtype}}

% upquote: in T1 the input ' is a right single quote, and a listing that sets
% f('x') with curly quotes is a SyntaxError when typed back in.
\IfFileExists{upquote.sty}{\usepackage{upquote}}{}
% And its twin: in T1 the pair -- is an en-dash ligature in the typewriter
% family too, so \code{--flag} would print a flag nobody can type. Scoped to
% tt*, so prose keeps its own dashes.
\DisableLigatures[-]{encoding = T1, family = tt*}

% ---------- Language ----------
% babel with a missing language file is FATAL, so each .ldf is probed first.
\IfFileExists{babel.sty}{%
  \IfFileExists{polish.ldf}{%
    \IfFileExists{british.ldf}{%
      \ifdefstring{\booklang}{pl}%
        {\usepackage[british,main=polish]{babel}}%
        {\usepackage[polish,main=british]{babel}}%
    }{\usepackage[polish]{babel}}%
  }{%
    \IfFileExists{british.ldf}{\usepackage[british]{babel}}{}%
  }%
}{}
% babel-polish makes " an active shorthand, which is a category-code accident
% waiting to happen inside listings. Polish quotation marks come from
% \enquote, which csquotes sets per language.
\ifdefstring{\booklang}{pl}{\AtBeginDocument{\shorthandoff{"}}}{}

% ---------- Core packages ----------
\usepackage{graphicx}
\usepackage{xcolor}
\usepackage{booktabs}
\usepackage{tabularx}
\usepackage{longtable}
\usepackage{array}
\usepackage{enumitem}
\usepackage{listings}
\usepackage[most]{tcolorbox}
\usepackage{fancyhdr}
\usepackage{titlesec}
\usepackage{caption}
\usepackage{imakeidx}
\IfFileExists{csquotes.sty}{\usepackage[autostyle=true]{csquotes}}{%
  \providecommand{\enquote}[1]{``##1''}}
\usepackage[hidelinks]{hyperref}

% ---------- Palette ----------
% The family's palette, so the three books read as a set. Two colours are
% this book's own: the lab (green) and the stop-and-think box (violet).
\definecolor{mblue}{HTML}{1F4E79}
\definecolor{mteal}{HTML}{0E7C7B}
\definecolor{mamber}{HTML}{B26A00}
\definecolor{mred}{HTML}{9B2C2C}
\definecolor{mviolet}{HTML}{7B4397}
\definecolor{mgreen}{HTML}{3B7A3B}
\definecolor{mgrey}{HTML}{5A5A5A}
\definecolor{mlight}{HTML}{F4F6F8}
\definecolor{mfaint}{HTML}{FBFCFD}
\definecolor{codebg}{HTML}{FAFAFA}
\definecolor{codeframe}{HTML}{D8DEE4}
\definecolor{kw}{HTML}{0B5394}
\definecolor{str}{HTML}{9B2C2C}
\definecolor{cmt}{HTML}{4F7A4F}
\definecolor{typ}{HTML}{0E7C7B}
\definecolor{dec}{HTML}{7B4397}

\hypersetup{
  colorlinks=true,
  linkcolor=mblue,
  urlcolor=mteal,
  citecolor=mblue,
  pdfauthor={Konrad Cinkusz}
}

% \ifpl{Polish}{English} -- for the handful of places where a string is
% structural rather than content. Robust, because a fragile \ifpl inside a
% moving argument fails on the next run in exactly one of the two editions.
\DeclareRobustCommand{\ifpl}[2]{\ifdefstring{\booklang}{pl}{#1}{#2}}
\makeatletter
\pdfstringdefDisableCommands{%
  \ifdefstring{\booklang}{pl}{\let\ifpl\@firstoftwo}{\let\ifpl\@secondoftwo}%
  % hyperref drops \color from a bookmark and keeps its ARGUMENT, so a title
  % carrying \code{} would bookmark as "mtealx". Gobble the colour name.
  \def\color#1{}%
}
\makeatother

% ---------- Language strings ----------
\input{lang/\booklang}

% The PDF's subject, keywords and language read \lblPdf... from the language
% file, so they cannot sit in the colour block above, which runs first. Two
% blocks is the ordering, not an oversight.
\hypersetup{
  pdfsubject={\lblPdfSubject},
  pdfkeywords={\lblPdfKeywords},
  pdflang={\lblPdfLang}
}

% ---------- Headings ----------
\titleformat{\chapter}[display]
  {\normalfont\huge\bfseries\color{mblue}\raggedright}
  {\normalfont\normalsize\color{mgrey}\MakeUppercase{\chaptertitlename}\ \thechapter}
  {8pt}{\Huge\raggedright}
\titlespacing*{\chapter}{0pt}{10pt}{28pt}
\titleformat{\section}{\normalfont\Large\bfseries\color{mblue}\raggedright}{\thesection}{0.8em}{}
\titleformat{\subsection}{\normalfont\large\bfseries\color{mgrey}\raggedright}{\thesubsection}{0.7em}{}

\setcounter{tocdepth}{1}
\setcounter{secnumdepth}{2}

\pagestyle{fancy}
\fancyhf{}
\fancyhead[L]{\small\nouppercase{\leftmark}}
\fancyhead[R]{\small\thepage}
\renewcommand{\headrulewidth}{0.4pt}
\fancypagestyle{plain}{\fancyhf{}\fancyfoot[R]{\small\thepage}%
  \renewcommand{\headrulewidth}{0pt}}
% After \pagestyle{fancy}, which installs its own \chaptermark.
\renewcommand{\chaptermark}[1]{\markboth{\thechapter.\ #1}{}}

% ---------- Mathematics ----------
% The handful of notations the book uses, as macros so both editions set the
% same symbol. The book writes \ln for the natural logarithm and \log_{10}
% where the base is ten; a bare \log is ambiguous between the two readerships
% and is not used.
\newcommand{\R}{\mathbb{R}}
\newcommand{\dd}{\,\mathrm{d}}
\newcommand{\abs}[1]{\left\lvert #1\right\rvert}
\newcommand{\iu}{\mathrm{i}}

% ============================================================
%  CODE LISTINGS
%  Every listing in a chapter is a FILE under code/, pulled in with \pyfile or
%  \pyregion, and CI runs it on the pinned interpreter. An in-source
%  \begin{python} block is allowed for a fragment of three or four lines that
%  shows a shape rather than runs.
% ============================================================
\lstdefinelanguage{pythonx}{
  language=Python,
  morekeywords={async,await,match,case,nonlocal,yield,type},
  morekeywords=[2]{
    orbit,iterate,logistic,rk4_step,euler_step,integrate,lorenz,henon,
    lyapunov,separation,escape_time,box_count
  },
  sensitive=true
}

\lstdefinestyle{bookcode}{
  basicstyle=\ttfamily\footnotesize,
  keywordstyle=\color{kw}\bfseries,
  keywordstyle=[2]\color{typ},
  stringstyle=\color{str},
  % Colour only: inconsolata has no italic, and an italic comment style is
  % silently substituted and logged on every build.
  commentstyle=\color{cmt},
  numbers=left,
  numberstyle=\ttfamily\tiny\color{mgrey},
  numbersep=8pt,
  backgroundcolor=\color{codebg},
  frame=single,
  rulecolor=\color{codeframe},
  framesep=6pt,
  breaklines=true,
  breakatwhitespace=false,
  postbreak=\mbox{\textcolor{mgrey}{$\hookrightarrow$}\space},
  showstringspaces=false,
  tabsize=4,
  columns=fullflexible,
  keepspaces=true,
  captionpos=b,
  aboveskip=1.2em,
  belowskip=1.2em,
  % Do NOT add a mapping for U+00A0: it aborts the run with a message that
  % names neither the character nor this line. Listings are ASCII, and
  % parity's C13 checks it.
  literate={ł}{{\l}}1 {Ł}{{\L}}1 {ą}{{\k a}}1 {ę}{{\k e}}1
           {ś}{{\'s}}1 {ć}{{\'c}}1 {ż}{{\.z}}1 {ź}{{\'z}}1 {ó}{{\'o}}1 {ń}{{\'n}}1
           {Ą}{{\k A}}1 {Ę}{{\k E}}1 {Ś}{{\'S}}1 {Ć}{{\'C}}1
           {Ż}{{\.Z}}1 {Ź}{{\'Z}}1 {Ó}{{\'O}}1 {Ń}{{\'N}}1
           {—}{{-{}-}}2 {–}{{-}}1 {‘}{{'}}1 {’}{{'}}1
           {“}{{"}}1 {”}{{"}}1 {°}{{\textdegree}}1
}
\lstset{style=bookcode}

\lstnewenvironment{python}[1][]
  {\lstset{language=pythonx,style=bookcode,#1}}
  {}
\lstnewenvironment{shellcmd}[1][]
  {\lstset{language=bash,style=bookcode,numbers=none,keywordstyle=\color{mteal}\bfseries,#1}}
  {}

% The path is printed above every file listing, in the form the reader opens
% in the editor. \nobreak keeps it on the page of the listing it names.
% The path must not be stranded at the foot of a page with its listing
% overleaf: \nobreak alone does not hold, because listings opens with its own
% vertical material. \needspace asks for room for the path and the first few
% lines of code, and turns the page first when there is not.
% \Needspace (capital N) measures \pagegoal-\pagetotal exactly; lower-case
% \needspace works through glue and penalties and was measured, by the
% Chapter 11 pass, turning pages early and leaving gaps. It must be used
% between paragraphs, which \listingpath's own \par guarantees. Chapters use
% \Needspace* too, so the fallback for a machine without the package defines
% both forms (the star swallowed with its argument).
\makeatletter
\IfFileExists{needspace.sty}{\usepackage{needspace}}{%
  \providecommand{\needspace}[1]{}%
  \providecommand{\Needspace}{\@ifstar\@gobble\@gobble}}
\makeatother
\newcommand{\listingpath}[1]{%
  \par\medskip\Needspace{6\baselineskip}\noindent
  {\ttfamily\footnotesize\color{mgrey}\detokenize{#1}}%
  \par\nobreak}

% Usage: \pyfile{code/ch05/cobweb.py}{Caption}{lst:label}
\newcommand{\pyfile}[3]{%
  \listingpath{#1}%
  \lstinputlisting[language=pythonx,style=bookcode,aboveskip=0.3em,
    caption={#2},label={#3}]{#1}%
}

% A named region of a file, delimited by `# --8<-- [start:name]` and
% `# --8<-- [end:name]`. Matched through listings' range prefix and suffix
% keys -- the other obvious form matches nothing and prints the WHOLE file
% without a warning (measured in the sibling book). A region name is a word,
% never a number. The line numbers printed are the file's own.
% Usage: \pyregion{code/ch05/cobweb.py}{step}{Caption}{lst:label}
\lstset{
  rangebeginprefix=\#\ --8<--\ [start:, rangebeginsuffix=],
  rangeendprefix=\#\ --8<--\ [end:, rangeendsuffix=],
  includerangemarker=false}
\newcommand{\pyregion}[4]{%
  \listingpath{#1}%
  \lstinputlisting[language=pythonx,style=bookcode,aboveskip=0.3em,
    linerange={#2-#2},caption={#3},label={#4}]{#1}%
}

% A program's output, written by code/measure/transcripts.py. There is no
% typed form and there must not be: \val{} does not expand in a verbatim
% body, and a transcript nobody ran looks exactly like one that was.
\newcommand{\transcript}[1]{%
  \par\smallskip
  \IfFileExists{figures/transcripts/#1.txt}{%
    \lstinputlisting[style=bookcode,numbers=none,basicstyle=\ttfamily\footnotesize,
      keywordstyle=\color{black},backgroundcolor=\color{white},
      language={}]{figures/transcripts/#1.txt}%
  }{%
    \begin{tcolorbox}[enhanced, sharp corners, width=\linewidth,
      colback=mlight, colframe=mgrey, boxrule=0.6pt,
      left=8pt, right=8pt, top=6pt, bottom=6pt]
      {\bfseries\color{mgrey}\small \lblTranscriptMissing}\\[2pt]
      {\ttfamily\scriptsize\color{mgrey} figures/transcripts/#1.txt}
    \end{tcolorbox}%
  }%
  \par\smallskip
}

% Inline literals. Underscores inside these must be written \_.
\newcommand{\code}[1]{\texttt{\small #1}}
\newcommand{\pkg}[1]{\texttt{\small\color{mblue}#1}}
\newcommand{\api}[1]{\texttt{\small\color{mteal}#1}}

% ============================================================
%  ADMONITIONS
%  A small, fixed vocabulary. `tinted' may be read in passing; `outlined' is
%  a block the reader must stop at. A third treatment would mean the first two
%  are not carrying their meaning.
% ============================================================
\tcbset{
  cfz tinted/.style={
    enhanced, breakable, sharp corners,
    lines before break=4,
    colback=mlight, boxrule=0pt, leftrule=3pt,
    left=8pt, right=8pt, top=6pt, bottom=6pt,
    before skip=13pt, after skip=13pt,
    fonttitle=\bfseries\small,
    coltitle=black, colbacktitle=mlight,
    attach title to upper={\par\smallskip},
  },
  cfz outlined/.style={
    enhanced, breakable, sharp corners,
    % A box broken straight after its title is a heading stranded at the foot
    % of a page with its question overleaf. Four lines is title plus three.
    lines before break=4,
    colback=white, boxrule=0.6pt,
    left=8pt, right=8pt, top=6pt, bottom=6pt,
    before skip=13pt, after skip=13pt,
    fonttitle=\bfseries\small,
    coltitle=black, colbacktitle=white,
    attach title to upper={\par\smallskip},
  },
}

% Read in passing.
\newtcolorbox{note}[1][]{cfz tinted, colframe=mblue, title={\lblNote}, #1}
\newtcolorbox{warning}[1][]{cfz tinted, colframe=mred, title={\lblWarning}, #1}
% Where you meet this in the world. If it cannot name a specific system,
% instrument, paper or event, it is prose and not this box.
\newtcolorbox{worldbox}[1][]{cfz tinted, colframe=mteal, title={\lblWorld}, #1}
% What is stated and not proved here, and where the proof lives.
\newtcolorbox{rigourbox}[1][]{cfz tinted, colframe=mgrey, title={\lblRigour}, #1}
% The companion library the reader builds, under code/src/chaoslab/.
\newtcolorbox{projectbox}[1][]{cfz tinted, colframe=mamber, title={\lblProject}, #1}

% Stop and read.
% The misconception, named AFTER the reader has been asked the question it
% gets wrong. Appendix B is the catalogue; a trap box is one entry in place.
\newtcolorbox{trapbox}[1][]{cfz outlined, colframe=mred, title={\lblTrap}, #1}
% A listing that has not been run on the pinned versions. Counted by
% `make debt'; nothing ships to a reader with one attached.
\newtcolorbox{verifybox}[1][]{cfz outlined, colframe=mgrey, title={\lblVerify}, #1}
\newtcolorbox{labbox}[1][]{cfz outlined, colframe=mgreen, title={\lblLab}, #1}

% ============================================================
%  THE METHOD --- questions before answers
%  ---------------------------------------------------------
%  The unit of teaching is a question the reader answers BEFORE reading the
%  answer. That is retrieval practice and errorful generation, the two
%  findings the programmed-learning format is defended by, and it is the one
%  thing the sibling mathematics book does that this one keeps whole.
%
%  \begin{think} ... \end{think}   the question, numbered per chapter, ending
%                                  with the instruction to answer before
%                                  reading on;
%  \reveal{short answer}           a centred box with the answer, which is the
%                                  first thing after the question -- the edge
%                                  the reader covers with a hand, or scrolls
%                                  to;
%  \begin{revealblock} ... \end{revealblock}  the same, for a worked answer.
%
%  tools/check_structure.py --think fails when a think box is not followed by
%  a reveal before the next think box or section, and when a reveal answers
%  nothing. That is a property of the source, so it is checked mechanically.
%
%  The instruction says "before you read on", never "before you turn over":
%  the two editions paginate differently, so the second is true in one of them.
% ============================================================
\newcounter{think}[chapter]
\renewcommand{\thethink}{\thechapter.\arabic{think}}
\newtcolorbox{thinkbox}[1][]{cfz outlined, colframe=mviolet, #1}
\NewDocumentEnvironment{think}{}{%
  \refstepcounter{think}%
  \begin{thinkbox}[title={\lblThink~\thethink}]%
}{%
  \par\smallskip
  {\raggedleft\itshape\small\color{mviolet}\lblThinkCue\par}%
  \end{thinkbox}%
}
\tcbset{
  cfz answer/.style={
    enhanced, sharp corners,
    colback=mfaint, colframe=mviolet!55, boxrule=0.5pt,
    before skip=4pt, after skip=10pt,
  },
}
\newcommand{\reveal}[1]{%
  \par\nobreak
  \begin{tcolorbox}[cfz answer, unbreakable,
    left=8pt, right=8pt, top=5pt, bottom=5pt,
    width=0.86\linewidth, center, halign=flush center]
    {\small\bfseries\color{mviolet}\lblAnswerMark}\enspace #1
  \end{tcolorbox}%
  \nobreak\noindent\ignorespaces
}
\NewDocumentEnvironment{revealblock}{}{%
  \par\nopagebreak
  \begin{tcolorbox}[cfz answer, breakable,
    left=10pt, right=10pt, top=6pt, bottom=6pt]
  {\small\bfseries\color{mviolet}\lblAnswerMark}\par\smallskip
}{%
  \end{tcolorbox}%
  \nopagebreak\par\noindent\ignorespaces
}

% ---- Outcomes on entry, and "Can you?" generated from them ----------------
% The self-check at the end of a chapter IS the list of outcomes at its
% start, one for one, because it is built from the same tokens: each
% \outcome both prints its item and appends a row to \cfz@canrows. A \loop
% inside a tabularx body does not survive alignment scanning, which is why
% the rows are accumulated at declaration time (the sibling book's finding).
\makeatletter
\newcommand{\cfz@canrows}{}
\newcommand{\cfz@ratebox}{%
  \fbox{\rule{0pt}{0.9ex}\rule{0.9ex}{0pt}}}
\newcommand{\cfz@rates}{%
  {\footnotesize\color{mgrey}%
   1\,\cfz@ratebox\ 2\,\cfz@ratebox\ 3\,\cfz@ratebox\ 4\,\cfz@ratebox\ 5\,\cfz@ratebox}}
\NewDocumentEnvironment{outcomes}{}{%
  \gdef\cfz@canrows{}%
  \begin{tcolorbox}[cfz tinted, colframe=mblue, title={\lblOutcomes}]%
  \lblOutcomesLead
  \begin{itemize}[leftmargin=1.4em, itemsep=2pt, topsep=4pt]%
}{%
  \end{itemize}%
  \end{tcolorbox}%
}
\newcommand{\outcome}[1]{%
  \item #1%
  \gappto\cfz@canrows{#1 & \cfz@rates \\ \addlinespace[3pt]}%
}
\newcommand{\canyou}{%
  \section*{\lblCanYou}%
  \phantomsection\addcontentsline{toc}{section}{\lblCanYou}%
  \noindent\lblCanYouIntro\par\medskip
  \ifdefempty{\cfz@canrows}{%
    \noindent{\color{mred}\lblNoOutcomes}\par}{%
  {\small
  \noindent\begin{tabularx}{\linewidth}{@{}>{\raggedright\arraybackslash}X r@{}}
    \toprule
    \cfz@canrows
    \bottomrule
  \end{tabularx}\par}}%
}
\makeatother

% ---- Summary -------------------------------------------------------------
\newtcolorbox{summarybox}[1][]{cfz tinted, colframe=mblue, title={\lblSummary}, #1}

% ---- Checkpoint questions, answered at the back --------------------------
% Answers are written AT THE QUESTION with \answerto{} and printed in
% Appendix A. They are carried there by a global token list rather than an
% auxiliary file, because an answer contains mathematics and \val{}, and
% material written to an .aux is expanded at shipout, where fragile commands
% break with errors naming neither the answer nor the chapter.
% \include is sequential within a run, so the list is complete by the time
% the appendix is typeset.
\makeatletter
\newcommand{\cfz@answers}{}
\newcommand{\cfz@lastansch}{0}
\NewDocumentEnvironment{checkpoint}{}{%
  \section*{\lblCheckpoint}%
  \phantomsection\addcontentsline{toc}{section}{\lblCheckpoint}%
  \noindent\lblCheckpointIntro\par
  \begin{enumerate}[label=\textbf{\arabic*.}, leftmargin=2em, itemsep=4pt]%
}{%
  \end{enumerate}%
}
\newcommand{\answerto}[1]{%
  \ifnum\value{chapter}=\cfz@lastansch\relax\else
    \xdef\cfz@lastansch{\arabic{chapter}}%
    \xappto\cfz@answers{\noexpand\cfz@anschapter{\thechapter}}%
  \fi
  \xappto\cfz@answers{\noexpand\cfz@ansitem{\arabic{enumi}}}%
  \gappto\cfz@answers{{#1}}%
}
\newcommand{\cfz@anschapter}[1]{%
  \par\medskip\noindent{\bfseries\color{mblue}\chaptername~#1}\par\nobreak\smallskip}
\newcommand{\cfz@ansitem}[2]{%
  \par\noindent\hangindent=2em\hangafter=1
  \makebox[2em][l]{\textbf{#1.}}#2\par}
\newcommand{\answersbody}{%
  \ifdefempty{\cfz@answers}{\noindent{\color{mred}\lblNoAnswers}\par}{%
    {\raggedright\cfz@answers}}%
}
\makeatother

% ============================================================
%  LABS
%  \begin{lab}{key}{title} ... \end{lab}
%
%  A lab is not a paragraph but a FILE the reader opens: a starter under
%  code/labs/chNN/<key>.py that must FAIL its test, a reference solution under
%  code/labs/chNN/solutions/, and test_<key>.py beside it. The box prints the
%  paths and the one command that runs the test. An environment rather than a
%  macro, because a macro argument cannot hold a verbatim fragment.
%  tools/check_structure.py --labs fails when a file is missing or the key's
%  chapter number disagrees with its chapter.
% ============================================================
\newcounter{lab}[chapter]
\renewcommand{\thelab}{\thechapter.\arabic{lab}}
\makeatletter
\newcommand{\labdir}{ch\two@digits{\value{chapter}}}
\newcommand{\listoflabs}{%
  \section*{\lblLabManifest}%
  \phantomsection\addcontentsline{toc}{section}{\lblLabManifest}%
  \begingroup\c@tocdepth=2\relax\@starttoc{lab}\endgroup%
}
\makeatother
\NewDocumentEnvironment{lab}{mm}{%
  \refstepcounter{lab}%
  \addcontentsline{lab}{subsection}{\protect\numberline{\thelab}\texttt{\protect\detokenize{#1}} --- #2}%
  \begin{labbox}[title={\lblLab~\thelab\ \dash{} #2}]%
}{%
  \par\medskip
  {\small\setlength{\parindent}{0pt}%
  \lblLabStarter: \texttt{code/labs/\labdir/\detokenize{#1}.py}\\
  \lblLabTest: \texttt{code/labs/\labdir/test\_\detokenize{#1}.py}\\
  \lblLabRun: \texttt{cd code \&\& uv run pytest -k \detokenize{#1}}\par}%
  \end{labbox}%
}

% ============================================================
%  FIGURES --- plots and diagrams
%  ---------------------------------------------------------
%  \plotfig{key}{caption}{one-line manifest description}
%    A PLOT, generated by code/plots/<chapter>.py from the same library the
%    listings use, into figures/plots/<lang>/<key>.pdf. Per language, because
%    an axis label in the wrong language is the half-translation that makes a
%    bilingual book look machine-made, and because the Polish edition owes a
%    decimal comma on its tick labels too. Build output, gitignored.
%
%  \mermaidfig{key}{caption}{one-line manifest description}
%    A DIAGRAM of an idea rather than of data, from
%    figures/mermaid/<lang>/<key>.mmd, rendered by `make diagrams'.
%
%  If the rendered PDF is absent either macro prints a visible marker in the
%  flow rather than a float, so a missing figure cannot be mistaken for a
%  layout choice. \phantomsection before \addcontentsline is load-bearing:
%  without it the manifest entry records whatever anchor preceded the figure,
%  which after a listing is a line that was never printed.
% ============================================================
\makeatletter
\newcommand{\listofplots}{%
  \section*{\lblPlotManifest}%
  \phantomsection\addcontentsline{toc}{section}{\lblPlotManifest}%
  \begingroup\c@tocdepth=2\relax\@starttoc{plt}\endgroup%
}
\newcommand{\listofdiagrams}{%
  \section*{\lblDiagramManifest}%
  \phantomsection\addcontentsline{toc}{section}{\lblDiagramManifest}%
  \begingroup\c@tocdepth=2\relax\@starttoc{dgm}\endgroup%
}
\makeatother

\newcommand{\figmissing}[3]{% #1 message, #2 path, #3 caption
  \par\medskip
  \begin{tcolorbox}[enhanced, breakable, width=\linewidth, sharp corners,
    colback=white, colframe=mgrey, boxrule=0.8pt,
    left=8pt, right=8pt, top=8pt, bottom=8pt]
    {\bfseries\color{mgrey}\small #1}\\[4pt]
    {\ttfamily\scriptsize\color{mgrey} #2}
  \end{tcolorbox}%
  \captionof{figure}{#3}%
  \par\medskip
}

\newcommand{\plotfig}[3]{%
  \phantomsection%
  \addcontentsline{plt}{subsection}{\protect\numberline{}\texttt{#1} --- #3}%
  \IfFileExists{figures/plots/\booklang/#1.pdf}{%
    \begin{figure}[htbp]
    \centering
    \includegraphics[width=\linewidth,height=0.45\textheight,keepaspectratio]{figures/plots/\booklang/#1.pdf}%
    \caption{#2}\label{fig:#1}
    \end{figure}%
  }{%
    \figmissing{\lblPlotNotRendered}{figures/plots/\booklang/#1.pdf}{#2}\label{fig:#1}%
  }%
}

\newcommand{\mermaidfig}[3]{%
  \phantomsection%
  \addcontentsline{dgm}{subsection}{\protect\numberline{}\texttt{#1.mmd} --- #3}%
  \IfFileExists{figures/diagrams/\booklang/#1.pdf}{%
    \begin{figure}[htbp]
    \centering
    \includegraphics[width=0.95\linewidth,height=0.42\textheight,keepaspectratio]{figures/diagrams/\booklang/#1.pdf}%
    \caption{#2}\label{fig:#1}
    \end{figure}%
  }{%
    \figmissing{\lblNotRendered}{figures/mermaid/\booklang/#1.mmd}{#2}\label{fig:#1}%
  }%
}

% ============================================================
%  APPENDIX B: the misreadings, generated by tools/gen_traps.py
%
%  \trapchapter{Chapter~\ref{ch:x}}{title}    a heading per chapter
%  \trapentry{n}{habit}{truth}{\ref{sec:x}}  one misreading
%
%  Set ragged right: an entry is mostly short sentences with \code{} spans in
%  them, and a justified line full of unbreakable runs is an overfull box
%  waiting for the edition with the longer words.
% ============================================================
\newcommand{\trapchapter}[2]{%
  \par\medskip
  {\large\bfseries\color{mblue}#1\quad\normalfont\itshape\color{mgrey}#2\par}%
  \nopagebreak\smallskip}
\newcommand{\trapentry}[4]{%
  \par\noindent
  \begingroup\raggedright\hangindent=2.4em\hangafter=1
  \makebox[2.4em][l]{\textbf{\color{mred}#1}}%
  \textbf{\enquote{#2}}\enspace #3\enspace{\color{mgrey}\S\,#4}\par
  \endgroup\smallskip}

% ============================================================
%  CHAPTER STUBS
% ============================================================
\newcommand{\chapterstub}[1]{%
  \begin{tcolorbox}[
    enhanced, breakable, width=\linewidth, sharp corners,
    colback=mlight, colframe=mred, boxrule=1pt,
    borderline={0.6pt}{3pt}{mred, dashed},
    left=12pt, right=12pt, top=12pt, bottom=12pt]
    {\bfseries\color{mred}\large \lblNotWritten}\\[8pt]
    \small\raggedright #1
  \end{tcolorbox}%
}

% ============================================================
%  COMPUTED VALUES
%  A number the reader cannot do in their head is computed by a script under
%  code/measure/, written to figures/values/<chapter>.tex as
%      \pyval{ch06.steps.apart}{34}
%  and reaches the page as \val{ch06.steps.apart}. The book contains
%  references, not digits; `make verify' fails when a committed value no
%  longer matches its script. \val applies the edition's decimal marker, which
%  is how the Polish edition gets its comma without anybody remembering.
% ============================================================
% The comma is BRACED: siunitx sets numbers in maths mode, where a bare comma
% is punctuation and gets a thin space after it ("9, 81"). {,} is an ordinary
% symbol, which is what a decimal marker is.
\newcommand{\bookdecimalmarker}{\ifpl{{,}}{.}}
\IfFileExists{siunitx.sty}{%
  \usepackage{siunitx}%
  \sisetup{
    group-digits = integer,
    group-minimum-digits = 5,
    group-separator = {\,},
    output-decimal-marker = {\bookdecimalmarker},
    exponent-product = \times,
    print-unity-mantissa = false
  }%
  \newcommand{\booknum}[1]{\num{##1}}%
}{%
  \newcommand{\booknum}[1]{##1}%
  \providecommand{\num}[1]{##1}%
}

\makeatletter
\newcommand{\pyvalues}{}
\newcommand{\pyval}[2]{\csgdef{pyval@#1}{#2}\listxadd{\pyvalues}{#1}}
\newcommand{\pyvaltext}[2]{%
  \csgdef{pyval@#1}{#2}\csgdef{pyvaltext@#1}{}\listxadd{\pyvalues}{#1}}
% \num on a non-numeric body is a FATAL siunitx error, so \val refuses a text
% value by name rather than letting siunitx fail on it.
\newcommand{\val}[1]{%
  \ifcsdef{pyval@#1}{%
    \ifcsdef{pyvaltext@#1}{%
      \PackageError{chaosbook}{Value `#1' is not a number}%
        {Use \string\valtext{#1} instead.}%
    }{}%
    \booknum{\csuse{pyval@#1}}%
  }{\textcolor{mred}{\textbf{[\lblValueMissing: #1]}}}%
}
\newcommand{\valtext}[1]{%
  \ifcsdef{pyval@#1}{\csuse{pyval@#1}}%
    {\textcolor{mred}{\textbf{[\lblValueMissing: #1]}}}%
}
\makeatother

% figures/values/all.tex is generated by `make numbers' and \inputs each
% value file. A build with no values still compiles; every \val{} then prints
% a visible marker. \valuesindex lets the one-chapter build point elsewhere.
\providecommand{\valuesindex}{figures/values/all}
\IfFileExists{\valuesindex.tex}{\input{\valuesindex}}{%
  \AtBeginDocument{\PackageWarning{chaosbook}{No computed values found. Run `make numbers'.}}}

% ============================================================
%  PINNED VERSIONS --- single source of truth
%  Change these here AND in code/pyproject.toml; tools/check_structure.py
%  --pins fails when the two disagree. A chapter never types a version.
% ============================================================
\newcommand{\bookedition}{0.1}
\newcommand{\bookyear}{2026}
\newcommand{\pyver}{3.14}              % the minor the book targets
\newcommand{\pypatch}{3.14.7}          % the exact interpreter the code ran on
\newcommand{\uvver}{0.12.13}           % uv
\newcommand{\numpyver}{2.5.3}          % numpy
\newcommand{\matplotlibver}{3.11.2}    % matplotlib
\newcommand{\pytestver}{9.1.1}         % pytest
\newcommand{\ruffver}{0.16.9}          % ruff
% The prose licence is a name with a version in it, not a quantity: it is a
% macro so that neither edition typesets it as a decimal to be localised.
\newcommand{\proselicence}{CC~BY-NC-SA~4.0}

% ============================================================
%  INDEX
% ============================================================
\apptocmd{\theindex}{\raggedcolumns\raggedright}{}{%
  \PackageWarning{chaosbook}{Could not patch theindex for ragged setting}}
\providecommand*\see[2]{}
\renewcommand*\see[2]{\emph{\lblIndexSee} #1}
\providecommand*\seealso[2]{}
\renewcommand*\seealso[2]{\emph{\lblIndexSeeAlso} #1}

\makeindex
; nothing else differs between them at this level.
%  Every user-visible string lives in lang/en.tex or lang/pl.tex, so a label
%  that appears in one edition and not the other is a missing macro rather
%  than a silent divergence.
%
%  Lineage. This book is the third in a family and takes one thing from each
%  of the other two:
%
%    * from Mathematics from Zero for the AI Engineer: the bilingual single
%      source (\booklang, lang/, one body.tex), the rule that every number is
%      computed and reaches the page as \val{}, and the programmed-learning
%      method -- outcomes on entry, a question before every answer, a
%      misconception elicited before it is named, a self-check generated from
%      the outcomes, and checkpoint answers collected at the back;
%    * from Python for .NET Engineers: every listing is a FILE that CI runs,
%      every lab is a starter that must fail and a solution that must pass,
%      and a console transcript is a computed number that happens to be
%      verbatim.
%
%  What is new here is \plotfig. A book about chaos cannot be taught in boxes
%  and arrows: the bifurcation diagram and the Lorenz attractor ARE the
%  argument, so the plots are generated by code/plots/ from the same library
%  the listings use, per language, and never drawn by hand.
% ============================================================

% \booklang must be set by the main file. Default to English rather than
% failing, so a stray % ============================================================
%  preamble.tex --- styling, environments, macros
%
%  Shared by BOTH editions. main-en.tex and main-pl.tex each define \booklang
%  before \input{preamble}; nothing else differs between them at this level.
%  Every user-visible string lives in lang/en.tex or lang/pl.tex, so a label
%  that appears in one edition and not the other is a missing macro rather
%  than a silent divergence.
%
%  Lineage. This book is the third in a family and takes one thing from each
%  of the other two:
%
%    * from Mathematics from Zero for the AI Engineer: the bilingual single
%      source (\booklang, lang/, one body.tex), the rule that every number is
%      computed and reaches the page as \val{}, and the programmed-learning
%      method -- outcomes on entry, a question before every answer, a
%      misconception elicited before it is named, a self-check generated from
%      the outcomes, and checkpoint answers collected at the back;
%    * from Python for .NET Engineers: every listing is a FILE that CI runs,
%      every lab is a starter that must fail and a solution that must pass,
%      and a console transcript is a computed number that happens to be
%      verbatim.
%
%  What is new here is \plotfig. A book about chaos cannot be taught in boxes
%  and arrows: the bifurcation diagram and the Lorenz attractor ARE the
%  argument, so the plots are generated by code/plots/ from the same library
%  the listings use, per language, and never drawn by hand.
% ============================================================

% \booklang must be set by the main file. Default to English rather than
% failing, so a stray \input{preamble} still compiles and says what it did.
\providecommand{\booklang}{en}

\usepackage{etoolbox}   % loaded first: the language switch below needs it

% ---------- Page geometry ----------
% ONE format: A4, 12pt, ONE-SIDED. The book is read on a screen, usually with
% an editor open beside it, so there is no recto, no verso, no gutter and no
% blank leaf before a chapter. The measure is about seventy-five characters
% of prose. At \footnotesize a listing line of 79 characters fits without
% wrapping, and 79 is the width every file under code/ is held to by
% tools/check_structure.py --lines.
\usepackage[
  a4paper,
  left=3.1cm, right=3.1cm, top=2.6cm, bottom=3.0cm,
  head=26pt, headsep=0.8cm, footskip=1.4cm
]{geometry}

% ---------- Encoding & fonts ----------
\usepackage[T1]{fontenc}
\usepackage[utf8]{inputenc}
\IfFileExists{lmodern.sty}{\usepackage{lmodern}}{}
% newtxtext takes its TS1 symbols from TeX Gyre Termes (Debian: tex-gyre).
% A machine with newtx and without tex-gyre dies on the copyright page. CI's
% full TeX Live is the reference; locally install tex-gyre as well.
\IfFileExists{newtxtext.sty}{\usepackage{newtxtext}}{}
% amsmath before newtxmath; amssymb only when newtxmath is absent, because
% loading both is a fatal clash (\Bbbk already defined) that is invisible on
% a machine without newtx.
\usepackage{amsmath}
\IfFileExists{newtxmath.sty}{\usepackage{newtxmath}}{\usepackage{amssymb}}
\IfFileExists{inconsolata.sty}{\usepackage[varqu,scaled=0.95]{inconsolata}}{}
% Without lmodern, microtype runs with font expansion off, which is what
% resolves a marginal overfull line -- so a container short of lmodern
% disagrees with CI about line breaks. Install lmodern locally too.
\IfFileExists{lmodern.sty}{\usepackage{microtype}}{\usepackage[expansion=false]{microtype}}

% upquote: in T1 the input ' is a right single quote, and a listing that sets
% f('x') with curly quotes is a SyntaxError when typed back in.
\IfFileExists{upquote.sty}{\usepackage{upquote}}{}
% And its twin: in T1 the pair -- is an en-dash ligature in the typewriter
% family too, so \code{--flag} would print a flag nobody can type. Scoped to
% tt*, so prose keeps its own dashes.
\DisableLigatures[-]{encoding = T1, family = tt*}

% ---------- Language ----------
% babel with a missing language file is FATAL, so each .ldf is probed first.
\IfFileExists{babel.sty}{%
  \IfFileExists{polish.ldf}{%
    \IfFileExists{british.ldf}{%
      \ifdefstring{\booklang}{pl}%
        {\usepackage[british,main=polish]{babel}}%
        {\usepackage[polish,main=british]{babel}}%
    }{\usepackage[polish]{babel}}%
  }{%
    \IfFileExists{british.ldf}{\usepackage[british]{babel}}{}%
  }%
}{}
% babel-polish makes " an active shorthand, which is a category-code accident
% waiting to happen inside listings. Polish quotation marks come from
% \enquote, which csquotes sets per language.
\ifdefstring{\booklang}{pl}{\AtBeginDocument{\shorthandoff{"}}}{}

% ---------- Core packages ----------
\usepackage{graphicx}
\usepackage{xcolor}
\usepackage{booktabs}
\usepackage{tabularx}
\usepackage{longtable}
\usepackage{array}
\usepackage{enumitem}
\usepackage{listings}
\usepackage[most]{tcolorbox}
\usepackage{fancyhdr}
\usepackage{titlesec}
\usepackage{caption}
\usepackage{imakeidx}
\IfFileExists{csquotes.sty}{\usepackage[autostyle=true]{csquotes}}{%
  \providecommand{\enquote}[1]{``##1''}}
\usepackage[hidelinks]{hyperref}

% ---------- Palette ----------
% The family's palette, so the three books read as a set. Two colours are
% this book's own: the lab (green) and the stop-and-think box (violet).
\definecolor{mblue}{HTML}{1F4E79}
\definecolor{mteal}{HTML}{0E7C7B}
\definecolor{mamber}{HTML}{B26A00}
\definecolor{mred}{HTML}{9B2C2C}
\definecolor{mviolet}{HTML}{7B4397}
\definecolor{mgreen}{HTML}{3B7A3B}
\definecolor{mgrey}{HTML}{5A5A5A}
\definecolor{mlight}{HTML}{F4F6F8}
\definecolor{mfaint}{HTML}{FBFCFD}
\definecolor{codebg}{HTML}{FAFAFA}
\definecolor{codeframe}{HTML}{D8DEE4}
\definecolor{kw}{HTML}{0B5394}
\definecolor{str}{HTML}{9B2C2C}
\definecolor{cmt}{HTML}{4F7A4F}
\definecolor{typ}{HTML}{0E7C7B}
\definecolor{dec}{HTML}{7B4397}

\hypersetup{
  colorlinks=true,
  linkcolor=mblue,
  urlcolor=mteal,
  citecolor=mblue,
  pdfauthor={Konrad Cinkusz}
}

% \ifpl{Polish}{English} -- for the handful of places where a string is
% structural rather than content. Robust, because a fragile \ifpl inside a
% moving argument fails on the next run in exactly one of the two editions.
\DeclareRobustCommand{\ifpl}[2]{\ifdefstring{\booklang}{pl}{#1}{#2}}
\makeatletter
\pdfstringdefDisableCommands{%
  \ifdefstring{\booklang}{pl}{\let\ifpl\@firstoftwo}{\let\ifpl\@secondoftwo}%
  % hyperref drops \color from a bookmark and keeps its ARGUMENT, so a title
  % carrying \code{} would bookmark as "mtealx". Gobble the colour name.
  \def\color#1{}%
}
\makeatother

% ---------- Language strings ----------
\input{lang/\booklang}

% The PDF's subject, keywords and language read \lblPdf... from the language
% file, so they cannot sit in the colour block above, which runs first. Two
% blocks is the ordering, not an oversight.
\hypersetup{
  pdfsubject={\lblPdfSubject},
  pdfkeywords={\lblPdfKeywords},
  pdflang={\lblPdfLang}
}

% ---------- Headings ----------
\titleformat{\chapter}[display]
  {\normalfont\huge\bfseries\color{mblue}\raggedright}
  {\normalfont\normalsize\color{mgrey}\MakeUppercase{\chaptertitlename}\ \thechapter}
  {8pt}{\Huge\raggedright}
\titlespacing*{\chapter}{0pt}{10pt}{28pt}
\titleformat{\section}{\normalfont\Large\bfseries\color{mblue}\raggedright}{\thesection}{0.8em}{}
\titleformat{\subsection}{\normalfont\large\bfseries\color{mgrey}\raggedright}{\thesubsection}{0.7em}{}

\setcounter{tocdepth}{1}
\setcounter{secnumdepth}{2}

\pagestyle{fancy}
\fancyhf{}
\fancyhead[L]{\small\nouppercase{\leftmark}}
\fancyhead[R]{\small\thepage}
\renewcommand{\headrulewidth}{0.4pt}
\fancypagestyle{plain}{\fancyhf{}\fancyfoot[R]{\small\thepage}%
  \renewcommand{\headrulewidth}{0pt}}
% After \pagestyle{fancy}, which installs its own \chaptermark.
\renewcommand{\chaptermark}[1]{\markboth{\thechapter.\ #1}{}}

% ---------- Mathematics ----------
% The handful of notations the book uses, as macros so both editions set the
% same symbol. The book writes \ln for the natural logarithm and \log_{10}
% where the base is ten; a bare \log is ambiguous between the two readerships
% and is not used.
\newcommand{\R}{\mathbb{R}}
\newcommand{\dd}{\,\mathrm{d}}
\newcommand{\abs}[1]{\left\lvert #1\right\rvert}
\newcommand{\iu}{\mathrm{i}}

% ============================================================
%  CODE LISTINGS
%  Every listing in a chapter is a FILE under code/, pulled in with \pyfile or
%  \pyregion, and CI runs it on the pinned interpreter. An in-source
%  \begin{python} block is allowed for a fragment of three or four lines that
%  shows a shape rather than runs.
% ============================================================
\lstdefinelanguage{pythonx}{
  language=Python,
  morekeywords={async,await,match,case,nonlocal,yield,type},
  morekeywords=[2]{
    orbit,iterate,logistic,rk4_step,euler_step,integrate,lorenz,henon,
    lyapunov,separation,escape_time,box_count
  },
  sensitive=true
}

\lstdefinestyle{bookcode}{
  basicstyle=\ttfamily\footnotesize,
  keywordstyle=\color{kw}\bfseries,
  keywordstyle=[2]\color{typ},
  stringstyle=\color{str},
  % Colour only: inconsolata has no italic, and an italic comment style is
  % silently substituted and logged on every build.
  commentstyle=\color{cmt},
  numbers=left,
  numberstyle=\ttfamily\tiny\color{mgrey},
  numbersep=8pt,
  backgroundcolor=\color{codebg},
  frame=single,
  rulecolor=\color{codeframe},
  framesep=6pt,
  breaklines=true,
  breakatwhitespace=false,
  postbreak=\mbox{\textcolor{mgrey}{$\hookrightarrow$}\space},
  showstringspaces=false,
  tabsize=4,
  columns=fullflexible,
  keepspaces=true,
  captionpos=b,
  aboveskip=1.2em,
  belowskip=1.2em,
  % Do NOT add a mapping for U+00A0: it aborts the run with a message that
  % names neither the character nor this line. Listings are ASCII, and
  % parity's C13 checks it.
  literate={ł}{{\l}}1 {Ł}{{\L}}1 {ą}{{\k a}}1 {ę}{{\k e}}1
           {ś}{{\'s}}1 {ć}{{\'c}}1 {ż}{{\.z}}1 {ź}{{\'z}}1 {ó}{{\'o}}1 {ń}{{\'n}}1
           {Ą}{{\k A}}1 {Ę}{{\k E}}1 {Ś}{{\'S}}1 {Ć}{{\'C}}1
           {Ż}{{\.Z}}1 {Ź}{{\'Z}}1 {Ó}{{\'O}}1 {Ń}{{\'N}}1
           {—}{{-{}-}}2 {–}{{-}}1 {‘}{{'}}1 {’}{{'}}1
           {“}{{"}}1 {”}{{"}}1 {°}{{\textdegree}}1
}
\lstset{style=bookcode}

\lstnewenvironment{python}[1][]
  {\lstset{language=pythonx,style=bookcode,#1}}
  {}
\lstnewenvironment{shellcmd}[1][]
  {\lstset{language=bash,style=bookcode,numbers=none,keywordstyle=\color{mteal}\bfseries,#1}}
  {}

% The path is printed above every file listing, in the form the reader opens
% in the editor. \nobreak keeps it on the page of the listing it names.
% The path must not be stranded at the foot of a page with its listing
% overleaf: \nobreak alone does not hold, because listings opens with its own
% vertical material. \needspace asks for room for the path and the first few
% lines of code, and turns the page first when there is not.
% \Needspace (capital N) measures \pagegoal-\pagetotal exactly; lower-case
% \needspace works through glue and penalties and was measured, by the
% Chapter 11 pass, turning pages early and leaving gaps. It must be used
% between paragraphs, which \listingpath's own \par guarantees. Chapters use
% \Needspace* too, so the fallback for a machine without the package defines
% both forms (the star swallowed with its argument).
\makeatletter
\IfFileExists{needspace.sty}{\usepackage{needspace}}{%
  \providecommand{\needspace}[1]{}%
  \providecommand{\Needspace}{\@ifstar\@gobble\@gobble}}
\makeatother
\newcommand{\listingpath}[1]{%
  \par\medskip\Needspace{6\baselineskip}\noindent
  {\ttfamily\footnotesize\color{mgrey}\detokenize{#1}}%
  \par\nobreak}

% Usage: \pyfile{code/ch05/cobweb.py}{Caption}{lst:label}
\newcommand{\pyfile}[3]{%
  \listingpath{#1}%
  \lstinputlisting[language=pythonx,style=bookcode,aboveskip=0.3em,
    caption={#2},label={#3}]{#1}%
}

% A named region of a file, delimited by `# --8<-- [start:name]` and
% `# --8<-- [end:name]`. Matched through listings' range prefix and suffix
% keys -- the other obvious form matches nothing and prints the WHOLE file
% without a warning (measured in the sibling book). A region name is a word,
% never a number. The line numbers printed are the file's own.
% Usage: \pyregion{code/ch05/cobweb.py}{step}{Caption}{lst:label}
\lstset{
  rangebeginprefix=\#\ --8<--\ [start:, rangebeginsuffix=],
  rangeendprefix=\#\ --8<--\ [end:, rangeendsuffix=],
  includerangemarker=false}
\newcommand{\pyregion}[4]{%
  \listingpath{#1}%
  \lstinputlisting[language=pythonx,style=bookcode,aboveskip=0.3em,
    linerange={#2-#2},caption={#3},label={#4}]{#1}%
}

% A program's output, written by code/measure/transcripts.py. There is no
% typed form and there must not be: \val{} does not expand in a verbatim
% body, and a transcript nobody ran looks exactly like one that was.
\newcommand{\transcript}[1]{%
  \par\smallskip
  \IfFileExists{figures/transcripts/#1.txt}{%
    \lstinputlisting[style=bookcode,numbers=none,basicstyle=\ttfamily\footnotesize,
      keywordstyle=\color{black},backgroundcolor=\color{white},
      language={}]{figures/transcripts/#1.txt}%
  }{%
    \begin{tcolorbox}[enhanced, sharp corners, width=\linewidth,
      colback=mlight, colframe=mgrey, boxrule=0.6pt,
      left=8pt, right=8pt, top=6pt, bottom=6pt]
      {\bfseries\color{mgrey}\small \lblTranscriptMissing}\\[2pt]
      {\ttfamily\scriptsize\color{mgrey} figures/transcripts/#1.txt}
    \end{tcolorbox}%
  }%
  \par\smallskip
}

% Inline literals. Underscores inside these must be written \_.
\newcommand{\code}[1]{\texttt{\small #1}}
\newcommand{\pkg}[1]{\texttt{\small\color{mblue}#1}}
\newcommand{\api}[1]{\texttt{\small\color{mteal}#1}}

% ============================================================
%  ADMONITIONS
%  A small, fixed vocabulary. `tinted' may be read in passing; `outlined' is
%  a block the reader must stop at. A third treatment would mean the first two
%  are not carrying their meaning.
% ============================================================
\tcbset{
  cfz tinted/.style={
    enhanced, breakable, sharp corners,
    lines before break=4,
    colback=mlight, boxrule=0pt, leftrule=3pt,
    left=8pt, right=8pt, top=6pt, bottom=6pt,
    before skip=13pt, after skip=13pt,
    fonttitle=\bfseries\small,
    coltitle=black, colbacktitle=mlight,
    attach title to upper={\par\smallskip},
  },
  cfz outlined/.style={
    enhanced, breakable, sharp corners,
    % A box broken straight after its title is a heading stranded at the foot
    % of a page with its question overleaf. Four lines is title plus three.
    lines before break=4,
    colback=white, boxrule=0.6pt,
    left=8pt, right=8pt, top=6pt, bottom=6pt,
    before skip=13pt, after skip=13pt,
    fonttitle=\bfseries\small,
    coltitle=black, colbacktitle=white,
    attach title to upper={\par\smallskip},
  },
}

% Read in passing.
\newtcolorbox{note}[1][]{cfz tinted, colframe=mblue, title={\lblNote}, #1}
\newtcolorbox{warning}[1][]{cfz tinted, colframe=mred, title={\lblWarning}, #1}
% Where you meet this in the world. If it cannot name a specific system,
% instrument, paper or event, it is prose and not this box.
\newtcolorbox{worldbox}[1][]{cfz tinted, colframe=mteal, title={\lblWorld}, #1}
% What is stated and not proved here, and where the proof lives.
\newtcolorbox{rigourbox}[1][]{cfz tinted, colframe=mgrey, title={\lblRigour}, #1}
% The companion library the reader builds, under code/src/chaoslab/.
\newtcolorbox{projectbox}[1][]{cfz tinted, colframe=mamber, title={\lblProject}, #1}

% Stop and read.
% The misconception, named AFTER the reader has been asked the question it
% gets wrong. Appendix B is the catalogue; a trap box is one entry in place.
\newtcolorbox{trapbox}[1][]{cfz outlined, colframe=mred, title={\lblTrap}, #1}
% A listing that has not been run on the pinned versions. Counted by
% `make debt'; nothing ships to a reader with one attached.
\newtcolorbox{verifybox}[1][]{cfz outlined, colframe=mgrey, title={\lblVerify}, #1}
\newtcolorbox{labbox}[1][]{cfz outlined, colframe=mgreen, title={\lblLab}, #1}

% ============================================================
%  THE METHOD --- questions before answers
%  ---------------------------------------------------------
%  The unit of teaching is a question the reader answers BEFORE reading the
%  answer. That is retrieval practice and errorful generation, the two
%  findings the programmed-learning format is defended by, and it is the one
%  thing the sibling mathematics book does that this one keeps whole.
%
%  \begin{think} ... \end{think}   the question, numbered per chapter, ending
%                                  with the instruction to answer before
%                                  reading on;
%  \reveal{short answer}           a centred box with the answer, which is the
%                                  first thing after the question -- the edge
%                                  the reader covers with a hand, or scrolls
%                                  to;
%  \begin{revealblock} ... \end{revealblock}  the same, for a worked answer.
%
%  tools/check_structure.py --think fails when a think box is not followed by
%  a reveal before the next think box or section, and when a reveal answers
%  nothing. That is a property of the source, so it is checked mechanically.
%
%  The instruction says "before you read on", never "before you turn over":
%  the two editions paginate differently, so the second is true in one of them.
% ============================================================
\newcounter{think}[chapter]
\renewcommand{\thethink}{\thechapter.\arabic{think}}
\newtcolorbox{thinkbox}[1][]{cfz outlined, colframe=mviolet, #1}
\NewDocumentEnvironment{think}{}{%
  \refstepcounter{think}%
  \begin{thinkbox}[title={\lblThink~\thethink}]%
}{%
  \par\smallskip
  {\raggedleft\itshape\small\color{mviolet}\lblThinkCue\par}%
  \end{thinkbox}%
}
\tcbset{
  cfz answer/.style={
    enhanced, sharp corners,
    colback=mfaint, colframe=mviolet!55, boxrule=0.5pt,
    before skip=4pt, after skip=10pt,
  },
}
\newcommand{\reveal}[1]{%
  \par\nobreak
  \begin{tcolorbox}[cfz answer, unbreakable,
    left=8pt, right=8pt, top=5pt, bottom=5pt,
    width=0.86\linewidth, center, halign=flush center]
    {\small\bfseries\color{mviolet}\lblAnswerMark}\enspace #1
  \end{tcolorbox}%
  \nobreak\noindent\ignorespaces
}
\NewDocumentEnvironment{revealblock}{}{%
  \par\nopagebreak
  \begin{tcolorbox}[cfz answer, breakable,
    left=10pt, right=10pt, top=6pt, bottom=6pt]
  {\small\bfseries\color{mviolet}\lblAnswerMark}\par\smallskip
}{%
  \end{tcolorbox}%
  \nopagebreak\par\noindent\ignorespaces
}

% ---- Outcomes on entry, and "Can you?" generated from them ----------------
% The self-check at the end of a chapter IS the list of outcomes at its
% start, one for one, because it is built from the same tokens: each
% \outcome both prints its item and appends a row to \cfz@canrows. A \loop
% inside a tabularx body does not survive alignment scanning, which is why
% the rows are accumulated at declaration time (the sibling book's finding).
\makeatletter
\newcommand{\cfz@canrows}{}
\newcommand{\cfz@ratebox}{%
  \fbox{\rule{0pt}{0.9ex}\rule{0.9ex}{0pt}}}
\newcommand{\cfz@rates}{%
  {\footnotesize\color{mgrey}%
   1\,\cfz@ratebox\ 2\,\cfz@ratebox\ 3\,\cfz@ratebox\ 4\,\cfz@ratebox\ 5\,\cfz@ratebox}}
\NewDocumentEnvironment{outcomes}{}{%
  \gdef\cfz@canrows{}%
  \begin{tcolorbox}[cfz tinted, colframe=mblue, title={\lblOutcomes}]%
  \lblOutcomesLead
  \begin{itemize}[leftmargin=1.4em, itemsep=2pt, topsep=4pt]%
}{%
  \end{itemize}%
  \end{tcolorbox}%
}
\newcommand{\outcome}[1]{%
  \item #1%
  \gappto\cfz@canrows{#1 & \cfz@rates \\ \addlinespace[3pt]}%
}
\newcommand{\canyou}{%
  \section*{\lblCanYou}%
  \phantomsection\addcontentsline{toc}{section}{\lblCanYou}%
  \noindent\lblCanYouIntro\par\medskip
  \ifdefempty{\cfz@canrows}{%
    \noindent{\color{mred}\lblNoOutcomes}\par}{%
  {\small
  \noindent\begin{tabularx}{\linewidth}{@{}>{\raggedright\arraybackslash}X r@{}}
    \toprule
    \cfz@canrows
    \bottomrule
  \end{tabularx}\par}}%
}
\makeatother

% ---- Summary -------------------------------------------------------------
\newtcolorbox{summarybox}[1][]{cfz tinted, colframe=mblue, title={\lblSummary}, #1}

% ---- Checkpoint questions, answered at the back --------------------------
% Answers are written AT THE QUESTION with \answerto{} and printed in
% Appendix A. They are carried there by a global token list rather than an
% auxiliary file, because an answer contains mathematics and \val{}, and
% material written to an .aux is expanded at shipout, where fragile commands
% break with errors naming neither the answer nor the chapter.
% \include is sequential within a run, so the list is complete by the time
% the appendix is typeset.
\makeatletter
\newcommand{\cfz@answers}{}
\newcommand{\cfz@lastansch}{0}
\NewDocumentEnvironment{checkpoint}{}{%
  \section*{\lblCheckpoint}%
  \phantomsection\addcontentsline{toc}{section}{\lblCheckpoint}%
  \noindent\lblCheckpointIntro\par
  \begin{enumerate}[label=\textbf{\arabic*.}, leftmargin=2em, itemsep=4pt]%
}{%
  \end{enumerate}%
}
\newcommand{\answerto}[1]{%
  \ifnum\value{chapter}=\cfz@lastansch\relax\else
    \xdef\cfz@lastansch{\arabic{chapter}}%
    \xappto\cfz@answers{\noexpand\cfz@anschapter{\thechapter}}%
  \fi
  \xappto\cfz@answers{\noexpand\cfz@ansitem{\arabic{enumi}}}%
  \gappto\cfz@answers{{#1}}%
}
\newcommand{\cfz@anschapter}[1]{%
  \par\medskip\noindent{\bfseries\color{mblue}\chaptername~#1}\par\nobreak\smallskip}
\newcommand{\cfz@ansitem}[2]{%
  \par\noindent\hangindent=2em\hangafter=1
  \makebox[2em][l]{\textbf{#1.}}#2\par}
\newcommand{\answersbody}{%
  \ifdefempty{\cfz@answers}{\noindent{\color{mred}\lblNoAnswers}\par}{%
    {\raggedright\cfz@answers}}%
}
\makeatother

% ============================================================
%  LABS
%  \begin{lab}{key}{title} ... \end{lab}
%
%  A lab is not a paragraph but a FILE the reader opens: a starter under
%  code/labs/chNN/<key>.py that must FAIL its test, a reference solution under
%  code/labs/chNN/solutions/, and test_<key>.py beside it. The box prints the
%  paths and the one command that runs the test. An environment rather than a
%  macro, because a macro argument cannot hold a verbatim fragment.
%  tools/check_structure.py --labs fails when a file is missing or the key's
%  chapter number disagrees with its chapter.
% ============================================================
\newcounter{lab}[chapter]
\renewcommand{\thelab}{\thechapter.\arabic{lab}}
\makeatletter
\newcommand{\labdir}{ch\two@digits{\value{chapter}}}
\newcommand{\listoflabs}{%
  \section*{\lblLabManifest}%
  \phantomsection\addcontentsline{toc}{section}{\lblLabManifest}%
  \begingroup\c@tocdepth=2\relax\@starttoc{lab}\endgroup%
}
\makeatother
\NewDocumentEnvironment{lab}{mm}{%
  \refstepcounter{lab}%
  \addcontentsline{lab}{subsection}{\protect\numberline{\thelab}\texttt{\protect\detokenize{#1}} --- #2}%
  \begin{labbox}[title={\lblLab~\thelab\ \dash{} #2}]%
}{%
  \par\medskip
  {\small\setlength{\parindent}{0pt}%
  \lblLabStarter: \texttt{code/labs/\labdir/\detokenize{#1}.py}\\
  \lblLabTest: \texttt{code/labs/\labdir/test\_\detokenize{#1}.py}\\
  \lblLabRun: \texttt{cd code \&\& uv run pytest -k \detokenize{#1}}\par}%
  \end{labbox}%
}

% ============================================================
%  FIGURES --- plots and diagrams
%  ---------------------------------------------------------
%  \plotfig{key}{caption}{one-line manifest description}
%    A PLOT, generated by code/plots/<chapter>.py from the same library the
%    listings use, into figures/plots/<lang>/<key>.pdf. Per language, because
%    an axis label in the wrong language is the half-translation that makes a
%    bilingual book look machine-made, and because the Polish edition owes a
%    decimal comma on its tick labels too. Build output, gitignored.
%
%  \mermaidfig{key}{caption}{one-line manifest description}
%    A DIAGRAM of an idea rather than of data, from
%    figures/mermaid/<lang>/<key>.mmd, rendered by `make diagrams'.
%
%  If the rendered PDF is absent either macro prints a visible marker in the
%  flow rather than a float, so a missing figure cannot be mistaken for a
%  layout choice. \phantomsection before \addcontentsline is load-bearing:
%  without it the manifest entry records whatever anchor preceded the figure,
%  which after a listing is a line that was never printed.
% ============================================================
\makeatletter
\newcommand{\listofplots}{%
  \section*{\lblPlotManifest}%
  \phantomsection\addcontentsline{toc}{section}{\lblPlotManifest}%
  \begingroup\c@tocdepth=2\relax\@starttoc{plt}\endgroup%
}
\newcommand{\listofdiagrams}{%
  \section*{\lblDiagramManifest}%
  \phantomsection\addcontentsline{toc}{section}{\lblDiagramManifest}%
  \begingroup\c@tocdepth=2\relax\@starttoc{dgm}\endgroup%
}
\makeatother

\newcommand{\figmissing}[3]{% #1 message, #2 path, #3 caption
  \par\medskip
  \begin{tcolorbox}[enhanced, breakable, width=\linewidth, sharp corners,
    colback=white, colframe=mgrey, boxrule=0.8pt,
    left=8pt, right=8pt, top=8pt, bottom=8pt]
    {\bfseries\color{mgrey}\small #1}\\[4pt]
    {\ttfamily\scriptsize\color{mgrey} #2}
  \end{tcolorbox}%
  \captionof{figure}{#3}%
  \par\medskip
}

\newcommand{\plotfig}[3]{%
  \phantomsection%
  \addcontentsline{plt}{subsection}{\protect\numberline{}\texttt{#1} --- #3}%
  \IfFileExists{figures/plots/\booklang/#1.pdf}{%
    \begin{figure}[htbp]
    \centering
    \includegraphics[width=\linewidth,height=0.45\textheight,keepaspectratio]{figures/plots/\booklang/#1.pdf}%
    \caption{#2}\label{fig:#1}
    \end{figure}%
  }{%
    \figmissing{\lblPlotNotRendered}{figures/plots/\booklang/#1.pdf}{#2}\label{fig:#1}%
  }%
}

\newcommand{\mermaidfig}[3]{%
  \phantomsection%
  \addcontentsline{dgm}{subsection}{\protect\numberline{}\texttt{#1.mmd} --- #3}%
  \IfFileExists{figures/diagrams/\booklang/#1.pdf}{%
    \begin{figure}[htbp]
    \centering
    \includegraphics[width=0.95\linewidth,height=0.42\textheight,keepaspectratio]{figures/diagrams/\booklang/#1.pdf}%
    \caption{#2}\label{fig:#1}
    \end{figure}%
  }{%
    \figmissing{\lblNotRendered}{figures/mermaid/\booklang/#1.mmd}{#2}\label{fig:#1}%
  }%
}

% ============================================================
%  APPENDIX B: the misreadings, generated by tools/gen_traps.py
%
%  \trapchapter{Chapter~\ref{ch:x}}{title}    a heading per chapter
%  \trapentry{n}{habit}{truth}{\ref{sec:x}}  one misreading
%
%  Set ragged right: an entry is mostly short sentences with \code{} spans in
%  them, and a justified line full of unbreakable runs is an overfull box
%  waiting for the edition with the longer words.
% ============================================================
\newcommand{\trapchapter}[2]{%
  \par\medskip
  {\large\bfseries\color{mblue}#1\quad\normalfont\itshape\color{mgrey}#2\par}%
  \nopagebreak\smallskip}
\newcommand{\trapentry}[4]{%
  \par\noindent
  \begingroup\raggedright\hangindent=2.4em\hangafter=1
  \makebox[2.4em][l]{\textbf{\color{mred}#1}}%
  \textbf{\enquote{#2}}\enspace #3\enspace{\color{mgrey}\S\,#4}\par
  \endgroup\smallskip}

% ============================================================
%  CHAPTER STUBS
% ============================================================
\newcommand{\chapterstub}[1]{%
  \begin{tcolorbox}[
    enhanced, breakable, width=\linewidth, sharp corners,
    colback=mlight, colframe=mred, boxrule=1pt,
    borderline={0.6pt}{3pt}{mred, dashed},
    left=12pt, right=12pt, top=12pt, bottom=12pt]
    {\bfseries\color{mred}\large \lblNotWritten}\\[8pt]
    \small\raggedright #1
  \end{tcolorbox}%
}

% ============================================================
%  COMPUTED VALUES
%  A number the reader cannot do in their head is computed by a script under
%  code/measure/, written to figures/values/<chapter>.tex as
%      \pyval{ch06.steps.apart}{34}
%  and reaches the page as \val{ch06.steps.apart}. The book contains
%  references, not digits; `make verify' fails when a committed value no
%  longer matches its script. \val applies the edition's decimal marker, which
%  is how the Polish edition gets its comma without anybody remembering.
% ============================================================
% The comma is BRACED: siunitx sets numbers in maths mode, where a bare comma
% is punctuation and gets a thin space after it ("9, 81"). {,} is an ordinary
% symbol, which is what a decimal marker is.
\newcommand{\bookdecimalmarker}{\ifpl{{,}}{.}}
\IfFileExists{siunitx.sty}{%
  \usepackage{siunitx}%
  \sisetup{
    group-digits = integer,
    group-minimum-digits = 5,
    group-separator = {\,},
    output-decimal-marker = {\bookdecimalmarker},
    exponent-product = \times,
    print-unity-mantissa = false
  }%
  \newcommand{\booknum}[1]{\num{##1}}%
}{%
  \newcommand{\booknum}[1]{##1}%
  \providecommand{\num}[1]{##1}%
}

\makeatletter
\newcommand{\pyvalues}{}
\newcommand{\pyval}[2]{\csgdef{pyval@#1}{#2}\listxadd{\pyvalues}{#1}}
\newcommand{\pyvaltext}[2]{%
  \csgdef{pyval@#1}{#2}\csgdef{pyvaltext@#1}{}\listxadd{\pyvalues}{#1}}
% \num on a non-numeric body is a FATAL siunitx error, so \val refuses a text
% value by name rather than letting siunitx fail on it.
\newcommand{\val}[1]{%
  \ifcsdef{pyval@#1}{%
    \ifcsdef{pyvaltext@#1}{%
      \PackageError{chaosbook}{Value `#1' is not a number}%
        {Use \string\valtext{#1} instead.}%
    }{}%
    \booknum{\csuse{pyval@#1}}%
  }{\textcolor{mred}{\textbf{[\lblValueMissing: #1]}}}%
}
\newcommand{\valtext}[1]{%
  \ifcsdef{pyval@#1}{\csuse{pyval@#1}}%
    {\textcolor{mred}{\textbf{[\lblValueMissing: #1]}}}%
}
\makeatother

% figures/values/all.tex is generated by `make numbers' and \inputs each
% value file. A build with no values still compiles; every \val{} then prints
% a visible marker. \valuesindex lets the one-chapter build point elsewhere.
\providecommand{\valuesindex}{figures/values/all}
\IfFileExists{\valuesindex.tex}{\input{\valuesindex}}{%
  \AtBeginDocument{\PackageWarning{chaosbook}{No computed values found. Run `make numbers'.}}}

% ============================================================
%  PINNED VERSIONS --- single source of truth
%  Change these here AND in code/pyproject.toml; tools/check_structure.py
%  --pins fails when the two disagree. A chapter never types a version.
% ============================================================
\newcommand{\bookedition}{0.1}
\newcommand{\bookyear}{2026}
\newcommand{\pyver}{3.14}              % the minor the book targets
\newcommand{\pypatch}{3.14.7}          % the exact interpreter the code ran on
\newcommand{\uvver}{0.12.13}           % uv
\newcommand{\numpyver}{2.5.3}          % numpy
\newcommand{\matplotlibver}{3.11.2}    % matplotlib
\newcommand{\pytestver}{9.1.1}         % pytest
\newcommand{\ruffver}{0.16.9}          % ruff
% The prose licence is a name with a version in it, not a quantity: it is a
% macro so that neither edition typesets it as a decimal to be localised.
\newcommand{\proselicence}{CC~BY-NC-SA~4.0}

% ============================================================
%  INDEX
% ============================================================
\apptocmd{\theindex}{\raggedcolumns\raggedright}{}{%
  \PackageWarning{chaosbook}{Could not patch theindex for ragged setting}}
\providecommand*\see[2]{}
\renewcommand*\see[2]{\emph{\lblIndexSee} #1}
\providecommand*\seealso[2]{}
\renewcommand*\seealso[2]{\emph{\lblIndexSeeAlso} #1}

\makeindex
 still compiles and says what it did.
\providecommand{\booklang}{en}

\usepackage{etoolbox}   % loaded first: the language switch below needs it

% ---------- Page geometry ----------
% ONE format: A4, 12pt, ONE-SIDED. The book is read on a screen, usually with
% an editor open beside it, so there is no recto, no verso, no gutter and no
% blank leaf before a chapter. The measure is about seventy-five characters
% of prose. At \footnotesize a listing line of 79 characters fits without
% wrapping, and 79 is the width every file under code/ is held to by
% tools/check_structure.py --lines.
\usepackage[
  a4paper,
  left=3.1cm, right=3.1cm, top=2.6cm, bottom=3.0cm,
  head=26pt, headsep=0.8cm, footskip=1.4cm
]{geometry}

% ---------- Encoding & fonts ----------
\usepackage[T1]{fontenc}
\usepackage[utf8]{inputenc}
\IfFileExists{lmodern.sty}{\usepackage{lmodern}}{}
% newtxtext takes its TS1 symbols from TeX Gyre Termes (Debian: tex-gyre).
% A machine with newtx and without tex-gyre dies on the copyright page. CI's
% full TeX Live is the reference; locally install tex-gyre as well.
\IfFileExists{newtxtext.sty}{\usepackage{newtxtext}}{}
% amsmath before newtxmath; amssymb only when newtxmath is absent, because
% loading both is a fatal clash (\Bbbk already defined) that is invisible on
% a machine without newtx.
\usepackage{amsmath}
\IfFileExists{newtxmath.sty}{\usepackage{newtxmath}}{\usepackage{amssymb}}
\IfFileExists{inconsolata.sty}{\usepackage[varqu,scaled=0.95]{inconsolata}}{}
% Without lmodern, microtype runs with font expansion off, which is what
% resolves a marginal overfull line -- so a container short of lmodern
% disagrees with CI about line breaks. Install lmodern locally too.
\IfFileExists{lmodern.sty}{\usepackage{microtype}}{\usepackage[expansion=false]{microtype}}

% upquote: in T1 the input ' is a right single quote, and a listing that sets
% f('x') with curly quotes is a SyntaxError when typed back in.
\IfFileExists{upquote.sty}{\usepackage{upquote}}{}
% And its twin: in T1 the pair -- is an en-dash ligature in the typewriter
% family too, so \code{--flag} would print a flag nobody can type. Scoped to
% tt*, so prose keeps its own dashes.
\DisableLigatures[-]{encoding = T1, family = tt*}

% ---------- Language ----------
% babel with a missing language file is FATAL, so each .ldf is probed first.
\IfFileExists{babel.sty}{%
  \IfFileExists{polish.ldf}{%
    \IfFileExists{british.ldf}{%
      \ifdefstring{\booklang}{pl}%
        {\usepackage[british,main=polish]{babel}}%
        {\usepackage[polish,main=british]{babel}}%
    }{\usepackage[polish]{babel}}%
  }{%
    \IfFileExists{british.ldf}{\usepackage[british]{babel}}{}%
  }%
}{}
% babel-polish makes " an active shorthand, which is a category-code accident
% waiting to happen inside listings. Polish quotation marks come from
% \enquote, which csquotes sets per language.
\ifdefstring{\booklang}{pl}{\AtBeginDocument{\shorthandoff{"}}}{}

% ---------- Core packages ----------
\usepackage{graphicx}
\usepackage{xcolor}
\usepackage{booktabs}
\usepackage{tabularx}
\usepackage{longtable}
\usepackage{array}
\usepackage{enumitem}
\usepackage{listings}
\usepackage[most]{tcolorbox}
\usepackage{fancyhdr}
\usepackage{titlesec}
\usepackage{caption}
\usepackage{imakeidx}
\IfFileExists{csquotes.sty}{\usepackage[autostyle=true]{csquotes}}{%
  \providecommand{\enquote}[1]{``##1''}}
\usepackage[hidelinks]{hyperref}

% ---------- Palette ----------
% The family's palette, so the three books read as a set. Two colours are
% this book's own: the lab (green) and the stop-and-think box (violet).
\definecolor{mblue}{HTML}{1F4E79}
\definecolor{mteal}{HTML}{0E7C7B}
\definecolor{mamber}{HTML}{B26A00}
\definecolor{mred}{HTML}{9B2C2C}
\definecolor{mviolet}{HTML}{7B4397}
\definecolor{mgreen}{HTML}{3B7A3B}
\definecolor{mgrey}{HTML}{5A5A5A}
\definecolor{mlight}{HTML}{F4F6F8}
\definecolor{mfaint}{HTML}{FBFCFD}
\definecolor{codebg}{HTML}{FAFAFA}
\definecolor{codeframe}{HTML}{D8DEE4}
\definecolor{kw}{HTML}{0B5394}
\definecolor{str}{HTML}{9B2C2C}
\definecolor{cmt}{HTML}{4F7A4F}
\definecolor{typ}{HTML}{0E7C7B}
\definecolor{dec}{HTML}{7B4397}

\hypersetup{
  colorlinks=true,
  linkcolor=mblue,
  urlcolor=mteal,
  citecolor=mblue,
  pdfauthor={Konrad Cinkusz}
}

% \ifpl{Polish}{English} -- for the handful of places where a string is
% structural rather than content. Robust, because a fragile \ifpl inside a
% moving argument fails on the next run in exactly one of the two editions.
\DeclareRobustCommand{\ifpl}[2]{\ifdefstring{\booklang}{pl}{#1}{#2}}
\makeatletter
\pdfstringdefDisableCommands{%
  \ifdefstring{\booklang}{pl}{\let\ifpl\@firstoftwo}{\let\ifpl\@secondoftwo}%
  % hyperref drops \color from a bookmark and keeps its ARGUMENT, so a title
  % carrying \code{} would bookmark as "mtealx". Gobble the colour name.
  \def\color#1{}%
}
\makeatother

% ---------- Language strings ----------
\input{lang/\booklang}

% The PDF's subject, keywords and language read \lblPdf... from the language
% file, so they cannot sit in the colour block above, which runs first. Two
% blocks is the ordering, not an oversight.
\hypersetup{
  pdfsubject={\lblPdfSubject},
  pdfkeywords={\lblPdfKeywords},
  pdflang={\lblPdfLang}
}

% ---------- Headings ----------
\titleformat{\chapter}[display]
  {\normalfont\huge\bfseries\color{mblue}\raggedright}
  {\normalfont\normalsize\color{mgrey}\MakeUppercase{\chaptertitlename}\ \thechapter}
  {8pt}{\Huge\raggedright}
\titlespacing*{\chapter}{0pt}{10pt}{28pt}
\titleformat{\section}{\normalfont\Large\bfseries\color{mblue}\raggedright}{\thesection}{0.8em}{}
\titleformat{\subsection}{\normalfont\large\bfseries\color{mgrey}\raggedright}{\thesubsection}{0.7em}{}

\setcounter{tocdepth}{1}
\setcounter{secnumdepth}{2}

\pagestyle{fancy}
\fancyhf{}
\fancyhead[L]{\small\nouppercase{\leftmark}}
\fancyhead[R]{\small\thepage}
\renewcommand{\headrulewidth}{0.4pt}
\fancypagestyle{plain}{\fancyhf{}\fancyfoot[R]{\small\thepage}%
  \renewcommand{\headrulewidth}{0pt}}
% After \pagestyle{fancy}, which installs its own \chaptermark.
\renewcommand{\chaptermark}[1]{\markboth{\thechapter.\ #1}{}}

% ---------- Mathematics ----------
% The handful of notations the book uses, as macros so both editions set the
% same symbol. The book writes \ln for the natural logarithm and \log_{10}
% where the base is ten; a bare \log is ambiguous between the two readerships
% and is not used.
\newcommand{\R}{\mathbb{R}}
\newcommand{\dd}{\,\mathrm{d}}
\newcommand{\abs}[1]{\left\lvert #1\right\rvert}
\newcommand{\iu}{\mathrm{i}}

% ============================================================
%  CODE LISTINGS
%  Every listing in a chapter is a FILE under code/, pulled in with \pyfile or
%  \pyregion, and CI runs it on the pinned interpreter. An in-source
%  \begin{python} block is allowed for a fragment of three or four lines that
%  shows a shape rather than runs.
% ============================================================
\lstdefinelanguage{pythonx}{
  language=Python,
  morekeywords={async,await,match,case,nonlocal,yield,type},
  morekeywords=[2]{
    orbit,iterate,logistic,rk4_step,euler_step,integrate,lorenz,henon,
    lyapunov,separation,escape_time,box_count
  },
  sensitive=true
}

\lstdefinestyle{bookcode}{
  basicstyle=\ttfamily\footnotesize,
  keywordstyle=\color{kw}\bfseries,
  keywordstyle=[2]\color{typ},
  stringstyle=\color{str},
  % Colour only: inconsolata has no italic, and an italic comment style is
  % silently substituted and logged on every build.
  commentstyle=\color{cmt},
  numbers=left,
  numberstyle=\ttfamily\tiny\color{mgrey},
  numbersep=8pt,
  backgroundcolor=\color{codebg},
  frame=single,
  rulecolor=\color{codeframe},
  framesep=6pt,
  breaklines=true,
  breakatwhitespace=false,
  postbreak=\mbox{\textcolor{mgrey}{$\hookrightarrow$}\space},
  showstringspaces=false,
  tabsize=4,
  columns=fullflexible,
  keepspaces=true,
  captionpos=b,
  aboveskip=1.2em,
  belowskip=1.2em,
  % Do NOT add a mapping for U+00A0: it aborts the run with a message that
  % names neither the character nor this line. Listings are ASCII, and
  % parity's C13 checks it.
  literate={ł}{{\l}}1 {Ł}{{\L}}1 {ą}{{\k a}}1 {ę}{{\k e}}1
           {ś}{{\'s}}1 {ć}{{\'c}}1 {ż}{{\.z}}1 {ź}{{\'z}}1 {ó}{{\'o}}1 {ń}{{\'n}}1
           {Ą}{{\k A}}1 {Ę}{{\k E}}1 {Ś}{{\'S}}1 {Ć}{{\'C}}1
           {Ż}{{\.Z}}1 {Ź}{{\'Z}}1 {Ó}{{\'O}}1 {Ń}{{\'N}}1
           {—}{{-{}-}}2 {–}{{-}}1 {‘}{{'}}1 {’}{{'}}1
           {“}{{"}}1 {”}{{"}}1 {°}{{\textdegree}}1
}
\lstset{style=bookcode}

\lstnewenvironment{python}[1][]
  {\lstset{language=pythonx,style=bookcode,#1}}
  {}
\lstnewenvironment{shellcmd}[1][]
  {\lstset{language=bash,style=bookcode,numbers=none,keywordstyle=\color{mteal}\bfseries,#1}}
  {}

% The path is printed above every file listing, in the form the reader opens
% in the editor. \nobreak keeps it on the page of the listing it names.
% The path must not be stranded at the foot of a page with its listing
% overleaf: \nobreak alone does not hold, because listings opens with its own
% vertical material. \needspace asks for room for the path and the first few
% lines of code, and turns the page first when there is not.
% \Needspace (capital N) measures \pagegoal-\pagetotal exactly; lower-case
% \needspace works through glue and penalties and was measured, by the
% Chapter 11 pass, turning pages early and leaving gaps. It must be used
% between paragraphs, which \listingpath's own \par guarantees. Chapters use
% \Needspace* too, so the fallback for a machine without the package defines
% both forms (the star swallowed with its argument).
\makeatletter
\IfFileExists{needspace.sty}{\usepackage{needspace}}{%
  \providecommand{\needspace}[1]{}%
  \providecommand{\Needspace}{\@ifstar\@gobble\@gobble}}
\makeatother
\newcommand{\listingpath}[1]{%
  \par\medskip\Needspace{6\baselineskip}\noindent
  {\ttfamily\footnotesize\color{mgrey}\detokenize{#1}}%
  \par\nobreak}

% Usage: \pyfile{code/ch05/cobweb.py}{Caption}{lst:label}
\newcommand{\pyfile}[3]{%
  \listingpath{#1}%
  \lstinputlisting[language=pythonx,style=bookcode,aboveskip=0.3em,
    caption={#2},label={#3}]{#1}%
}

% A named region of a file, delimited by `# --8<-- [start:name]` and
% `# --8<-- [end:name]`. Matched through listings' range prefix and suffix
% keys -- the other obvious form matches nothing and prints the WHOLE file
% without a warning (measured in the sibling book). A region name is a word,
% never a number. The line numbers printed are the file's own.
% Usage: \pyregion{code/ch05/cobweb.py}{step}{Caption}{lst:label}
\lstset{
  rangebeginprefix=\#\ --8<--\ [start:, rangebeginsuffix=],
  rangeendprefix=\#\ --8<--\ [end:, rangeendsuffix=],
  includerangemarker=false}
\newcommand{\pyregion}[4]{%
  \listingpath{#1}%
  \lstinputlisting[language=pythonx,style=bookcode,aboveskip=0.3em,
    linerange={#2-#2},caption={#3},label={#4}]{#1}%
}

% A program's output, written by code/measure/transcripts.py. There is no
% typed form and there must not be: \val{} does not expand in a verbatim
% body, and a transcript nobody ran looks exactly like one that was.
\newcommand{\transcript}[1]{%
  \par\smallskip
  \IfFileExists{figures/transcripts/#1.txt}{%
    \lstinputlisting[style=bookcode,numbers=none,basicstyle=\ttfamily\footnotesize,
      keywordstyle=\color{black},backgroundcolor=\color{white},
      language={}]{figures/transcripts/#1.txt}%
  }{%
    \begin{tcolorbox}[enhanced, sharp corners, width=\linewidth,
      colback=mlight, colframe=mgrey, boxrule=0.6pt,
      left=8pt, right=8pt, top=6pt, bottom=6pt]
      {\bfseries\color{mgrey}\small \lblTranscriptMissing}\\[2pt]
      {\ttfamily\scriptsize\color{mgrey} figures/transcripts/#1.txt}
    \end{tcolorbox}%
  }%
  \par\smallskip
}

% Inline literals. Underscores inside these must be written \_.
\newcommand{\code}[1]{\texttt{\small #1}}
\newcommand{\pkg}[1]{\texttt{\small\color{mblue}#1}}
\newcommand{\api}[1]{\texttt{\small\color{mteal}#1}}

% ============================================================
%  ADMONITIONS
%  A small, fixed vocabulary. `tinted' may be read in passing; `outlined' is
%  a block the reader must stop at. A third treatment would mean the first two
%  are not carrying their meaning.
% ============================================================
\tcbset{
  cfz tinted/.style={
    enhanced, breakable, sharp corners,
    lines before break=4,
    colback=mlight, boxrule=0pt, leftrule=3pt,
    left=8pt, right=8pt, top=6pt, bottom=6pt,
    before skip=13pt, after skip=13pt,
    fonttitle=\bfseries\small,
    coltitle=black, colbacktitle=mlight,
    attach title to upper={\par\smallskip},
  },
  cfz outlined/.style={
    enhanced, breakable, sharp corners,
    % A box broken straight after its title is a heading stranded at the foot
    % of a page with its question overleaf. Four lines is title plus three.
    lines before break=4,
    colback=white, boxrule=0.6pt,
    left=8pt, right=8pt, top=6pt, bottom=6pt,
    before skip=13pt, after skip=13pt,
    fonttitle=\bfseries\small,
    coltitle=black, colbacktitle=white,
    attach title to upper={\par\smallskip},
  },
}

% Read in passing.
\newtcolorbox{note}[1][]{cfz tinted, colframe=mblue, title={\lblNote}, #1}
\newtcolorbox{warning}[1][]{cfz tinted, colframe=mred, title={\lblWarning}, #1}
% Where you meet this in the world. If it cannot name a specific system,
% instrument, paper or event, it is prose and not this box.
\newtcolorbox{worldbox}[1][]{cfz tinted, colframe=mteal, title={\lblWorld}, #1}
% What is stated and not proved here, and where the proof lives.
\newtcolorbox{rigourbox}[1][]{cfz tinted, colframe=mgrey, title={\lblRigour}, #1}
% The companion library the reader builds, under code/src/chaoslab/.
\newtcolorbox{projectbox}[1][]{cfz tinted, colframe=mamber, title={\lblProject}, #1}

% Stop and read.
% The misconception, named AFTER the reader has been asked the question it
% gets wrong. Appendix B is the catalogue; a trap box is one entry in place.
\newtcolorbox{trapbox}[1][]{cfz outlined, colframe=mred, title={\lblTrap}, #1}
% A listing that has not been run on the pinned versions. Counted by
% `make debt'; nothing ships to a reader with one attached.
\newtcolorbox{verifybox}[1][]{cfz outlined, colframe=mgrey, title={\lblVerify}, #1}
\newtcolorbox{labbox}[1][]{cfz outlined, colframe=mgreen, title={\lblLab}, #1}

% ============================================================
%  THE METHOD --- questions before answers
%  ---------------------------------------------------------
%  The unit of teaching is a question the reader answers BEFORE reading the
%  answer. That is retrieval practice and errorful generation, the two
%  findings the programmed-learning format is defended by, and it is the one
%  thing the sibling mathematics book does that this one keeps whole.
%
%  \begin{think} ... \end{think}   the question, numbered per chapter, ending
%                                  with the instruction to answer before
%                                  reading on;
%  \reveal{short answer}           a centred box with the answer, which is the
%                                  first thing after the question -- the edge
%                                  the reader covers with a hand, or scrolls
%                                  to;
%  \begin{revealblock} ... \end{revealblock}  the same, for a worked answer.
%
%  tools/check_structure.py --think fails when a think box is not followed by
%  a reveal before the next think box or section, and when a reveal answers
%  nothing. That is a property of the source, so it is checked mechanically.
%
%  The instruction says "before you read on", never "before you turn over":
%  the two editions paginate differently, so the second is true in one of them.
% ============================================================
\newcounter{think}[chapter]
\renewcommand{\thethink}{\thechapter.\arabic{think}}
\newtcolorbox{thinkbox}[1][]{cfz outlined, colframe=mviolet, #1}
\NewDocumentEnvironment{think}{}{%
  \refstepcounter{think}%
  \begin{thinkbox}[title={\lblThink~\thethink}]%
}{%
  \par\smallskip
  {\raggedleft\itshape\small\color{mviolet}\lblThinkCue\par}%
  \end{thinkbox}%
}
\tcbset{
  cfz answer/.style={
    enhanced, sharp corners,
    colback=mfaint, colframe=mviolet!55, boxrule=0.5pt,
    before skip=4pt, after skip=10pt,
  },
}
\newcommand{\reveal}[1]{%
  \par\nobreak
  \begin{tcolorbox}[cfz answer, unbreakable,
    left=8pt, right=8pt, top=5pt, bottom=5pt,
    width=0.86\linewidth, center, halign=flush center]
    {\small\bfseries\color{mviolet}\lblAnswerMark}\enspace #1
  \end{tcolorbox}%
  \nobreak\noindent\ignorespaces
}
\NewDocumentEnvironment{revealblock}{}{%
  \par\nopagebreak
  \begin{tcolorbox}[cfz answer, breakable,
    left=10pt, right=10pt, top=6pt, bottom=6pt]
  {\small\bfseries\color{mviolet}\lblAnswerMark}\par\smallskip
}{%
  \end{tcolorbox}%
  \nopagebreak\par\noindent\ignorespaces
}

% ---- Outcomes on entry, and "Can you?" generated from them ----------------
% The self-check at the end of a chapter IS the list of outcomes at its
% start, one for one, because it is built from the same tokens: each
% \outcome both prints its item and appends a row to \cfz@canrows. A \loop
% inside a tabularx body does not survive alignment scanning, which is why
% the rows are accumulated at declaration time (the sibling book's finding).
\makeatletter
\newcommand{\cfz@canrows}{}
\newcommand{\cfz@ratebox}{%
  \fbox{\rule{0pt}{0.9ex}\rule{0.9ex}{0pt}}}
\newcommand{\cfz@rates}{%
  {\footnotesize\color{mgrey}%
   1\,\cfz@ratebox\ 2\,\cfz@ratebox\ 3\,\cfz@ratebox\ 4\,\cfz@ratebox\ 5\,\cfz@ratebox}}
\NewDocumentEnvironment{outcomes}{}{%
  \gdef\cfz@canrows{}%
  \begin{tcolorbox}[cfz tinted, colframe=mblue, title={\lblOutcomes}]%
  \lblOutcomesLead
  \begin{itemize}[leftmargin=1.4em, itemsep=2pt, topsep=4pt]%
}{%
  \end{itemize}%
  \end{tcolorbox}%
}
\newcommand{\outcome}[1]{%
  \item #1%
  \gappto\cfz@canrows{#1 & \cfz@rates \\ \addlinespace[3pt]}%
}
\newcommand{\canyou}{%
  \section*{\lblCanYou}%
  \phantomsection\addcontentsline{toc}{section}{\lblCanYou}%
  \noindent\lblCanYouIntro\par\medskip
  \ifdefempty{\cfz@canrows}{%
    \noindent{\color{mred}\lblNoOutcomes}\par}{%
  {\small
  \noindent\begin{tabularx}{\linewidth}{@{}>{\raggedright\arraybackslash}X r@{}}
    \toprule
    \cfz@canrows
    \bottomrule
  \end{tabularx}\par}}%
}
\makeatother

% ---- Summary -------------------------------------------------------------
\newtcolorbox{summarybox}[1][]{cfz tinted, colframe=mblue, title={\lblSummary}, #1}

% ---- Checkpoint questions, answered at the back --------------------------
% Answers are written AT THE QUESTION with \answerto{} and printed in
% Appendix A. They are carried there by a global token list rather than an
% auxiliary file, because an answer contains mathematics and \val{}, and
% material written to an .aux is expanded at shipout, where fragile commands
% break with errors naming neither the answer nor the chapter.
% \include is sequential within a run, so the list is complete by the time
% the appendix is typeset.
\makeatletter
\newcommand{\cfz@answers}{}
\newcommand{\cfz@lastansch}{0}
\NewDocumentEnvironment{checkpoint}{}{%
  \section*{\lblCheckpoint}%
  \phantomsection\addcontentsline{toc}{section}{\lblCheckpoint}%
  \noindent\lblCheckpointIntro\par
  \begin{enumerate}[label=\textbf{\arabic*.}, leftmargin=2em, itemsep=4pt]%
}{%
  \end{enumerate}%
}
\newcommand{\answerto}[1]{%
  \ifnum\value{chapter}=\cfz@lastansch\relax\else
    \xdef\cfz@lastansch{\arabic{chapter}}%
    \xappto\cfz@answers{\noexpand\cfz@anschapter{\thechapter}}%
  \fi
  \xappto\cfz@answers{\noexpand\cfz@ansitem{\arabic{enumi}}}%
  \gappto\cfz@answers{{#1}}%
}
\newcommand{\cfz@anschapter}[1]{%
  \par\medskip\noindent{\bfseries\color{mblue}\chaptername~#1}\par\nobreak\smallskip}
\newcommand{\cfz@ansitem}[2]{%
  \par\noindent\hangindent=2em\hangafter=1
  \makebox[2em][l]{\textbf{#1.}}#2\par}
\newcommand{\answersbody}{%
  \ifdefempty{\cfz@answers}{\noindent{\color{mred}\lblNoAnswers}\par}{%
    {\raggedright\cfz@answers}}%
}
\makeatother

% ============================================================
%  LABS
%  \begin{lab}{key}{title} ... \end{lab}
%
%  A lab is not a paragraph but a FILE the reader opens: a starter under
%  code/labs/chNN/<key>.py that must FAIL its test, a reference solution under
%  code/labs/chNN/solutions/, and test_<key>.py beside it. The box prints the
%  paths and the one command that runs the test. An environment rather than a
%  macro, because a macro argument cannot hold a verbatim fragment.
%  tools/check_structure.py --labs fails when a file is missing or the key's
%  chapter number disagrees with its chapter.
% ============================================================
\newcounter{lab}[chapter]
\renewcommand{\thelab}{\thechapter.\arabic{lab}}
\makeatletter
\newcommand{\labdir}{ch\two@digits{\value{chapter}}}
\newcommand{\listoflabs}{%
  \section*{\lblLabManifest}%
  \phantomsection\addcontentsline{toc}{section}{\lblLabManifest}%
  \begingroup\c@tocdepth=2\relax\@starttoc{lab}\endgroup%
}
\makeatother
\NewDocumentEnvironment{lab}{mm}{%
  \refstepcounter{lab}%
  \addcontentsline{lab}{subsection}{\protect\numberline{\thelab}\texttt{\protect\detokenize{#1}} --- #2}%
  \begin{labbox}[title={\lblLab~\thelab\ \dash{} #2}]%
}{%
  \par\medskip
  {\small\setlength{\parindent}{0pt}%
  \lblLabStarter: \texttt{code/labs/\labdir/\detokenize{#1}.py}\\
  \lblLabTest: \texttt{code/labs/\labdir/test\_\detokenize{#1}.py}\\
  \lblLabRun: \texttt{cd code \&\& uv run pytest -k \detokenize{#1}}\par}%
  \end{labbox}%
}

% ============================================================
%  FIGURES --- plots and diagrams
%  ---------------------------------------------------------
%  \plotfig{key}{caption}{one-line manifest description}
%    A PLOT, generated by code/plots/<chapter>.py from the same library the
%    listings use, into figures/plots/<lang>/<key>.pdf. Per language, because
%    an axis label in the wrong language is the half-translation that makes a
%    bilingual book look machine-made, and because the Polish edition owes a
%    decimal comma on its tick labels too. Build output, gitignored.
%
%  \mermaidfig{key}{caption}{one-line manifest description}
%    A DIAGRAM of an idea rather than of data, from
%    figures/mermaid/<lang>/<key>.mmd, rendered by `make diagrams'.
%
%  If the rendered PDF is absent either macro prints a visible marker in the
%  flow rather than a float, so a missing figure cannot be mistaken for a
%  layout choice. \phantomsection before \addcontentsline is load-bearing:
%  without it the manifest entry records whatever anchor preceded the figure,
%  which after a listing is a line that was never printed.
% ============================================================
\makeatletter
\newcommand{\listofplots}{%
  \section*{\lblPlotManifest}%
  \phantomsection\addcontentsline{toc}{section}{\lblPlotManifest}%
  \begingroup\c@tocdepth=2\relax\@starttoc{plt}\endgroup%
}
\newcommand{\listofdiagrams}{%
  \section*{\lblDiagramManifest}%
  \phantomsection\addcontentsline{toc}{section}{\lblDiagramManifest}%
  \begingroup\c@tocdepth=2\relax\@starttoc{dgm}\endgroup%
}
\makeatother

\newcommand{\figmissing}[3]{% #1 message, #2 path, #3 caption
  \par\medskip
  \begin{tcolorbox}[enhanced, breakable, width=\linewidth, sharp corners,
    colback=white, colframe=mgrey, boxrule=0.8pt,
    left=8pt, right=8pt, top=8pt, bottom=8pt]
    {\bfseries\color{mgrey}\small #1}\\[4pt]
    {\ttfamily\scriptsize\color{mgrey} #2}
  \end{tcolorbox}%
  \captionof{figure}{#3}%
  \par\medskip
}

\newcommand{\plotfig}[3]{%
  \phantomsection%
  \addcontentsline{plt}{subsection}{\protect\numberline{}\texttt{#1} --- #3}%
  \IfFileExists{figures/plots/\booklang/#1.pdf}{%
    \begin{figure}[htbp]
    \centering
    \includegraphics[width=\linewidth,height=0.45\textheight,keepaspectratio]{figures/plots/\booklang/#1.pdf}%
    \caption{#2}\label{fig:#1}
    \end{figure}%
  }{%
    \figmissing{\lblPlotNotRendered}{figures/plots/\booklang/#1.pdf}{#2}\label{fig:#1}%
  }%
}

\newcommand{\mermaidfig}[3]{%
  \phantomsection%
  \addcontentsline{dgm}{subsection}{\protect\numberline{}\texttt{#1.mmd} --- #3}%
  \IfFileExists{figures/diagrams/\booklang/#1.pdf}{%
    \begin{figure}[htbp]
    \centering
    \includegraphics[width=0.95\linewidth,height=0.42\textheight,keepaspectratio]{figures/diagrams/\booklang/#1.pdf}%
    \caption{#2}\label{fig:#1}
    \end{figure}%
  }{%
    \figmissing{\lblNotRendered}{figures/mermaid/\booklang/#1.mmd}{#2}\label{fig:#1}%
  }%
}

% ============================================================
%  APPENDIX B: the misreadings, generated by tools/gen_traps.py
%
%  \trapchapter{Chapter~\ref{ch:x}}{title}    a heading per chapter
%  \trapentry{n}{habit}{truth}{\ref{sec:x}}  one misreading
%
%  Set ragged right: an entry is mostly short sentences with \code{} spans in
%  them, and a justified line full of unbreakable runs is an overfull box
%  waiting for the edition with the longer words.
% ============================================================
\newcommand{\trapchapter}[2]{%
  \par\medskip
  {\large\bfseries\color{mblue}#1\quad\normalfont\itshape\color{mgrey}#2\par}%
  \nopagebreak\smallskip}
\newcommand{\trapentry}[4]{%
  \par\noindent
  \begingroup\raggedright\hangindent=2.4em\hangafter=1
  \makebox[2.4em][l]{\textbf{\color{mred}#1}}%
  \textbf{\enquote{#2}}\enspace #3\enspace{\color{mgrey}\S\,#4}\par
  \endgroup\smallskip}

% ============================================================
%  CHAPTER STUBS
% ============================================================
\newcommand{\chapterstub}[1]{%
  \begin{tcolorbox}[
    enhanced, breakable, width=\linewidth, sharp corners,
    colback=mlight, colframe=mred, boxrule=1pt,
    borderline={0.6pt}{3pt}{mred, dashed},
    left=12pt, right=12pt, top=12pt, bottom=12pt]
    {\bfseries\color{mred}\large \lblNotWritten}\\[8pt]
    \small\raggedright #1
  \end{tcolorbox}%
}

% ============================================================
%  COMPUTED VALUES
%  A number the reader cannot do in their head is computed by a script under
%  code/measure/, written to figures/values/<chapter>.tex as
%      \pyval{ch06.steps.apart}{34}
%  and reaches the page as \val{ch06.steps.apart}. The book contains
%  references, not digits; `make verify' fails when a committed value no
%  longer matches its script. \val applies the edition's decimal marker, which
%  is how the Polish edition gets its comma without anybody remembering.
% ============================================================
% The comma is BRACED: siunitx sets numbers in maths mode, where a bare comma
% is punctuation and gets a thin space after it ("9, 81"). {,} is an ordinary
% symbol, which is what a decimal marker is.
\newcommand{\bookdecimalmarker}{\ifpl{{,}}{.}}
\IfFileExists{siunitx.sty}{%
  \usepackage{siunitx}%
  \sisetup{
    group-digits = integer,
    group-minimum-digits = 5,
    group-separator = {\,},
    output-decimal-marker = {\bookdecimalmarker},
    exponent-product = \times,
    print-unity-mantissa = false
  }%
  \newcommand{\booknum}[1]{\num{##1}}%
}{%
  \newcommand{\booknum}[1]{##1}%
  \providecommand{\num}[1]{##1}%
}

\makeatletter
\newcommand{\pyvalues}{}
\newcommand{\pyval}[2]{\csgdef{pyval@#1}{#2}\listxadd{\pyvalues}{#1}}
\newcommand{\pyvaltext}[2]{%
  \csgdef{pyval@#1}{#2}\csgdef{pyvaltext@#1}{}\listxadd{\pyvalues}{#1}}
% \num on a non-numeric body is a FATAL siunitx error, so \val refuses a text
% value by name rather than letting siunitx fail on it.
\newcommand{\val}[1]{%
  \ifcsdef{pyval@#1}{%
    \ifcsdef{pyvaltext@#1}{%
      \PackageError{chaosbook}{Value `#1' is not a number}%
        {Use \string\valtext{#1} instead.}%
    }{}%
    \booknum{\csuse{pyval@#1}}%
  }{\textcolor{mred}{\textbf{[\lblValueMissing: #1]}}}%
}
\newcommand{\valtext}[1]{%
  \ifcsdef{pyval@#1}{\csuse{pyval@#1}}%
    {\textcolor{mred}{\textbf{[\lblValueMissing: #1]}}}%
}
\makeatother

% figures/values/all.tex is generated by `make numbers' and \inputs each
% value file. A build with no values still compiles; every \val{} then prints
% a visible marker. \valuesindex lets the one-chapter build point elsewhere.
\providecommand{\valuesindex}{figures/values/all}
\IfFileExists{\valuesindex.tex}{\input{\valuesindex}}{%
  \AtBeginDocument{\PackageWarning{chaosbook}{No computed values found. Run `make numbers'.}}}

% ============================================================
%  PINNED VERSIONS --- single source of truth
%  Change these here AND in code/pyproject.toml; tools/check_structure.py
%  --pins fails when the two disagree. A chapter never types a version.
% ============================================================
\newcommand{\bookedition}{0.1}
\newcommand{\bookyear}{2026}
\newcommand{\pyver}{3.14}              % the minor the book targets
\newcommand{\pypatch}{3.14.7}          % the exact interpreter the code ran on
\newcommand{\uvver}{0.12.13}           % uv
\newcommand{\numpyver}{2.5.3}          % numpy
\newcommand{\matplotlibver}{3.11.2}    % matplotlib
\newcommand{\pytestver}{9.1.1}         % pytest
\newcommand{\ruffver}{0.16.9}          % ruff
% The prose licence is a name with a version in it, not a quantity: it is a
% macro so that neither edition typesets it as a decimal to be localised.
\newcommand{\proselicence}{CC~BY-NC-SA~4.0}

% ============================================================
%  INDEX
% ============================================================
\apptocmd{\theindex}{\raggedcolumns\raggedright}{}{%
  \PackageWarning{chaosbook}{Could not patch theindex for ragged setting}}
\providecommand*\see[2]{}
\renewcommand*\see[2]{\emph{\lblIndexSee} #1}
\providecommand*\seealso[2]{}
\renewcommand*\seealso[2]{\emph{\lblIndexSeeAlso} #1}

\makeindex
 still compiles and says what it did.
\providecommand{\booklang}{en}

\usepackage{etoolbox}   % loaded first: the language switch below needs it

% ---------- Page geometry ----------
% ONE format: A4, 12pt, ONE-SIDED. The book is read on a screen, usually with
% an editor open beside it, so there is no recto, no verso, no gutter and no
% blank leaf before a chapter. The measure is about seventy-five characters
% of prose. At \footnotesize a listing line of 79 characters fits without
% wrapping, and 79 is the width every file under code/ is held to by
% tools/check_structure.py --lines.
\usepackage[
  a4paper,
  left=3.1cm, right=3.1cm, top=2.6cm, bottom=3.0cm,
  head=26pt, headsep=0.8cm, footskip=1.4cm
]{geometry}

% ---------- Encoding & fonts ----------
\usepackage[T1]{fontenc}
\usepackage[utf8]{inputenc}
\IfFileExists{lmodern.sty}{\usepackage{lmodern}}{}
% newtxtext takes its TS1 symbols from TeX Gyre Termes (Debian: tex-gyre).
% A machine with newtx and without tex-gyre dies on the copyright page. CI's
% full TeX Live is the reference; locally install tex-gyre as well.
\IfFileExists{newtxtext.sty}{\usepackage{newtxtext}}{}
% amsmath before newtxmath; amssymb only when newtxmath is absent, because
% loading both is a fatal clash (\Bbbk already defined) that is invisible on
% a machine without newtx.
\usepackage{amsmath}
\IfFileExists{newtxmath.sty}{\usepackage{newtxmath}}{\usepackage{amssymb}}
\IfFileExists{inconsolata.sty}{\usepackage[varqu,scaled=0.95]{inconsolata}}{}
% Without lmodern, microtype runs with font expansion off, which is what
% resolves a marginal overfull line -- so a container short of lmodern
% disagrees with CI about line breaks. Install lmodern locally too.
\IfFileExists{lmodern.sty}{\usepackage{microtype}}{\usepackage[expansion=false]{microtype}}

% upquote: in T1 the input ' is a right single quote, and a listing that sets
% f('x') with curly quotes is a SyntaxError when typed back in.
\IfFileExists{upquote.sty}{\usepackage{upquote}}{}
% And its twin: in T1 the pair -- is an en-dash ligature in the typewriter
% family too, so \code{--flag} would print a flag nobody can type. Scoped to
% tt*, so prose keeps its own dashes.
\DisableLigatures[-]{encoding = T1, family = tt*}

% ---------- Language ----------
% babel with a missing language file is FATAL, so each .ldf is probed first.
\IfFileExists{babel.sty}{%
  \IfFileExists{polish.ldf}{%
    \IfFileExists{british.ldf}{%
      \ifdefstring{\booklang}{pl}%
        {\usepackage[british,main=polish]{babel}}%
        {\usepackage[polish,main=british]{babel}}%
    }{\usepackage[polish]{babel}}%
  }{%
    \IfFileExists{british.ldf}{\usepackage[british]{babel}}{}%
  }%
}{}
% babel-polish makes " an active shorthand, which is a category-code accident
% waiting to happen inside listings. Polish quotation marks come from
% \enquote, which csquotes sets per language.
\ifdefstring{\booklang}{pl}{\AtBeginDocument{\shorthandoff{"}}}{}

% ---------- Core packages ----------
\usepackage{graphicx}
\usepackage{xcolor}
\usepackage{booktabs}
\usepackage{tabularx}
\usepackage{longtable}
\usepackage{array}
\usepackage{enumitem}
\usepackage{listings}
\usepackage[most]{tcolorbox}
\usepackage{fancyhdr}
\usepackage{titlesec}
\usepackage{caption}
\usepackage{imakeidx}
\IfFileExists{csquotes.sty}{\usepackage[autostyle=true]{csquotes}}{%
  \providecommand{\enquote}[1]{``##1''}}
\usepackage[hidelinks]{hyperref}

% ---------- Palette ----------
% The family's palette, so the three books read as a set. Two colours are
% this book's own: the lab (green) and the stop-and-think box (violet).
\definecolor{mblue}{HTML}{1F4E79}
\definecolor{mteal}{HTML}{0E7C7B}
\definecolor{mamber}{HTML}{B26A00}
\definecolor{mred}{HTML}{9B2C2C}
\definecolor{mviolet}{HTML}{7B4397}
\definecolor{mgreen}{HTML}{3B7A3B}
\definecolor{mgrey}{HTML}{5A5A5A}
\definecolor{mlight}{HTML}{F4F6F8}
\definecolor{mfaint}{HTML}{FBFCFD}
\definecolor{codebg}{HTML}{FAFAFA}
\definecolor{codeframe}{HTML}{D8DEE4}
\definecolor{kw}{HTML}{0B5394}
\definecolor{str}{HTML}{9B2C2C}
\definecolor{cmt}{HTML}{4F7A4F}
\definecolor{typ}{HTML}{0E7C7B}
\definecolor{dec}{HTML}{7B4397}

\hypersetup{
  colorlinks=true,
  linkcolor=mblue,
  urlcolor=mteal,
  citecolor=mblue,
  pdfauthor={Konrad Cinkusz}
}

% \ifpl{Polish}{English} -- for the handful of places where a string is
% structural rather than content. Robust, because a fragile \ifpl inside a
% moving argument fails on the next run in exactly one of the two editions.
\DeclareRobustCommand{\ifpl}[2]{\ifdefstring{\booklang}{pl}{#1}{#2}}
\makeatletter
\pdfstringdefDisableCommands{%
  \ifdefstring{\booklang}{pl}{\let\ifpl\@firstoftwo}{\let\ifpl\@secondoftwo}%
  % hyperref drops \color from a bookmark and keeps its ARGUMENT, so a title
  % carrying \code{} would bookmark as "mtealx". Gobble the colour name.
  \def\color#1{}%
}
\makeatother

% ---------- Language strings ----------
\input{lang/\booklang}

% The PDF's subject, keywords and language read \lblPdf... from the language
% file, so they cannot sit in the colour block above, which runs first. Two
% blocks is the ordering, not an oversight.
\hypersetup{
  pdfsubject={\lblPdfSubject},
  pdfkeywords={\lblPdfKeywords},
  pdflang={\lblPdfLang}
}

% ---------- Headings ----------
\titleformat{\chapter}[display]
  {\normalfont\huge\bfseries\color{mblue}\raggedright}
  {\normalfont\normalsize\color{mgrey}\MakeUppercase{\chaptertitlename}\ \thechapter}
  {8pt}{\Huge\raggedright}
\titlespacing*{\chapter}{0pt}{10pt}{28pt}
\titleformat{\section}{\normalfont\Large\bfseries\color{mblue}\raggedright}{\thesection}{0.8em}{}
\titleformat{\subsection}{\normalfont\large\bfseries\color{mgrey}\raggedright}{\thesubsection}{0.7em}{}

\setcounter{tocdepth}{1}
\setcounter{secnumdepth}{2}

\pagestyle{fancy}
\fancyhf{}
\fancyhead[L]{\small\nouppercase{\leftmark}}
\fancyhead[R]{\small\thepage}
\renewcommand{\headrulewidth}{0.4pt}
\fancypagestyle{plain}{\fancyhf{}\fancyfoot[R]{\small\thepage}%
  \renewcommand{\headrulewidth}{0pt}}
% After \pagestyle{fancy}, which installs its own \chaptermark.
\renewcommand{\chaptermark}[1]{\markboth{\thechapter.\ #1}{}}

% ---------- Mathematics ----------
% The handful of notations the book uses, as macros so both editions set the
% same symbol. The book writes \ln for the natural logarithm and \log_{10}
% where the base is ten; a bare \log is ambiguous between the two readerships
% and is not used.
\newcommand{\R}{\mathbb{R}}
\newcommand{\dd}{\,\mathrm{d}}
\newcommand{\abs}[1]{\left\lvert #1\right\rvert}
\newcommand{\iu}{\mathrm{i}}

% ============================================================
%  CODE LISTINGS
%  Every listing in a chapter is a FILE under code/, pulled in with \pyfile or
%  \pyregion, and CI runs it on the pinned interpreter. An in-source
%  \begin{python} block is allowed for a fragment of three or four lines that
%  shows a shape rather than runs.
% ============================================================
\lstdefinelanguage{pythonx}{
  language=Python,
  morekeywords={async,await,match,case,nonlocal,yield,type},
  morekeywords=[2]{
    orbit,iterate,logistic,rk4_step,euler_step,integrate,lorenz,henon,
    lyapunov,separation,escape_time,box_count
  },
  sensitive=true
}

\lstdefinestyle{bookcode}{
  basicstyle=\ttfamily\footnotesize,
  keywordstyle=\color{kw}\bfseries,
  keywordstyle=[2]\color{typ},
  stringstyle=\color{str},
  % Colour only: inconsolata has no italic, and an italic comment style is
  % silently substituted and logged on every build.
  commentstyle=\color{cmt},
  numbers=left,
  numberstyle=\ttfamily\tiny\color{mgrey},
  numbersep=8pt,
  backgroundcolor=\color{codebg},
  frame=single,
  rulecolor=\color{codeframe},
  framesep=6pt,
  breaklines=true,
  breakatwhitespace=false,
  postbreak=\mbox{\textcolor{mgrey}{$\hookrightarrow$}\space},
  showstringspaces=false,
  tabsize=4,
  columns=fullflexible,
  keepspaces=true,
  captionpos=b,
  aboveskip=1.2em,
  belowskip=1.2em,
  % Do NOT add a mapping for U+00A0: it aborts the run with a message that
  % names neither the character nor this line. Listings are ASCII, and
  % parity's C13 checks it.
  literate={ł}{{\l}}1 {Ł}{{\L}}1 {ą}{{\k a}}1 {ę}{{\k e}}1
           {ś}{{\'s}}1 {ć}{{\'c}}1 {ż}{{\.z}}1 {ź}{{\'z}}1 {ó}{{\'o}}1 {ń}{{\'n}}1
           {Ą}{{\k A}}1 {Ę}{{\k E}}1 {Ś}{{\'S}}1 {Ć}{{\'C}}1
           {Ż}{{\.Z}}1 {Ź}{{\'Z}}1 {Ó}{{\'O}}1 {Ń}{{\'N}}1
           {—}{{-{}-}}2 {–}{{-}}1 {‘}{{'}}1 {’}{{'}}1
           {“}{{"}}1 {”}{{"}}1 {°}{{\textdegree}}1
}
\lstset{style=bookcode}

\lstnewenvironment{python}[1][]
  {\lstset{language=pythonx,style=bookcode,#1}}
  {}
\lstnewenvironment{shellcmd}[1][]
  {\lstset{language=bash,style=bookcode,numbers=none,keywordstyle=\color{mteal}\bfseries,#1}}
  {}

% The path is printed above every file listing, in the form the reader opens
% in the editor. \nobreak keeps it on the page of the listing it names.
% The path must not be stranded at the foot of a page with its listing
% overleaf: \nobreak alone does not hold, because listings opens with its own
% vertical material. \needspace asks for room for the path and the first few
% lines of code, and turns the page first when there is not.
% \Needspace (capital N) measures \pagegoal-\pagetotal exactly; lower-case
% \needspace works through glue and penalties and was measured, by the
% Chapter 11 pass, turning pages early and leaving gaps. It must be used
% between paragraphs, which \listingpath's own \par guarantees. Chapters use
% \Needspace* too, so the fallback for a machine without the package defines
% both forms (the star swallowed with its argument).
\makeatletter
\IfFileExists{needspace.sty}{\usepackage{needspace}}{%
  \providecommand{\needspace}[1]{}%
  \providecommand{\Needspace}{\@ifstar\@gobble\@gobble}}
\makeatother
\newcommand{\listingpath}[1]{%
  \par\medskip\Needspace{6\baselineskip}\noindent
  {\ttfamily\footnotesize\color{mgrey}\detokenize{#1}}%
  \par\nobreak}

% Usage: \pyfile{code/ch05/cobweb.py}{Caption}{lst:label}
\newcommand{\pyfile}[3]{%
  \listingpath{#1}%
  \lstinputlisting[language=pythonx,style=bookcode,aboveskip=0.3em,
    caption={#2},label={#3}]{#1}%
}

% A named region of a file, delimited by `# --8<-- [start:name]` and
% `# --8<-- [end:name]`. Matched through listings' range prefix and suffix
% keys -- the other obvious form matches nothing and prints the WHOLE file
% without a warning (measured in the sibling book). A region name is a word,
% never a number. The line numbers printed are the file's own.
% Usage: \pyregion{code/ch05/cobweb.py}{step}{Caption}{lst:label}
\lstset{
  rangebeginprefix=\#\ --8<--\ [start:, rangebeginsuffix=],
  rangeendprefix=\#\ --8<--\ [end:, rangeendsuffix=],
  includerangemarker=false}
\newcommand{\pyregion}[4]{%
  \listingpath{#1}%
  \lstinputlisting[language=pythonx,style=bookcode,aboveskip=0.3em,
    linerange={#2-#2},caption={#3},label={#4}]{#1}%
}

% A program's output, written by code/measure/transcripts.py. There is no
% typed form and there must not be: \val{} does not expand in a verbatim
% body, and a transcript nobody ran looks exactly like one that was.
\newcommand{\transcript}[1]{%
  \par\smallskip
  \IfFileExists{figures/transcripts/#1.txt}{%
    \lstinputlisting[style=bookcode,numbers=none,basicstyle=\ttfamily\footnotesize,
      keywordstyle=\color{black},backgroundcolor=\color{white},
      language={}]{figures/transcripts/#1.txt}%
  }{%
    \begin{tcolorbox}[enhanced, sharp corners, width=\linewidth,
      colback=mlight, colframe=mgrey, boxrule=0.6pt,
      left=8pt, right=8pt, top=6pt, bottom=6pt]
      {\bfseries\color{mgrey}\small \lblTranscriptMissing}\\[2pt]
      {\ttfamily\scriptsize\color{mgrey} figures/transcripts/#1.txt}
    \end{tcolorbox}%
  }%
  \par\smallskip
}

% Inline literals. Underscores inside these must be written \_.
\newcommand{\code}[1]{\texttt{\small #1}}
\newcommand{\pkg}[1]{\texttt{\small\color{mblue}#1}}
\newcommand{\api}[1]{\texttt{\small\color{mteal}#1}}

% ============================================================
%  ADMONITIONS
%  A small, fixed vocabulary. `tinted' may be read in passing; `outlined' is
%  a block the reader must stop at. A third treatment would mean the first two
%  are not carrying their meaning.
% ============================================================
\tcbset{
  cfz tinted/.style={
    enhanced, breakable, sharp corners,
    lines before break=4,
    colback=mlight, boxrule=0pt, leftrule=3pt,
    left=8pt, right=8pt, top=6pt, bottom=6pt,
    before skip=13pt, after skip=13pt,
    fonttitle=\bfseries\small,
    coltitle=black, colbacktitle=mlight,
    attach title to upper={\par\smallskip},
  },
  cfz outlined/.style={
    enhanced, breakable, sharp corners,
    % A box broken straight after its title is a heading stranded at the foot
    % of a page with its question overleaf. Four lines is title plus three.
    lines before break=4,
    colback=white, boxrule=0.6pt,
    left=8pt, right=8pt, top=6pt, bottom=6pt,
    before skip=13pt, after skip=13pt,
    fonttitle=\bfseries\small,
    coltitle=black, colbacktitle=white,
    attach title to upper={\par\smallskip},
  },
}

% Read in passing.
\newtcolorbox{note}[1][]{cfz tinted, colframe=mblue, title={\lblNote}, #1}
\newtcolorbox{warning}[1][]{cfz tinted, colframe=mred, title={\lblWarning}, #1}
% Where you meet this in the world. If it cannot name a specific system,
% instrument, paper or event, it is prose and not this box.
\newtcolorbox{worldbox}[1][]{cfz tinted, colframe=mteal, title={\lblWorld}, #1}
% What is stated and not proved here, and where the proof lives.
\newtcolorbox{rigourbox}[1][]{cfz tinted, colframe=mgrey, title={\lblRigour}, #1}
% The companion library the reader builds, under code/src/chaoslab/.
\newtcolorbox{projectbox}[1][]{cfz tinted, colframe=mamber, title={\lblProject}, #1}

% Stop and read.
% The misconception, named AFTER the reader has been asked the question it
% gets wrong. Appendix B is the catalogue; a trap box is one entry in place.
\newtcolorbox{trapbox}[1][]{cfz outlined, colframe=mred, title={\lblTrap}, #1}
% A listing that has not been run on the pinned versions. Counted by
% `make debt'; nothing ships to a reader with one attached.
\newtcolorbox{verifybox}[1][]{cfz outlined, colframe=mgrey, title={\lblVerify}, #1}
\newtcolorbox{labbox}[1][]{cfz outlined, colframe=mgreen, title={\lblLab}, #1}

% ============================================================
%  THE METHOD --- questions before answers
%  ---------------------------------------------------------
%  The unit of teaching is a question the reader answers BEFORE reading the
%  answer. That is retrieval practice and errorful generation, the two
%  findings the programmed-learning format is defended by, and it is the one
%  thing the sibling mathematics book does that this one keeps whole.
%
%  \begin{think} ... \end{think}   the question, numbered per chapter, ending
%                                  with the instruction to answer before
%                                  reading on;
%  \reveal{short answer}           a centred box with the answer, which is the
%                                  first thing after the question -- the edge
%                                  the reader covers with a hand, or scrolls
%                                  to;
%  \begin{revealblock} ... \end{revealblock}  the same, for a worked answer.
%
%  tools/check_structure.py --think fails when a think box is not followed by
%  a reveal before the next think box or section, and when a reveal answers
%  nothing. That is a property of the source, so it is checked mechanically.
%
%  The instruction says "before you read on", never "before you turn over":
%  the two editions paginate differently, so the second is true in one of them.
% ============================================================
\newcounter{think}[chapter]
\renewcommand{\thethink}{\thechapter.\arabic{think}}
\newtcolorbox{thinkbox}[1][]{cfz outlined, colframe=mviolet, #1}
\NewDocumentEnvironment{think}{}{%
  \refstepcounter{think}%
  \begin{thinkbox}[title={\lblThink~\thethink}]%
}{%
  \par\smallskip
  {\raggedleft\itshape\small\color{mviolet}\lblThinkCue\par}%
  \end{thinkbox}%
}
\tcbset{
  cfz answer/.style={
    enhanced, sharp corners,
    colback=mfaint, colframe=mviolet!55, boxrule=0.5pt,
    before skip=4pt, after skip=10pt,
  },
}
\newcommand{\reveal}[1]{%
  \par\nobreak
  \begin{tcolorbox}[cfz answer, unbreakable,
    left=8pt, right=8pt, top=5pt, bottom=5pt,
    width=0.86\linewidth, center, halign=flush center]
    {\small\bfseries\color{mviolet}\lblAnswerMark}\enspace #1
  \end{tcolorbox}%
  \nobreak\noindent\ignorespaces
}
\NewDocumentEnvironment{revealblock}{}{%
  \par\nopagebreak
  \begin{tcolorbox}[cfz answer, breakable,
    left=10pt, right=10pt, top=6pt, bottom=6pt]
  {\small\bfseries\color{mviolet}\lblAnswerMark}\par\smallskip
}{%
  \end{tcolorbox}%
  \nopagebreak\par\noindent\ignorespaces
}

% ---- Outcomes on entry, and "Can you?" generated from them ----------------
% The self-check at the end of a chapter IS the list of outcomes at its
% start, one for one, because it is built from the same tokens: each
% \outcome both prints its item and appends a row to \cfz@canrows. A \loop
% inside a tabularx body does not survive alignment scanning, which is why
% the rows are accumulated at declaration time (the sibling book's finding).
\makeatletter
\newcommand{\cfz@canrows}{}
\newcommand{\cfz@ratebox}{%
  \fbox{\rule{0pt}{0.9ex}\rule{0.9ex}{0pt}}}
\newcommand{\cfz@rates}{%
  {\footnotesize\color{mgrey}%
   1\,\cfz@ratebox\ 2\,\cfz@ratebox\ 3\,\cfz@ratebox\ 4\,\cfz@ratebox\ 5\,\cfz@ratebox}}
\NewDocumentEnvironment{outcomes}{}{%
  \gdef\cfz@canrows{}%
  \begin{tcolorbox}[cfz tinted, colframe=mblue, title={\lblOutcomes}]%
  \lblOutcomesLead
  \begin{itemize}[leftmargin=1.4em, itemsep=2pt, topsep=4pt]%
}{%
  \end{itemize}%
  \end{tcolorbox}%
}
\newcommand{\outcome}[1]{%
  \item #1%
  \gappto\cfz@canrows{#1 & \cfz@rates \\ \addlinespace[3pt]}%
}
\newcommand{\canyou}{%
  \section*{\lblCanYou}%
  \phantomsection\addcontentsline{toc}{section}{\lblCanYou}%
  \noindent\lblCanYouIntro\par\medskip
  \ifdefempty{\cfz@canrows}{%
    \noindent{\color{mred}\lblNoOutcomes}\par}{%
  {\small
  \noindent\begin{tabularx}{\linewidth}{@{}>{\raggedright\arraybackslash}X r@{}}
    \toprule
    \cfz@canrows
    \bottomrule
  \end{tabularx}\par}}%
}
\makeatother

% ---- Summary -------------------------------------------------------------
\newtcolorbox{summarybox}[1][]{cfz tinted, colframe=mblue, title={\lblSummary}, #1}

% ---- Checkpoint questions, answered at the back --------------------------
% Answers are written AT THE QUESTION with \answerto{} and printed in
% Appendix A. They are carried there by a global token list rather than an
% auxiliary file, because an answer contains mathematics and \val{}, and
% material written to an .aux is expanded at shipout, where fragile commands
% break with errors naming neither the answer nor the chapter.
% \include is sequential within a run, so the list is complete by the time
% the appendix is typeset.
\makeatletter
\newcommand{\cfz@answers}{}
\newcommand{\cfz@lastansch}{0}
\NewDocumentEnvironment{checkpoint}{}{%
  \section*{\lblCheckpoint}%
  \phantomsection\addcontentsline{toc}{section}{\lblCheckpoint}%
  \noindent\lblCheckpointIntro\par
  \begin{enumerate}[label=\textbf{\arabic*.}, leftmargin=2em, itemsep=4pt]%
}{%
  \end{enumerate}%
}
\newcommand{\answerto}[1]{%
  \ifnum\value{chapter}=\cfz@lastansch\relax\else
    \xdef\cfz@lastansch{\arabic{chapter}}%
    \xappto\cfz@answers{\noexpand\cfz@anschapter{\thechapter}}%
  \fi
  \xappto\cfz@answers{\noexpand\cfz@ansitem{\arabic{enumi}}}%
  \gappto\cfz@answers{{#1}}%
}
\newcommand{\cfz@anschapter}[1]{%
  \par\medskip\noindent{\bfseries\color{mblue}\chaptername~#1}\par\nobreak\smallskip}
\newcommand{\cfz@ansitem}[2]{%
  \par\noindent\hangindent=2em\hangafter=1
  \makebox[2em][l]{\textbf{#1.}}#2\par}
\newcommand{\answersbody}{%
  \ifdefempty{\cfz@answers}{\noindent{\color{mred}\lblNoAnswers}\par}{%
    {\raggedright\cfz@answers}}%
}
\makeatother

% ============================================================
%  LABS
%  \begin{lab}{key}{title} ... \end{lab}
%
%  A lab is not a paragraph but a FILE the reader opens: a starter under
%  code/labs/chNN/<key>.py that must FAIL its test, a reference solution under
%  code/labs/chNN/solutions/, and test_<key>.py beside it. The box prints the
%  paths and the one command that runs the test. An environment rather than a
%  macro, because a macro argument cannot hold a verbatim fragment.
%  tools/check_structure.py --labs fails when a file is missing or the key's
%  chapter number disagrees with its chapter.
% ============================================================
\newcounter{lab}[chapter]
\renewcommand{\thelab}{\thechapter.\arabic{lab}}
\makeatletter
\newcommand{\labdir}{ch\two@digits{\value{chapter}}}
\newcommand{\listoflabs}{%
  \section*{\lblLabManifest}%
  \phantomsection\addcontentsline{toc}{section}{\lblLabManifest}%
  \begingroup\c@tocdepth=2\relax\@starttoc{lab}\endgroup%
}
\makeatother
\NewDocumentEnvironment{lab}{mm}{%
  \refstepcounter{lab}%
  \addcontentsline{lab}{subsection}{\protect\numberline{\thelab}\texttt{\protect\detokenize{#1}} --- #2}%
  \begin{labbox}[title={\lblLab~\thelab\ \dash{} #2}]%
}{%
  \par\medskip
  {\small\setlength{\parindent}{0pt}%
  \lblLabStarter: \texttt{code/labs/\labdir/\detokenize{#1}.py}\\
  \lblLabTest: \texttt{code/labs/\labdir/test\_\detokenize{#1}.py}\\
  \lblLabRun: \texttt{cd code \&\& uv run pytest -k \detokenize{#1}}\par}%
  \end{labbox}%
}

% ============================================================
%  FIGURES --- plots and diagrams
%  ---------------------------------------------------------
%  \plotfig{key}{caption}{one-line manifest description}
%    A PLOT, generated by code/plots/<chapter>.py from the same library the
%    listings use, into figures/plots/<lang>/<key>.pdf. Per language, because
%    an axis label in the wrong language is the half-translation that makes a
%    bilingual book look machine-made, and because the Polish edition owes a
%    decimal comma on its tick labels too. Build output, gitignored.
%
%  \mermaidfig{key}{caption}{one-line manifest description}
%    A DIAGRAM of an idea rather than of data, from
%    figures/mermaid/<lang>/<key>.mmd, rendered by `make diagrams'.
%
%  If the rendered PDF is absent either macro prints a visible marker in the
%  flow rather than a float, so a missing figure cannot be mistaken for a
%  layout choice. \phantomsection before \addcontentsline is load-bearing:
%  without it the manifest entry records whatever anchor preceded the figure,
%  which after a listing is a line that was never printed.
% ============================================================
\makeatletter
\newcommand{\listofplots}{%
  \section*{\lblPlotManifest}%
  \phantomsection\addcontentsline{toc}{section}{\lblPlotManifest}%
  \begingroup\c@tocdepth=2\relax\@starttoc{plt}\endgroup%
}
\newcommand{\listofdiagrams}{%
  \section*{\lblDiagramManifest}%
  \phantomsection\addcontentsline{toc}{section}{\lblDiagramManifest}%
  \begingroup\c@tocdepth=2\relax\@starttoc{dgm}\endgroup%
}
\makeatother

\newcommand{\figmissing}[3]{% #1 message, #2 path, #3 caption
  \par\medskip
  \begin{tcolorbox}[enhanced, breakable, width=\linewidth, sharp corners,
    colback=white, colframe=mgrey, boxrule=0.8pt,
    left=8pt, right=8pt, top=8pt, bottom=8pt]
    {\bfseries\color{mgrey}\small #1}\\[4pt]
    {\ttfamily\scriptsize\color{mgrey} #2}
  \end{tcolorbox}%
  \captionof{figure}{#3}%
  \par\medskip
}

\newcommand{\plotfig}[3]{%
  \phantomsection%
  \addcontentsline{plt}{subsection}{\protect\numberline{}\texttt{#1} --- #3}%
  \IfFileExists{figures/plots/\booklang/#1.pdf}{%
    \begin{figure}[htbp]
    \centering
    \includegraphics[width=\linewidth,height=0.45\textheight,keepaspectratio]{figures/plots/\booklang/#1.pdf}%
    \caption{#2}\label{fig:#1}
    \end{figure}%
  }{%
    \figmissing{\lblPlotNotRendered}{figures/plots/\booklang/#1.pdf}{#2}\label{fig:#1}%
  }%
}

\newcommand{\mermaidfig}[3]{%
  \phantomsection%
  \addcontentsline{dgm}{subsection}{\protect\numberline{}\texttt{#1.mmd} --- #3}%
  \IfFileExists{figures/diagrams/\booklang/#1.pdf}{%
    \begin{figure}[htbp]
    \centering
    \includegraphics[width=0.95\linewidth,height=0.42\textheight,keepaspectratio]{figures/diagrams/\booklang/#1.pdf}%
    \caption{#2}\label{fig:#1}
    \end{figure}%
  }{%
    \figmissing{\lblNotRendered}{figures/mermaid/\booklang/#1.mmd}{#2}\label{fig:#1}%
  }%
}

% ============================================================
%  APPENDIX B: the misreadings, generated by tools/gen_traps.py
%
%  \trapchapter{Chapter~\ref{ch:x}}{title}    a heading per chapter
%  \trapentry{n}{habit}{truth}{\ref{sec:x}}  one misreading
%
%  Set ragged right: an entry is mostly short sentences with \code{} spans in
%  them, and a justified line full of unbreakable runs is an overfull box
%  waiting for the edition with the longer words.
% ============================================================
\newcommand{\trapchapter}[2]{%
  \par\medskip
  {\large\bfseries\color{mblue}#1\quad\normalfont\itshape\color{mgrey}#2\par}%
  \nopagebreak\smallskip}
\newcommand{\trapentry}[4]{%
  \par\noindent
  \begingroup\raggedright\hangindent=2.4em\hangafter=1
  \makebox[2.4em][l]{\textbf{\color{mred}#1}}%
  \textbf{\enquote{#2}}\enspace #3\enspace{\color{mgrey}\S\,#4}\par
  \endgroup\smallskip}

% ============================================================
%  CHAPTER STUBS
% ============================================================
\newcommand{\chapterstub}[1]{%
  \begin{tcolorbox}[
    enhanced, breakable, width=\linewidth, sharp corners,
    colback=mlight, colframe=mred, boxrule=1pt,
    borderline={0.6pt}{3pt}{mred, dashed},
    left=12pt, right=12pt, top=12pt, bottom=12pt]
    {\bfseries\color{mred}\large \lblNotWritten}\\[8pt]
    \small\raggedright #1
  \end{tcolorbox}%
}

% ============================================================
%  COMPUTED VALUES
%  A number the reader cannot do in their head is computed by a script under
%  code/measure/, written to figures/values/<chapter>.tex as
%      \pyval{ch06.steps.apart}{34}
%  and reaches the page as \val{ch06.steps.apart}. The book contains
%  references, not digits; `make verify' fails when a committed value no
%  longer matches its script. \val applies the edition's decimal marker, which
%  is how the Polish edition gets its comma without anybody remembering.
% ============================================================
% The comma is BRACED: siunitx sets numbers in maths mode, where a bare comma
% is punctuation and gets a thin space after it ("9, 81"). {,} is an ordinary
% symbol, which is what a decimal marker is.
\newcommand{\bookdecimalmarker}{\ifpl{{,}}{.}}
\IfFileExists{siunitx.sty}{%
  \usepackage{siunitx}%
  \sisetup{
    group-digits = integer,
    group-minimum-digits = 5,
    group-separator = {\,},
    output-decimal-marker = {\bookdecimalmarker},
    exponent-product = \times,
    print-unity-mantissa = false
  }%
  \newcommand{\booknum}[1]{\num{##1}}%
}{%
  \newcommand{\booknum}[1]{##1}%
  \providecommand{\num}[1]{##1}%
}

\makeatletter
\newcommand{\pyvalues}{}
\newcommand{\pyval}[2]{\csgdef{pyval@#1}{#2}\listxadd{\pyvalues}{#1}}
\newcommand{\pyvaltext}[2]{%
  \csgdef{pyval@#1}{#2}\csgdef{pyvaltext@#1}{}\listxadd{\pyvalues}{#1}}
% \num on a non-numeric body is a FATAL siunitx error, so \val refuses a text
% value by name rather than letting siunitx fail on it.
\newcommand{\val}[1]{%
  \ifcsdef{pyval@#1}{%
    \ifcsdef{pyvaltext@#1}{%
      \PackageError{chaosbook}{Value `#1' is not a number}%
        {Use \string\valtext{#1} instead.}%
    }{}%
    \booknum{\csuse{pyval@#1}}%
  }{\textcolor{mred}{\textbf{[\lblValueMissing: #1]}}}%
}
\newcommand{\valtext}[1]{%
  \ifcsdef{pyval@#1}{\csuse{pyval@#1}}%
    {\textcolor{mred}{\textbf{[\lblValueMissing: #1]}}}%
}
\makeatother

% figures/values/all.tex is generated by `make numbers' and \inputs each
% value file. A build with no values still compiles; every \val{} then prints
% a visible marker. \valuesindex lets the one-chapter build point elsewhere.
\providecommand{\valuesindex}{figures/values/all}
\IfFileExists{\valuesindex.tex}{\input{\valuesindex}}{%
  \AtBeginDocument{\PackageWarning{chaosbook}{No computed values found. Run `make numbers'.}}}

% ============================================================
%  PINNED VERSIONS --- single source of truth
%  Change these here AND in code/pyproject.toml; tools/check_structure.py
%  --pins fails when the two disagree. A chapter never types a version.
% ============================================================
\newcommand{\bookedition}{0.1}
\newcommand{\bookyear}{2026}
\newcommand{\pyver}{3.14}              % the minor the book targets
\newcommand{\pypatch}{3.14.7}          % the exact interpreter the code ran on
\newcommand{\uvver}{0.12.13}           % uv
\newcommand{\numpyver}{2.5.3}          % numpy
\newcommand{\matplotlibver}{3.11.2}    % matplotlib
\newcommand{\pytestver}{9.1.1}         % pytest
\newcommand{\ruffver}{0.16.9}          % ruff
% The prose licence is a name with a version in it, not a quantity: it is a
% macro so that neither edition typesets it as a decimal to be localised.
\newcommand{\proselicence}{CC~BY-NC-SA~4.0}

% ============================================================
%  INDEX
% ============================================================
\apptocmd{\theindex}{\raggedcolumns\raggedright}{}{%
  \PackageWarning{chaosbook}{Could not patch theindex for ragged setting}}
\providecommand*\see[2]{}
\renewcommand*\see[2]{\emph{\lblIndexSee} #1}
\providecommand*\seealso[2]{}
\renewcommand*\seealso[2]{\emph{\lblIndexSeeAlso} #1}

\makeindex
 still compiles and says what it did.
\providecommand{\booklang}{en}

\usepackage{etoolbox}   % loaded first: the language switch below needs it

% ---------- Page geometry ----------
% ONE format: A4, 12pt, ONE-SIDED. The book is read on a screen, usually with
% an editor open beside it, so there is no recto, no verso, no gutter and no
% blank leaf before a chapter. The measure is about seventy-five characters
% of prose. At \footnotesize a listing line of 79 characters fits without
% wrapping, and 79 is the width every file under code/ is held to by
% tools/check_structure.py --lines.
\usepackage[
  a4paper,
  left=3.1cm, right=3.1cm, top=2.6cm, bottom=3.0cm,
  head=26pt, headsep=0.8cm, footskip=1.4cm
]{geometry}

% ---------- Encoding & fonts ----------
\usepackage[T1]{fontenc}
\usepackage[utf8]{inputenc}
\IfFileExists{lmodern.sty}{\usepackage{lmodern}}{}
% newtxtext takes its TS1 symbols from TeX Gyre Termes (Debian: tex-gyre).
% A machine with newtx and without tex-gyre dies on the copyright page. CI's
% full TeX Live is the reference; locally install tex-gyre as well.
\IfFileExists{newtxtext.sty}{\usepackage{newtxtext}}{}
% amsmath before newtxmath; amssymb only when newtxmath is absent, because
% loading both is a fatal clash (\Bbbk already defined) that is invisible on
% a machine without newtx.
\usepackage{amsmath}
\IfFileExists{newtxmath.sty}{\usepackage{newtxmath}}{\usepackage{amssymb}}
\IfFileExists{inconsolata.sty}{\usepackage[varqu,scaled=0.95]{inconsolata}}{}
% Without lmodern, microtype runs with font expansion off, which is what
% resolves a marginal overfull line -- so a container short of lmodern
% disagrees with CI about line breaks. Install lmodern locally too.
\IfFileExists{lmodern.sty}{\usepackage{microtype}}{\usepackage[expansion=false]{microtype}}

% upquote: in T1 the input ' is a right single quote, and a listing that sets
% f('x') with curly quotes is a SyntaxError when typed back in.
\IfFileExists{upquote.sty}{\usepackage{upquote}}{}
% And its twin: in T1 the pair -- is an en-dash ligature in the typewriter
% family too, so \code{--flag} would print a flag nobody can type. Scoped to
% tt*, so prose keeps its own dashes.
\DisableLigatures[-]{encoding = T1, family = tt*}

% ---------- Language ----------
% babel with a missing language file is FATAL, so each .ldf is probed first.
\IfFileExists{babel.sty}{%
  \IfFileExists{polish.ldf}{%
    \IfFileExists{british.ldf}{%
      \ifdefstring{\booklang}{pl}%
        {\usepackage[british,main=polish]{babel}}%
        {\usepackage[polish,main=british]{babel}}%
    }{\usepackage[polish]{babel}}%
  }{%
    \IfFileExists{british.ldf}{\usepackage[british]{babel}}{}%
  }%
}{}
% babel-polish makes " an active shorthand, which is a category-code accident
% waiting to happen inside listings. Polish quotation marks come from
% \enquote, which csquotes sets per language.
\ifdefstring{\booklang}{pl}{\AtBeginDocument{\shorthandoff{"}}}{}

% ---------- Core packages ----------
\usepackage{graphicx}
\usepackage{xcolor}
\usepackage{booktabs}
\usepackage{tabularx}
\usepackage{longtable}
\usepackage{array}
\usepackage{enumitem}
\usepackage{listings}
\usepackage[most]{tcolorbox}
\usepackage{fancyhdr}
\usepackage{titlesec}
\usepackage{caption}
\usepackage{imakeidx}
\IfFileExists{csquotes.sty}{\usepackage[autostyle=true]{csquotes}}{%
  \providecommand{\enquote}[1]{``##1''}}
\usepackage[hidelinks]{hyperref}

% ---------- Palette ----------
% The family's palette, so the three books read as a set. Two colours are
% this book's own: the lab (green) and the stop-and-think box (violet).
\definecolor{mblue}{HTML}{1F4E79}
\definecolor{mteal}{HTML}{0E7C7B}
\definecolor{mamber}{HTML}{B26A00}
\definecolor{mred}{HTML}{9B2C2C}
\definecolor{mviolet}{HTML}{7B4397}
\definecolor{mgreen}{HTML}{3B7A3B}
\definecolor{mgrey}{HTML}{5A5A5A}
\definecolor{mlight}{HTML}{F4F6F8}
\definecolor{mfaint}{HTML}{FBFCFD}
\definecolor{codebg}{HTML}{FAFAFA}
\definecolor{codeframe}{HTML}{D8DEE4}
\definecolor{kw}{HTML}{0B5394}
\definecolor{str}{HTML}{9B2C2C}
\definecolor{cmt}{HTML}{4F7A4F}
\definecolor{typ}{HTML}{0E7C7B}
\definecolor{dec}{HTML}{7B4397}

\hypersetup{
  colorlinks=true,
  linkcolor=mblue,
  urlcolor=mteal,
  citecolor=mblue,
  pdfauthor={Konrad Cinkusz}
}

% \ifpl{Polish}{English} -- for the handful of places where a string is
% structural rather than content. Robust, because a fragile \ifpl inside a
% moving argument fails on the next run in exactly one of the two editions.
\DeclareRobustCommand{\ifpl}[2]{\ifdefstring{\booklang}{pl}{#1}{#2}}
\makeatletter
\pdfstringdefDisableCommands{%
  \ifdefstring{\booklang}{pl}{\let\ifpl\@firstoftwo}{\let\ifpl\@secondoftwo}%
  % hyperref drops \color from a bookmark and keeps its ARGUMENT, so a title
  % carrying \code{} would bookmark as "mtealx". Gobble the colour name.
  \def\color#1{}%
}
\makeatother

% ---------- Language strings ----------
\input{lang/\booklang}

% The PDF's subject, keywords and language read \lblPdf... from the language
% file, so they cannot sit in the colour block above, which runs first. Two
% blocks is the ordering, not an oversight.
\hypersetup{
  pdfsubject={\lblPdfSubject},
  pdfkeywords={\lblPdfKeywords},
  pdflang={\lblPdfLang}
}

% ---------- Headings ----------
\titleformat{\chapter}[display]
  {\normalfont\huge\bfseries\color{mblue}\raggedright}
  {\normalfont\normalsize\color{mgrey}\MakeUppercase{\chaptertitlename}\ \thechapter}
  {8pt}{\Huge\raggedright}
\titlespacing*{\chapter}{0pt}{10pt}{28pt}
\titleformat{\section}{\normalfont\Large\bfseries\color{mblue}\raggedright}{\thesection}{0.8em}{}
\titleformat{\subsection}{\normalfont\large\bfseries\color{mgrey}\raggedright}{\thesubsection}{0.7em}{}

\setcounter{tocdepth}{1}
\setcounter{secnumdepth}{2}

\pagestyle{fancy}
\fancyhf{}
\fancyhead[L]{\small\nouppercase{\leftmark}}
\fancyhead[R]{\small\thepage}
\renewcommand{\headrulewidth}{0.4pt}
\fancypagestyle{plain}{\fancyhf{}\fancyfoot[R]{\small\thepage}%
  \renewcommand{\headrulewidth}{0pt}}
% After \pagestyle{fancy}, which installs its own \chaptermark.
\renewcommand{\chaptermark}[1]{\markboth{\thechapter.\ #1}{}}

% ---------- Mathematics ----------
% The handful of notations the book uses, as macros so both editions set the
% same symbol. The book writes \ln for the natural logarithm and \log_{10}
% where the base is ten; a bare \log is ambiguous between the two readerships
% and is not used.
\newcommand{\R}{\mathbb{R}}
\newcommand{\dd}{\,\mathrm{d}}
\newcommand{\abs}[1]{\left\lvert #1\right\rvert}
\newcommand{\iu}{\mathrm{i}}

% ============================================================
%  CODE LISTINGS
%  Every listing in a chapter is a FILE under code/, pulled in with \pyfile or
%  \pyregion, and CI runs it on the pinned interpreter. An in-source
%  \begin{python} block is allowed for a fragment of three or four lines that
%  shows a shape rather than runs.
% ============================================================
\lstdefinelanguage{pythonx}{
  language=Python,
  morekeywords={async,await,match,case,nonlocal,yield,type},
  morekeywords=[2]{
    orbit,iterate,logistic,rk4_step,euler_step,integrate,lorenz,henon,
    lyapunov,separation,escape_time,box_count
  },
  sensitive=true
}

\lstdefinestyle{bookcode}{
  basicstyle=\ttfamily\footnotesize,
  keywordstyle=\color{kw}\bfseries,
  keywordstyle=[2]\color{typ},
  stringstyle=\color{str},
  % Colour only: inconsolata has no italic, and an italic comment style is
  % silently substituted and logged on every build.
  commentstyle=\color{cmt},
  numbers=left,
  numberstyle=\ttfamily\tiny\color{mgrey},
  numbersep=8pt,
  backgroundcolor=\color{codebg},
  frame=single,
  rulecolor=\color{codeframe},
  framesep=6pt,
  breaklines=true,
  breakatwhitespace=false,
  postbreak=\mbox{\textcolor{mgrey}{$\hookrightarrow$}\space},
  showstringspaces=false,
  tabsize=4,
  columns=fullflexible,
  keepspaces=true,
  captionpos=b,
  aboveskip=1.2em,
  belowskip=1.2em,
  % Do NOT add a mapping for U+00A0: it aborts the run with a message that
  % names neither the character nor this line. Listings are ASCII, and
  % parity's C13 checks it.
  literate={ł}{{\l}}1 {Ł}{{\L}}1 {ą}{{\k a}}1 {ę}{{\k e}}1
           {ś}{{\'s}}1 {ć}{{\'c}}1 {ż}{{\.z}}1 {ź}{{\'z}}1 {ó}{{\'o}}1 {ń}{{\'n}}1
           {Ą}{{\k A}}1 {Ę}{{\k E}}1 {Ś}{{\'S}}1 {Ć}{{\'C}}1
           {Ż}{{\.Z}}1 {Ź}{{\'Z}}1 {Ó}{{\'O}}1 {Ń}{{\'N}}1
           {—}{{-{}-}}2 {–}{{-}}1 {‘}{{'}}1 {’}{{'}}1
           {“}{{"}}1 {”}{{"}}1 {°}{{\textdegree}}1
}
\lstset{style=bookcode}

\lstnewenvironment{python}[1][]
  {\lstset{language=pythonx,style=bookcode,#1}}
  {}
\lstnewenvironment{shellcmd}[1][]
  {\lstset{language=bash,style=bookcode,numbers=none,keywordstyle=\color{mteal}\bfseries,#1}}
  {}

% The path is printed above every file listing, in the form the reader opens
% in the editor. \nobreak keeps it on the page of the listing it names.
% The path must not be stranded at the foot of a page with its listing
% overleaf: \nobreak alone does not hold, because listings opens with its own
% vertical material. \needspace asks for room for the path and the first few
% lines of code, and turns the page first when there is not.
% \Needspace (capital N) measures \pagegoal-\pagetotal exactly; lower-case
% \needspace works through glue and penalties and was measured, by the
% Chapter 11 pass, turning pages early and leaving gaps. It must be used
% between paragraphs, which \listingpath's own \par guarantees. Chapters use
% \Needspace* too, so the fallback for a machine without the package defines
% both forms (the star swallowed with its argument).
\makeatletter
\IfFileExists{needspace.sty}{\usepackage{needspace}}{%
  \providecommand{\needspace}[1]{}%
  \providecommand{\Needspace}{\@ifstar\@gobble\@gobble}}
\makeatother
\newcommand{\listingpath}[1]{%
  \par\medskip\Needspace{6\baselineskip}\noindent
  {\ttfamily\footnotesize\color{mgrey}\detokenize{#1}}%
  \par\nobreak}

% Usage: \pyfile{code/ch05/cobweb.py}{Caption}{lst:label}
\newcommand{\pyfile}[3]{%
  \listingpath{#1}%
  \lstinputlisting[language=pythonx,style=bookcode,aboveskip=0.3em,
    caption={#2},label={#3}]{#1}%
}

% A named region of a file, delimited by `# --8<-- [start:name]` and
% `# --8<-- [end:name]`. Matched through listings' range prefix and suffix
% keys -- the other obvious form matches nothing and prints the WHOLE file
% without a warning (measured in the sibling book). A region name is a word,
% never a number. The line numbers printed are the file's own.
% Usage: \pyregion{code/ch05/cobweb.py}{step}{Caption}{lst:label}
\lstset{
  rangebeginprefix=\#\ --8<--\ [start:, rangebeginsuffix=],
  rangeendprefix=\#\ --8<--\ [end:, rangeendsuffix=],
  includerangemarker=false}
\newcommand{\pyregion}[4]{%
  \listingpath{#1}%
  \lstinputlisting[language=pythonx,style=bookcode,aboveskip=0.3em,
    linerange={#2-#2},caption={#3},label={#4}]{#1}%
}

% A program's output, written by code/measure/transcripts.py. There is no
% typed form and there must not be: \val{} does not expand in a verbatim
% body, and a transcript nobody ran looks exactly like one that was.
\newcommand{\transcript}[1]{%
  \par\smallskip
  \IfFileExists{figures/transcripts/#1.txt}{%
    \lstinputlisting[style=bookcode,numbers=none,basicstyle=\ttfamily\footnotesize,
      keywordstyle=\color{black},backgroundcolor=\color{white},
      language={}]{figures/transcripts/#1.txt}%
  }{%
    \begin{tcolorbox}[enhanced, sharp corners, width=\linewidth,
      colback=mlight, colframe=mgrey, boxrule=0.6pt,
      left=8pt, right=8pt, top=6pt, bottom=6pt]
      {\bfseries\color{mgrey}\small \lblTranscriptMissing}\\[2pt]
      {\ttfamily\scriptsize\color{mgrey} figures/transcripts/#1.txt}
    \end{tcolorbox}%
  }%
  \par\smallskip
}

% Inline literals. Underscores inside these must be written \_.
\newcommand{\code}[1]{\texttt{\small #1}}
\newcommand{\pkg}[1]{\texttt{\small\color{mblue}#1}}
\newcommand{\api}[1]{\texttt{\small\color{mteal}#1}}

% ============================================================
%  ADMONITIONS
%  A small, fixed vocabulary. `tinted' may be read in passing; `outlined' is
%  a block the reader must stop at. A third treatment would mean the first two
%  are not carrying their meaning.
% ============================================================
\tcbset{
  cfz tinted/.style={
    enhanced, breakable, sharp corners,
    lines before break=4,
    colback=mlight, boxrule=0pt, leftrule=3pt,
    left=8pt, right=8pt, top=6pt, bottom=6pt,
    before skip=13pt, after skip=13pt,
    fonttitle=\bfseries\small,
    coltitle=black, colbacktitle=mlight,
    attach title to upper={\par\smallskip},
  },
  cfz outlined/.style={
    enhanced, breakable, sharp corners,
    % A box broken straight after its title is a heading stranded at the foot
    % of a page with its question overleaf. Four lines is title plus three.
    lines before break=4,
    colback=white, boxrule=0.6pt,
    left=8pt, right=8pt, top=6pt, bottom=6pt,
    before skip=13pt, after skip=13pt,
    fonttitle=\bfseries\small,
    coltitle=black, colbacktitle=white,
    attach title to upper={\par\smallskip},
  },
}

% Read in passing.
\newtcolorbox{note}[1][]{cfz tinted, colframe=mblue, title={\lblNote}, #1}
\newtcolorbox{warning}[1][]{cfz tinted, colframe=mred, title={\lblWarning}, #1}
% Where you meet this in the world. If it cannot name a specific system,
% instrument, paper or event, it is prose and not this box.
\newtcolorbox{worldbox}[1][]{cfz tinted, colframe=mteal, title={\lblWorld}, #1}
% What is stated and not proved here, and where the proof lives.
\newtcolorbox{rigourbox}[1][]{cfz tinted, colframe=mgrey, title={\lblRigour}, #1}
% The companion library the reader builds, under code/src/chaoslab/.
\newtcolorbox{projectbox}[1][]{cfz tinted, colframe=mamber, title={\lblProject}, #1}

% Stop and read.
% The misconception, named AFTER the reader has been asked the question it
% gets wrong. Appendix B is the catalogue; a trap box is one entry in place.
\newtcolorbox{trapbox}[1][]{cfz outlined, colframe=mred, title={\lblTrap}, #1}
% A listing that has not been run on the pinned versions. Counted by
% `make debt'; nothing ships to a reader with one attached.
\newtcolorbox{verifybox}[1][]{cfz outlined, colframe=mgrey, title={\lblVerify}, #1}
\newtcolorbox{labbox}[1][]{cfz outlined, colframe=mgreen, title={\lblLab}, #1}

% ============================================================
%  THE METHOD --- questions before answers
%  ---------------------------------------------------------
%  The unit of teaching is a question the reader answers BEFORE reading the
%  answer. That is retrieval practice and errorful generation, the two
%  findings the programmed-learning format is defended by, and it is the one
%  thing the sibling mathematics book does that this one keeps whole.
%
%  \begin{think} ... \end{think}   the question, numbered per chapter, ending
%                                  with the instruction to answer before
%                                  reading on;
%  \reveal{short answer}           a centred box with the answer, which is the
%                                  first thing after the question -- the edge
%                                  the reader covers with a hand, or scrolls
%                                  to;
%  \begin{revealblock} ... \end{revealblock}  the same, for a worked answer.
%
%  tools/check_structure.py --think fails when a think box is not followed by
%  a reveal before the next think box or section, and when a reveal answers
%  nothing. That is a property of the source, so it is checked mechanically.
%
%  The instruction says "before you read on", never "before you turn over":
%  the two editions paginate differently, so the second is true in one of them.
% ============================================================
\newcounter{think}[chapter]
\renewcommand{\thethink}{\thechapter.\arabic{think}}
\newtcolorbox{thinkbox}[1][]{cfz outlined, colframe=mviolet, #1}
\NewDocumentEnvironment{think}{}{%
  \refstepcounter{think}%
  \begin{thinkbox}[title={\lblThink~\thethink}]%
}{%
  \par\smallskip
  {\raggedleft\itshape\small\color{mviolet}\lblThinkCue\par}%
  \end{thinkbox}%
}
\tcbset{
  cfz answer/.style={
    enhanced, sharp corners,
    colback=mfaint, colframe=mviolet!55, boxrule=0.5pt,
    before skip=4pt, after skip=10pt,
  },
}
\newcommand{\reveal}[1]{%
  \par\nobreak
  \begin{tcolorbox}[cfz answer, unbreakable,
    left=8pt, right=8pt, top=5pt, bottom=5pt,
    width=0.86\linewidth, center, halign=flush center]
    {\small\bfseries\color{mviolet}\lblAnswerMark}\enspace #1
  \end{tcolorbox}%
  \nobreak\noindent\ignorespaces
}
\NewDocumentEnvironment{revealblock}{}{%
  \par\nopagebreak
  \begin{tcolorbox}[cfz answer, breakable,
    left=10pt, right=10pt, top=6pt, bottom=6pt]
  {\small\bfseries\color{mviolet}\lblAnswerMark}\par\smallskip
}{%
  \end{tcolorbox}%
  \nopagebreak\par\noindent\ignorespaces
}

% ---- Outcomes on entry, and "Can you?" generated from them ----------------
% The self-check at the end of a chapter IS the list of outcomes at its
% start, one for one, because it is built from the same tokens: each
% \outcome both prints its item and appends a row to \cfz@canrows. A \loop
% inside a tabularx body does not survive alignment scanning, which is why
% the rows are accumulated at declaration time (the sibling book's finding).
\makeatletter
\newcommand{\cfz@canrows}{}
\newcommand{\cfz@ratebox}{%
  \fbox{\rule{0pt}{0.9ex}\rule{0.9ex}{0pt}}}
\newcommand{\cfz@rates}{%
  {\footnotesize\color{mgrey}%
   1\,\cfz@ratebox\ 2\,\cfz@ratebox\ 3\,\cfz@ratebox\ 4\,\cfz@ratebox\ 5\,\cfz@ratebox}}
\NewDocumentEnvironment{outcomes}{}{%
  \gdef\cfz@canrows{}%
  \begin{tcolorbox}[cfz tinted, colframe=mblue, title={\lblOutcomes}]%
  \lblOutcomesLead
  \begin{itemize}[leftmargin=1.4em, itemsep=2pt, topsep=4pt]%
}{%
  \end{itemize}%
  \end{tcolorbox}%
}
\newcommand{\outcome}[1]{%
  \item #1%
  \gappto\cfz@canrows{#1 & \cfz@rates \\ \addlinespace[3pt]}%
}
\newcommand{\canyou}{%
  \section*{\lblCanYou}%
  \phantomsection\addcontentsline{toc}{section}{\lblCanYou}%
  \noindent\lblCanYouIntro\par\medskip
  \ifdefempty{\cfz@canrows}{%
    \noindent{\color{mred}\lblNoOutcomes}\par}{%
  {\small
  \noindent\begin{tabularx}{\linewidth}{@{}>{\raggedright\arraybackslash}X r@{}}
    \toprule
    \cfz@canrows
    \bottomrule
  \end{tabularx}\par}}%
}
\makeatother

% ---- Summary -------------------------------------------------------------
\newtcolorbox{summarybox}[1][]{cfz tinted, colframe=mblue, title={\lblSummary}, #1}

% ---- Checkpoint questions, answered at the back --------------------------
% Answers are written AT THE QUESTION with \answerto{} and printed in
% Appendix A. They are carried there by a global token list rather than an
% auxiliary file, because an answer contains mathematics and \val{}, and
% material written to an .aux is expanded at shipout, where fragile commands
% break with errors naming neither the answer nor the chapter.
% \include is sequential within a run, so the list is complete by the time
% the appendix is typeset.
\makeatletter
\newcommand{\cfz@answers}{}
\newcommand{\cfz@lastansch}{0}
\NewDocumentEnvironment{checkpoint}{}{%
  \section*{\lblCheckpoint}%
  \phantomsection\addcontentsline{toc}{section}{\lblCheckpoint}%
  \noindent\lblCheckpointIntro\par
  \begin{enumerate}[label=\textbf{\arabic*.}, leftmargin=2em, itemsep=4pt]%
}{%
  \end{enumerate}%
}
\newcommand{\answerto}[1]{%
  \ifnum\value{chapter}=\cfz@lastansch\relax\else
    \xdef\cfz@lastansch{\arabic{chapter}}%
    \xappto\cfz@answers{\noexpand\cfz@anschapter{\thechapter}}%
  \fi
  \xappto\cfz@answers{\noexpand\cfz@ansitem{\arabic{enumi}}}%
  \gappto\cfz@answers{{#1}}%
}
\newcommand{\cfz@anschapter}[1]{%
  \par\medskip\noindent{\bfseries\color{mblue}\chaptername~#1}\par\nobreak\smallskip}
\newcommand{\cfz@ansitem}[2]{%
  \par\noindent\hangindent=2em\hangafter=1
  \makebox[2em][l]{\textbf{#1.}}#2\par}
\newcommand{\answersbody}{%
  \ifdefempty{\cfz@answers}{\noindent{\color{mred}\lblNoAnswers}\par}{%
    {\raggedright\cfz@answers}}%
}
\makeatother

% ============================================================
%  LABS
%  \begin{lab}{key}{title} ... \end{lab}
%
%  A lab is not a paragraph but a FILE the reader opens: a starter under
%  code/labs/chNN/<key>.py that must FAIL its test, a reference solution under
%  code/labs/chNN/solutions/, and test_<key>.py beside it. The box prints the
%  paths and the one command that runs the test. An environment rather than a
%  macro, because a macro argument cannot hold a verbatim fragment.
%  tools/check_structure.py --labs fails when a file is missing or the key's
%  chapter number disagrees with its chapter.
% ============================================================
\newcounter{lab}[chapter]
\renewcommand{\thelab}{\thechapter.\arabic{lab}}
\makeatletter
\newcommand{\labdir}{ch\two@digits{\value{chapter}}}
\newcommand{\listoflabs}{%
  \section*{\lblLabManifest}%
  \phantomsection\addcontentsline{toc}{section}{\lblLabManifest}%
  \begingroup\c@tocdepth=2\relax\@starttoc{lab}\endgroup%
}
\makeatother
\NewDocumentEnvironment{lab}{mm}{%
  \refstepcounter{lab}%
  \addcontentsline{lab}{subsection}{\protect\numberline{\thelab}\texttt{\protect\detokenize{#1}} --- #2}%
  \begin{labbox}[title={\lblLab~\thelab\ \dash{} #2}]%
}{%
  \par\medskip
  {\small\setlength{\parindent}{0pt}%
  \lblLabStarter: \texttt{code/labs/\labdir/\detokenize{#1}.py}\\
  \lblLabTest: \texttt{code/labs/\labdir/test\_\detokenize{#1}.py}\\
  \lblLabRun: \texttt{cd code \&\& uv run pytest -k \detokenize{#1}}\par}%
  \end{labbox}%
}

% ============================================================
%  FIGURES --- plots and diagrams
%  ---------------------------------------------------------
%  \plotfig{key}{caption}{one-line manifest description}
%    A PLOT, generated by code/plots/<chapter>.py from the same library the
%    listings use, into figures/plots/<lang>/<key>.pdf. Per language, because
%    an axis label in the wrong language is the half-translation that makes a
%    bilingual book look machine-made, and because the Polish edition owes a
%    decimal comma on its tick labels too. Build output, gitignored.
%
%  \mermaidfig{key}{caption}{one-line manifest description}
%    A DIAGRAM of an idea rather than of data, from
%    figures/mermaid/<lang>/<key>.mmd, rendered by `make diagrams'.
%
%  If the rendered PDF is absent either macro prints a visible marker in the
%  flow rather than a float, so a missing figure cannot be mistaken for a
%  layout choice. \phantomsection before \addcontentsline is load-bearing:
%  without it the manifest entry records whatever anchor preceded the figure,
%  which after a listing is a line that was never printed.
% ============================================================
\makeatletter
\newcommand{\listofplots}{%
  \section*{\lblPlotManifest}%
  \phantomsection\addcontentsline{toc}{section}{\lblPlotManifest}%
  \begingroup\c@tocdepth=2\relax\@starttoc{plt}\endgroup%
}
\newcommand{\listofdiagrams}{%
  \section*{\lblDiagramManifest}%
  \phantomsection\addcontentsline{toc}{section}{\lblDiagramManifest}%
  \begingroup\c@tocdepth=2\relax\@starttoc{dgm}\endgroup%
}
\makeatother

\newcommand{\figmissing}[3]{% #1 message, #2 path, #3 caption
  \par\medskip
  \begin{tcolorbox}[enhanced, breakable, width=\linewidth, sharp corners,
    colback=white, colframe=mgrey, boxrule=0.8pt,
    left=8pt, right=8pt, top=8pt, bottom=8pt]
    {\bfseries\color{mgrey}\small #1}\\[4pt]
    {\ttfamily\scriptsize\color{mgrey} #2}
  \end{tcolorbox}%
  \captionof{figure}{#3}%
  \par\medskip
}

\newcommand{\plotfig}[3]{%
  \phantomsection%
  \addcontentsline{plt}{subsection}{\protect\numberline{}\texttt{#1} --- #3}%
  \IfFileExists{figures/plots/\booklang/#1.pdf}{%
    \begin{figure}[htbp]
    \centering
    \includegraphics[width=\linewidth,height=0.45\textheight,keepaspectratio]{figures/plots/\booklang/#1.pdf}%
    \caption{#2}\label{fig:#1}
    \end{figure}%
  }{%
    \figmissing{\lblPlotNotRendered}{figures/plots/\booklang/#1.pdf}{#2}\label{fig:#1}%
  }%
}

\newcommand{\mermaidfig}[3]{%
  \phantomsection%
  \addcontentsline{dgm}{subsection}{\protect\numberline{}\texttt{#1.mmd} --- #3}%
  \IfFileExists{figures/diagrams/\booklang/#1.pdf}{%
    \begin{figure}[htbp]
    \centering
    \includegraphics[width=0.95\linewidth,height=0.42\textheight,keepaspectratio]{figures/diagrams/\booklang/#1.pdf}%
    \caption{#2}\label{fig:#1}
    \end{figure}%
  }{%
    \figmissing{\lblNotRendered}{figures/mermaid/\booklang/#1.mmd}{#2}\label{fig:#1}%
  }%
}

% ============================================================
%  APPENDIX B: the misreadings, generated by tools/gen_traps.py
%
%  \trapchapter{Chapter~\ref{ch:x}}{title}    a heading per chapter
%  \trapentry{n}{habit}{truth}{\ref{sec:x}}  one misreading
%
%  Set ragged right: an entry is mostly short sentences with \code{} spans in
%  them, and a justified line full of unbreakable runs is an overfull box
%  waiting for the edition with the longer words.
% ============================================================
\newcommand{\trapchapter}[2]{%
  \par\medskip
  {\large\bfseries\color{mblue}#1\quad\normalfont\itshape\color{mgrey}#2\par}%
  \nopagebreak\smallskip}
\newcommand{\trapentry}[4]{%
  \par\noindent
  \begingroup\raggedright\hangindent=2.4em\hangafter=1
  \makebox[2.4em][l]{\textbf{\color{mred}#1}}%
  \textbf{\enquote{#2}}\enspace #3\enspace{\color{mgrey}\S\,#4}\par
  \endgroup\smallskip}

% ============================================================
%  CHAPTER STUBS
% ============================================================
\newcommand{\chapterstub}[1]{%
  \begin{tcolorbox}[
    enhanced, breakable, width=\linewidth, sharp corners,
    colback=mlight, colframe=mred, boxrule=1pt,
    borderline={0.6pt}{3pt}{mred, dashed},
    left=12pt, right=12pt, top=12pt, bottom=12pt]
    {\bfseries\color{mred}\large \lblNotWritten}\\[8pt]
    \small\raggedright #1
  \end{tcolorbox}%
}

% ============================================================
%  COMPUTED VALUES
%  A number the reader cannot do in their head is computed by a script under
%  code/measure/, written to figures/values/<chapter>.tex as
%      \pyval{ch06.steps.apart}{34}
%  and reaches the page as \val{ch06.steps.apart}. The book contains
%  references, not digits; `make verify' fails when a committed value no
%  longer matches its script. \val applies the edition's decimal marker, which
%  is how the Polish edition gets its comma without anybody remembering.
% ============================================================
% The comma is BRACED: siunitx sets numbers in maths mode, where a bare comma
% is punctuation and gets a thin space after it ("9, 81"). {,} is an ordinary
% symbol, which is what a decimal marker is.
\newcommand{\bookdecimalmarker}{\ifpl{{,}}{.}}
\IfFileExists{siunitx.sty}{%
  \usepackage{siunitx}%
  \sisetup{
    group-digits = integer,
    group-minimum-digits = 5,
    group-separator = {\,},
    output-decimal-marker = {\bookdecimalmarker},
    exponent-product = \times,
    print-unity-mantissa = false
  }%
  \newcommand{\booknum}[1]{\num{##1}}%
}{%
  \newcommand{\booknum}[1]{##1}%
  \providecommand{\num}[1]{##1}%
}

\makeatletter
\newcommand{\pyvalues}{}
\newcommand{\pyval}[2]{\csgdef{pyval@#1}{#2}\listxadd{\pyvalues}{#1}}
\newcommand{\pyvaltext}[2]{%
  \csgdef{pyval@#1}{#2}\csgdef{pyvaltext@#1}{}\listxadd{\pyvalues}{#1}}
% \num on a non-numeric body is a FATAL siunitx error, so \val refuses a text
% value by name rather than letting siunitx fail on it.
\newcommand{\val}[1]{%
  \ifcsdef{pyval@#1}{%
    \ifcsdef{pyvaltext@#1}{%
      \PackageError{chaosbook}{Value `#1' is not a number}%
        {Use \string\valtext{#1} instead.}%
    }{}%
    \booknum{\csuse{pyval@#1}}%
  }{\textcolor{mred}{\textbf{[\lblValueMissing: #1]}}}%
}
\newcommand{\valtext}[1]{%
  \ifcsdef{pyval@#1}{\csuse{pyval@#1}}%
    {\textcolor{mred}{\textbf{[\lblValueMissing: #1]}}}%
}
\makeatother

% figures/values/all.tex is generated by `make numbers' and \inputs each
% value file. A build with no values still compiles; every \val{} then prints
% a visible marker. \valuesindex lets the one-chapter build point elsewhere.
\providecommand{\valuesindex}{figures/values/all}
\IfFileExists{\valuesindex.tex}{\input{\valuesindex}}{%
  \AtBeginDocument{\PackageWarning{chaosbook}{No computed values found. Run `make numbers'.}}}

% ============================================================
%  PINNED VERSIONS --- single source of truth
%  Change these here AND in code/pyproject.toml; tools/check_structure.py
%  --pins fails when the two disagree. A chapter never types a version.
% ============================================================
\newcommand{\bookedition}{0.1}
\newcommand{\bookyear}{2026}
\newcommand{\pyver}{3.14}              % the minor the book targets
\newcommand{\pypatch}{3.14.7}          % the exact interpreter the code ran on
\newcommand{\uvver}{0.12.13}           % uv
\newcommand{\numpyver}{2.5.3}          % numpy
\newcommand{\matplotlibver}{3.11.2}    % matplotlib
\newcommand{\pytestver}{9.1.1}         % pytest
\newcommand{\ruffver}{0.16.9}          % ruff
% The prose licence is a name with a version in it, not a quantity: it is a
% macro so that neither edition typesets it as a decimal to be localised.
\newcommand{\proselicence}{CC~BY-NC-SA~4.0}

% ============================================================
%  INDEX
% ============================================================
\apptocmd{\theindex}{\raggedcolumns\raggedright}{}{%
  \PackageWarning{chaosbook}{Could not patch theindex for ragged setting}}
\providecommand*\see[2]{}
\renewcommand*\see[2]{\emph{\lblIndexSee} #1}
\providecommand*\seealso[2]{}
\renewcommand*\seealso[2]{\emph{\lblIndexSeeAlso} #1}

\makeindex
